% !TEX program = lualatex
% ======================================================================
%  Quantification in a Grammar without Variable Names
%  Scope, Storage and Indices after Inclusiveness
%  Keio Linguistics Colloquium, 7-8 November 2026 -- the seminar volume
%
%  One file: the four lectures, the exercises (in an appendix), the glossary
%  and the references; the Keio announcement is reproduced at the end.
%
%  Build:   latexmk
%           (lualatex -> bibtex -> lualatex x2; gtbook, gtling and minsyn
%           in tex/, found through latexmkrc; keio.bib, keio-slides.bib and
%           glossa.bst in this folder; keio-extra.bib is written by this
%           file.)
% ======================================================================

\documentclass[
  layout=tufte,           % wide outer margin; footnotes become sidenotes
  twoside=false, open=any,
  font=pagella,
  chapterprefix=false,
  numbers=noendperiod,    % no full stop after section numbers, appendix or not
  headings=optiontohead,  % \chapter[head=...]
  parskip=half-,          % no 1em reserve at the end of paragraphs
  colorlinks=true,
  examples=gb4e,
  minsyn=true,            % workspaces, boxes, Self-Merge trees (loads forest)
  bib=natbib, bibstyle=glossa, bibfile=keio,
  greek=false,
]{gtbook}

\usepackage{pdfpages}
\usepackage{etoc}
\usepackage{adjustbox}
\usetikzlibrary{shapes.misc,fit,backgrounds,positioning,calc,arrows.meta,decorations.pathreplacing}
\tcbuselibrary{skins,breakable}

\makeatletter
% ======================================================================
%  1. Colours
% ======================================================================
% The slides' colblanc names, so that diagrams carry over unchanged.
\colorlet{colblancaccent}{gtaccent}
\colorlet{colblancalert}{gtalert}
\colorlet{colblancghost}{gtghost}
\colorlet{colblancexbg}{gtexbg}
\colorlet{colblancbg}{gtbg}
\colorlet{colblanctext}{gttext}
% minsyn workspaces in the volume's palette, as on the slides
\colorlet{minsyn@wsbg}{gtexbg}
\colorlet{minsyn@wsfr}{gtaccent}
\tcbset{minsyn@wsbox/.append style={left=5pt,right=5pt,top=8pt,bottom=5pt,
  before skip=6pt, after skip=4pt, coltitle=gttext,
  boxed title style={colback=white, colframe=minsyn@wsfr, arc=3pt, boxrule=0.6pt,
    top=1.5pt, bottom=1.5pt, left=5pt, right=5pt}}}

% ======================================================================
%  2. Page: the margin, wide material, tables
% ======================================================================
% the text column is narrow: let TeX stretch a hard paragraph rather than
% push a line into the margin
\setlength\emergencystretch{1.5em}
\newlength\keio@overhang
\newlength\keio@pagewidth   % the text column, as set by the class
\AtBeginDocument{\setlength\keio@overhang{\dimexpr\marginparsep+\marginparwidth\relax}%
  \setlength\keio@pagewidth{\textwidth}}
% wide: text width plus the margin (one-sided, so the margin is on the right)
\newenvironment{wide}
  {\par\begin{adjustwidth}{0pt}{-\keio@overhang}}
  {\end{adjustwidth}\par}
% A tabularx at full text width runs over the margin automatically;
% one already inside a narrower or wider measure is left alone.
\newcolumntype{L}{>{\raggedright\arraybackslash}X}
\newcolumntype{P}[1]{>{\raggedright\arraybackslash}p{#1}}
\let\keio@tabularx\tabularx
\newif\ifkeio@widetab
\def\tabularx#1{%
  \ifkeio@widetab\expandafter\@firstoftwo\else\expandafter\@secondoftwo\fi
  {\keio@tabularx{\dimexpr#1+\keio@overhang\relax}}{\keio@tabularx{#1}}}
\BeforeBeginEnvironment{tabularx}{%
  \ifdim\linewidth=\keio@pagewidth
    \keio@widetabtrue\leavevmode\hbox to\linewidth\bgroup
  \else\keio@widetabfalse\fi}
\AfterEndEnvironment{tabularx}{\ifkeio@widetab\hss\egroup\fi\keio@widetabfalse}
\setlist{itemsep=1pt,topsep=3pt,parsep=0pt}
\renewcommand*{\arraystretch}{1.15}

% ======================================================================
%  3. Lectures (chapters), sections, running heads
% ======================================================================
\renewcommand\thesection{\arabic{section}}      % §3, not §1.3
\setcounter{tocdepth}{\sectiontocdepth}
\setcounter{secnumdepth}{\subsectionnumdepth}
\renewcaptionname{british}{\chaptername}{Lecture}
\RedeclareSectionCommand[beforeskip=0pt, afterskip=1.4\baselineskip,
  font=\normalfont\bfseries\LARGE]{chapter}
\RedeclareSectionCommand[beforeskip=-2.4ex plus -.6ex minus -.3ex,
  afterskip=1ex plus .2ex, font=\normalfont\bfseries\Large]{section}
\RedeclareSectionCommand[beforeskip=-1.8ex plus -.4ex minus -.2ex,
  afterskip=.6ex plus .1ex, font=\normalfont\bfseries\large]{subsection}
% Lecture heading: a small-caps label above the title, and the numeral, large
% and pale, hanging in the margin.
\renewcommand\chapterlinesformat[3]{%
  \Ifstr{#1}{chapter}{%
    \IfArgIsEmpty{#2}{#3}{%
      \makebox[0pt][l]{\hspace*{\dimexpr\textwidth+\marginparsep\relax}%
        \raisebox{-0.62\height}[0pt][0pt]{\fontsize{96}{96}\selectfont
          \color{gtaccent!22}\thechapter}}%
      {\normalsize\mdseries\scshape\color{gtaccent}\letterspace{\chapapp~\thechapter}\par}%
      \vspace{0.4ex}#3}}%
    {\@hangfrom{#2}{#3}}}
\providecommand\letterspace[1]{#1}
\renewcommand*\chaptermarkformat{\chapapp~\thechapter\ \textperiodcentered\ }
\automark{chapter}
\clearpairofpagestyles
\ihead{{\scshape Tsoulas} \textperiodcentered{} {\scshape Keio, 7--8 November 2026}}
\ohead{\headmark}
\cfoot*{\pagemark}
\addtokomafont{pagehead}{\upshape}

% label namespaces per lecture: \label{ex:x} in Lecture 2 is H2:ex:x
\def\keio@pfx{}
\AddToHook{begindocument/end}{%
  \let\keio@label\label
  \let\keio@ref\ref
  \renewcommand\label[1]{\keio@label{\keio@pfx#1}}%
  \renewcommand\ref[1]{\keio@ref{\keio@fullname{#1}}}}
\def\keio@fullname#1{\keio@fullname@#1\@empty\@empty\@empty\keio@nil}
\def\keio@fullname@#1#2#3#4\keio@nil{%
  \ifx H#1\ifx :#3\else\keio@pfx\fi\else\keio@pfx\fi#1#2#3#4}
% cross-references to other lectures:  \Hx{1}{ex:desideratum} -> Lecture 1 (38)
%                                      \Hs{1}{sec:copy}       -> Lecture 1, §3.2
\newcommand{\Hx}[2]{Lecture~#1~(\keio@ref{H#1:#2})}
\newcommand{\Hs}[2]{Lecture~#1, \S\keio@ref{H#1:#2}}
\newcommand{\Hp}[1]{Lecture~#1}

% ---------------------------------------------------------------- progress strip
% The four lectures as a strip of cells; the current one is filled.
\def\keio@ptitle#1{\ifcase#1\or A grammar that cannot write \textit{x}\or
  Quantification on the clausal spine\or Scope in the box theory\or
  Indices after Inclusiveness\fi}
\def\keio@pday#1{\ifcase#1\or Saturday\or Saturday\or Sunday\or Sunday\fi}
\newcommand\keio@progress[1]{%
  \par\noindent\makebox[\textwidth][l]{\begin{tikzpicture}[x=\dimexpr(\textwidth+\keio@overhang)/4\relax]
    \foreach \i in {1,...,4}{%
      \pgfmathtruncatemacro\cur{\i==#1 ? 1 : 0}%
      \ifnum\cur=1
        \node[anchor=north west, minimum width=\dimexpr(\textwidth+\keio@overhang)/4-3pt\relax,
          minimum height=11mm, fill=gtaccent, text=white, inner sep=3pt,
          text width=\dimexpr(\textwidth+\keio@overhang)/4-9pt\relax, align=left,
          font=\sffamily\scriptsize] at (\i-1,0)
          {\textbf{\i}\enspace\keio@pday{\i}\\[1pt]\keio@ptitle{\i}};
      \else
        \node[anchor=north west, minimum width=\dimexpr(\textwidth+\keio@overhang)/4-3pt\relax,
          minimum height=11mm, draw=gtghost!60, text=gtghost, inner sep=3pt,
          text width=\dimexpr(\textwidth+\keio@overhang)/4-9pt\relax, align=left,
          font=\sffamily\scriptsize] at (\i-1,0)
          {\textbf{\i}\enspace\keio@pday{\i}\\[1pt]\keio@ptitle{\i}};
      \fi}
  \end{tikzpicture}}\par}

% \keiolecture{n}{title}{short title for the head}{session line}
\newcommand{\keiolecture}[4]{%
  \cleardoublepage
  \setcounter{chapter}{\numexpr#1-1\relax}%
  \gdef\keio@pfx{H#1:}%
  \chapter[head={#3}]{#2}%
  \vspace*{-0.6\baselineskip}%
  {\small\itshape #4\par}%
  \vspace{0.8\baselineskip}%
  \keio@progress{#1}%
  \vspace{0.4\baselineskip}%
  \setcounter{section}{-1}%
  \setcounter{exx}{0}%
  \setcounter{figure}{0}%
}
% \keioepigraph{text}{source}
\newcommand\keioepigraph[2]{%
  \par\vspace{0.6\baselineskip}%
  {\noindent\hfill\begin{minipage}{0.86\linewidth}\raggedleft\itshape
    #1\par\vspace{0.3ex}{\upshape\small\color{gtghost}#2}\end{minipage}\par}%
  \vspace{0.6\baselineskip}}
% \lecturemap: a contents of this lecture, set in the margin beside the opening
\newcommand\lecturemap{%
  \marginnote{\begin{minipage}{\marginparwidth}\raggedright\footnotesize\sffamily
    {\color{gtaccent}\textsc{In this lecture}}\par\smallskip
    \etocsetnexttocdepth{section}%
    \etocsettocstyle{}{}%
    \etocsetstyle{section}{}{}%
      {{\noindent\hangindent1em\rightskip1.6em plus 1fil\parfillskip-1.6em
        \etocnumber\enspace\etocname\nobreak\hfill\nobreak
        \makebox[1.6em][r]{\etocpage}\par}\smallskip}{}%
    \localtableofcontents
  \end{minipage}}}

% ======================================================================
%  4. The note fields of keio.bib
% ======================================================================
% keio.bib carries note fields of the form \bibchk, \bibchk[...] and
% \bibGT{...}, left over from checking the references. Here each prints
% nothing and removes the full stop that the bibliography style puts after it.
\newcommand{\keio@unskip}{\ifhmode\unskip\fi}
\newcommand{\bibchk}[1][]{\keio@unskip\@gobble}
\newcommand{\bibGT}[1]{\keio@unskip\@gobble}

% Inside the boxes a footnote cannot float to the margin; set it there
% directly instead (same numbering, same look).
\newcommand\keio@boxnotes{\RenewDocumentCommand\footnote{o m}{\sidenote[][0pt]{##2}}}

% ======================================================================
%  5. Boxes
% ======================================================================
% excursus: the slides' optional excursions (other frameworks, logics,
% extra trees); [title]. excursus* runs over the margin.
\tcbset{keio@excursus/.style={enhanced jigsaw, breakable,
    colback=gtexbg!55, colframe=gtghost!70, boxrule=0pt, leftrule=2.4pt,
    arc=0pt, outer arc=0pt, left=8pt, right=8pt, top=5pt, bottom=6pt,
    fontupper=\small, detach title,
    before upper*={\setlength\parskip{0.45\baselineskip}\keio@boxnotes\keio@exhook
      \noindent{\sffamily\footnotesize\color{gtghost}\textsc{Excursus}}\enspace\tcbtitle\par\smallskip},
    fonttitle=\bfseries\small, coltitle=gttext,
    before skip=0.9\baselineskip, after skip=0.9\baselineskip, parbox=false}}
\let\keio@exhook\relax
\newtcolorbox{excursus}[1][]{keio@excursus, title={#1}}
\newtcolorbox{excursus*}[1][]{keio@excursus, title={#1},
  grow to right by=\keio@overhang,
  after upper={\keio@printwidenotes}, code={\let\keio@exhook\keio@widenotes}}
% A wide box covers the margin, so its notes are gathered at its foot,
% numbered in the sidenote sequence.
\newtoks\keio@wnotes
\newcommand\keio@widenotes{\global\keio@wnotes{}%
  \RenewDocumentCommand\footnote{o m}{\textsuperscript{\thesidenote}%
    \edef\keio@n{\thesidenote}\stepcounter{sidenote}%
    \global\keio@wnotes\expandafter{\the\expandafter\keio@wnotes
      \expandafter\keio@wnoteitem\expandafter{\keio@n}{##2}}}}
\newcommand\keio@wnoteitem[2]{\par\noindent{\footnotesize\textsuperscript{#1}\,#2\par}}
\newcommand\keio@printwidenotes{%
  \edef\keio@tmp{\the\keio@wnotes}%
  \ifx\keio@tmp\@empty\else
    \par\smallskip\noindent{\color{gtghost}\rule{4em}{0.4pt}}\par\the\keio@wnotes
  \fi\global\keio@wnotes{}}
% exercise: the exercises of the seminar, collected in the appendix.
% \begin{exercise}{key}{title}{n}{sec:label} ... \end{exercise}: numbered
% Exercise 1, 2, ...; the foot of the box links back to Lecture n, \S sec:label,
% where the exercise belongs. In the text, \toexercise{key} puts a pointer in
% the margin (running text only) and \exref{key} gives 'Exercise 3'; both link.
\newcounter{keioexercise}
\tcbset{keio@exercise/.style={enhanced jigsaw, breakable, colback=white,
  frame hidden, borderline={0.7pt}{0pt}{gtaccent, dash pattern=on 3pt off 2pt},
  arc=2pt, left=8pt, right=8pt, top=6pt, bottom=6pt, fontupper=\small,
  detach title, before upper*={\setlength\parskip{0.45\baselineskip}\keio@boxnotes
    \noindent{\sffamily\footnotesize\color{gtaccent}\textsc{Exercise
    \thekeioexercise}}\enspace\tcbtitle\par\smallskip},
  fonttitle=\bfseries\small, coltitle=gttext,
  before skip=0.9\baselineskip, after skip=1.2\baselineskip, parbox=false}}
\NewDocumentEnvironment{exercise}{m m m m}
  {\par\refstepcounter{keioexercise}\keio@label{XR:#1}%
   \def\keio@exback{\hyperref[H#3:#4]{Lecture~#3, \S\keio@ref*{H#3:#4}}}%
   \begin{tcolorbox}[keio@exercise, title={#2}]}
  {\par\smallskip\noindent{\sffamily\footnotesize\color{gtghost}%
     $\blacktriangleleft$\enspace Where it belongs in the text: \keio@exback}%
   \end{tcolorbox}}
\newcommand\exref[1]{\hyperref[XR:#1]{Exercise~\keio@ref*{XR:#1}}}
\newcommand\toexercise[1]{\marginpar{\raggedright\sffamily\scriptsize
  \hyperref[XR:#1]{\color{gtaccent}$\blacktriangleright$\enspace Exercise~\keio@ref*{XR:#1}}%
  \\{\color{gtghost}in the Appendix}}}
% the appendix: its own label space, its own place in the lists of notes
\newcommand\keioappendix{%
  \cleardoublepage\appendix
  \gdef\keio@pfx{XA:}%
  \chapter[head={The exercises}]{The exercises}\label{keio:exercises}%
  \setcounter{section}{0}\setcounter{exx}{0}}
% the sentences of an exercise, reprinted with their numbers in the text:
% \exagain{1}{ex:someone} gives the label (1) of Lecture 1, linked
\newcommand\exagain[2]{\exi{{(\hyperref[H#1:#2]{\keio@ref*{H#1:#2}})}}}
% the reading asked about, under an item of the judgement sheet
\newcommand\keioreading[1]{\par\nopagebreak{\small\textit{The reading asked about:} #1\par}}

% keybox: the central statements of the argument (hypotheses, the
% desideratum, the proposal's clauses, the conjectures)
\newtcolorbox{keybox}[1][]{enhanced jigsaw, breakable, colback=gtaccent!4,
  frame hidden, borderline west={2pt}{0pt}{gtaccent}, arc=0pt,
  left=4pt, right=6pt, top=1pt, bottom=3pt, before upper={\keio@boxnotes},
  before skip=0.6\baselineskip, after skip=0.6\baselineskip, parbox=false, #1}
% questions: the questions for the discussion at the end of each lecture
\newtcolorbox{questions}[1][Questions for the discussion]{enhanced jigsaw,
  breakable, colback=white, colframe=gtaccent, boxrule=0.6pt, arc=2pt,
  left=8pt, right=8pt, top=6pt, bottom=6pt, fontupper=\small,
  before upper={\keio@boxnotes}, title={#1}, fonttitle=\sffamily\bfseries\small, colbacktitle=gtaccent,
  coltitle=white, toptitle=1pt, bottomtitle=1pt,
  before skip=0.9\baselineskip, after skip=0.9\baselineskip, parbox=false}

% ======================================================================
%  6. Figures: trees, derivations and diagrams
% ======================================================================
% \begin{keiofig}{label}{caption}: in the text column, caption in the margin.
% \begin{keiofig*}{label}{caption}: over the margin, caption beneath.
% Figures do not float: they stay where the argument needs them.
% (\captionof swaps in its own \label, so the lecture prefix is added here.)
\captionsetup{font={small}, labelfont={sf,bf,color=gtaccent}, labelsep=quad,
  justification=raggedright, singlelinecheck=false}
\renewcommand\thefigure{\thechapter.\arabic{figure}}
% The caption hangs beside the figure. If it is taller than the figure it is
% raised, but no higher than the top of the figure (so it cannot run into the
% running head); whatever is still left over is kept free below the figure,
% so that the caption never runs past the foot of the text block.
\newsavebox\keio@figbox \newsavebox\keio@capbox
\newlength\keio@capraise \newlength\keio@caprest
\NewDocumentEnvironment{keiofig}{m m +b}
  {\par\medskip\noindent
   \sbox\keio@figbox{\begin{minipage}{\linewidth}\centering #3\end{minipage}}%
   \sbox\keio@capbox{\begin{minipage}[t]{\marginparwidth}\captionsetup{font={footnotesize}}%
     \captionof{figure}{#2}\keio@label{\keio@pfx#1}\end{minipage}}%
   % overshoot of the caption below the figure
   \setlength\keio@caprest{\dimexpr\dp\keio@capbox-\dp\keio@figbox\relax}%
   \ifdim\keio@caprest<\z@\setlength\keio@caprest{\z@}\fi
   % raise by at most the room between the caption's first line and the figure's top
   \setlength\keio@capraise{\dimexpr\ht\keio@figbox-\ht\keio@capbox\relax}%
   \ifdim\keio@capraise<\z@\setlength\keio@capraise{\z@}\fi
   \ifdim\keio@capraise>\keio@caprest\setlength\keio@capraise{\keio@caprest}\fi
   \addtolength\keio@caprest{-\keio@capraise}%
   \marginnote{\usebox\keio@capbox}[-\keio@capraise]%
   \usebox\keio@figbox\rule[-\dimexpr\dp\keio@figbox+\keio@caprest\relax]{0pt}{0pt}%
   \par\medskip}{}
% the wide figure and its caption are one unbreakable block
\NewDocumentEnvironment{keiofig*}{m m +b}
  {\par\medskip\begin{adjustwidth}{0pt}{-\keio@overhang}%
   \begin{minipage}{\linewidth}\centering #3\par\vspace{2pt}%
   \begin{minipage}{0.82\linewidth}\captionof{figure}{#2}\keio@label{\keio@pfx#1}\end{minipage}%
   \end{minipage}\end{adjustwidth}\par\medskip}{}
% side by side, as the slides' twocol: \begin{twocol}{0.48} ... \midcol ... \end{twocol}
\newlength\keio@leftwd \newlength\keio@rightwd
\newenvironment{twocol}[1]
  {\par\noindent\setlength\keio@leftwd{#1\linewidth}%
   \setlength\keio@rightwd{\dimexpr0.96\linewidth-\keio@leftwd\relax}%
   \begin{minipage}[t]{\keio@leftwd}}
  {\end{minipage}\par}
\newcommand\midcol{\end{minipage}\hfill\begin{minipage}[t]{\keio@rightwd}}

% Trees: the slides' compact styles, sized for the page
\forestset{
  slim/.style={for tree={font=\small, l sep=8pt, s sep=10pt, inner sep=2pt}},
  xslim/.style={for tree={font=\footnotesize, l sep=6pt, s sep=7pt, inner sep=1.5pt}},
  spine/.style={for tree={font=\footnotesize, l sep=3pt, s sep=6pt, inner sep=1pt}},
  tiny/.style={for tree={font=\scriptsize, l sep=2pt, s sep=5pt, inner sep=1pt}},
  flat/.style={for tree={l=0pt, l sep=5pt}},
  ty/.style={font=\footnotesize, text=gtghost},
  default preamble={for tree={edge={gtaccent!55, line width=0.5pt}}},
}
% continue minsyn's W numbering across separate derivation panels
\newcommand{\WSsetstep}[1]{\setcounter{minsyn@step}{#1}}
\newcommand{\WSright}[1][]{\begin{center}$\Longrightarrow$\\[-2pt]{\scriptsize #1}\end{center}}
% workspaces in the book: the set wraps, at \allowbreak in the items
\RenewDocumentCommand{\WSeq}{m}{\par{\centering $W_{\arabic{minsyn@step}} = \{\,#1\,\}$\par}}
\RenewDocumentCommand{\WSeqTree}{m m}{%
  \par{\centering $W_{\arabic{minsyn@step}} = \{\,#1\,\}$\par}%
  \vspace{2pt}{\color{minsyn@wsfr}\hrule height 0.6pt}\vspace{6pt}%
  \par{\centering\begin{forest} #2 \end{forest}\par}}

% ======================================================================
%  7. Linguistic shorthands (as in the handouts and the slides)
% ======================================================================
\newcommand{\uth}{u\ensuremath{\theta}}
\newcommand{\uthr}{\ensuremath{\langle}u\ensuremath{\theta\rangle}}
\newcommand{\vs}{\textit{v}*}
\makeatletter
% boxed element in running text and derivation tables: a light, text-size
% version of minsyn's \BOX (same colour), so that it does not open up the line
\newcommand{\bx}[1]{\tikz[baseline=(keiobx.base)]\node[draw=minsyn@box,
  line width=0.6pt, inner sep=1pt, rounded corners=1pt] (keiobx) {#1};}
\makeatother
\newcommand{\BX}[1]{\BOX{\mathrm{#1}}}
\newcommand{\lbr}[1]{[\textsubscript{#1}\,}
\newcommand{\src}[1]{{\unskip\nobreak\hfil\penalty50\hskip2em\hbox{}\nobreak\hfil
  \mbox{#1}\parfillskip=0pt\finalhyphendemerits=0\par}}
\newcommand{\mstar}{\llap{*}}           % judgement hung in the margin of a glossed line
\newcommand{\mqstar}{\llap{?*}}
\newcommand{\lab}[1]{{\footnotesize\scshape #1}}
\let\eachwordone=\itshape
% the slides' biblatex citation commands, for safety
\providecommand\textcite{\citet}
\providecommand\parencite{\citep}
\newcommand\posscitet[1]{\citeauthor{#1}'s \citeyearpar{#1}}

% ======================================================================
%  8. Bibliography: one list at the end
% ======================================================================
\bibpunct[: ]{(}{)}{; }{a}{}{,}
\setlength{\bibsep}{1.5pt plus 1pt}
\setlength{\bibhang}{1.5em}
\renewcommand{\bibsection}{\addchap{References}}
\renewcommand{\bibfont}{\small}
\newcommand\keioreferences{\bibliography{keio-extra,keio,keio-slides}}
\providecommand{\doi}[1]{\url{https://doi.org/#1}}
% the release: change these for a new version (title page and PDF metadata)
\newcommand\keioversion{1.0}
\newcommand\keioversiondate{November 2026}
\newcommand\keiodoi{10.5281/zenodo.24495824}
\AtBeginDocument{\urlstyle{rm}\hypersetup{
  pdftitle={Quantification in a grammar without variable names: scope, storage and indices after Inclusiveness},
  pdfsubject={Lecture notes, version \keioversion, \keioversiondate, doi:\keiodoi},
  pdfauthor={George Tsoulas}}}
\makeatother

% ======================================================================
%  9. Macros the lectures need
% ======================================================================
\makeatletter
% ---- from Lecture 2
% =====================================================================
% Lecture 2 -- Quantification on the clausal spine
% Day 1, Lecture 2 (Saturday 7 November 2026, 15:45-17:45)
% The handout H2, with everything from the slide deck S2 folded in.
% =====================================================================

% lecture-local macros (to hoist)
% \pIIneedspace{len}: start a new page unless len of space is left
% (needspace's test), so that a figure keeps its caption, and a box
% that opens with a picture does not leave its title stranded.
\makeatletter
\providecommand\pIIneedspace[1]{\par\penalty-100\begingroup
  \setlength\dimen@{#1}\dimen@ii\pagegoal\advance\dimen@ii-\pagetotal
  \ifdim\dimen@>\dimen@ii\ifdim\dimen@ii>\z@\vfil\fi\break\fi\endgroup}
\makeatother

% ---- from Lecture 3
% =====================================================================
% Lecture 3 -- Scope in the box theory
% Day 2, Lecture 3 (Sunday 8 November 2026, 13:30-15:30)
% The handout H3 with the slide deck S3 folded in: every tree, workspace
% derivation, table and diagram of the deck, its slide-only examples, its
% optional frames (as excursus boxes) and the two refinements made after
% the handout was printed (as `revised' text).
% =====================================================================

% lecture-local macros (to hoist)
% minsyn's boxed-workspace equations set the set on one unbreakable line;
% as in the slides, let them wrap at the \allowbreak points in the items.
\makeatletter
\RenewDocumentCommand{\WSeqBox}{m m}{%
  \par{\centering $W^{\arabic{minsyn@step}}_{\minsyn@boxsub{#1}} = \{\,#2\,\}$\par}}
\RenewDocumentCommand{\WSeqBoxTree}{m m m}{%
  \par{\centering $W^{\arabic{minsyn@step}}_{\minsyn@boxsub{#1}} = \{\,#2\,\}$\par}%
  \vspace{2pt}{\color{minsyn@wsfr}\hrule height 0.6pt}\vspace{6pt}%
  \par{\centering\begin{forest} #3 \end{forest}\par}}
\makeatother
\makeatother

% Entries cited in the volume that are in neither keio.bib nor keio-slides.bib.
\begin{filecontents*}[overwrite]{keio-extra.bib}
@incollection{Fox2003,
  author = {Fox, Danny},
  title = {On logical form},
  booktitle = {Minimalist syntax},
  editor = {Hendrick, Randall},
  pages = {82--123},
  publisher = {Blackwell},
  address = {Oxford},
  year = {2003}}

@incollection{KawashimaKitahara1993,
  author = {Kawashima, Ruriko and Kitahara, Hisatsugu},
  title = {On the distribution and interpretation of subjects and their numeral classifiers},
  booktitle = {Proceedings of {S}emantics and {L}inguistic {T}heory 3},
  editor = {Lahiri, Utpal and Wyner, Adam Zachary},
  pages = {97--116},
  publisher = {Cornell University, DMLL Publications},
  address = {Ithaca, NY},
  year = {1993}}

@incollection{Kitahara1993,
  author = {Kitahara, Hisatsugu},
  title = {Numeral classifier phrases inside {DP} and the specificity effect},
  booktitle = {Japanese/{K}orean linguistics 3},
  editor = {Choi, Soonja},
  pages = {171--186},
  publisher = {CSLI Publications},
  address = {Stanford, CA},
  year = {1993}}

@misc{PenkaZeijlstra2005,
  author = {Penka, Doris and Zeijlstra, Hedde},
  title = {Negative indefinites in {D}utch and {G}erman},
  howpublished = {Paper presented at the 20th Comparative Germanic Syntax Workshop, Tilburg},
  year = {2005},
  note = {\bibchk[venue]}}

@article{Quine1960,
  author = {Quine, Willard Van Orman},
  title = {Variables explained away},
  journal = {Proceedings of the American Philosophical Society},
  volume = {104},
  number = {3},
  pages = {343--347},
  year = {1960}}

@article{PollardYoo1998,
  author = {Pollard, Carl and Yoo, Eun Jung},
  title = {A unified theory of scope for quantifiers and wh-phrases},
  journal = {Journal of Linguistics},
  volume = {34},
  number = {2},
  pages = {415--445},
  year = {1998}}

@book{GinzburgSag2000,
  author = {Ginzburg, Jonathan and Sag, Ivan A.},
  title = {Interrogative investigations: {T}he form, meaning, and use of {E}nglish interrogatives},
  publisher = {CSLI Publications},
  address = {Stanford, CA},
  year = {2000}}

@book{Nunes2004,
  author = {Nunes, Jairo},
  title = {Linearization of chains and sideward movement},
  publisher = {MIT Press},
  address = {Cambridge, MA},
  year = {2004}}

@phdthesis{Hendriks1993,
  author = {Hendriks, Herman},
  title = {Studied flexibility: {C}ategories and types in syntax and semantics},
  school = {University of Amsterdam},
  year = {1993}}

@article{PietroskiEtAl2009,
  author = {Pietroski, Paul and Lidz, Jeffrey and Hunter, Tim and Halberda, Justin},
  title = {The meaning of `most': {S}emantics, numerosity and psychology},
  journal = {Mind \& Language},
  volume = {24},
  number = {5},
  pages = {554--585},
  year = {2009}}

@article{Jayaseelan2001,
  author = {Jayaseelan, K. A.},
  title = {Questions and question-word incorporating quantifiers in {M}alayalam},
  journal = {Syntax},
  volume = {4},
  number = {2},
  pages = {63--93},
  year = {2001}}

@article{Shimoyama2011,
  author = {Shimoyama, Junko},
  title = {Japanese indeterminate negative polarity items and their scope},
  journal = {Journal of Semantics},
  volume = {28},
  number = {4},
  pages = {413--450},
  year = {2011},
  note = {\bibchk[volume and pages]}}

@article{Kitahara2024,
  author = {Kitahara, Hisatsugu},
  title = {A review and reanalysis of some recent developments from {GK} to {MC}, and beyond},
  journal = {Reports of the Keio Institute of Cultural and Linguistic Studies},
  volume = {55},
  pages = {199--214},
  year = {2024}}

@book{Kuno1987,
  author = {Kuno, Susumu},
  title = {Functional syntax: {A}naphora, discourse and empathy},
  publisher = {University of Chicago Press},
  address = {Chicago},
  year = {1987}}

% The entries below correct placeholder titles, or complete entries, in keio.bib; this
% file is read first, so BibTeX keeps these and warns about the repeated keys.
\end{filecontents*}

\begin{document}

% ======================================================================
%  Front matter
% ======================================================================
\frontmatter
% ---------------------------------------------------------------- cover
% Keio's announcement, as issued, is the cover.
%\includepdf[pages=1]{2026110708B.pdf}

% ---------------------------------------------------------------- title page
\begin{titlepage}
\begin{tikzpicture}[remember picture, overlay, every node/.style={inner sep=0pt}]
  % a rule down the left edge of the page, in the accent colour
  \fill[gtaccent] ([xshift=14mm]current page.north west) rectangle
    ([xshift=16mm]current page.south west);
  % running label
  \node[anchor=north west, font=\sffamily\footnotesize, text=gtaccent]
    at ([xshift=24mm, yshift=-22mm]current page.north west)
    {\textsc{Keio Linguistics Colloquium} \enspace\textperiodcentered\enspace
     \textsc{7--8 November 2026}};
  % title block
  \node[anchor=north west, align=left, text width=150mm]
    (title) at ([xshift=24mm, yshift=-44mm]current page.north west)
    {{\fontsize{30}{34}\selectfont\bfseries Quantification in a Grammar\\
      without Variable Names\par}\vspace{5mm}
     {\Large\itshape Scope, Storage and Indices after Inclusiveness\par}};
  \node[anchor=north west, align=left] (author) at ([yshift=-14mm]title.south west)
    {{\large\scshape George Tsoulas}\\[2pt]{\normalsize University of York}};
  % the generalisation, as a fading ladder: what natural language pronounces
  \begin{scope}[shift={([xshift=24mm, yshift=104mm]current page.south west)}]
    \foreach \word/\how/\shade [count=\i] in
      {restriction/always/100, force/often/62, scope site/sometimes/32, variable/never/9}{
      \node[anchor=base west, font=\fontsize{30}{30}\selectfont\itshape,
        text=gttext!\shade] at (0, {-(\i-1)*15mm}) {\word};
      \node[anchor=base west, font=\sffamily\small, text=gtaccent]
        at (68mm, {-(\i-1)*15mm}) {\how\ pronounced};
    }
    \draw[gtaccent!40, line width=0.4pt] (64mm, 9mm) -- (64mm, -50mm);
  \end{scope}
  % foot of the page
  \node[anchor=south west, align=left, font=\small]
    at ([xshift=24mm, yshift=22mm]current page.south west)
    {Keio University\\The Keio Institute of Cultural and Linguistic Studies\\
     Lecture Notes\\[6pt]
     {\color{gtaccent}Version \keioversion\enspace\textperiodcentered\enspace \keioversiondate}\\[2pt]
     {\footnotesize\href{https://doi.org/\keiodoi}{\texttt{doi:\keiodoi}}}};
\end{tikzpicture}
\end{titlepage}

% ---------------------------------------------------------------- about
\addchap{About these lectures}

This document is the text of the four lectures given at Keio University on 7 and 8 November 2026, under the title of the Keio announcement. It is written to be read as well as heard. Each lecture corresponds to one session, and carries, beside the argument, the trees and the derivations, the tables, the examples there is never time for in a lecture, and excursions for readers from other frameworks, so that it can be read on its own, before the lectures or after them. The order of the argument is the order of the two afternoons. The text departs from the lectures as given in one respect. The comparison with other frameworks, which the announcement puts as the question whether a monocyclic grammar with a store is a monostratal one under another name, ran through the lectures as their thesis. It is a question about where the proposal stands and not the proposal, and it is gathered here into a section of its own near the end of Lecture~4 (\Hs{4}{sec:neighbours}), where it is put in the more usual terms of single-cycle syntax.

The question is the one the announcement poses. Natural language seems to quantify almost exclusively by operator and variable, and the grammar most of us assume has no way of writing the variable. Lecture~1 sets out the architecture and states the problem; Lecture~2 follows force along the clausal spine; Lecture~3 develops a box theory of Quantifier Raising and puts it to the test of locality; Lecture~4 asks what an index can be in a grammar that cannot add one. The overview that follows the contents gives the argument in six statements.

\minisec{Conventions}

Examples are numbered afresh in each lecture. A reference to an example in another lecture takes the form Lecture~1~(38), and to a section, Lecture~3, \S4.2. Japanese is romanised in Kunrei-shiki throughout (\textit{syootaisita}, \textit{mittu-zutu}, \textit{zibun}, \textit{soitu}), Korean in Yale. Quantificational particles are glossed by their own names (\gl{mo}, \gl{ka}, \gl{demo}, \gl{zutu}; Korean \gl{nka}, \gl{na}, \gl{to}, \gl{ssik}) rather than by a translation, so as not to prejudge what they mean, which is, after all, one of the things at issue. In the derivations \vs\ marks the phase head of the verbal domain; in trees reproduced from the literature the same head is written \textit{v}. Judgements are given as reported in the sources cited. Some of them are delicate, and some of the Japanese examples have been constructed after the literature's descriptions rather than copied; native speakers who disagree with any of them should tell me, since the argument is meant to turn on the data and not on my reading of it.

\minisec{The apparatus}

The outer margin carries notes, figure captions and, at the opening of each lecture, a map of its sections. Two kinds of box recur in the text, and they mean the same thing throughout.

\begin{keybox}
\small\noindent\textbf{Central statements} are set against a red rule: the hypotheses, the desideratum, the clauses of the proposal and the conjectures, the sentences on which the argument turns.
\end{keybox}

\begin{excursus}[For readers from other frameworks]
Excursions are set in grey. They connect the argument to HPSG, LFG, CCG, variable-free and dynamic semantics, and to the logics that quantify without variables, and they can be skipped without loss of the thread. Lecture~4 draws them together (\Hs{4}{sec:neighbours}).
\end{excursus}

Six exercises go with the text: judgements on three sentences before any theory is on the table, the grid of Lecture~1, a judgement sheet, the hierarchy of Lecture~2, the worksheet of Lecture~3, and a choice between three exits in Lecture~4. They are collected in an appendix, so that they can be done at the reader's own pace, or used in a class. A pointer in the margin, \textsf{\scriptsize\color{gtaccent}$\blacktriangleright$\,Exercise}, marks the place in the text where each belongs, and links to it; the exercise links back.

The trees and workspace derivations follow the conventions of the \texttt{minsyn} package: a workspace is a numbered panel $W_n$, listing the objects the computation can see; an element in a box, as in \BOX{\mathrm{QP}}, is set aside in the box of Lecture~3; a ringed node is the root of a Self-Merge, and a dashed node a copy, or a lower occurrence; a shaded node is a phase head.

\minisec{Acknowledgements}

I am grateful to Keio University for the invitation to give these lectures, and to the colleagues at Keio who made them possible, in particular Hisa Kitahara who made the arrangements and made extremely useful comments.  The present version has not yet absorbed all the comments and questions from these sessions, it is still a work in progress. For the work on which much of this rests I thank David Adger, Kook-Hee Gil, Steve Harlow, Murat Kural, Norman Yeo among many other individuals and audiences that I cannot list in full here. Earlier versions of Lecture~3 were presented at the International Circle for Korean Linguistics 2026 in Sheffield and at a NESS (North East Syntax Seminar) meeting in Leeds, and I thank both audiences for their questions.  Part of Lecture~4 was presented at Peking University prior to the Keio sessions.  I thank that audience and especially Freddy Hu and Victor Hu for the invitation, arrangements and comments on that part of the work. 



% ---------------------------------------------------------------- contents
\cleardoublepage
\tableofcontents

% ======================================================================
%  The argument at a glance (front matter, after the contents)
% ======================================================================
\addchap{The argument at a glance}

\begin{wide}
\centering
\begin{tikzpicture}[
    x=1mm, y=1mm,
    q/.style={draw=gtaccent, line width=0.8pt, fill=gtaccent!5, rounded corners=2pt,
      align=center, inner sep=5pt, text width=150mm},
    ans/.style={draw=gtghost!70, line width=0.5pt, rounded corners=2pt, fill=white,
      align=left, inner sep=4pt, text width=46mm, anchor=north, font=\small},
    head/.style={font=\sffamily\bfseries\small, text=gtaccent},
    arr/.style={-{Stealth[length=2.2mm]}, gtaccent!70, line width=0.6pt},
  ]
  \node[q] (Q) at (80,0) {\small\textsc{The question}\\[2pt]
    {\large How does a grammar that cannot write \textit{x} quantify?}\\[2pt]
    {\small\itshape Natural language quantifies by operator and variable; Inclusiveness forbids indices,
     the copy theory leaves a quantifier where the semantics wants a variable,
     and there is no level of LF at which either could be repaired.}};
  \node[ans] (A1) at (27,-24) {\parbox[t][35mm][t]{46mm}{\raggedright{\sffamily\bfseries\color{gtaccent}1\enspace Movement}\\[2pt]
    Copies, repaired by Trace Conversion.\\[4pt]
    {\itshape It breaks Inclusiveness: the definite and the variable were never in the workspace.}\\[4pt]
    {\sffamily\footnotesize\color{gtghost}Lecture~1, \S\ref{H1:sec:problem}; Lecture~3, \S\ref{H3:sec:wh-to-QR}}}};
  \node[ans] (A2) at (80,-24) {\parbox[t][35mm][t]{46mm}{\raggedright{\sffamily\bfseries\color{gtaccent}2\enspace Storage}\\[2pt]
    Cooper storage, then the box.\\[4pt]
    {\itshape It works, and its locality is inherited from phases, where a semantic store has to stipulate it.}\\[4pt]
    {\sffamily\footnotesize\color{gtghost}Lecture 2; Lecture 3}}};
  \node[ans] (A3) at (133,-24) {\parbox[t][35mm][t]{46mm}{\raggedright{\sffamily\bfseries\color{gtaccent}3\enspace Relations over tokens}\\[2pt]
    Indices, read as Form Copy.\\[4pt]
    {\itshape It works, and the box already relies on such a relation: its \uth\ is an index, honestly described.}\\[4pt]
    {\sffamily\footnotesize\color{gtghost}Lecture 4}}};
  \foreach \a in {A1,A2,A3} \draw[arr] (Q.south) -- (\a.north);
  \node[q, anchor=north, text width=118mm] (T) at (80,-72) {\small\textsc{The answer}\\[2pt]
    {\large The variable is not an item\\ but a relation between tokens.}\\[2pt]
    {\small\itshape Scope is storage in the syntax, and its locality is the locality of phases.}};
  \draw[arr] (A2.south) -- (A2.south |- T.north);
  \draw[arr] (A3.south) -- (A3.south |- T.north);
  \draw[arr, gtghost!60, dashed] (A1.south) -- node[pos=0.5, right=1mm,
    font=\sffamily\scriptsize, text=gtghost]{set aside} (A1.south |- T.north);
\end{tikzpicture}
\end{wide}

\medskip
\minisec{The argument in six statements}

\begin{enumerate}[label={\sffamily\bfseries\color{gtaccent}\arabic*}, leftmargin=1.6em, itemsep=4pt]
\item \textit{The generalisation} (Lecture~1, \S\ref{H1:sec:spot}). The restriction of a quantifier is always pronounced, its force often, its scope site sometimes, and its variable never (\Hx{1}{ex:generalisation}).
\item \textit{The desideratum} (Lecture~1, \S\ref{H1:sec:desideratum}). Force must be separable from the restriction: the grammar must fix force and scope at a position other than the one at which the restriction is merged and pronounced, and relate the two without adding anything to the workspace (\Hx{1}{ex:desideratum}).
\item \textit{Two localities} (Lecture~2, \S\ref{H2:sec:localities}). Domain locality counts edges; path locality counts interveners. The two cut orthogonally, and which of them governs a dependency tells a syntactic store from a semantic one (\Hx{2}{ex:wedge}).
\item \textit{The box theory of QR} (Lecture~3, \S\ref{H3:sec:boxQR}). A quantifier is boxed by Self-Merge, the verb's internal $\theta$-feature stands in as the variable, and the quantifier discharges it at a phase edge: it is pronounced at its Case site, the highest relation the sensorimotor interface can see, and interpreted where it discharges (\Hx{3}{ex:proposal}).
\item \textit{Inclusiveness revised} (Lecture~4, \S\ref{H4:sec:revised}). The computation operates on tokens; it may relate them, by Merge, Agree and Form Copy, but it may not create new values (\Hx{4}{ex:incl-revised}).
\item \textit{The answer} (Lecture~4, \S\ref{H4:sec:conclusion}). The variable is not an item but a relation between tokens; scope is storage in the syntax, and its locality is the locality of phases (\Hx{1}{ex:answer}). How the proposal stands to its neighbours in other frameworks, Cooper storage first among them, is the subject of Lecture~4, \S\ref{H4:sec:neighbours}.
\end{enumerate}

\minisec{The two afternoons}

\noindent
\begin{tabularx}{\linewidth}{@{}p{2.3cm}p{2.1cm}L@{}}
\toprule
\textbf{Day 1} & Saturday & 7 November 2026\\
\midrule
13:30--15:30 & Lecture 1 & A grammar that cannot write \textit{x}: the minimalist architecture, and what makes a QP a QP\\
15:45--17:45 & Lecture 2 & Quantification on the clausal spine: unselective binding, Japanese and Korean particles, and a parameter hierarchy\\
18:00--18:30 & & Discussion\\
\midrule
\textbf{Day 2} & Sunday & 8 November 2026\\
\midrule
13:30--15:30 & Lecture 3 & Scope in the box theory: from Quantifier Raising to storage, with locality, rigidity and scrambling in Japanese\\
15:45--17:45 & Lecture 4 & Indices after Inclusiveness: Form Copy, bound-variable anaphora, and the architecture revisited\\
18:00--18:30 & & Discussion\\
\bottomrule
\end{tabularx}

\mainmatter

% ======================================================================
%  PART 1
% ======================================================================
\keiolecture{1}{A grammar that cannot write \emph{x}}{A grammar that cannot write \emph{x}}{Day 1, Lecture 1 \textperiodcentered{} Saturday 7 November 2026, 13:30--15:30}
\keioepigraph{In natural language, as far as I can tell, the restriction of a quantifier is always pronounced, its force often, its scope site sometimes, and its variable never.}{}
\lecturemap

% =====================================================================
\section{Aims and plan}\label{sec:aims}

The question this seminar asks is easy to state and, I think, surprisingly hard to answer: how does a grammar that cannot write \emph{x} quantify? Quantification in natural language is, as far as anyone has been able to establish, operator--variable quantification. A quantifier binds a position, and the sentence is understood as a claim about the values that position can take. The grammar that most of us assume, however, has no way of writing the variable. Inclusiveness forbids the computation to add indices; the copy theory of movement leaves a copy of the quantifier exactly where the semantics wants a variable; and there is no level of Logical Form, after the fact, at which either could be repaired. Something has to give, and over the four lectures I will weigh three candidates for the source of the variable: movement, storage, and relations between tokens.

It is worth keeping the three components of the problem separate from the start, since each of the three answers abandons, or reinterprets, a different one of them. Movement keeps the copies and repairs them, by converting the lower copy into something that contains a variable (this is Trace Conversion, \S\ref{sec:copy}). Storage keeps Inclusiveness and gives up the idea that the variable is something movement leaves behind. Relations over tokens keep both, and say that the variable was never an object in the first place. Which of the three survives is, in a sense, the question of the seminar.

This first lecture has two aims. The first is to set out, once and in one vocabulary, the architecture that the rest of the seminar assumes, so that when the box arrives in Lecture~3 it can be tested against conditions everyone has agreed to in advance. That aim is architectural: nothing in \S\ref{sec:architecture} is mine, and no result there depends on my own proposals. The second is to state the problem precisely and to argue, from the internal syntax of quantified nominals, for a desideratum that the remaining three lectures must meet: quantificational force must be separable from the restriction, and the grammar must be able to fix force and scope at a position other than the one at which the restriction is merged and pronounced. This is the first step of the argument proper, and the desideratum (\ref{ex:desideratum}) is stated so as to be deliberately neutral among the three answers.

It is worth turning the question around before we start. The obvious puzzle is how a grammar without variables manages operator--variable quantification at all. The less obvious one, and to my mind the more interesting, is its inverse: why does natural language do so little besides operator and variable, given how many other ways of quantifying logicians have devised, and why is the one ingredient that the logician writes into every formula, the variable, the one ingredient that no language pronounces? The obvious question invites a technical repair; the inverse question invites an explanation. I return to this inversion in \S\ref{sec:inversion}, with the logics that quantify without variables on the table, and the seminar as a whole can be read as an attempt to answer it.

The plan is as follows. \S\ref{sec:vote} asks for judgements on three sentences before any theory is on the table. \S\ref{sec:architecture} sets out the architecture in five ideas, each tagged with the lecture where it matters; for readers who work in other frameworks, it is accompanied by excursions on HPSG, LFG, CCG and variable-free semantics, set off in boxes that can be skipped without loss to the argument. \S\ref{sec:problem} states the problem: what Quantifier Raising (QR from now on) was for, why the copy theory and Inclusiveness leave it without a variable, and what the three candidate answers are. \S\ref{sec:spot} looks for the variable in six sentences, and finds it in none of them. \S\ref{sec:nominal} turns to the nominal extended projection and asks what makes a QP a QP. \S\ref{sec:desideratum} states the desideratum, and \S\ref{sec:readings} collects questions for the discussion and the readings.

A note on the glosses of \textit{mo} and \textit{ka}. There are several theories of both particles, some of them developed at Keio, and glossing \textit{mo} as `every' or as `also' would take a side before the argument has started. The glosses are therefore deliberately uninformative.

% =====================================================================
\section{Three sentences before any theory}\label{sec:vote}

Before any theory is on the table, consider three sentences, and one question about each of them: can the books vary with the readers? That is, can the sentence be true in a situation in which every book was read by someone, but not by the same someone in each case?\toexercise{vote}

\begin{exe}
\ex \label{ex:someone} Someone read every book.
\ex \label{ex:dareka-dono}
\gll Dare-ka-ga dono hon-mo yonda.\\
     who-\gl{ka}-\gl{nom} which book-\gl{mo} read.\gl{pst}\\
\glt `Someone read every book.'
\ex \label{ex:dono-dareka}
\gll Dono hon-mo dare-ka-ga yonda.\\
     which book-\gl{mo} who-\gl{ka}-\gl{nom} read.\gl{pst}\\
\glt `Every book, someone read.'
\end{exe}

Scope judgements given out of the blue are unreliable, and out of context the surface reading is usually preferred, so the varying reading needs a concrete situation to be true in. Here is one. The Mita reading circle drew up a list of twelve novels for the summer. By September every novel on the list had been read by at least one member (Hanako read \textit{Kokoro}, Taroo read \textit{Sansiroo}, Kenta read \textit{Yukiguni}), but nobody had read all twelve. The convenor is now reporting to the department. Is (\ref{ex:someone}) a true report of what happened? Is (\ref{ex:dareka-dono})? Is (\ref{ex:dono-dareka})? The appendix puts the three questions as an exercise (\exref{vote}), to be answered before anything else is read, and Lecture~3 returns to them (\Hs{3}{sec:rigid}).

The answers reported in the literature are, I think, the ones most speakers will give. (\ref{ex:someone}) allows both readings: the surface one, on which a single person read all the books, and the inverse one, on which the books vary with the readers. Out of the blue many speakers prefer the first, but the second is there, and with a less tendentious pair, such as \textit{An engineer inspected every plane}, it is the more natural of the two. (\ref{ex:dareka-dono}) has only the surface reading. If it is true, one person read all twelve novels, and the convenor's report is false \citep{Kuroda1970,Kuno1973,Hoji1985}. (\ref{ex:dono-dareka}), with the object scrambled to the front, has both. Note already that (\ref{ex:dareka-dono}) and (\ref{ex:dono-dareka}) contain the same words. \citet{Hoji1985} generalised the pattern roughly as follows (my paraphrase):

\begin{exe}
\ex \label{ex:hoji-gen} In Japanese, a quantified phrase $\alpha$ can take scope over a quantified phrase $\beta$ only if $\alpha$ c-commands $\beta$ or a copy of $\beta$ (in Hoji's terms, a trace).
\end{exe}

(\ref{ex:hoji-gen}) is a paraphrase, and a deliberately simplified one. Hoji's own formulation is stated in terms of traces, and of a condition on scope that he derived within Lasnik and Saito's version of the government-and-binding framework; the version given here is the one we will need. It is, in any case, a generalisation over judgements, not yet an explanation. What it does make plain is that the pattern can be read off the structures, once copies are counted, as in Figure~\ref{fig:vote-trees}.

\begin{keiofig*}{fig:vote-trees}{(\ref{ex:dareka-dono}) and (\ref{ex:dono-dareka}) as structures. On the left, the subject c-commands the object, and the object c-commands neither the subject nor a copy of it: $\exists > \forall$ only. On the right, the scrambled object c-commands the subject, and the subject c-commands a copy of the object: both readings. The trees are deliberately pre-theoretical (no claim yet about what scrambling is); the lower copies are drawn dashed.}
\begin{twocol}{0.47}
\centering\small (\ref{ex:dareka-dono}) \textit{dare-ka-ga dono hon-mo yonda}\par\smallskip
\begin{forest}
[TP, SOtree, spine
  [DP, SOmoved [{\LI{dare-ka-ga}}, SOlexical]]
  [\vs P
    [DP, SOtrace [{$\langle$\LI{dare-ka-ga}$\rangle$}]]
    [VP [DP [{\LI{dono hon-mo}}, SOlexical]] [V, SOhead [{\LI{yonda}}, SOlexical]]]]]
\end{forest}
\midcol
\centering\small (\ref{ex:dono-dareka}) \textit{dono hon-mo dare-ka-ga yonda}\par\smallskip
\begin{forest}
[TP, SOtree, spine
  [DP, SOmoved [{\LI{dono hon-mo}}, SOlexical]]
  [TP
    [DP [{\LI{dare-ka-ga}}, SOlexical]]
    [\vs P
      [DP, SOtrace [{$\langle$\LI{dare-ka-ga}$\rangle$}]]
      [VP [DP, SOtrace [{$\langle$\LI{dono hon-mo}$\rangle$}]] [V, SOhead [{\LI{yonda}}, SOlexical]]]]]]
\end{forest}
\end{twocol}
\end{keiofig*}

Whether the copy of the subject in Spec-\vs P matters for (\ref{ex:dareka-dono}), that is, whether the object could take scope over it without taking scope over the subject's higher copy, is a further question, which I leave open here.

Three things follow even before we have a theory, and they frame everything else. First, scope is at least partly a matter of syntax: (\ref{ex:dareka-dono}) and (\ref{ex:dono-dareka}) contain exactly the same words and differ only in their order, yet only one of them is ambiguous. Second, the availability of the inverse reading is not a property of the quantifiers involved. \textit{Dono hon-mo} is the same universal in (\ref{ex:dareka-dono}) and (\ref{ex:dono-dareka}), and \textit{dare-ka} the same existential; what changes is where they are. This is the point semanticists may want to resist, and it is worth seeing what resisting it really means: a lexical account would have to say that \textit{dono hon-mo} is lexically unable to take inverse scope, and then explain (\ref{ex:dono-dareka}) by something else. Third, whatever English does to the object of (\ref{ex:someone}) to let it take scope over the subject, Japanese does not do to the object of (\ref{ex:dareka-dono}), or does it only when the word order already puts the object on top. So the first question to ask is not a semantic one, \emph{what do the quantifiers mean?}, but a syntactic one: \emph{what is the syntax allowed to do?} It is the third point that turns the problem into a question about operations.

A caution about judgements of this kind. A sheet filled in at a lecture, or a judgement given while reading, is not an experiment, and judgements of scope are notoriously delicate, nowhere more so than in Japanese, where a whole research programme (\citealt{Hoji2015} and the Language Faculty Science work that followed) has gone into separating what the grammar rules out from what is merely hard to get. I will treat readers' answers as data, not decoration, but as pilot data. The judgement sheet in the appendix (\exref{sheet}) is a slightly more careful version of the same exercise; Lecture~2 explains its design (\Hs{2}{sec:localities}), and Lecture~3 tests it (\Hs{3}{sec:wedge}). If a speaker accepts the varying reading of (\ref{ex:dareka-dono}), that is not noise to be explained away; it is something the theory will have to say something about. Note, finally, that the accusative \textit{-o} is not pronounced before \textit{mo} in (\ref{ex:dareka-dono}) and (\ref{ex:dono-dareka}). That small fact, that the force particle and the case particle compete for one position, will matter when we get to Case and externalisation in Lecture~3.

% =====================================================================
\section{The architecture: five ideas}\label{sec:architecture}

What follows sets out the architecture the seminar assumes, in five ideas. For each I give a short statement, say what it does, and tag the lecture where it matters most. None of it is mine. The point is to fix a vocabulary in which the box, when it arrives in Lecture~3, can be tested, and in which those who do not work in this framework can see exactly what is being assumed. I lean throughout on \citet{Chomsky2021} and on \textit{Merge and the Strong Minimalist Thesis} (\citealt{ChomskyEtAl2023}; MSMT from now on), and its authors will, I hope, correct me where my paraphrase is loose. The tags are there so that the architecture does not read as a detour: every idea is a promissory note for a later lecture.

\begin{center}\small
\begin{tabular}{@{}l P{4.6cm} P{3.1cm}@{}}
\toprule
& Idea & Where it matters\\
\midrule
(a) & MERGE in a workspace, and the conditions on it & Lectures 3 and 4\\
(b) & Inclusiveness and No Tampering & \S\ref{sec:problem}; Lecture 4\\
(c) & The duality of semantics & Lecture 3\\
(d) & Cyclicity without LF & Lecture 3\\
(e) & Spines and externalisation & Lectures 1 and 2; rigidity in Lecture 3\\
\bottomrule
\end{tabular}
\end{center}

\subsection{(a) MERGE in a workspace (Lectures 3 and 4)}\label{sec:merge}

The computation operates on a workspace (WS), the collection of syntactic objects available to it at a given point. There is one structure-building operation, MERGE, which takes two objects and the workspace and returns a new workspace:

\begin{exe}
\ex \label{ex:merge} \textit{MERGE in a workspace} \citep{Chomsky2019,Chomsky2021,ChomskyEtAl2023}:\\[3pt]
$\textsf{MERGE}(P, Q, \mathit{WS}) = \mathit{WS}' = [\,\{P, Q\},\ X_1, \dots, X_n\,]$
\end{exe}

Here $\mathit{WS}'$ contains the new set $\{P, Q\}$ together with every member of WS other than $P$ and $Q$. External and Internal Merge are not two operations but two cases of one. In External Merge (EM), $P$ and $Q$ are both members of WS; in Internal Merge (IM), one is a term of the other, and the result contains two \emph{inscriptions} of the same object. That is all the copy theory of movement amounts to. The term `inscription' is Chomsky's, and it is chosen to avoid both `copy' and `occurrence', each of which carries theory with it. Two inscriptions of one object is what IM returns; whether they are to count as copies or as repetitions is decided later, by Form Copy (\ref{ex:formcopy}).

It is worth watching this happen once, slowly, on the first sentence of \S\ref{sec:vote}, (\ref{ex:someone}). Figure~\ref{fig:merge-deriv} gives the derivation as a sequence of workspaces. Each application of MERGE removes two objects from the workspace and adds one, the set it forms; nothing else changes. The labels (QP, VP, \vs P) are added for readability only, because we read trees faster than sets: MERGE returns sets, and the workspace notation above each tree is the real object. The step from $W_4$ to $W_5$ is the one the duality of semantics, idea (c), cares about: the external argument is externally merged in Spec-\vs P, its $\theta$-position. Then T is merged, and a single object remains in the workspace.

\begin{keiofig*}{fig:merge-deriv}{The derivation of (\ref{ex:someone}) in a workspace, $W_1$--$W_7$. $W_2$--$W_6$ are External Merge; $W_7$ is Internal Merge, with \textit{someone} a term of TP. Labels are for readability only, and QP is kept as a unit from $W_4$ on, for space.}
\begin{derivation}
\begin{minipage}[t]{0.31\linewidth}\vspace{0pt}
\begin{workspace}[Select]
  \scriptsize\WSeq{\LI{every},\allowbreak\ \LI{book},\allowbreak\ \LI{read},\allowbreak\ \vs,\allowbreak\ \LI{someone},\allowbreak\ \text{T}}
\end{workspace}
\end{minipage}\hfill
\begin{minipage}[t]{0.31\linewidth}\vspace{0pt}
\begin{workspace}[EM]
  \scriptsize\WSeqTree{\{\LI{every},\LI{book}\},\allowbreak\ \LI{read},\allowbreak\ \vs,\allowbreak\ \LI{someone},\allowbreak\ \text{T}}
    {[QP, SOtree, xslim [Q, SOhead [{\LI{every}}, SOlexical]] [NP [{\LI{book}}, SOlexical]]]}
\end{workspace}
\end{minipage}\hfill
\begin{minipage}[t]{0.31\linewidth}\vspace{0pt}
\begin{workspace}[EM]
  \scriptsize\WSeqTree{\{\LI{read},\{\LI{every},\LI{book}\}\},\allowbreak\ \vs,\allowbreak\ \LI{someone},\allowbreak\ \text{T}}
    {[VP, SOtree, xslim [V, SOhead [{\LI{read}}, SOlexical]] [QP [Q, SOhead [{\LI{every}}, SOlexical]] [NP [{\LI{book}}, SOlexical]]]]}
\end{workspace}
\end{minipage}

\medskip
\begin{minipage}[t]{0.31\linewidth}\vspace{0pt}
\begin{workspace}[EM]
  \scriptsize\WSeqTree{\{\vs,\{\LI{read},\mathrm{QP}\}\},\allowbreak\ \LI{someone},\allowbreak\ \text{T}}
    {[\vs P, SOtree, tiny [\vs, SOhead] [VP [V, SOhead [{\LI{read}}, SOlexical]] [QP [{\LI{every book}}, SOlexical]]]]}
\end{workspace}
\end{minipage}\hfill
\begin{minipage}[t]{0.31\linewidth}\vspace{0pt}
\begin{workspace}[EM: $\theta$-position]
  \scriptsize\WSeqTree{\{\LI{someone},\allowbreak\{\vs,\{\LI{read},\mathrm{QP}\}\}\},\allowbreak\ \text{T}}
    {[\vs P, SOtree, tiny [DP [{\LI{someone}}, SOlexical]] [\vs$'$ [\vs, SOhead] [VP [V, SOhead [{\LI{read}}, SOlexical]] [QP [{\LI{every book}}, SOlexical]]]]]}
\end{workspace}
\end{minipage}\hfill
\begin{minipage}[t]{0.31\linewidth}\vspace{0pt}
\begin{workspace}[EM]
  \scriptsize\WSeqTree{\{\text{T},\allowbreak\{\LI{someone},\allowbreak\{\vs,\allowbreak\{\LI{read},\mathrm{QP}\}\}\}\}}
    {[TP, SOtree, tiny [T, SOhead] [\vs P [DP [{\LI{someone}}, SOlexical]] [\vs$'$ [\vs, SOhead] [VP [V, SOhead [{\LI{read}}, SOlexical]] [QP [{\LI{every book}}, SOlexical]]]]]]}
\end{workspace}
\end{minipage}

\medskip
\begin{twocol}{0.46}\vspace{0pt}
\begin{workspace}[IM: \textit{someone} is a term of TP]
  \scriptsize\WSeqTree{\{\LI{someone},\allowbreak\{\text{T},\allowbreak\{\LI{someone},\allowbreak\{\vs,\allowbreak\{\LI{read},\mathrm{QP}\}\}\}\}\}}
    {[TP, SOtree, tiny [DP, SOmoved [{\LI{someone}}, SOlexical]] [T$'$ [T, SOhead] [\vs P [DP, SOtrace [{$\langle$\LI{someone}$\rangle$}]] [\vs$'$ [\vs, SOhead] [VP [{\LI{read every book}}, SOlexical]]]]]]}
\end{workspace}
\midcol\vspace{0pt}
\small\raggedright
\vspace*{1.2em}
MERGE(\textit{someone}, TP, WS): one of the two objects is a term of the other. The result contains two inscriptions of one object, \textit{someone}; the set on the left shows it twice because it \emph{is} there twice. The dashed node is a drawing convention for the lower inscription: there is no trace, no index, and no new object.
\end{twocol}
\end{derivation}
\end{keiofig*}

The last step deserves a moment, particularly for the semanticists. At $W_7$ nothing has been copied, moved or left behind. The lower inscription of \textit{someone} is not a variable, not a trace and not a pronoun: it is the DP \textit{someone}, and which of its two inscriptions is pronounced is a matter for externalisation, idea (e). That is exactly the problem of \S\ref{sec:problem}, seen before any quantifier has moved.

MERGE is constrained, not by stipulations specific to it, but by general conditions on computation, which I list in the form in which we will use them:

\begin{exe}
\ex \label{ex:conditions}
\begin{xlist}
\ex \label{ex:RR} \textit{Resource Restriction / Minimal Yield.} An operation adds to the workspace as few new accessible objects as possible; MERGE adds exactly one, the set it forms. This is what rules out Parallel Merge and Sideward Movement, both of which yield two accessible inscriptions where one would do.
\ex \label{ex:determinacy} \textit{Determinacy.} An operation applies only to a uniquely determined target. If the search for an object of the required description returns two candidates, the operation cannot apply \citep{GotoIshii2019,Chomsky2021}.
\ex \label{ex:minsearch} \textit{Minimal Search.} MERGE, Agree and labelling find their targets by the least costly search: the first, structurally closest candidate is found, and the search stops.
\ex \label{ex:formcopy} \textit{Form Copy (FC).} Identical inscriptions are construed as copies of one object, rather than as repetitions of two, only under certain structural conditions: identity, c-command (with minimality), and co-occurrence within a single phase \citep{Chomsky2021,ChomskyEtAl2023}.
\end{xlist}
\end{exe}

Three remarks. First, (\ref{ex:RR}) is what will make Self-Merge, the operation on which my version of the box depends, a real test case: whatever MERGE($\alpha$, $\alpha$) produces, it must add no more than one accessible object. Second, (\ref{ex:determinacy}) is in obvious tension with optionality. If boxing is optional, or if a boxed quantifier may be discharged at either of two phases, something must make the choice determinate, or else the ambiguity of (\ref{ex:someone}) must come from somewhere other than the narrow syntax. These are among the hard questions of Lecture~3 (\Hs{3}{sec:objections}), and I would rather they were found here than there. Third, (\ref{ex:formcopy}) is, as I will argue in Lecture~4, already a theory of indices: what FC computes is which tokens count as the same, which is precisely the information an index carries. Form Copy is the idea the whole of Lecture~4 rests on, and whether co-phasehood belongs in its definition is exactly what the most recent work on it has been revising (\citealt{KitaharaSeely2025}).

What (\ref{ex:RR}) rules out is easiest to see in a picture (Figure~\ref{fig:minyield}). Under Internal Merge the lower inscription of $\alpha$ is inside the new object, and so it is not independently accessible: one new accessible object, as required. Under Parallel Merge, $\gamma$ is merged with $\alpha$ and then, now a term of $\{\alpha,\gamma\}$, merged again with $\beta$, which leaves two accessible objects where one would do. Sideward movement \citep{Nunes2004} is the same in the relevant respect; `Parallel Merge' is Citko's \citeyearpar{Citko2005} name for the structure just described. Minimal Yield, put in one line, says that of all the ways of forming a new object from $\gamma$, MERGE is allowed only the one that leaves the fewest objects accessible to later operations. Both of the excluded operations are things the box will have to avoid in Lecture~3.

\begin{keiofig*}{fig:minyield}{What (\ref{ex:RR}) rules out. Left: Internal Merge yields one new accessible object; the lower inscription of $\alpha$ is inside it. Right: Parallel Merge yields two, $\{\alpha,\gamma\}$ and $\{\beta,\gamma\}$, with $\gamma$ a term of both.}
\begin{twocol}{0.48}
\centering\small\textbf{Internal Merge}: one new accessible object\par\medskip
\begin{tikzpicture}[every node/.style={font=\small}, e/.style={gtaccent!60, line width=0.5pt}]
  \node (r) at (0,1.6) {$\{\alpha, \{\beta,\alpha\}\}$};
  \node (a1) at (-1.0,0.7) {$\alpha$};
  \node (ba) at (1.0,0.7) {$\{\beta,\alpha\}$};
  \node (b) at (0.5,-0.2) {$\beta$};
  \node (a2) at (1.6,-0.2) {$\alpha$};
  \draw[e] (r)--(a1); \draw[e] (r)--(ba); \draw[e] (ba)--(b); \draw[e] (ba)--(a2);
\end{tikzpicture}
\midcol
\centering\small\textbf{Parallel Merge}: two accessible objects\par\medskip
\begin{tikzpicture}[every node/.style={font=\small}, e/.style={gtaccent!60, line width=0.5pt}]
  \node (x) at (-1.2,1.6) {$\{\alpha,\gamma\}$};
  \node (y) at (1.2,1.6) {$\{\beta,\gamma\}$};
  \node (al) at (-2.0,0.7) {$\alpha$};
  \node (g) at (0,0.7) {$\gamma$};
  \node (be) at (2.0,0.7) {$\beta$};
  \draw[e] (x)--(al); \draw[e] (x)--(g); \draw[e] (y)--(g); \draw[e] (y)--(be);
\end{tikzpicture}
\end{twocol}
\end{keiofig*}

The structure on the right of Figure~\ref{fig:minyield} has a name of its own,
and it is worth defining it here, because it will come back as one of the
answers to the question of what a copy is. A \emph{multidominant} phrase marker
is one in which a single node may have more than one mother: a directed acyclic
graph rather than a tree. Displacement is then the remerge of one object in a
second position, with no copy and no trace
\citep{Gartner2002,Citko2005,deVries2009,Johnson2012}; in the Parallel Merge
structure, $\gamma$ is one object, dominated by both $\{\alpha,\gamma\}$ and
$\{\beta,\gamma\}$, and in Figure~\ref{fig:multidominance} a QP is dominated by
both TP and VP. Under (\ref{ex:RR}) Parallel Merge is not an available
\emph{operation}. As a \emph{representation} of movement, however,
multidominance is worth having on the table, since it is one of three ways of
saying what the lower copy is (\S\ref{sec:copy}), and it still has to say which
occurrence is read as a variable, and how. The question \emph{what is a copy?}
will in fact be asked three times in this lecture, and receive three answers: two
inscriptions of one object (MERGE), one node with two mothers (multidominance),
and a lower copy converted into a definite description (Trace Conversion).

\begin{keiofig}{fig:multidominance}{Multidominance: QP has two mothers, TP and VP. No copy, no trace; one object in two positions.}
\begin{tikzpicture}[every node/.style={font=\small}, e/.style={gtaccent!60, line width=0.5pt}]
  \node (tp) at (0,2.2) {TP};
  \node (vp) at (0.9,1.3) {VP};
  \node (q) at (-0.6,0.4) {QP};
  \node (v) at (1.6,0.4) {V};
  \node (s) at (-1.4,1.3) {DP};
  \draw[e] (tp)--(s); \draw[e] (tp)--(vp); \draw[e] (vp)--(v); \draw[e] (vp)--(q); \draw[e] (tp)--(q);
\end{tikzpicture}
\end{keiofig}

\subsection{(b) Inclusiveness and No Tampering (\texorpdfstring{\S\ref{sec:problem}}{section 3}; Lecture 4)}\label{sec:inclusiveness}

\begin{exe}
\ex \label{ex:inclusiveness} \textit{Inclusiveness.} The objects the computation forms are made up of the features of the lexical items that enter it; nothing is added in the course of the computation, `in particular, no indices, bar levels' and the like \citep[228]{Chomsky1995keio}.
\ex \label{ex:NTC} \textit{No Tampering.} MERGE(X, Y) leaves X and Y unchanged \citep{Chomsky2005,Chomsky2008keio}.
\end{exe}

Together they yield the copy theory. A moved element cannot leave behind a
trace, since a trace is a new object, usually an indexed one, and Inclusiveness
forbids both. Nor can it be altered in its original position, since No Tampering
forbids that. What remains in the original position is the element itself, a
second inscription of one and the same object. The copy theory, in other words,
is not an extra assumption: it is what is left once traces and tampering have
been removed, and giving it up would mean giving up (\ref{ex:inclusiveness}) or
(\ref{ex:NTC}). That matters for what follows, since none of the answers I
consider gives up either.

Now set beside (\ref{ex:inclusiveness}) the rules the textbook semantics uses to interpret movement \citep{HeimKratzer1998}:

\begin{exe}
\ex \label{ex:PA}
\begin{xlist}
\ex \label{ex:PA-a} \textit{Predicate Abstraction:}\quad $\sem{i\ \beta}^{g} = \lambda x \in D_e.\ \sem{\beta}^{g[i \mapsto x]}$
\ex \label{ex:PA-b} \textit{Traces and Pronouns:}\quad $\sem{t_i}^{g} = g(i)$
\end{xlist}
\end{exe}

On this treatment a moved quantifier leaves a trace bearing an index, and a copy
of that index, adjoined just below the moved phrase, triggers
$\lambda$-abstraction over the trace's position (Figure~\ref{fig:PA-tree}). The
index does real work twice: on the binder, where it tells (\ref{ex:PA-a}) which
value of the assignment to shift, and on the trace, where (\ref{ex:PA-b})
interprets it. Both are exactly what (\ref{ex:inclusiveness}) forbids. Heim and
Kratzer are, moreover, quite explicit that the index on the binder is a
syntactic object, adjoined by movement itself; that is precisely an addition to
the workspace. What the structure delivers,
$\sem{\text{TP}} = \sem{\textit{every book}}(\lambda x.\ \text{someone read }
x)$, is right; how it delivers it is the problem. This is the problem of
\S\ref{sec:problem} in miniature, and it is also where Lecture~4 will end up: if
Inclusiveness is right, the index in (\ref{ex:PA}) has to be something the
grammar computes, not something it writes.

\begin{keiofig}{fig:PA-tree}{The textbook structure for the inverse reading of
    (\ref{ex:someone}). The two framed nodes are the two places where the index
    is consulted: the binder, by (\ref{ex:PA-a}), and the trace, by
    (\ref{ex:PA-b}).}
\begin{forest}
[TP, SOtree, xslim
  [QP, SOmoved [{\LI{every book}}, SOlexical]]
  [{$1$}, draw=gtaccent, rounded corners=2pt
    [TP
      [DP [{\LI{someone}}, SOlexical]]
      [VP [V, SOhead [{\LI{read}}, SOlexical]] [{$t_1$}, draw=gtaccent, rounded corners=2pt]]]]]
\end{forest}
\end{keiofig}

\subsection{(c) The duality of semantics (Lecture 3)}\label{sec:duality}

\begin{exe}
\ex \label{ex:duality} External Merge yields argument structure; Internal Merge yields scopal and discourse-related properties \citep{Chomsky2004,Chomsky2008keio}.
\end{exe}

The duality ties the two cases of MERGE to the two broad kinds of semantic
information. EM into $\theta$-positions builds the predicate--argument skeleton;
IM to the edge builds operator--variable structures and the
information-structural layer. QR is its textbook instance: a QP is externally
merged in its $\theta$-position and internally merged at the edge where it takes
scope. MSMT's restriction that $\theta$-positions are filled only by EM (so that
nothing moves into a $\theta$-position, and there is no movement theory of
control) is the same idea seen from the other side. I flag it now because the
box will put pressure on it. In my version of the box the quantifier discharges
a $\theta$-feature of the verb at its scope position, which looks like
$\theta$-assignment away from the $\theta$-position, by External Merge but at a phase edge, and I owe the reader an answer
to the question of which gives. It is better to announce the debt openly than to
have it found; Lecture~3 offers a way of paying it (\Hs{3}{sec:objections}).

\subsection{(d) Cyclicity without LF (Lecture 3)}\label{sec:cyclicity}

The architecture we inherited is the inverted Y of Figure~\ref{fig:Y-phases}, on
the left. There is one syntactic derivation, and at Spell-Out it forks to the
two interfaces. Scope lives on the LF branch, and the covert component is where
QR has lived since \citet{May1977,May1985}. The question for the seminar is
whether that LF branch exists as a separate stratum at all. I assume, with
\citet{Chomsky2004}, that it does not, and that there is a single cycle. The
phase model on the right replaces the single fork by many:

\begin{exe}
\ex \label{ex:transfer}
\begin{xlist}
\ex Transfer applies at each phase (\vs P, CP), sending the complement of the phase head to both interfaces at once.
\ex \label{ex:PIC} \textit{Phase Impenetrability (PIC).} In a phase with head H, the domain of H is not accessible to operations outside the phase; only H and its edge are \citep{Chomsky2000keio,Chomsky2001}.
\end{xlist}
\end{exe}

\begin{keiofig*}{fig:Y-phases}{Left: the inverted Y, with scope on the covert LF branch. Right: cyclic Transfer, in which each phase sends its complement to both interfaces, and only the edge remains accessible. One fork becomes many, and there is no separate LF stratum.}
\begin{twocol}{0.46}
\centering
\begin{tikzpicture}[every node/.style={font=\small}, >=Stealth, e/.style={gtaccent, line width=0.6pt}]
  \node (lex) at (0,3.2) {Lexicon / numeration};
  \node (syn) at (0,2.2) {Syntax: External / Internal Merge};
  \node[draw=gtaccent, rounded corners=2pt, inner sep=3pt] (so) at (0,1.2) {Spell-Out};
  \node (pf) at (-1.8,0) {PF};
  \node (lf) at (1.8,0) {LF};
  \node[text=gtghost, font=\footnotesize] at (-1.8,-0.45) {sensorimotor};
  \node[text=gtghost, font=\footnotesize] at (1.8,-0.45) {conceptual--intentional};
  \draw[e,->] (lex)--(syn); \draw[e,->] (syn)--(so);
  \draw[e,->] (so)--(pf) node[midway, left, text=gtghost] {\scriptsize overt};
  \draw[e,->] (so)--(lf) node[midway, right, text=gtghost] {\scriptsize covert};
\end{tikzpicture}
\midcol
\centering
\begin{tikzpicture}[every node/.style={font=\small}, >=Stealth, e/.style={gtaccent, line width=0.6pt}]
  \node[draw=gtaccent, rounded corners=2pt, inner sep=4pt, fill=gtexbg] (vp) at (0,0) {\vs P phase: complement VP};
  \node[draw=gtaccent, rounded corners=2pt, inner sep=4pt, fill=gtexbg] (cp) at (0,1.5) {CP phase: complement TP};
  \node (sm1) at (3.2,0.25) {SM};
  \node (ci1) at (3.2,-0.3) {C--I};
  \node (sm2) at (3.2,1.75) {SM};
  \node (ci2) at (3.2,1.2) {C--I};
  \draw[e,->] (vp.east)--(sm1); \draw[e,->] (vp.east)--(ci1);
  \draw[e,->] (cp.east)--(sm2); \draw[e,->] (cp.east)--(ci2);
  \draw[e,->, dashed] (vp.north)--(cp.south) node[midway, right, text=gtghost] {\scriptsize builds to};
  \node[text=gtghost] at (3.2,2.25) {\scriptsize Transfer};
  \node[text=gtghost] at (0,-0.6) {\scriptsize edge = escape hatch (accessible)};
\end{tikzpicture}
\end{twocol}
\end{keiofig*}

There is, then, no level of LF in the old sense: no point after Spell-Out at
which the derivation carries on applying operations invisible to the
sensorimotor interface. What used to be called covert movement has to be
rethought. The standard move is the single-output model: movement is always
movement in the one syntax, and `covert' names the case in which the lower copy
is the one pronounced \citep{GroatONeil1996,Pesetsky2000,Bobaljik2002}. The idea
that the derivation reaches the interfaces piece by piece, rather than in one or
two packages, is older than phases, and it is worth crediting
a tradition with roots at Keio: Epstein, Groat, Kawashima \& Kitahara
\citeyearpar{EpsteinEtAl1998}, who derived c-command from the derivation itself,
and \citet{EpsteinSeely2002} on rule applications as cycles in a level-free
syntax.

(\ref{ex:transfer}) is also the model the box lives in: set an element aside at
a phase and consult it later. That is storage, not a covert component, and I say
so now so that the proposal of Lecture~3 can be recognised as a consequence of
(\ref{ex:transfer}) rather than as a new architecture. The single-output model
is the one the box will compete with. On it, QR is movement plus low
pronunciation; on the box the quantifier never moves at all, and the low position at
which the quantifier is pronounced is where its Case relation is. The question
this raises, and it is a fair one, is whether a single-cycle grammar that stores
quantifiers and consults them at later phases differs in anything but name from
the grammars that have stored quantifiers all along. Lecture~3 compares the box
with Cooper's store on exactly that point (\Hs{3}{sec:box-cooper}), and
Lecture~4 widens the comparison to the other frameworks
(\Hs{4}{sec:neighbours}).

\subsection{(e) Spines and externalisation (Lectures 1 and 2; rigidity in Lecture 3)}\label{sec:spines}

\begin{exe}
\ex \label{ex:spines}
\begin{xlist}
\ex \label{ex:EP} \textit{Extended projections.} N $<$ n $<$ Div $<$ \# $<$ D ($<$ Q); V $<$ \textit{v} $<$ Asp $<$ T $<$ C (\citealt{Grimshaw2005}; for the nominal, \citealt{Borer2005}).
\ex \label{ex:scopepositions} \textit{Scope positions.} RefP $>$ CP $>$ AgrSP $>$ DistP $>$ ShareP $>$ NegP $>$ AgrOP $>$ VP \citep{BeghelliStowell1997}.
\ex \label{ex:variation} \textit{Variation.} The computational system is uniform; variation lies in the features of functional heads and in externalisation \citep{BerwickChomsky2016,Roberts2019}.
\end{xlist}
\end{exe}

The two spines of (\ref{ex:EP}) have the same shape (Figure~\ref{fig:spines}).
The nominal one is Borer's: n categorises, Div divides (realised by a classifier
or by plural morphology), \# counts, and D determines; the Q above D is the
Q-layer of \S\ref{sec:Qlayer}. I draw heads only; the specifiers, where the action will be, come in \S\ref{sec:nominal} and in Lecture~2. The question of
\S\ref{sec:nominal} is where on the nominal spine quantificational force lives;
the question of Lecture~2 is where on the clausal spine it can be fixed.

\begin{keiofig}{fig:spines}{Two spines, one shape (\textit{n}P and AspP omitted). Left: the nominal extended projection, after \citet{Borer2005}, with the Q-layer on top. Right: the clausal one.}
\begin{forest}
[QP, SOtree, tiny, flat [Q, SOhead] [DP [D, SOhead] [\#P [\#, SOhead] [DivP [Div, SOhead] [NP [N, SOhead]]]]]]
\end{forest}\hspace{2.2em}%
\begin{forest}
[CP, SOtree, tiny, flat [C, SOhead] [TP [T, SOhead] [\textit{v}P [\textit{v}, SOhead] [VP [V, SOhead]]]]]
\end{forest}
\end{keiofig}

There are, broadly, two ways of thinking about where scope comes from. On one,
there are scope positions: designated heads on the clausal spine whose
specifiers host quantifiers of particular kinds (distributive universals in
DistP, group-denoting indefinites in ShareP and RefP, negative quantifiers in
NegP, and so on, in Beghelli and Stowell's version). On the other, scope is
given by an operation, such as QR, storage, type-shifting or continuations,
which can in principle target any position of the right type. The box is a
theory of the second kind, but its discharge sites, the phase edges, sit among
the positions of the first. The two pictures are not, in any case,
exclusive, since a theory of the second kind may well predict the positions of
the first; and the particles of Japanese and Korean, in Lecture~2, are where they
meet. Beghelli and Stowell's spine is also the one I needed some years ago for
floating quantifiers as scope markers (\S\ref{sec:FQ}), and it is the one the
box will let us dispense with.

(\ref{ex:variation}) is the working hypothesis about variation on which the
seminar leans hardest: the languages we will look at differ less in what their
syntax computes than in which steps of the computation they pronounce. English
pronounces quantificational force on the restriction (\textit{every}), Japanese
on a particle (\textit{mo}), Mandarin on the clausal spine (\textit{dou}), and
Japanese questions pronounce the scope site itself (\textit{ka}, \textit{no}).
Table~\ref{tab:externalisation} is the whole typological bet in one picture:
each row pronounces a different step of one computation. Lecture~2 builds a
parameter hierarchy on exactly this, and Lecture~3 argues that a language which fixes force at retrieval, and
pronounces that step, is scope-rigid. The claim runs in one direction only.
Pronounce the retrieving operator, and the order in which scope is fixed can be
read off an order the derivation already has; but a language may be rigid for
other reasons, as German is, and as Japanese is with universals that carry no
particle, so rigidity does not entail pronounced retrieval (\Hs{3}{sec:rigid}).

\begin{center}\small
\begin{tabularx}{\linewidth}{@{}l L l l@{}}
\toprule
Language & Force is pronounced\ldots & on & example\\
\midrule
English & on the restriction & D & \textit{every} book\\
Japanese & on a particle attached to the restriction & Q-particle & dono hon-\textit{mo}\\
Mandarin & on the clausal spine & Adv/Dist head & \textit{dou}\\
Japanese questions & at the scope site itself & C & dare-ga kita \textit{no}\\
\bottomrule
\end{tabularx}
\captionof{table}{Where force is pronounced: one computation, four externalisations.}\label{tab:externalisation}
\end{center}

\subsection{One more word, for Lecture~3}\label{sec:box-preview}

One more piece of vocabulary, which Lecture~3 will use: the \emph{box}
(\citealt{Chomsky2024}, from the 2023 Keio--EMU lectures; see now
\citealt{KitaharaSeely2025}). Informally, an element that has been internally
merged is set aside in a box, written \BX{X}: segregated from the ongoing
derivation, positionally neutral, but still consulted for instructions by phase
heads at later cycles. Four ideas are bundled here. The element is
\emph{stored}, set aside from the structure being built; it is
\emph{segregated}, in that its box is neither a phase edge nor a scope position;
it remains \emph{accessible}, consulted rather than re-merged, so that nothing moves; and it is eventually \emph{retrieved}, discharged later, which is where its scope and force are fixed. (So in Chomsky's version; in mine, as Lecture~3 will say, discharge is itself an External Merge of the boxed element.)

I will argue that the box is, to a first approximation, a syntactic Cooper
store, and that QR can be rebuilt on it without covert movement and without
adding a variable to the workspace. I will not say more here, except to flag one
difference. In Chomsky's box nothing but Internal Merge to the edge puts an element in a box; in mine the element is boxed by a separate application of MERGE, Self-Merge, which I treat as a degenerate case of External Merge. Lecture~3 states the
difference openly, and argues for it (\Hs{3}{sec:box}).

\begin{excursus*}[For readers from other frameworks: a map]
  Each of the five ideas has a counterpart elsewhere, and it is worth saying so
  now, since how far the counterparts are notational variants is a question the
  lectures come back to (\Hs{3}{sec:box-cooper}; \Hs{4}{sec:neighbours}). HPSG
  and LFG have no derivation, and so no Inclusiveness in this sense. HPSG long
  kept quantifiers in a store and retrieved them at a sign of the right type
  \citep{Cooper1983,PollardSag1994}, and now more often leaves scope
  underspecified; LFG, with glue semantics, gets scope from the order of steps
  in a proof. CCG gets
  scope without a store, from the combinatorics of the derivation itself
  \citep{Steedman2012}, and variable-free semantics does without variables
  altogether \citep{Jacobson1999}. Between them they show what a grammar that
  takes the ban on indices literally might look like. The table is a map of the
  four excursions that follow. The HPSG and LFG rows are the ones the
  comparison with storage in Lecture~3 is about; the CCG and variable-free rows are
  the ones the inversion of \S\ref{sec:inversion} is about.

\medskip

\begin{tabularx}{\linewidth}{@{}P{2.3cm}LLLL@{}}
\toprule
& Derivation? & Variable? & Scope fixed by & Locality from\\
\midrule
Minimalism (QR) & yes & a copy, converted & a position (IM) & phases, economy\\
HPSG & no: constraints & indexed placeholder & retrieval from a store, or left underspecified & stated on retrieval, or on resolution\\
LFG + glue & no & none: proof terms & order of proof steps & added: a path constraint, or modes of the logic\\
CCG & yes: combinatory & none & the combinatorics & the combinatory rules\\
Variable-free semantics & -- & none: pronouns are identity maps & the \textsf{z} and \textsf{g} combinators & type-driven\\
\bottomrule
\end{tabularx}
\end{excursus*}

\begin{excursus}[HPSG: a store in the sign, and underspecification]
  In the HPSG of \citet[ch.~8]{PollardSag1994} every sign carries a feature
  \textsc{qstore}, the set of quantifiers not yet scoped, and a feature
  \textsc{retrieved}, the list of those scoped at that node. Two principles do
  the work. By the Quantifier Inheritance Principle, the \textsc{qstore} of a
  phrase is the union of its daughters' \textsc{qstore}s minus whatever is
  \textsc{retrieved} there; by the Scope Principle, the \textsc{quants} list of
  a phrase is its \textsc{retrieved} list prepended to the \textsc{quants} of
  its head daughter, and order in that list \emph{is} scope. Retrieval is
  possible at any phrasal node of the right content type. This is Cooper
  storage, written into the feature structure. The first sentence of \S\ref{sec:vote}, (\ref{ex:someone}), comes
  out as follows, on its surface reading; retrieving in the order
  $\langle\fbox{\scriptsize 2},\fbox{\scriptsize 1}\rangle$ would give
  $\forall > \exists$:
\[\left[\begin{array}{@{}l@{\ \ }l@{}}
  \textsc{phon} & \langle \textit{someone read every book}\rangle\\[2pt]
  \textsc{cont} & \left[\begin{array}{@{}l@{\ \ }l@{}}
      \textsc{quants} & \langle \fbox{\scriptsize 1}\,\exists x,\ \fbox{\scriptsize 2}\,\forall y\rangle\\
      \textsc{nucleus} & \mathit{read}(x,y)
    \end{array}\right]\\[8pt]
  \textsc{qstore} & \{\ \}\\
  \textsc{retrieved} & \langle \fbox{\scriptsize 1}, \fbox{\scriptsize 2}\rangle
\end{array}\right]\]
Later versions tie retrieval to lexical heads \citep{PollardYoo1998,GinzburgSag2000}. The point for Lecture~3 is that HPSG's store has no locality of its own: the clause-boundedness of retrieval is stipulated, or follows from the content type of the retrieval site. The box inherits it from phases, and that is the whole difference I will press.

Minimal Recursion Semantics \citep{CopestakeEtAl2005}, the semantics of the LinGO grammars (including the Japanese grammar Jacy), takes the opposite course and does not resolve scope in the grammar at all. A sentence denotes a flat \emph{bag} of elementary predications, each carrying a \emph{handle}, a label such as $h_1$ or $h_5$; a quantifier has two \emph{holes}, one for its restriction and one for its body, to be filled by handles; and a few constraints say which fillings are allowed:
\[\begin{array}{@{}l@{}}
  h_0\\
  h_1{:}\ \mathit{some}(y, h_2, h_3)\qquad h_4{:}\ \mathit{person}(y)\\
  h_5{:}\ \mathit{every}(x, h_6, h_7)\qquad h_8{:}\ \mathit{book}(x)\\
  h_9{:}\ \mathit{read}(y, x)\\
  h_2 =_q h_4,\quad h_6 =_q h_8
\end{array}\]
Read it as follows. $h_0$ is the \emph{top}, the handle of whatever ends up outermost. The line $h_1{:}\ \mathit{some}(y, h_2, h_3)$ says that the predication labelled $h_1$ is \textit{some}, binding $y$, with its restriction whatever fills the hole $h_2$ and its body whatever fills the hole $h_3$. The constraint $h_2 =_q h_4$, `equal modulo quantifiers', says that the restriction of \textit{some} is \textit{person}, possibly with other quantifiers plugged in between. The restrictions are thus fixed, and the bodies, $h_3$ and $h_7$, are left open: nothing says which quantifier fills the top, or whose body \textit{read} ends up in, and that is exactly where scope is left undecided. A \emph{resolution} plugs every hole with a handle so that the result is a tree with $h_0$ at the top, and there are two. On the first, $h_0 = h_1$, $h_3 = h_5$, $h_7 = h_9$: \textit{some} is on top, the body of \textit{some} is \textit{every}, and the body of \textit{every} is \textit{read}, which is $\exists > \forall$. On the second, $h_0 = h_5$, $h_7 = h_1$, $h_3 = h_9$: \textit{every} is on top, the body of \textit{every} is \textit{some}, and the body of \textit{some} is \textit{read}, which is $\forall > \exists$. The variable is written, twice: $x$ in \textit{every} and in \textit{read}, $y$ in \textit{some} and in \textit{read}. The grammar simply does not say which resolution holds; that is left to whatever consumes the representation. This is the opposite of the box: since the syntax never fixes scope, no syntactic locality can constrain it, and the rigidity of Japanese (\ref{ex:dareka-dono}) would have to be stated as an extra constraint on resolutions, added to a representation that is by design indifferent to scope.
\end{excursus}

\begin{excursus*}[LFG: scope from proofs (glue semantics)]
In LFG with glue semantics \citep{DalrympleEtAl1997,Dalrymple2001} the syntax is the f-structure, and meaning is assembled by deduction. Each word contributes a meaning constructor, a premise of linear logic, to be used exactly once, and a \emph{proof} that the premises yield the meaning of the sentence is a meaning for it. For (\ref{ex:someone}):

\medskip
\begin{twocol}{0.36}
\centering
$\left[\begin{array}{@{}l@{\ \ }l@{}}
  \textsc{pred} & \text{`read}\langle\textsc{subj},\textsc{obj}\rangle\text{'}\\
  \textsc{subj} & s{:}\left[\textsc{pred}\ \text{`someone'}\right]\\
  \textsc{obj}  & o{:}\left[\begin{array}{@{}l@{\ }l@{}}\textsc{pred} & \text{`book'}\\ \textsc{spec} & \text{`every'}\end{array}\right]
\end{array}\right]_f$
\midcol
\footnotesize
\begin{align*}
  \textit{read}: &\ \lambda x\lambda y.\mathit{read}(y,x)\ :\ o_\sigma \multimap (s_\sigma \multimap f_\sigma)\\
  \textit{someone}: &\ \lambda P.\exists y[\mathit{person}(y)\wedge P(y)]\ :\ \forall H.(s_\sigma \multimap H)\multimap H\\
  \textit{every book}: &\ \lambda P.\forall x[\mathit{book}(x)\to P(x)]\ :\ \forall H.(o_\sigma \multimap H)\multimap H
\end{align*}
\end{twocol}

\medskip
In words: $s_\sigma$, $o_\sigma$ and $f_\sigma$ are the meaning slots of the subject, the object and the clause, and $\multimap$ reads `consume the resource on the left to produce the one on the right', each premise being used exactly once. \textit{Read} consumes an object meaning, then a subject meaning, and yields a clause meaning. \textit{Every book} consumes any meaning $H$ that still has an object-shaped gap in it, $o_\sigma \multimap H$, and returns $H$ with the gap filled by the quantifier: it can take any clause-sized meaning as its scope, as long as that meaning is still waiting for an object. The universal over $H$ is where scope flexibility lives; the quantifier is tied not to a position but to a gap.

Scope ambiguity is then a matter of alternative proofs from one set of premises. On the first, suppose, for the sake of argument, that the object meaning is some $x$ ($x{:}o_\sigma$); combine it with \textit{read} and with \textit{someone}, which gives `someone read $x$', a clause meaning, $\exists y[\mathit{person}(y)\wedge\mathit{read}(y,x)]{:}f_\sigma$; now withdraw the supposition ($\multimap$-introduction). What remains is a clause with an object-shaped gap, $o_\sigma\multimap f_\sigma$, and that is exactly what \textit{every book} consumes, with $H = f_\sigma$; the result is $\forall x[\mathit{book}(x)\to\exists y[\mathit{person}(y)\wedge\mathit{read}(y,x)]]$, the inverse reading. On the second, suppose instead a subject meaning $y$ ($y{:}s_\sigma$); combine it with \textit{read} and with \textit{every book}, `$y$ read every book', $\forall x[\mathit{book}(x)\to\mathit{read}(y,x)]{:}f_\sigma$; withdraw the supposition, leaving a clause with a subject-shaped gap, $s_\sigma\multimap f_\sigma$, and feed it to \textit{someone}, which yields $\exists y[\mathit{person}(y)\wedge\forall x[\mathit{book}(x)\to\mathit{read}(y,x)]]$, the surface reading. The quantifier whose gap is filled last takes widest scope. There is no store, no movement and no LF; no variable of the syntax is ever bound, and the $\lambda$-terms are proof terms.

Two things about this are worth carrying forward. The analogue of a variable is
the \emph{hypothesis} that is assumed and then discharged: a step of the proof,
not a part of the sentence. That is, I think, the nicest analogue in any
framework of what I will claim in Lecture~4, that the variable is a discharged
relation rather than an item. And locality, why a quantifier cannot scope out of
an island, is not given by the logic: glue is, in the words of
\citet[147]{Gotham2021}, `towards the unconstrained end of the spectrum'. Two
ways of constraining it have been worked out. One restricts the path along
which the scope level $H$ may be found. The other makes the logic multi-modal,
so that an island is a mode which the premise of a quantifier cannot cross, and
the constraint then follows from the rules of inference
\citep{Gotham2021}. On the first, locality is stated; on the second
it is built into the logic, at the price of putting a good deal of syntax
there. \Hs{4}{sec:nb-locality} sets both beside the box.
\end{excursus*}

\begin{excursus*}[CCG: scope from the combinatorics]
CCG \citep{Steedman2000,Steedman2012} has a derivation, but no movement and no store. The combinatory rules we need are the following, where the combinators are Curry's \textbf{B} (composition) and \textbf{T} (type-raising; Smullyan's `thrush'); there are no variables in the rules, only functions and their arguments:

\medskip
\begin{twocol}{0.5}
\footnotesize
\begin{tabular}{@{}ll@{}}
\toprule
Forward application ($>$) & $X/Y\ \ Y \Rightarrow X$\\
Backward application ($<$) & $Y\ \ X\backslash Y \Rightarrow X$\\
Forward composition ($>\mathbf{B}$) & $X/Y\ \ Y/Z \Rightarrow X/Z$\\
Forward type-raising ($>\mathbf{T}$) & $X \Rightarrow T/(T\backslash X)$\\
Backward type-raising ($<\mathbf{T}$) & $X \Rightarrow T\backslash(T/X)$\\
\bottomrule
\end{tabular}
\midcol
\footnotesize
\begin{tabular}{@{}ll@{}}
\toprule
\textit{read} & $(S\backslash NP)/NP : \lambda x\lambda y.\mathit{read}'xy$\\
\textit{every book} & $NP : \lambda p.\forall x[\mathit{book}'x \to p\,x]$\\
\textit{someone} & $NP : \mathit{skolem}'\,\mathit{person}'$\\
\bottomrule
\end{tabular}
\end{twocol}

\medskip
In words: a category $X/Y$ is a function that wants a $Y$ on its right and gives an $X$; $X\backslash Y$ wants its $Y$ on the left. Application feeds a function its argument. Composition, $\mathbf{B}$, lets a function that wants a $Y$ combine with a function that will \emph{produce} a $Y$ once it gets a $Z$: the two are chained into one function still waiting for the $Z$. Type-raising, $\mathbf{T}$, turns an argument into a function that takes its own functor as argument, so that a subject becomes the thing that applies rather than the thing applied to.

\citet{Steedman2012} draws the line not between quantifiers and non-quantifiers but between \emph{true} quantifiers, the universals, and \emph{indefinites}, which are Skolem terms. A Skolem term stands for a particular but as yet unnamed individual. When it is \emph{specified} inside the scope of a universal it becomes a function of that universal's variable, a possibly different person for each book; specified outside, it is a constant, one person for all the books. The sentence of \S\ref{sec:vote} has two derivations:

\medskip
\begin{twocol}{0.5}
\centering\small\textbf{Application only}
\[\begin{array}{@{}c@{\ \ }c@{\ \ }c@{}}
  \textit{someone} & \textit{read} & \textit{every book}\\
  NP & (S\backslash NP)/NP & NP\\
  \cline{2-3}
  & \multicolumn{2}{c}{S\backslash NP \quad {}^{>}}\\
  \cline{1-3}
  \multicolumn{3}{c}{S \quad {}^{<}}
\end{array}\]
\midcol
\centering\small\textbf{Composition, then the object}
\[\begin{array}{@{}c@{\ \ }c@{\ \ }c@{}}
  \textit{someone} & \textit{read} & \textit{every book}\\
  NP & (S\backslash NP)/NP & NP\\
  \cline{1-1}
  S/(S\backslash NP)\ {}^{>\mathbf{T}} & & \\
  \cline{1-2}
  \multicolumn{2}{c}{S/NP \quad {}^{>\mathbf{B}}} & \\
  \cline{1-3}
  \multicolumn{3}{c}{S \quad {}^{>}}
\end{array}\]
\end{twocol}

\medskip On the left, \textit{read} takes \textit{every book} as its argument, and the subject then takes the result; the object has the category of an $NP$ but the meaning of a generalised quantifier, so it applies to the verb's open argument rather than filling it. On the right, the subject is first type-raised into a function that wants a verb phrase; it then composes with \textit{read} into one function, `someone read', still waiting for an object, and that non-standard constituent takes \textit{every book} as its argument. The constituent is a genuine one (it coordinates, and it is a possible intonational phrase), and that is Steedman's main argument for the combinators. One has to be careful,
though: in Steedman's system the scope ambiguity of (\ref{ex:someone}) is not a
matter of these two derivations but of \emph{when} the Skolem term for
\textit{someone} is specified. The inverse reading does not come from moving
\textit{every book}; it comes from specifying the Skolem term inside the scope
of the universal, as a function of $x$. So CCG has no variable in the syntax, no
store and no movement; what it has instead is a lexical division between
universals and indefinites, and derivational constituency. The family
resemblance to \S\ref{sec:indefinites}, indefinites as restrictions without
force, and to Reinhart's choice functions, is worth noting.
\end{excursus*}

\begin{excursus*}[Variable-free semantics: pronouns as identity maps]
  In the variable-free semantics of \citet{Jacobson1999} there are no variables,
  no assignment functions and no indices. A pronoun denotes the identity
  function on individuals, $\sem{\textit{he}} = \lambda x.x$, of type
  $\langle e,e\rangle$, and an expression containing an unbound pronoun simply
  denotes a function from individuals; that is all `containing a free variable'
  amounts to. Two combinators do the rest. The Geach rule \textsf{g} passes the
  dependence up, and the rule \textsf{z} binds:
\begin{align*}
\mathsf{g}(f) &= \lambda V_{\langle e,a\rangle}\lambda x.\ f(V(x)) && \text{for $f$ of type $\langle a,b\rangle$}\\
\mathsf{z}(f) &= \lambda V_{\langle e,a\rangle}\lambda x.\ f(V(x))(x) && \text{for $f$ of type $\langle a,\langle e,b\rangle\rangle$}
\end{align*}
In words, \textsf{g} says that a function $f$ that wants an $a$ may take instead something that \emph{would be} an $a$ if it were given an individual, and the need for that individual is handed up to the whole; \textsf{z} says that a two-place verb may be shifted so that its \emph{subject} also supplies the individual its object was waiting for. So \textit{his mother} and \textit{loves his mother} come out as follows, the first of type $\langle e,e\rangle$:
\begin{align*}
\textit{his mother}: &\quad \mathsf{g}(\sem{\textit{mother-of}})(\lambda x.x) = \lambda x.\mathit{mother}(x)\\
\textit{loves his mother}: &\quad \mathsf{z}(\sem{\textit{loves}})(\lambda x.\mathit{mother}(x)) = \lambda x.\mathit{loves}(\mathit{mother}(x))(x)
\end{align*}
The first is not a mother but the function `mother of', still waiting for an individual; the second is the property of loving one's own mother. Binding has happened, and no pronoun was ever indexed. Binding is thus a property of the \emph{verb}, shifted by \textsf{z}, not a relation between an antecedent and an indexed pronoun:
\begin{align*}
\sem{\textit{Every man loves his mother}} &= \sem{\textit{every man}}\bigl(\lambda x.\mathit{loves}(\mathit{mother}(x))(x)\bigr)
\end{align*}

Jacobson's is the cleanest existing answer to the inverse question of
\S\ref{sec:aims}: a semantics that never writes $x$. What it buys is
considerable. It takes Inclusiveness literally, since nothing in the semantics
corresponds to an index; it is directly compositional, with every constituent
assigned a meaning and no LF; and Fine's point (\S\ref{sec:pronouns}) is built
in, since there is no such thing as `the same variable', only the same function,
applied. What it owes is, for our purposes, as instructive. Quantifier scope
still needs an operation, Hendriks-style type-shifting or lifting
\citep{Hendriks1993}, which is QR by another name, or a store. The particles of
Japanese and Korean pronounce something, and one would like to know what, if
there is no index to pronounce. A pronoun that resists binding, like
\textit{kare} (Lecture~4), must be lexically special, since the \textsf{z} rule
binds from the subject, and crossover is something it has to say something extra
about.

The moral for the seminar is that logic and semantics both know how to quantify
without variables. The question is whether the \emph{syntax} of natural language
does, and that is what the pronounced particles of Lecture~2 let us test. The
inversion of \S\ref{sec:inversion} is not that variable-free systems are
impossible; it is that natural language, which could have used one, appears not
to, except in the one respect that it never pronounces the variable.
\end{excursus*}

% =====================================================================
\section{The problem: where does the variable come from?}\label{sec:problem}

This section has three aims. First, to recall what QR was for, and the two jobs
it did at once; second, to show why, in the architecture of
\S\ref{sec:architecture}, it has no variable to bind, since the copy theory
removes the trace and the repair adds to the workspace; and third, to set out
the three places the variable might come from instead, which are the three
answers the rest of the seminar weighs. The section is historical in shape, but
the point is architectural. Nothing in it is a criticism of May or of Fox on
their own terms; it is an audit of what each assumption costs once
(\ref{ex:inclusiveness}) and (\ref{ex:NTC}) are taken literally.

\subsection{QR and its two jobs}\label{sec:QR}

\begin{exe}
\ex \label{ex:QR} \textit{Quantifier Raising} \citep{May1977,May1985}: a quantified phrase adjoins to a clausal projection (TP, or \textit{v}P), leaving behind a trace that is interpreted as a variable bound by it.
\end{exe}

QR was meant to solve two problems at once. The first is a type mismatch. If \textit{read} denotes a relation of type $\langle e,\langle e,t\rangle\rangle$ and \textit{every book} a generalised quantifier of type $\langle\langle e,t\rangle,t\rangle$, then [\textit{read} [\textit{every book}]] cannot be composed by function application. Once the object has raised, its sister is a derived predicate of type $\langle e,t\rangle$, and the quantifier can take it as its argument (Figure~\ref{fig:type-mismatch}). The trace is of type $e$, so the verb composes with it; the binder index turns the clause into a function of type $\langle e,t\rangle$; and both steps need the index.

\begin{keiofig*}{fig:type-mismatch}{QR's first job. Left: before QR, no rule composes V with QP. Right: after QR, the trace is of type $e$, and the index turns its sister into a predicate of type $\langle e,t\rangle$, which the quantifier takes as its argument. Types are in grey; the tree on the right is Heim and Kratzer's.}
\begin{twocol}{0.42}
\centering\small Before QR\par\medskip
\begin{forest}
[{VP \quad \textcolor{gtalert}{?}}, SOtree, slim, flat
  [{V\\ \ghost{$\langle e,\langle e,t\rangle\rangle$}}, SOhead [{\LI{read}}, SOlexical]]
  [{QP\\ \ghost{$\langle\langle e,t\rangle,t\rangle$}} [{\LI{every book}}, SOlexical]]]
\end{forest}
\midcol
\centering\small After QR\par\medskip
\begin{forest}
[{TP \quad \ghost{$t$}}, SOtree, slim, flat
  [{QP\\ \ghost{$\langle\langle e,t\rangle,t\rangle$}}, SOmoved [{\LI{every book}}, SOlexical]]
  [{$\lambda_1$ \quad \ghost{$\langle e,t\rangle$}}
    [{TP \quad \ghost{$t$}}
      [DP [{\LI{someone}}, SOlexical]]
      [VP [V, SOhead [{\LI{read}}, SOlexical]] [{$t_1$ \quad \ghost{$e$}}, SOtrace]]]]]
\end{forest}
\end{twocol}
\end{keiofig*}

The second problem is scope. The inverse reading of (\ref{ex:someone}) is the one on which \textit{every book} outranks \textit{someone}, and on the assumption that semantic scope is c-command at the level where interpretation takes place, \textit{every book} must c-command \textit{someone} there, although it does not in the surface string:

\begin{exe}
\ex \label{ex:QR-LFs}
\begin{xlist}
\ex \label{ex:QR-vP} \lbr{TP} someone \lbr{\textit{v}P} every book [1 \lbr{\textit{v}P} $t_{\textit{someone}}$ read $t_1$ ]]]] \src{$\exists > \forall$}
\ex \label{ex:QR-TP} \lbr{TP} every book [1 \lbr{TP} someone read $t_1$ ]]] \src{$\forall > \exists$}
\end{xlist}
\end{exe}

In (\ref{ex:QR-vP}) the object raises only to \textit{v}P, enough to solve the type problem and no more \citep{Johnson2000}; in (\ref{ex:QR-TP}) it raises to TP, above the subject (Figure~\ref{fig:two-LFs}). The difference between the two readings is a difference in landing site, which is to say in the syntax; the semantics is the same rule applied twice. Note, too, that in (\ref{ex:QR-vP}) the subject's own copy in Spec-\textit{v}P has been drawn as a trace, which is of course the same problem over again.

\begin{keiofig*}{fig:two-LFs}{QR's second job: two landing sites, two scopes. Left, (\ref{ex:QR-vP}): the object adjoins to \textit{v}P, below the subject. Right, (\ref{ex:QR-TP}): it adjoins to TP, above it.}
\begin{twocol}{0.5}
\centering
\begin{forest}
[TP, SOtree, xslim, flat
  [{\LI{someone}}, SOmoved]
  [\textit{v}P
    [{\LI{every book}}, SOmoved]
    [{$1$} [\textit{v}P [{$t_{\textit{someone}}$}, SOtrace] [VP [{\LI{read}}, SOlexical] [{$t_1$}, SOtrace]]]]]]
\end{forest}
\midcol
\centering
\begin{forest}
[TP, SOtree, xslim, flat
  [{\LI{every book}}, SOmoved]
  [{$1$} [TP [{\LI{someone}}, SOlexical] [VP [{\LI{read}}, SOlexical] [{$t_1$}, SOtrace]]]]]
\end{forest}
\end{twocol}
\end{keiofig*}

Both jobs, then, rest on one idea: that the sentence contains a position, the
trace, which is read as a variable. It is worth keeping the two jobs apart
nonetheless, because they need not be done by the same thing. The first has
purely semantic solutions, by type-shifting \citep{Hendriks1993}, and the box of
Lecture~3 will do the second syntactically while leaving the first to the token of
type $e$ that fission supplies. The point for now is only that QR did both.

QR has two further properties, and both will matter in Lecture~3. It is clause-bound, by and large (\citealt{May1977}; \citealt{Cecchetto2004}; see \citealt{FarkasGiannakidou1996} on the exceptions), where \textit{wh}-movement is not. In (\ref{ex:clausebound-a}) the universal cannot take scope over the matrix subject, while in (\ref{ex:clausebound-b}) the \textit{wh}-phrase crosses the same clause boundary without difficulty:

\begin{exe}
\ex \label{ex:clausebound}
\begin{xlist}
\ex \label{ex:clausebound-a} Someone said that every book was read. \src{$\exists > \forall$; *$\forall > \exists$}
\ex \label{ex:clausebound-b} Which book did someone say that Mary read? \src{(long \textit{wh}-movement)}
\end{xlist}
\end{exe}

And it is subject to an economy condition: it applies only if it makes a semantic difference \citep{Fox2000}. The evidence is ellipsis parallelism, which blocks inverse scope in the first conjunct when QR would be vacuous in the second:

\begin{exe}
\ex \label{ex:economy} A boy admires every teacher, and Mary does too. \src{*$\forall > \exists$}
\end{exe}

Clause-boundedness can also be derived without a dedicated rule, by taking QR to be A-movement \citep{Hornstein1995}. Lecture~3 has to derive clause-boundedness in its own terms, from the discharge of a $\theta$-feature at the phase (\Hs{3}{sec:boxQR}); economy is another matter.

On economy, I do not think that Scope Economy should be restated inside the derivation at all. Both derivations of (\ref{ex:someone}) are generated, and the preference for the surface reading is left to processing; what remains owed is an account of ellipsis parallelism of the kind in (\ref{ex:economy}).

\subsection{The copy theory removes the trace}\label{sec:copy}

Now recall \S\ref{sec:inclusiveness}. If movement leaves a copy rather than a trace, (\ref{ex:QR-TP}) is really (\ref{ex:copy-LF}):

\begin{exe}
\ex \label{ex:copy-LF} \lbr{TP} every book [1 \lbr{TP} someone read every book ]]]
\end{exe}

The lower copy is a quantifier, not a variable (Figure~\ref{fig:copy-tree}). Interpreted as it stands it gives either a type clash, which is the first of QR's jobs over again, or, if the clash is somehow resolved, the wrong meaning: `every book is such that someone read every book'. The binder index of (\ref{ex:PA-a}) has nothing to bind, and the index itself was never in the workspace. Students are often shown the lower copy already struck through; the point here is that nothing in the syntax strikes it through.

\begin{keiofig}{fig:copy-tree}{The copy theory: the lower copy of \textit{every book} is drawn in full, because it is there in full.}
\begin{forest}
[TP, SOtree, xslim
  [QP, SOmoved [{\LI{every book}}, SOlexical]]
  [TP [DP [{\LI{someone}}, SOlexical]] [VP [V, SOhead [{\LI{read}}, SOlexical]] [QP, SOtrace [{\LI{every book}}, SOlexical]]]]]
\end{forest}
\end{keiofig}

\citet{Fox2002,Fox2003} proposed that the lower copy is converted, before interpretation, into something that contains a variable:\footnote{The citation is sometimes given as Fox (1999), which is about reconstruction and the interpretation of chains; the two rules in (\ref{ex:TC}) are stated in \citet{Fox2002,Fox2003}.}

\begin{exe}
\ex \label{ex:TC} \textit{Trace Conversion} \citep{Fox2002,Fox2003}:
\begin{xlist}
\ex \textit{Variable Insertion:}\\ (Det) Pred $\;\rightarrow\;$ (Det) [Pred \mbox{$\lambda y\,(y = x)$}]
\ex \textit{Determiner Replacement:}\\ (Det) [Pred \mbox{$\lambda y\,(y = x)$}] $\;\rightarrow\;$ \textit{the} [Pred \mbox{$\lambda y\,(y = x)$}]
\end{xlist}
\ex \label{ex:TC-result} [ every book [ $\lambda x$ [ someone read\\ \hspace*{1.5em}[ the [ book \mbox{$\lambda y\,(y = x)$} ]]]]]\\
`for every book $x$, someone read the book identical to $x$'
\end{exe}

(\ref{ex:TC-result}) has the right truth conditions, it carries the right presupposition (that $x$ is a book), and the restriction left in the lower position is what Fox needed for reconstruction effects with Condition~C and for antecedent-contained deletion. It is a good solution to the problem it was designed for: the right answer, I think, to the question of how a copy can be read as a variable. The objection is not that it fails, but that what it adds was never in the workspace. Figure~\ref{fig:TC-picture} shows what the two rules do to the lower copy, and Figure~\ref{fig:TC-workspace} states the cost with the workspace.

\begin{keiofig*}{fig:TC-picture}{Trace Conversion in a picture. Left: the lower copy, as MERGE left it. Right: after Variable Insertion and Determiner Replacement; the framed nodes mark the added material.}
\begin{twocol}{0.42}
\centering\small The lower copy, as MERGE left it\par\medskip
\begin{forest}
[QP, SOtree, slim, SOtrace [Q, SOhead [{\LI{every}}, SOlexical]] [NP [{\LI{book}}, SOlexical]]]
\end{forest}
\midcol
\centering\small After Variable Insertion and Determiner Replacement\par\medskip
\begin{forest}
[DP, SOtree, slim
  [{D}, SOhead, draw=gtaccent, rounded corners=2pt [{\LI{the}}, SOlexical]]
  [NP [{\LI{book}}, SOlexical] [{$\lambda y\,(y = x)$}, draw=gtaccent, rounded corners=2pt]]]
\end{forest}
\end{twocol}
\end{keiofig*}

The cost lies elsewhere. Variable Insertion adds a variable and an identity
predicate; Determiner Replacement adds a definite determiner. None of the three was in the workspace. And the quantifier \textit{every}, which was in the workspace, has
been removed, which No Tampering also forbids; this half of the cost is usually
overlooked. Trace Conversion is, of course, a solution to the tension between
the copy theory and the need for an operator--variable chain at the point of
interpretation, but it is not a solution to the problem posed by Inclusiveness,
unless we are prepared to assume that the workspace contained an unused definite
determiner, and that this determiner survived long after the workspace had
become the input to semantic interpretation. Nor does it help to say that the
conversion happens `at LF', and so escapes Inclusiveness: in the architecture of
\S\ref{sec:cyclicity} there is no LF, and Transfer hands the phase complement to
the conceptual--intentional interface as it is.

\begin{keiofig*}{fig:TC-workspace}{The cost, stated with the workspace: three additions and one deletion.}
\begin{twocol}{0.44}\vspace{0pt}
\begin{workspace}[what MERGE delivered to C--I]
  \scriptsize\WSeq{\{\mathrm{QP},\allowbreak\{\LI{someone},\allowbreak\{\text{T},\allowbreak\{\LI{someone},\allowbreak\{\vs,\{\LI{read},\mathrm{QP}\}\}\}\}\}\}}
  \par\smallskip\footnotesize where $\mathrm{QP} = \{\LI{every},\LI{book}\}$, the same object twice.
\end{workspace}
\midcol\vspace{0pt}
\small\raggedright
\vspace*{0.6em}
What (\ref{ex:TC-result}) needs the lower QP to be:\\[0.3em]
\quad$\{\LI{the},\{\LI{book},\ \lambda y\,(y=x)\}\}$\\[0.6em]
\footnotesize
\begin{tabular}{@{}ll@{}}
\toprule
\LI{the} & $\notin$ any workspace of the derivation\\
$\lambda y\,(y=x)$ & $\notin$ any workspace of the derivation\\
$x$ & $\notin$ any workspace of the derivation\\
\LI{every} & $\in$ the workspace, but gone from (\ref{ex:TC-result})\\
\bottomrule
\end{tabular}
\end{twocol}
\end{keiofig*}

The semantic versions of the idea (\citealt{Sauerland2004};
\citealt{Pasternak2020}) let the semantics, rather than the syntax, read the
lower copy as a definite description,
$\iota y\,[\mathit{book}(y) \wedge y = x]$. This is good as far as it goes, but
it concedes the point: the syntax delivers no variable, and the semantics has to
manufacture one out of a full copy of the quantifier. It is only fair to add
that Sauerland's point is that the restriction of the lower copy is semantically
active, as Condition~C reconstruction shows, and that any theory has to capture
this; the box will do so by interpreting the restriction at the token (Lecture~3).
Multidominance (\S\ref{sec:merge}; e.g.\ \citealt{Johnson2012}) changes what a
copy is, one object with two mothers, so that there is nothing to convert. But
it still has to say which occurrence is read as a variable, and how; the
question has moved from the syntax to the linearisation and the interpretation
of the graph. Johnson's own concern is in fact with why \textit{wh}-movement and
QR are pronounced differently, with the restriction shared between the two
positions. I take it, then, that movement plus copies does not by itself answer
the question of where the variable comes from, and that the first of the three
answers is in trouble before the seminar has properly started.

\subsection{When do indices matter?}\label{sec:indices-matter}

One might reply that the indices in (\ref{ex:PA}) and (\ref{ex:QR-LFs}) are
harmless notation, a way of writing down relations the grammar computes anyway.
\citet[114]{Chomsky2000keio} draws exactly this line: some indices are innocuous
notation, while others function within the computation, and his example of the
second kind is the percolation of indices. The question, then, is not whether we
write subscripts but whether the computation consults them. The index that
triggers Predicate Abstraction in (\ref{ex:PA-a}) is consulted: without it, the
rule does not know which value of the assignment to shift. Keep this distinction
in mind, because it will return to confront the box in Lecture~4. This is also where
honesty about my own proposal has to start: the $\theta$-feature on which my
version of the box depends must not turn out to be a percolating index of the
forbidden kind. The way out I will take is that it does not percolate at all,
and that the phase head finds the undischarged token by Minimal Search
(\Hs{4}{sec:percolation}).

\subsection{Pronouns are not variables}\label{sec:pronouns}

The natural reply is that language has variables after all: pronouns. But
pronouns are not variables in the logician's sense. They carry
$\varphi$-features and presuppositions (\textit{she} presupposes a female
referent), they can be referential or E-type as well as bound, and whether a
given pronoun is bound depends on its structural relation to a quantifier, not
on the pronoun alone. Japanese even has a pronoun, \textit{kare}, which famously
resists bound-variable readings altogether \citep{SaitoHoji1983,Hoji1991}; we
will meet it in Lecture~4 (\Hs{4}{sec:testbed}). In (\ref{ex:kare}) \textit{kare}
cannot be read as bound by the universal:

\begin{exe}
\ex \label{ex:kare}
\gll Dono gakusei-mo kare-no sensei-o sonkeisite iru.\\
     which student-\gl{mo} he-\gl{gen} teacher-\gl{acc} respect-\gl{te} be.\gl{npst}\\
\glt `Every student respects his teacher.' \src{(\textit{kare} not bound)}
\end{exe}

If a language had a pronounced variable, \textit{kare} would not be it: it is
the pronoun that refuses to be one. Resumptive pronouns, which might look like
pronounced variables, are pronouns too, with their $\varphi$-features intact.

There is a deeper reason, which philosophers will recognise as Fine's antinomy
of the variable \citep{Fine2003,Fine2007}. The semantic role of $x$ in `$x > 0$'
is the same as that of $y$ in `$y > 0$'; yet `$x > y$' and `$x > x$' differ in
meaning. So the difference cannot lie in anything intrinsic to the two variables
taken one at a time. If variables had meanings of their own, variables with the
same meaning could be substituted for one another without change of meaning, and
they cannot; if they had no meaning, open formulas would mean nothing, and they
do. Fine's resolution is that the meaning lies in the relation between their
occurrences, in whether they are, in his term, \emph{coordinated}. This is
semantic relationism: the view that the semantic relations between expressions
need not be reducible to their intrinsic semantic features. Fine's own examples
are as much about names (\textit{Cicero} and \textit{Tully}) as about variables,
but the variable case is the one that bears on us. Sameness of variable is a
relation. This is the thought Lecture~4 tries to make syntactic: an index is not a
feature of a token but a relation between tokens. A grammar that cannot write
$x$ may yet be able to compute the relation that the occurrences of $x$ bear to
one another, and Form Copy (\ref{ex:formcopy}) is a candidate for exactly that
computation.

\subsection{Kuroda's two legacies}\label{sec:kuroda}

It is fitting, in lectures given in Tokyo, to note that both halves of the problem were seen early
in work on Japanese. \citet{Kuroda1965} named the \emph{indeterminate} pronouns:
one root (\textit{dare}, \textit{nani}, \textit{doko}), several forces, the
force supplied by a particle (Table~\ref{tab:indeterminates}). That is the first
clear demonstration I know of that force and restriction come apart in grammar:
the root is the restriction, the particle the force.

\begin{center}\small
\begin{tabularx}{\linewidth}{@{}l L L L l@{}}
\toprule
root & + \textit{ka} & + \textit{mo} & + \textit{demo} & \ldots\textit{no?} (C)\\
\midrule
\textit{dare} & \textit{dare-ka} `someone' & \textit{dare-mo} `everyone / no one' & \textit{dare-demo} `anyone' & `who?'\\
\textit{nani} & \textit{nani-ka} `something' & \textit{nani-mo} `nothing' (with negation) & \textit{nan-demo} `anything' & `what?'\\
\textit{doko} & \textit{doko-ka} `somewhere' & \textit{doko-mo} `everywhere' & \textit{doko-demo} `anywhere' & `where?'\\
\bottomrule
\end{tabularx}
\captionof{table}{Kuroda's indeterminates: one root, several forces.}\label{tab:indeterminates}
\end{center}

The glosses in the table are deliberately rough. Whether \textit{dare-mo} is
`everyone' or `no one' depends on accent and on negation, and \textit{nani-mo}
is largely restricted to negative contexts; the table is there only to show the
split, and Lecture~2 does the indeterminates properly (\Hs{2}{sec:indeterminates}).

\citet{Kuroda1972}, drawing on Brentano and Marty, distinguished the
\emph{categorical} judgement, which picks out a subject and predicates something
of it, from the \emph{thetic} judgement, which simply presents a state of
affairs, and found the distinction in \textit{wa} and \textit{ga}:

\begin{exe}
\ex \label{ex:wa-ga}
\begin{xlist}
\ex \label{ex:wa}
\gll Inu-wa hasiru.\\
     dog-\gl{top} run.\gl{npst}\\
\glt `Dogs run.' \src{(categorical, generic)}
\ex \label{ex:ga}
\gll Inu-ga hasitte iru.\\
     dog-\gl{nom} run-\gl{te} be.\gl{npst}\\
\glt `A dog is running / Some dogs are running.' \src{(thetic)}
\end{xlist}
\end{exe}

The categorical form is the Aristotelian one, \textit{S is P}: a restriction and
something said of it. Frege's achievement was to abandon that form for the
operator and the variable. Natural language did not follow him all the way.
Every natural-language determiner is conservative
\citep{BarwiseCooper1981,KeenanStavi1986}:

\begin{keybox}
\begin{exe}
\ex \label{ex:conservativity} \textit{Conservativity:}\quad $Q(A)(B) \;\Leftrightarrow\; Q(A)(A \cap B)$
\end{exe}
\end{keybox}

That is, \textit{every student smokes} is equivalent to \textit{every student is
  a student who smokes}, and \textit{most students smoke} to \textit{most
  students are students who smoke}. A determiner that looked only at $B$, or at
$B - A$, would not be conservative, and none is attested. The restriction is
obligatory and privileged: the determiner lives on its first argument, and
natural-language quantification is restricted quantification, always. Japanese,
then, wears both of Kuroda's legacies on its sleeve: the split between force and
restriction in its indeterminates, and the split between restrictor and nuclear
scope in \textit{wa} and \textit{ga}. We will need both. The second connects
directly with Lecture~2, since adverbs of quantification and existential closure
sit on exactly this divide \citep{Diesing1992keio}, and one of the questions for
the discussion asks whether \textit{wa} marks the edge of the restrictor
structurally.

\subsection{The inversion}\label{sec:inversion}

We can now put the inversion promised in \S\ref{sec:aims}. Logic offers natural
language several ways of quantifying without variables, and a way of quantifying
without a restriction (Table~\ref{tab:logics}). As far as I can tell, no natural
language uses any of them in pure form.

\begin{center}\small
\begin{tabularx}{\linewidth}{@{}P{4.4cm} L L@{}}
\toprule
System & How it quantifies & Used by natural language?\\
\midrule
First-order logic \citep{Frege1879} & operator, variable, unrestricted domain & operator--variable: yes; unrestricted: no (\ref{ex:conservativity})\\
Generalised quantifiers (\citealt{Montague1973}; \citealt{BarwiseCooper1981}) & a relation between two sets; restricted & yes: this is the one\\
$\lambda$-calculus \citep{Church1941} & bound variables, names immaterial ($\alpha$-conversion) & as the metalanguage of semantics\\
Combinatory logic (\citealt{Schonfinkel1924}; \citealt{CurryFeys1958}) & no variables: combinators & no (but see CCG)\\
Predicate-functor logic \citep{Quine1960,Quine1971} & no variables: functors on predicates & no\\
Variable-free semantics \citep{Jacobson1999} & no variables: pronouns as identity maps & arguably, in the semantics\\
\bottomrule
\end{tabularx}
\captionof{table}{What logic offers, and what language takes.}\label{tab:logics}
\end{center}

Logicians will know all of this; the point is to have the alternatives named before I claim that language passes them by. The excursion that follows makes the rows explicit, and can be skipped; the conclusion comes after it.

\begin{excursus*}[Five logics, and what each does with the variable]
  \textit{(i) Frege: the variable, and unrestricted quantification.} In Frege's
  notation \citep{Frege1879} \textit{every student left} is $\forall x\,[\mathit{Student}(x) \to \mathit{Left}(x)]$: for every object $x$, if $x$ is a student then $x$ left. The quantifier ranges over everything, and the restriction is smuggled in by the connective, so that
  `every student left' and `everything is either not a student or left' are the
  same formula. The variable occurs three times. Its occurrences mean nothing on
  their own (Fine); what carries the meaning is that they are \emph{the same}
  variable, so that the two predicates are evaluated of the same object. The
  variable, in other words, is the device that identifies argument positions
  across predicates. Natural language does not quantify unrestrictedly: there is
  no determiner \textit{every} that takes a sentence and ranges over the
  universe, and the closest thing, \textit{everything}, is \textit{every} plus
  \textit{thing}, a restriction after all. What natural language keeps from
  Frege is the identity across positions, the variable; what it refuses is the
  unrestricted domain. Frege's notation bundles the two, and conservativity
  (\ref{ex:conservativity}) is the formal statement of the refusal.

\textit{(ii) Generalised quantifiers, or restricted quantification.} On the treatment of \citet{Montague1973} and \citet{BarwiseCooper1981}, determiners are relations between sets:
\[\sem{\textit{every}} = \lambda A\,\lambda B.\ A \subseteq B \qquad \sem{\textit{some}} = \lambda A\,\lambda B.\ A \cap B \neq \emptyset \qquad \sem{\textit{most}} = \lambda A\,\lambda B.\ |A \cap B| > |A - B|\]
Read: a determiner takes two sets, the restriction $A$ (the students) and the predicate $B$ (the leavers), and says how they are related. \textit{Every}: $A$ is included in $B$. \textit{Some}: the two overlap. \textit{Most}: the part of $A$ inside $B$ outnumbers the part outside. A determiner is of type $\langle\langle e,t\rangle,\langle\langle e,t\rangle,t\rangle\rangle$, and a QP,
of type $\langle\langle e,t\rangle,t\rangle$, is a set of sets. The restriction
is the first argument, and conservativity says the determiner never looks beyond
it. \textit{Most} is not first-order definable at all, so that a language with
only Frege's quantifiers could not say it; natural language says it easily. This
is the row natural language actually uses. But where is the variable? In the
metalanguage: $A \subseteq B$ is $\forall x[x \in A \to x \in B]$. The notation
of generalised quantifiers hides the variable; it does not remove it. Whether
the hidden variable corresponds to anything in the sentence is exactly where the
theories divide: for QR it does, and it is the trace; for the box it does, and
it is a token of type $e$; for variable-free semantics it does not. The GQ
treatment, then, is not variable-free but variable-hiding, and that is why the
type-$e$ position in the object matters so much in Lecture~3.

\textit{(iii) The $\lambda$-calculus, $\alpha$-conversion, and nameless
  variables.} In the $\lambda$-calculus \citep{Church1941} bound variables can
be renamed freely:
$\lambda x.\,\mathit{read}(x) =_\alpha \lambda y.\,\mathit{read}(y)$. Here $\lambda x.\,\mathit{read}(x)$ is `the function that, given an $x$, returns $\mathit{read}(x)$': the $\lambda$ marks which position is the argument, and renaming $x$ as $y$ throughout changes nothing. Whatever the variable contributes, it is not its name; it is the link between the binder
and the position. De Bruijn's `nameless dummies' \citep{deBruijn1972} make the
point explicit, replacing each variable by its \emph{distance} from its binder,
the number of $\lambda$s between the occurrence and the $\lambda$ that binds it:
$\lambda x.\lambda y.\ x\,y$ becomes $\lambda\,\lambda\,2\,1$, where $1$ means `bound by the nearest $\lambda$' and $2$ `bound by the one above it'. There are then no names, no clashes, and no $\alpha$-conversion. A de Bruijn index is neither a
name nor a feature of the occurrence: it is a structural relation between two
positions in the term. Replace `number of $\lambda$s' by `number of phase
heads', or by `c-command plus minimality', and the index becomes something a
grammar could compute from the structure rather than write into it. That is, to
a first approximation, the proposal of Lecture~4: Form Copy computes which tokens
are the same, and the $\theta$-relation of the box is an index that carries no
value. The analogy is not exact, since de Bruijn indices are still written; but
what they write is a relation, which is just what Chomsky's innocuous indices
are allowed to be.

\textit{(iv) Combinatory logic, or quantification without variables.} With the combinators of \citet{Schonfinkel1924} and \citet{CurryFeys1958},
\[\mathbf{I}\,x = x \qquad \mathbf{K}\,x\,y = x \qquad \mathbf{S}\,x\,y\,z = x\,z\,(y\,z) \qquad \mathbf{B}\,x\,y\,z = x\,(y\,z),\]
in words: $\mathbf{I}$ returns its argument; $\mathbf{K}$ takes two arguments and returns the first, throwing the second away; $\mathbf{S}$ takes three, feeds the third to the first and to the second, and applies the one result to the other; $\mathbf{B}$ composes, applying $y$ to $z$ and then $x$ to the result. Every $\lambda$-term can be rewritten with $\mathbf{S}$ and $\mathbf{K}$ alone: $\lambda x.x = \mathbf{I}$; $\lambda x.M = \mathbf{K}M$ if $x$ does not occur in $M$; and $\lambda x.MN = \mathbf{S}(\lambda x.M)(\lambda x.N)$. Read from right to left, the three clauses are an algorithm: push the $\lambda$ inwards through every application with $\mathbf{S}$, discard it with $\mathbf{K}$ wherever the variable is absent, and cancel it with $\mathbf{I}$ where the variable is all that is left; at the end no $\lambda$ remains, and with it no variable. One example, step by step. $\lambda x.\,f(g\,x)$, `the function that applies $f$ to the result of $g$', is $\mathbf{B}\,f\,g$, and by the rules it is $\mathbf{S}(\mathbf{K}f)\,g$: the body $f(g\,x)$ is an application, so the $\lambda$ goes through it with $\mathbf{S}$; $x$ does not occur in $f$, so $\lambda x.f = \mathbf{K}f$; and $\lambda x.\,g\,x$ behaves exactly as $g$ does. Give the result an argument $x$ and reduce:
\[\mathbf{S}(\mathbf{K}f)\,g\,x \ =\ (\mathbf{K}f\,x)\,(g\,x) \ =\ f\,(g\,x)\]
$\mathbf{S}$ feeds $x$ to both $\mathbf{K}f$ and $g$; $\mathbf{K}f\,x$ throws $x$ away and returns $f$; so the result is $f$ applied to $g\,x$, and the function has been written without ever naming the argument that $f$ and $g$ are applied \emph{to}. The variable has been explained away, its argument position encoded in the shape of the combinator. Sch\"onfinkel's motive was exactly ours: to get rid of the variable, which he regarded, roughly, as a mere auxiliary device of notation, and of its bookkeeping. CCG takes this seriously, since \textbf{B} and \textbf{T} are its rules (\S\ref{sec:box-preview}, excursion). The evidence of \S\ref{sec:spot} points the other way, however: language pronounces restriction and force and leaves the identification of positions silent, which is what a variable, not a combinator, would do.

\textit{(v) Predicate-functor logic.} Quine \citeyearpar{Quine1960,Quine1971} replaces the binding of variables by a handful of functors that operate on predicates: intersection and complement; \textbf{Inv} and \textbf{inv}, which permute argument places; \textbf{Ref}, which identifies two places; \textbf{Pad}, which adds a vacuous place; and \textbf{Der}, which existentially closes the first place (the names vary between the two papers; these are the later ones). With $R$ for `$x$ loves $y$': $\mathbf{Inv}\,R$ is `$y$ loves $x$', that is, `is loved by'; $\mathbf{Ref}\,R$ is `$x$ loves $x$', `loves oneself'; $\mathbf{Pad}$ lets predicates of different arity be intersected place by place; and $\mathbf{Der}\,F$ is `something is $F$'. Quantification is \textbf{Der}, and co-reference across predicates is \textbf{Ref}:
\[\exists x\,[\mathit{F}x \wedge \mathit{G}x] \ \leadsto\ \mathbf{Der}\,(\mathit{F} \cap \mathit{G}) \qquad\qquad \exists x\,\mathit{R}xx \ \leadsto\ \mathbf{Der}\,\mathbf{Ref}\,\mathit{R}\]
In words: intersect the two predicates place by place (the things that are both $F$ and $G$), then close the place: something is both. And: identify the two places of $R$ (loves oneself), then close: someone loves themselves. No variable is ever written. `Variables explained away', as Quine put it: what the variable did, identify positions across predicates, is done by \textbf{Ref} and the permutation functors, and what the quantifier did is done by \textbf{Der}. The identification of positions is still there; it has become an operation on predicates rather than a shared name. No natural language has a morpheme that does what \textbf{Ref} does and nothing else, and reflexives and reciprocals, the nearest candidates, carry $\varphi$-features. That is `pronouns are not variables' seen from the other side. The moral is the same as for the combinators: eliminating the name does not eliminate the identification of positions, which simply moves into the operations.
\end{excursus*}

Quantification in natural language, then, is always restricted, as (\ref{ex:conservativity}) says, and it always relates a restriction to a nuclear scope through an identity between positions, which is to say through a variable. It uses none of the variable-free systems in pure form, and it does not quantify unrestrictedly. Yet the variable, which the logician writes into every formula, is the one ingredient that is never pronounced (we will see this in \S\ref{sec:spot}). If it is never pronounced, it is unlikely to be a lexical item. The bet of this seminar is that it is not an object at all, but a relation the grammar computes. This is the answer in its most general form, and the three answers that follow are three places to look for that relation; the first, as we shall see, does not find one.

\subsection{Three answers}\label{sec:three-answers}

There are, I think, three places the variable could come from, and the seminar takes them in turn (Table~\ref{tab:three-answers}).

\begin{center}\small
\begin{tabularx}{\linewidth}{@{}P{2.4cm}P{3.6cm}P{1.0cm}L@{}}
\toprule
Answer & Where the variable comes from & Lectures & What I will argue\\
\midrule
Movement & Copies, plus Trace Conversion & 1, 3 & It breaks Inclusiveness: the definite and the variable were never in the workspace\\
Storage & Cooper storage, then the box & 2, 3 & It works, and its locality is inherited from phases, where a semantic store has to stipulate it\\
Relations over tokens & Indices, Form Copy & 4 & It works, and the box already relies on such a relation\\
\bottomrule
\end{tabularx}
\captionof{table}{Three answers to the question of where the variable comes from: a map of the seminar.}\label{tab:three-answers}
\end{center}

The answer the four lectures defend is argued in Lectures~3 and~4; here it is only stated, so that it is clear where the argument is going:

\begin{keybox}
\begin{exe}
\ex \label{ex:answer} \textit{The answer.} The variable is not an item but a relation between tokens. Scope is storage in the syntax, and its locality is the locality of phases.
\end{exe}
\end{keybox}

The second sentence is the result of Lecture~3, and the first that of Lecture~4. How the answer stands to the other ways of doing the same job, Cooper storage first among them, is a separate matter, taken up where the box is introduced (\Hs{3}{sec:box-cooper}) and at the end (\Hs{4}{sec:neighbours}).

% =====================================================================
\section{Spot the variable}\label{sec:spot}

Here is a simple way of making the problem concrete. Take each of the sentences
in (\ref{ex:spot}) and ask four things of it: what its \emph{restriction} is
(what the quantification is over), what its \emph{force} is (what kind of
quantification it is), where its \emph{scope site} is (where its scope is
fixed), and where its \emph{variable} is (the position it binds). For each of
the four, ask whether it is pronounced, and if so, by which
morpheme.\toexercise{spot}

\begin{exe}
\ex \label{ex:spot}
\begin{xlist}
\ex \label{ex:spot-every} Every student left.
\ex \label{ex:spot-sannin}
\gll Gakusei-ga san-nin kita.\\
     student-\gl{nom} three-\gl{cl} come.\gl{pst}\\
\glt `Three students came.'
\ex \label{ex:spot-dareka}
\gll Dare-ka-ga kita.\\
     who-\gl{ka}-\gl{nom} come.\gl{pst}\\
\glt `Someone came.'
\ex \label{ex:spot-dono}
\gll Dono gakusei-mo kita.\\
     which student-\gl{mo} come.\gl{pst}\\
\glt `Every student came.'
\ex \label{ex:spot-zutu}
\gll Kodomo-tati-ga ringo-o mittu-zutu tabeta.\\
     child-\gl{pl}-\gl{nom} apple-\gl{acc} three.\gl{cl}-\gl{zutu} eat.\gl{pst}\\
\glt `The children ate three apples each.'
\ex \label{ex:spot-who}
\gll Dare-ga kita no?\\
     who-\gl{nom} come.\gl{pst} \gl{q}\\
\glt `Who came?'
\end{xlist}
\end{exe}

Disagreement is to be expected, particularly over the force of (\ref{ex:spot-sannin}), and the disagreement is part of the lesson. The appendix sets this out as an exercise, with a blank grid (\exref{spot}); the answers follow here, one sentence at a time.

In (\ref{ex:spot-every}) the restriction, \textit{student}, is pronounced inside the QP, and so is the force, \textit{every}, syncretic with the determiner. The scope site is not pronounced, and neither is the variable. If the subject is merged in Spec-\vs P and raised to Spec-TP, its lower copy marks the position the variable ought to occupy, which is exactly where the problem of \S\ref{sec:problem} arises (Figure~\ref{fig:spot-every}). Even the simplest English sentence, in other words, already poses the problem.

\begin{keiofig}{fig:spot-every}{(\ref{ex:spot-every}): the lower copy of the subject marks where the variable ought to be.}
\begin{forest}
[TP, SOtree, xslim, flat
  [QP, SOmoved [{\LI{every student}}, SOlexical]]
  [\textit{v}P [QP, SOtrace [{$\langle$\LI{every student}$\rangle$}]] [VP [{\LI{left}}, SOlexical]]]]
\end{forest}
\end{keiofig}

In (\ref{ex:spot-sannin}) the restriction, \textit{gakusei}, is pronounced, and so is the numeral classifier \textit{san-nin}; but the numeral contributes cardinality, a measure, not force. This is where most people asked to fill in the grid put \textit{san-nin} under force, and the reply is simple: remove it, and the sentence keeps its existential force (\textit{Gakusei-ga kita} `A student came / Students came'). The existential force of the sentence comes from the clause (existential closure, in the sense of \citealt{Diesing1992keio}, to which Lecture~2 returns), not from anything pronounced. The scope site is not pronounced, unless the position of the floated numeral is taken to mark it (\citealt{Tsoulas2003}; \S\ref{sec:FQ}), and the variable is not pronounced either.

In (\ref{ex:spot-dareka}) the restriction is \textit{dare}, `person', so to speak, and the force is \textit{ka}, pronounced and existential. The scope site is wherever \textit{ka} attaches, here the edge of the subject nominal itself, and the variable is silent. (\ref{ex:spot-dono}) is the same with \textit{mo} in place of \textit{ka}: \textit{dono gakusei} is the restriction, \textit{mo} the force, now universal, and scope is fixed at \textit{mo}'s attachment site.

(\ref{ex:spot-zutu}) is the instructive one. Distributive force is pronounced, by \textit{-zutu}, but not on the restriction it distributes over. It is attached to the share, \textit{mittu} `three', while the plural over which distribution takes place, \textit{kodomo-tati} (the sorting key, in the terms of \citealt{Choe1987}), carries nothing beyond the plural (Figure~\ref{fig:zutu}). Force is pronounced on a different argument from its restriction. The scope site is not pronounced, and neither is the variable, which ranges over the children one by one. Of the six, this is the case that shows most clearly that force and restriction come apart across arguments; the nominal data of \S\ref{sec:nominal} will show them coming apart within a single nominal. Choe's terms, sorting key and distributive share, are the ones Lecture~2 uses for \textit{-zutu} and Korean \textit{-ssik}.

\begin{keiofig}{fig:zutu}{(\ref{ex:spot-zutu}): distributive force on the share, distributing over the sorting key.}
\begin{tikzpicture}[every node/.style={font=\footnotesize}, >=Stealth]
  \node (key) at (0,0) {\textit{kodomo-tati-ga}};
  \node (obj) at (1.8,0) {\textit{ringo-o}};
  \node (share) at (3.4,0) {\textit{mittu-\textbf{zutu}}};
  \node (v) at (4.8,0) {\textit{tabeta}};
  \node[text=gtghost] at (0,-0.45) {\scriptsize sorting key};
  \node[text=gtghost] at (3.4,-0.45) {\scriptsize share + force};
  \draw[->, gtaccent, line width=0.6pt] (share.north) to[bend right=35] node[midway, above, text=gtghost] {\scriptsize distributes over} (key.north);
\end{tikzpicture}
\end{keiofig}

(\ref{ex:spot-who}) is the only one of the six in which the scope site itself is pronounced, by a morpheme of its own at a distance from the restriction. The restriction, \textit{dare}, stays in situ; the force, interrogative, is supplied by \textit{no} at C, which is also where the question takes its scope. The variable is, once again, silent. Putting the six together (Table~\ref{tab:spot-filled}):

\begin{keybox}
\begin{exe}
\ex \label{ex:generalisation} The restriction is always pronounced; force often; the scope site sometimes; the variable never.
\end{exe}
\end{keybox}

\begin{center}\small
\begin{tabularx}{\linewidth}{@{}P{3.6cm}LLLl@{}}
\toprule
& Restriction & Force & Scope site & Variable\\
\midrule
(\ref{ex:spot-every}) \textit{every student left} & \jok\ \textit{student} & \jok\ \textit{every} & -- & --\\
(\ref{ex:spot-sannin}) \textit{gakusei-ga san-nin kita} & \jok\ \textit{gakusei} & -- (closure) & -- (or \textit{san-nin}?) & --\\
(\ref{ex:spot-dareka}) \textit{dare-ka-ga kita} & \jok\ \textit{dare} & \jok\ \textit{ka} & \jok\ at \textit{ka} & --\\
(\ref{ex:spot-dono}) \textit{dono gakusei-mo kita} & \jok\ \textit{dono gakusei} & \jok\ \textit{mo} & \jok\ at \textit{mo} & --\\
(\ref{ex:spot-zutu}) \textit{kodomo-tati-ga \ldots mittu-zutu tabeta} & \jok\ \textit{kodomo-tati} & \jok\ \textit{zutu} (on the share) & -- & --\\
(\ref{ex:spot-who}) \textit{dare-ga kita no} & \jok\ \textit{dare} & \jok\ \textit{no} & \jok\ at C & --\\
\bottomrule
\end{tabularx}
\captionof{table}{The grid, filled in.}\label{tab:spot-filled}
\end{center}

Where a reader's grid disagrees with Table~\ref{tab:spot-filled}, it is worth asking why, and in one place the disagreement is genuine. The scope-site column for (\ref{ex:spot-dareka}) and (\ref{ex:spot-dono}) records the attachment site of the particle; whether that is the scope site, or whether scope is fixed higher, is a real question, and Lecture~2 is largely about it.

(\ref{ex:generalisation}) is an empirical generalisation, resting here on six sentences from two languages and, in Lecture~2, on Korean, Mandarin and the Salish languages besides, and I offer it as
such. I would very much like to hear of a language in which some morpheme does
nothing but mark the bound position. Resumptive pronouns are the obvious
candidate, but, as \S\ref{sec:pronouns} noted, they are pronouns, with their
$\varphi$-features, and their distribution is not that of a bound position left by movement (in many languages they are what rescues an island).

The generalisation should, however, be stated with some care about what `the variable' is. What is never pronounced, on the view I will defend, is the foot of the quantifier's own chain: the position the quantifier is related to, which in the box theory is the verb's own token. Bound pronouns, minimal pronouns and \textit{wh}-resumptives (in Vata, Swedish and Irish, for instance) are then not counterexamples to be dismissed but cases to be discussed, and I return to them at the end of Lecture~4: each pronounces a token related to the operator, and none pronounces the relation itself.

If (\ref{ex:generalisation}) holds, then the one ingredient of the logician's formula that natural language never pronounces is also the one the grammar has no means of writing. That is either a coincidence or a clue. I will treat it as a clue: the variable is not an item but a relation.

% =====================================================================
\section{What makes a QP a QP: the nominal spine}\label{sec:nominal}

This section asks the question on which the nominal half of the argument turns:
what makes a QP a QP? There are two broad answers. On the first, the lexicalist
one, the quantifier is a lexical item of a particular semantic type, a
determiner of type
$\langle\langle e,t\rangle,\langle\langle e,t\rangle,t\rangle\rangle$ in the
sense of \citet{BarwiseCooper1981}, and a QP is simply a phrase headed by one.
On the second, the configurational one, quantificational force is a feature,
[Q], that may be realised on any of several heads of the nominal extended
projection (\#, Div/Cl, D, or a Q head above D) or on a particle outside the nominal
altogether, and a QP is a configuration in which such a feature composes with a
restriction. I will argue for the second. For concreteness and ease of
reference, call it the Q-configuration hypothesis:

\begin{keybox}
\begin{exe}
\ex \label{ex:Qconfig} \textit{The Q-configuration hypothesis.} A QP is any nominal extended projection in which a [Q] feature is associated with the D-region and composes with the projection's predicative core, its restriction. Where [Q] is pronounced (its locus of externalisation) and where its force is fixed (its locus of interpretation) are in principle independent, and their match or mismatch is the main axis of cross-linguistic variation in quantification.
\end{exe}
\end{keybox}

The last clause of (\ref{ex:Qconfig}) is the one to watch: it is what the hierarchy of Lecture~2 and the rigidity claim of Lecture~3 are built on. The argument for it is cumulative. None of the domains below proves (\ref{ex:Qconfig}) on its own; each forces one piece of the split between force and restriction, and by the end the lexicalist alternative has, I think, very little left to stand on. Let us take them in turn, in the order of Table~\ref{tab:cumulative}.

\begin{center}\small
\begin{tabularx}{\linewidth}{@{}l P{4.8cm} L@{}}
\toprule
& Domain & The piece of the split it forces\\
\midrule
\S\ref{sec:Qlayer} & The DP hypothesis and the Q-layer & force above D, or two heads where English pronounces one\\
\S\ref{sec:FQ} & Floating quantifiers, outside and inside DP & force pronounced away from its restriction\\
\S\ref{sec:numerals} & Numerals, classifiers, count and mass & measure and division are not force\\
\S\ref{sec:indefinites} & (In)definites & restriction without force\\
\S\ref{sec:bare} & Bare nominals & restriction plus index\\
\S\ref{sec:split} & Split scope and decomposition & force and restriction scoping apart in German and English; force assembled\\
\S\ref{sec:jp-universals} & Two Japanese universals & one language choosing both options\\
\bottomrule
\end{tabularx}
\captionof{table}{The cumulative argument for (\ref{ex:Qconfig}).}\label{tab:cumulative}
\end{center}

\subsection{The DP hypothesis and the Q-layer}\label{sec:Qlayer}

\begin{exe}
\ex \label{ex:Qlayer}
\begin{xlist}
\ex all the students / both the students / *every the students / *the every student
\ex
\gll oli i fitites\\
     all.\gl{pl} the students\\
\glt `all the students' \src{(Greek)}
\ex
\gll o kathe fititis\\
     the every student\\
\glt `each (and every) student' \src{(Greek)}
\end{xlist}
\ex \label{ex:statimcets}
\gll t\'akem i smelhm\'ulhats-a\\
     all \gl{det} women-\gl{det}\\
\glt `all the women' \src{(St'\'at'imcets; \citealt{Matthewson2001})}
\end{exe}

The DP hypothesis \citep{Szabolcsi1983,Szabolcsi1994,Abney1987} gave the nominal a functional spine, and quantifiers that co-occur with determiners suggested a Q head above D (\citealt{Shlonsky1991} on Hebrew \textit{kol}; \citealt{Giusti1991,Giusti1997}). \citet{Matthewson2001} went further. From St'\'at'imcets, where quantifiers combine overtly with full DPs as in (\ref{ex:statimcets}), she argued that quantification proceeds in two steps in every language: D supplies an individual or a domain, and Q quantifies over its parts. English \textit{every} would then conceal a covert D. On Barwise and Cooper's treatment \textit{every} is a single item doing both jobs; on Matthewson's it is two heads, of which English pronounces one. Greek sits interestingly in between: \textit{kathe} does without the article, as in English, but \textit{o kathe} stacks the two, with a reading that is, as far as I can tell, more individuating, `each and every'. Notice that the structure Figure~\ref{fig:QoverD} assigns to \textit{o kathe} puts D above Q, the opposite of the English and St'\'at'imcets order; whether that is right, or whether \textit{o} is a D that has raised over Q, is one of the things that remain to be settled.

\begin{keiofig*}{fig:QoverD}{Q over D, D over Q. Left: English and St'\'at'imcets, with Q above a full DP. Right: Greek \textit{o kathe}, with the article above Q.}
\begin{twocol}{0.48}
\centering\small English, St'\'at'imcets\par\medskip
\begin{forest}
[QP, SOtree, spine [Q, SOhead [{\LI{all} / \LI{t\'akem}}, SOlexical]] [DP [D, SOhead [{\LI{the} / \LI{i\ldots-a}}, SOlexical]] [NP [{\LI{students}}, SOlexical]]]]
\end{forest}
\midcol
\centering\small Greek \textit{o kathe}\par\medskip
\begin{forest}
[DP, SOtree, spine [D, SOhead [{\LI{o}}, SOlexical]] [QP [Q, SOhead [{\LI{kathe}}, SOlexical]] [NP [{\LI{fititis}}, SOlexical]]]]
\end{forest}
\end{twocol}
\end{keiofig*}

The question to ask, then, is whether the Q-layer is a single head or a span across \#, Div, D and Q on which [Q] may be realised. Figure~\ref{fig:Qsites} marks the places where force has been argued to be realised: on Q (\textit{all}, \textit{t\'akem}), on D (\textit{every}, on the fused view), on \# (numerals, on some analyses), on Div or Cl (classifier quantifiers), or outside the nominal altogether, on a particle (\textit{mo}, \textit{ka}, \textit{dou}). The figure is deliberately noncommittal about which head is the `real' one; that is the head-or-span question. On (\ref{ex:Qconfig}) these are one configuration, externalised at different heads; on the lexicalist view they are different lexical items, of different types. (\ref{ex:Qlayer})--(\ref{ex:statimcets}) suggest at least that the answer differs by item and by language, which is what (\ref{ex:Qconfig}) leads us to expect.

\begin{keiofig}{fig:Qsites}{Where [Q] can live: the nominal spine of Figure~\ref{fig:spines}, with the candidate sites for [Q] marked; [Q] may also be realised on a particle outside the nominal (\textit{mo}, \textit{ka}, \textit{dou}).}
\begin{forest}
[{QP}, SOtree, spine
  [{Q \quad \textcolor{gtaccent}{[Q]}}, SOhead]
  [DP [{D \quad \textcolor{gtaccent}{[Q]}}, SOhead] [\#P [{\# \quad \textcolor{gtaccent}{[Q]}}, SOhead] [DivP [{Div/Cl \quad \textcolor{gtaccent}{[Q]}}, SOhead] [NP [N, SOhead]]]]]]
\end{forest}
\end{keiofig}

\subsection{Floating quantifiers, outside and inside DP}\label{sec:FQ}

\begin{exe}
\ex \label{ex:FQ}
\begin{xlist}
\ex \label{ex:FQ-a} All the students have left.
\ex \label{ex:FQ-b} The students have all left.
\end{xlist}
\end{exe}

\citet{Sportiche1988} derived (\ref{ex:FQ-b}) from the structure of (\ref{ex:FQ-a}): the DP moves and strands the quantifier, so that \lbr{QP} \textit{all} \lbr{DP} \textit{the students}]] yields [\textit{the students}] \textit{have} \lbr{QP} \textit{all} $\langle$\textit{the students}$\rangle$] \textit{left} (Figure~\ref{fig:stranding}). The floated quantifier marks a position the subject has passed through, which is what made it the classic evidence for VP-internal subjects. Adverbial analyses \citep{Bobaljik2003} and mixed ones \citep{Fitzpatrick2006} followed. The stranding analysis makes the point that matters here most directly, but the adverbial one makes it too: on it, an adverb \textit{all} quantifies over a subject with which it does not form a constituent. On either analysis, force is pronounced away from its restriction.

\begin{keiofig*}{fig:stranding}{Stranding \citep{Sportiche1988}. Left: the QP in Spec-\textit{v}P, with the DP inside it, as in (\ref{ex:FQ-a}). Right: the DP raised to Spec-TP, stranding \textit{all}, as in (\ref{ex:FQ-b}).}
\begin{twocol}{0.48}
\centering
\begin{forest}
[TP, SOtree, spine, flat
  [{}, no edge]
  [T$'$ [T, SOhead [{\LI{have}}, SOlexical]]
    [\textit{v}P
      [QP [Q, SOhead [{\LI{all}}, SOlexical]] [DP [{\LI{the students}}, SOlexical]]]
      [VP [V, SOhead [{\LI{left}}, SOlexical]]]]]]
\end{forest}
\midcol
\centering
\begin{forest}
[TP, SOtree, spine, flat
  [DP, SOmoved [{\LI{the students}}, SOlexical]]
  [T$'$ [T, SOhead [{\LI{have}}, SOlexical]]
    [\textit{v}P
      [QP [Q, SOhead [{\LI{all}}, SOlexical]] [DP, SOtrace [{$\langle$\LI{the students}$\rangle$}]]]
      [VP [V, SOhead [{\LI{left}}, SOlexical]]]]]]
\end{forest}
\end{twocol}
\end{keiofig*}

Japanese numeral quantifiers float too:

\begin{exe}
\ex \label{ex:FQ-jp}
\begin{xlist}
\ex
\gll San-nin-no gakusei-ga kita.\\
     three-\gl{cl}-\gl{gen} student-\gl{nom} come.\gl{pst}\\
\glt `Three students came.'
\ex
\gll Gakusei-ga san-nin kita.\\
     student-\gl{nom} three-\gl{cl} come.\gl{pst}\\
\glt `Three students came.'
\end{xlist}
\end{exe}

The floated numeral is subject to a locality condition with its host (\citealt{Miyagawa1989}; on numeral classifier phrases see also \citealt{KawashimaKitahara1993} and \citealt{Kitahara1993}), and on \posscitet{Nakanishi2007} analysis it is a measure function that can measure events as well as individuals, which is why floated numerals allow readings that prenominal ones do not. I argued some time ago \citep{Tsoulas2003} that floating quantifiers neither float nor quantify. They are overt scope markers: the DP carries an unvalued scope feature, valued by long-distance Agree with a scopal head, and the valuing site is spelled out as the `floated' element, so that scope freezes there. Korean shows the effect with a scrambled object (Figure~\ref{fig:scope-marker}):

\begin{exe}
\ex \label{ex:FQ-kor}
\gll Chayk-ul haksayng-i se-kwen sa-ss-ta.\\
     book-\gl{acc} student-\gl{nom} three-\gl{cl} buy-\gl{pst}-\gl{decl}\\
\glt `A student bought three books.' \src{(the numeral marks the object's scope; \citealt{Tsoulas2003})}
\end{exe}

\begin{keiofig}{fig:scope-marker}{Floating quantifiers as overt scope markers \citep{Tsoulas2003}: the scope feature of the scrambled object is valued by Agree at the position of \textit{se-kwen}, where scope freezes.}
\begin{tikzpicture}[every node/.style={font=\footnotesize}, >=Stealth]
  \node (dp) at (0,0) {[\textit{chayk-ul}]$_{[u\mathrm{Scope}]}$};
  \node (subj) at (1.95,0) {\textit{haksayng-i}};
  \node (fq) at (3.4,0) {\textbf{\textit{se-kwen}}};
  \node (v) at (4.6,0) {\textit{sa-ss-ta}};
  \draw[->, gtaccent, line width=0.6pt] (fq.north) to[bend right=25] node[midway, above, text=gtghost] {\scriptsize Agree: scope valued here} (dp.north);
\end{tikzpicture}
\end{keiofig}

In retrospect, that paper was the first precedent for the box that I know of in my own work. Scope is syntactic, set by Agree, and spelled out at the valuing site: the scope feature approximates the boxing feature, the marker approximates the externalised particle, and freezing approximates rigidity. What it lacks is storage, and it needs the Beghelli--Stowell spine of (\ref{ex:scopepositions}), which the box dispenses with. Lecture~3 returns to it among the precedents (\Hs{3}{sec:precedents}).

The float can also happen inside the nominal. \citet{Green1971} observed that quantificational expressions seem unable to serve as possessors of singular head nouns. \citet{TsoulasWoods2019} showed, with data from a range of English dialects, that such possessives are not simply unacceptable: they are possible under restrictions, and only in predicative position, which is why we called them uniquely predicational possessives (UPPs). Universal quantifiers are generally acceptable in them, existential and intermediate ones generally not.

\begin{exe}
\ex \label{ex:UPP}
\begin{xlist}
\ex \label{ex:UPP-a} Norman is both of our friend(s).
\ex \label{ex:UPP-b} Getting a fair deal for millennials is in all our interests. \src{(\textit{The Guardian}, March 2016)}
\ex \label{ex:UPP-c} Is this both of our decision?! \src{(US English)}
\ex \label{ex:UPP-d} It is all of our responsibility to help girls everywhere defy the odds. \src{(Canadian English)}
\ex[*]{\label{ex:UPP-e} Jim is every of our friend(s).}
\ex[*]{\label{ex:UPP-f} The Skoda is both of our car(s).}
\end{xlist}
\end{exe}

For the speakers who accept it, (\ref{ex:UPP-a}) does not have the pragmatically implausible reading on which the two of us have exactly two friends between us and they happen to be one person, Norman; it means that Norman is a friend of each of us. Likewise (\ref{ex:UPP-b}), a headline, does not say that all the interests a group of people have are served by a fair deal for millennials, but that one interest, which each of us has, is. The possessor is quantified, the quantifier distributes over the possessors, and the head noun is understood once per possessor. Three restrictions hold. The quantifiers that occur are the three that float in English, \textit{all}, \textit{both} and \textit{each}, together with \textit{none} and \textit{neither}; \textit{every}, which does not float, is out, as in (\ref{ex:UPP-e}), and so, by and large, are \textit{some}, \textit{most}, \textit{many} and \textit{few}. The head noun must be relational (\textit{friend}, \textit{decision}, \textit{responsibility}, \textit{border}), so that (\ref{ex:UPP-f}) is out whatever the context. And the construction is possible only in predicative position: as an argument, \textit{both of our friends} denotes two friends and nothing else, and \textit{*both of our friend is coming} is simply ungrammatical. The plural \textit{-s} on the head noun, finally, is obligatory for some speakers, optional for others and banned for still others, which is why it is bracketed throughout.\footnote{The attested examples are from \citet{TsoulasWoods2019}, where (\ref{ex:UPP-c}) and (\ref{ex:UPP-d}) were found with Google in March 2017 and labelled by the dialect of their writers; the British National Corpus has the same construction (\textit{which is all our responsibility}, text~KGR~697). The squib's analysis builds on \citeauthor{Kayne1993}'s (\citeyear{Kayne1993}, \citeyear{Kayne1994}) treatment of English possessives: the quantifier originates with the possessor and is stranded, which is why only the floating quantifiers occur.}

For present purposes the point is simple. The quantifier and the restriction it quantifies over, the possessors, are pronounced apart within a single nominal, and where the construction can appear is fixed by the syntax, not by the lexical quantifier. Floating, outside DP and inside it, is the first piece of evidence that force can be pronounced away from its restriction.

\subsection{Numerals, classifiers and the count--mass divide}\label{sec:numerals}

On the simplest analysis a numeral is a modifier, of type $\langle\langle e,t\rangle,\langle e,t\rangle\rangle$:

\begin{exe}
\ex \label{ex:three} $\sem{\textit{three}} = \lambda P\,\lambda x\,[\,P(x) \wedge |x| = 3\,]$
\end{exe}

Existential force then comes from elsewhere: from D, or from existential closure. On the more articulated analysis of \citet{IoninMatushansky2006} numerals are heads, and complex numerals are built by syntax. Either way, numerals carry measure, not force, which is what (\ref{ex:spot-sannin}) already suggested. Classifiers make the same point one level down. On \posscitet{Chierchia1998} Nominal Mapping Parameter, nouns in classifier languages denote kinds and need a classifier to be counted; on \posscitet{Borer2005} syntactic version, every noun stem is mass, a Div head (realised as a classifier or as plural morphology) divides it, and a \# head counts (Figure~\ref{fig:borer}). `Count noun' then names a structure, not a lexical class. (\ref{ex:CL-c}) shows the classifier requirement: without the classifier, the numeral cannot combine with the noun.

\begin{exe}
\ex \label{ex:CL}
\begin{xlist}
\ex
\gll hon san-satu\\
     book three-\gl{cl}\\
\glt `three books'
\ex
\gll san-satu-no hon\\
     three-\gl{cl}-\gl{gen} book\\
\glt `three books'
\ex \label{ex:CL-c}
\gll \mstar san hon\\
     three book\\
\end{xlist}
\end{exe}

\begin{keiofig}{fig:borer}{Three heads, three jobs: \# counts, Div divides, N names. Borer's structure, with English and Japanese exponents side by side (word order set aside).}
\begin{forest}
[\#P, SOtree, xslim, flat [{\#}, SOhead [{\LI{three} / \LI{san}}, SOlexical]] [DivP [Div, SOhead [{\LI{-s} / \LI{-satu}}, SOlexical]] [NP [{\LI{book} / \LI{hon}}, SOlexical]]]]
\end{forest}
\end{keiofig}

Greek has no classifiers, but it pluralises mass nouns, with a reading of abundance rather than of division:

\begin{exe}
\ex \label{ex:waters}
\gll Ta nera plimirisan to ipogio.\\
     the waters flood.\gl{pst}.3\gl{pl} the basement\\
\glt `Water flooded the basement (and lots of it).'
\end{exe}

In the number/mass paper \citep{Tsoulas2009} I took this to show that number is a syntactic head in its own right, which can combine with a mass noun and measure rather than divide. Division, cardinality and force, then, are three different things, realised on three different heads, and only the last is quantification in the sense that concerns us.

\subsection{(In)definites: restriction without force}\label{sec:indefinites}

\begin{exe}
\ex \label{ex:texas} If a friend of mine from Texas had died in the fire, I would have heard about it. \src{\citep{FodorSag1982}}
\ex \label{ex:rumour} Each teacher overheard the rumour that a student of mine had been called before the dean.
\end{exe}

In (\ref{ex:texas}) the indefinite takes scope outside the conditional clause, which QR could not leave; in (\ref{ex:rumour}) the disputed intermediate reading, on which the student varies with the teachers, is the one the literature has argued about ever since \citep{Farkas1981,Abusch1994}. Fodor and Sag took indefinites to be ambiguous between quantificational and referential uses. \citet{Reinhart1997} and \citet{Winter1997} instead treated them as choice functions, existentially closed at any height without movement, so that islands are simply irrelevant; \citet{Kratzer1998} argued that the closure is contextual, and the wide readings pseudo-scope. Behind all of these stand \citet{Heim1982} and \citet{Kamp1981}: an indefinite is a restricted variable, and its force comes from elsewhere. \citet{Kratzer2005} traces the dependence of indefinites on outside operators from Japanese to Salish. The indefinite is thus the limiting case of (\ref{ex:Qconfig}), a restriction with no force of its own, whose apparent scope is wherever force is supplied.

One consequence will matter in Lecture~3: a force-less indefinite has no reason to be boxed, and its exceptional scope is not storage scope. The box, in other words, is not a theory of indefinite scope, and should not be. Korean \textit{nwukwu-nka} `someone' is not a counterexample, since there the force is on the particle, and the indefinite is not force-less.

\subsection{Bare nominals: restriction plus index}\label{sec:bare}

\begin{exe}
\ex \label{ex:bare}
\gll Gakusei-ga kita.\\
     student-\gl{nom} come.\gl{pst}\\
\glt `A student came / Students came / The student(s) came.'
\end{exe}

A Japanese bare noun has no article and no number, and its force and definiteness come entirely from the clause and the context \citep{Carlson1977,Longobardi1994,Chierchia1998,Dayal2004keio}. Alexopoulou, Folli \& Tsoulas \citeyearpar{AlexopoulouFolliTsoulas2013} argued, largely from Greek, that bare nominals are best analysed as a restriction plus an index. That characterisation is the right one, I think, but notice what it commits us to: the bare nominal contributes an index that the clause binds. Inclusiveness forbids indices. So the question of Lecture~4, what an index can be in a grammar that cannot add one, arises already inside the nominal. Lecture~4 gives the general form of the answer, an index as a relation between tokens; I do not work the bare nominal through in these lectures, and I note the debt.

\subsection{English is only quieter: split scope and decomposition}\label{sec:split}

\begin{exe}
\ex \label{ex:keine}
\gll Du musst keine Krawatte anziehen.\\
     you must no tie put.on\\
\glt `You don't have to wear a tie.' \src{($\neg >$ \textit{must} $> \exists$; German; \citealt{PenkaZeijlstra2005})}
\end{exe}

The salient reading of (\ref{ex:keine}) puts negation above the modal and the indefinite below it: \textit{keine} splits, and its two halves take scope on either side of an intervening operator (\citealt{Jacobs1980}; \citealt{Rullmann1995} for Dutch \textit{geen}; \citealt{Penka2011}). On the analysis Penka develops, \textit{kein} is an indefinite carrying an uninterpretable negative feature, checked under clause-bound Agree by an abstract negation on the clausal spine (Figure~\ref{fig:split}). That is a force-less restriction low, force high, and a local syntactic relation between them. It is, in effect, the store-and-retrieve architecture discovered from inside a Germanic nominal, and it is the reason the Japanese and Korean facts of Lecture~2 should not look exotic: German does covertly what Korean \textit{amwu-to} does overtly.

\begin{keiofig}{fig:split}{Split scope, schematically: \textit{keine} carries [\textit{u}Neg], checked by Agree with an abstract negation above the modal. (In Penka's analysis the negation is an operator in the specifier of a functional projection above the modal.)}
\begin{forest}
[NegP, SOtree, xslim, tikz+={\draw[->, gtaccent, line width=0.6pt] (keinlex.south west) .. controls +(-3.4,-0.3) and +(-1.0,-2.2) .. node[pos=0.55, left, text=gtghost, font=\scriptsize] {Agree} (neg.south west);}
  [{Op$_\neg$ [\textit{i}Neg]}, name=neg]
  [ModP [Mod, SOhead [{\LI{musst}}, SOlexical]]
    [\textit{v}P [DP [{\LI{keine}$_{[u\mathrm{Neg}]}$ \LI{Krawatte}}, SOlexical, name=keinlex]] [V, SOhead [{\LI{anziehen}}, SOlexical]]]]]
\end{forest}
\end{keiofig}

\begin{exe}
\ex \label{ex:most} Most students passed.
\end{exe}

\citet{Hackl2009} argued that proportional \textit{most} is \textit{many} plus the superlative \textit{-est}, assembled in the syntax (Figure~\ref{fig:most}), and that the decomposition explains why \textit{most} and \textit{more than half}, though truth-conditionally close, are verified differently; the experimental work on how speakers verify \textit{most} \citep{PietroskiEtAl2009} is what gives the decomposition its teeth. Even a flagship determiner, then, is built from a measure and a degree operator rather than listed as a primitive of type $\langle\langle e,t\rangle,\langle\langle e,t\rangle,t\rangle\rangle$.

\begin{keiofig}{fig:most}{\textit{Most} as \textit{many} plus \textit{-est} \citep{Hackl2009}; a simplification of his structure.}
\begin{forest}
[QP, SOtree, xslim [DegP [Deg, SOhead [{\LI{-est}}, SOlexical]] [AP [A, SOhead [{\LI{many}}, SOlexical]]]] [NP [{\LI{students}}, SOlexical]]]
\end{forest}
\end{keiofig}

\subsection{Two Japanese universals}\label{sec:jp-universals}

\begin{exe}
\ex \label{ex:universals}
\begin{xlist}
\ex \label{ex:subete}
\gll Subete-no gakusei-ga kita.\\
     all-\gl{gen} student-\gl{nom} come.\gl{pst}\\
\glt `All the students came.'
\ex \label{ex:dono-gakusei}
\gll Dono gakusei-mo kita.\\
     which student-\gl{mo} come.\gl{pst}\\
\glt `Every student came.'
\ex \label{ex:motun}
\gll Motun haksayng-i o-ass-ta.\\
     every student-\gl{nom} come-\gl{pst}-\gl{decl}\\
\glt `Every student came.' \src{(Korean)}
\end{xlist}
\end{exe}

Japanese has a universal whose force is carried by a nominal modifier, \textit{subete-no} N (beside the nominal universals \textit{zen'in} and \textit{minna}), beside one whose force is carried by a particle, \textit{dono} N-\textit{mo}; Korean has \textit{motun} N beside the indeterminate series (\textit{nwukwu-to}, \textit{nwukwu-na}). One language, then, chooses both options, and this is the cleanest test of (\ref{ex:Qconfig}) available (the pair that makes the test is \Hx{2}{ex:mineshima}); it is why the judgement sheet in the appendix (\exref{sheet}) has items on both universals. On (\ref{ex:Qconfig}) the two Japanese universals should differ precisely where it matters. Where force is a particle, its scope should be pinned to the particle's attachment site; where it is syncretic with the nominal, the universal should behave more like English \textit{every}. \citet{Mineshima2005}, in a talk to the Semantics Research Group at the University of Tokyo, argued for exactly such a difference: \textit{subete-no} N can take scope outside a complex NP that contains it, while an indeterminate-\textit{mo} phrase cannot. One would like to know whether the wide reading of \textit{subete-no} N is genuine scope or the cumulative, plural reading that \textit{subete-no} N, like English \textit{all}, also makes available; that is one of the things the judgement sheet asks. If Mineshima is right, the two universals land on different branches of the hierarchy we build in Lecture~2, and the contrast is evidence for the rigidity claim of Lecture~3.

\subsection{Taking stock}\label{sec:stock}

Each domain has forced a piece of the split (Table~\ref{tab:stock}). The lexicalist view can describe each of them, but only by multiplying lexical items and their types; (\ref{ex:Qconfig}) describes all of them as one configuration, externalised in different ways.

\begin{center}\small
\begin{tabularx}{\linewidth}{@{}P{4.2cm} L@{}}
\toprule
Floating quantifiers & pronounce force away from its restriction, outside the nominal and inside it\\
Numerals and classifiers & show that measure and division are not force\\
Indefinites and bare nominals & are restrictions without force (plus, for bare nominals, an index)\\
Split scope & shows force and restriction taking scope separately even in a language with fused determiners\\
Decomposition & shows force being assembled\\
The two Japanese universals & show the same language choosing both options, with different scope behaviour\\
\bottomrule
\end{tabularx}
\captionof{table}{Taking stock: one configuration, externalised in different ways.}\label{tab:stock}
\end{center}

If the lexicalist wants to push back, the best place to do it is decomposition. One can accept that \textit{most} is built in the syntax and still say that the result is an item of the generalised-quantifier type. The reply is that the type is then a property of the configuration, and that is what (\ref{ex:Qconfig}) says.

% =====================================================================
\section{The desideratum}\label{sec:desideratum}

We can now state what the remaining three lectures must deliver.

\begin{keybox}
\begin{exe}
\ex \label{ex:desideratum} \textit{Desideratum.} Quantificational force must be separable from the restriction. The grammar must be able to introduce force, and to fix scope, at a position other than the one at which the restriction is merged and pronounced; and it must relate the two without adding to the workspace anything that was not already in it.
\end{exe}
\end{keybox}

A QP, in other words, is a configuration, not a lexical class. In an earlier formulation I stated the requirement differently: QPs had to be built so that, under Internal Merge, their original copy could be interpreted as a variable. That formulation was wrong in a precise way. It presupposes the very answer, movement and copies, that Lecture~3 argues against, and it puts the burden in the wrong place. Restated as separability, the requirement is neutral among the three answers of \S\ref{sec:three-answers}, and it is what the data of \S\S\ref{sec:spot}--\ref{sec:nominal} actually demand.

It is worth seeing how each answer fares against the two clauses of (\ref{ex:desideratum}), which serve from here on as a scoreboard (Table~\ref{tab:scoreboard}). Movement satisfies the first: QR is precisely a way of separating the position where force takes scope from the position where the restriction is pronounced. It fails the second, since Trace Conversion adds a determiner and a variable. Storage can satisfy both, provided that the store is a syntactic object and that its locality follows from something independent. Cooper storage satisfies the first clause in the semantics; the box, I will argue, satisfies it in the syntax, with its locality inherited from phases. Relations over tokens are what the second clause requires: whatever links a force pronounced in one place to a restriction pronounced in another must be a relation between tokens already present, not a new object. Whether such relations can also do the work of the first clause is what Lecture~4 asks.

\begin{center}\small
\renewcommand{\arraystretch}{1.25}
\begin{tabularx}{\linewidth}{@{}P{2.5cm}LL@{}}
\toprule
& \textbf{Clause 1:} force and scope fixed away from the restriction & \textbf{Clause 2:} nothing added to the workspace\\
\midrule
\textbf{Movement} & \jok\ QR separates the position where force takes scope from the position where the restriction is pronounced & \jbad\ Trace Conversion adds a determiner and a variable\\
\textbf{Storage} & \jok\ Cooper storage, in the semantics; the box, I will argue, in the syntax & \jok\ provided the store is a syntactic object and its locality follows from something independent (phases)\\
\textbf{Relations over tokens} & (\jok) whatever links a force pronounced in one place to a restriction pronounced in another & \jok\ by definition: a relation between tokens already present, not a new object\\
\bottomrule
\end{tabularx}
\captionof{table}{The three answers against the two clauses of (\ref{ex:desideratum}).}\label{tab:scoreboard}
\end{center}

The rest of the seminar follows from this. Lecture~2 shows that the clausal spine offers a series of places where force can be fixed, with the particles of Japanese and Korean as the paradigm. Lecture~3 develops storage, first as Cooper had it and then as the box. Lecture~4 asks what the relation between force and restriction is, and argues that it is an index, honestly described. If the desideratum is right, the question we started with, where does the variable come from, becomes a different and more tractable one: what relates a force pronounced in one place to a restriction pronounced in another? That is a question the syntax can answer.

% =====================================================================
\section{Questions and readings}\label{sec:readings}

\begin{questions}
\begin{enumerate}
\item Is (\ref{ex:generalisation}) true? Is there a language that pronounces the variable, and are resumptive pronouns a counterexample?
\item If scope is partly syntax, as (\ref{ex:dareka-dono})--(\ref{ex:dono-dareka}) suggest, is it wholly syntax? What would a purely semantic account of the Japanese contrast look like, and what would it have to stipulate?
\item Is conservativity a property of determiners or of the syntax of restriction? Could (\ref{ex:conservativity}) follow from (\ref{ex:Qconfig}), on which [Q] always composes with the projection's predicative core first?
\item Does Kuroda's categorical/thetic distinction correspond to a structural position, with \textit{wa} marking the edge of the restrictor? What would that predict for adverbs of quantification (Lecture~2)?
\item Is the Q-layer a head or a span, and what could decide the question in Japanese, which has no articles?
\item Which of the five ideas of \S\ref{sec:architecture} would you give up first, and what would the box look like without it?
\end{enumerate}
\end{questions}

\subsection*{Key readings}

\citet{Chomsky2021}; \citet{ChomskyEtAl2023}; \citet{May1985}; \citet{Fox2002}; \citet{Fine2007}. These are the five texts the four lectures keep returning to.

\subsection*{Further readings}

\citet{Chomsky1995keio,Chomsky2000keio,Chomsky2008keio}; \citet{EpsteinEtAl1998}; \citet{Bobaljik2002}; \citet{HeimKratzer1998}; \citet{Grimshaw2005}; \citet{BarwiseCooper1981}; \citet{Kuroda1972}; \citet{Matthewson2001}; \citet{Hackl2009}; \citet{Penka2011}; \citet{AlexopoulouFolliTsoulas2013}; \citet{TsoulasWoods2019}. For the other frameworks: \citet{PollardSag1994}; \citet{CopestakeEtAl2005}; \citet{Dalrymple2001}; \citet{Steedman2012}; \citet{Jacobson1999}.

\subsection*{A textbook, for the overview}

\citet{Szabolcsi2010} is the best single map of the semantic side, for those who want one.


% ======================================================================
%  PART 2
% ======================================================================
\keiolecture{2}{Quantification on the clausal spine}{Quantification on the clausal spine}%
  {Day 1, Lecture 2 \textperiodcentered{} Saturday 7 November 2026, 15:45--17:45}
\keioepigraph{English uses the clausal spine as much as Japanese or Mandarin does,\\ and simply leaves it silent.}{}
\lecturemap

% =====================================================================
\section{Aims and plan}\label{sec:aims}

Lecture~1 ended on a desideratum (\Hx{1}{ex:desideratum}): quantificational force must be separable from the restriction; the grammar must be able to fix force and scope somewhere other than where the restriction is merged and pronounced; and it must relate the two without adding to the workspace anything that was not already in it. Lecture~2 asks where that somewhere is. The answer I will defend is that the clausal spine offers a series of places at which force can be fixed, and that the languages of East Asia, Japanese and Korean above all, pronounce the elements that sit there. The desideratum is the thread this lecture takes up from Lecture~1; the externalisation hypothesis (\Hx{1}{ex:variation}) and the nominal spine of \Hs{1}{sec:nominal} return at the end. The data are new.

This lecture has three aims. First, to show, with the particles of Japanese and Korean as the paradigm case, that force can be assigned at a distance from the restriction, on the clause rather than in the nominal. Second, to separate two kinds of locality, domain and path, and to design the two-by-two on which Lecture~3 tests them; the judgement sheet in the appendix (\exref{sheet}) gives readers its items. Third, to end on a single parameter hierarchy for both the nominal and the clausal domain, built up one question at a time; readers who would rather build it for themselves first will find a blank tree in the appendix (\exref{hierarchy}).

The discussion is organised around three problems, which any theory of quantification on the clause has to solve:

\begin{exe}
\ex \label{ex:problems}
\begin{xlist}
\ex \label{ex:restrictor-problem} \textit{The restrictor problem.} How is the restriction of a clausal quantifier assembled, when it is not the sister of the quantifier?
\ex \label{ex:binding-problem} \textit{The binding problem.} What does an adverb of quantification bind, and how many things at once?
\ex \label{ex:force-problem} \textit{The force-assignment problem.} If force can be assigned at a distance, on the spine, what fixes its scope, and what enforces its locality?
\end{xlist}
\end{exe}

Every proposal we will meet in this lecture, Hamblin alternatives, unselective binding, Diesing's mapping, Mandarin \textit{dou}, is an answer to one or more of these, and the hierarchy at the end asks how the answers pattern across languages. (\ref{ex:restrictor-problem}) is the problem that Lewis's \textit{if}-clauses and the Japanese \textit{-te-mo} clauses of \S\ref{sec:tripartite} solve overtly; (\ref{ex:binding-problem}) is Lewis's and Heim's problem; and (\ref{ex:force-problem}) is the one the box answers in Lecture~3, with $\theta$-discharge at the phase.

The obvious question is why some languages put force on the clause at all. Let me turn it around. The clausal spine is there in every language, with its operators and its closure positions, as \citet{BeghelliStowell1997} argued for English from the fixed scope of \textit{each} and \textit{every} relative to negation and to group-denoting indefinites (\Hx{1}{ex:scopepositions}); the more puzzling question may be why English uses it so little for quantification, placing force on the determiner almost everywhere and leaving the positions on the spine silent. I will suggest by the end that English does use them, and simply does not pronounce them. The answer will come at the last step of the hierarchy (\S\ref{sec:hierarchy}): English answers `no' to the clausal parameter not because it lacks the positions, which Beghelli and Stowell showed it has, but because it does not pronounce them.

The plan is as follows. \S\ref{sec:indeterminates} sets out the indeterminate system of Japanese and Korean and its semantics in terms of alternatives. \S\ref{sec:distance} looks at force at a distance: an indeterminate inside an island, a particle outside it. \S\ref{sec:unselective} turns to adverbs of quantification and unselective binding, with Japanese \textit{-te-mo} as the pronounced case. \S\ref{sec:clause} surveys languages that put force on the clause, and states the spine. \S\ref{sec:localities} previews the two kinds of locality on which the argument of Lecture~3 turns, and \S\ref{sec:hierarchy} is the hierarchy. \S\ref{sec:readings} has questions and readings. The boxed excursions (on the particles as join and meet, on the typology of indefinites, on LF pied-piping, and on the semantic frameworks: discourse representation, dynamic binding, situations and focus) are for readers from those traditions, and can be skipped without loss to the argument.

A short judgement sheet goes with this lecture (\exref{sheet}).\toexercise{sheet} Readers who would rather test the argument of Lecture~3 against their own judgements than against the literature's should fill it in before reading \S\ref{sec:localities}. Its design is the two-by-two of \S\ref{sec:wedge}: same clause against across a clause boundary, with and without a focus intervener on the path, with items in each of the two cells that decide, for the association of \textit{ka} and \textit{mo} with an indeterminate and for the licensing of \textit{nani-mo} separately, and, beside them, items on the two Japanese universals of \Hs{1}{sec:jp-universals}.

% =====================================================================
\section{The indeterminate system}\label{sec:indeterminates}

Let me start with some basic facts. Japanese has a series of words, which \citet{Kuroda1965} called \emph{indeterminates}, that are neither interrogative nor indefinite nor universal on their own, but become each of these in combination with a particle:

\begin{exe}
\ex \label{ex:paradigm} The Japanese indeterminate paradigm:
\end{exe}

\begin{center}\footnotesize
\begin{tabularx}{\linewidth}{@{}P{2.0cm}LLLLP{1.7cm}@{}}
\toprule
Indeterminate & + \textit{ka} & + \textit{mo}, with negation & + \textit{mo}, affirmative & + \textit{demo} & In a question (\ldots{} \textit{ka}/\textit{no})\\
\midrule
\textit{dare} `who' & \textit{dare-ka} `someone' & \textit{dare-mo} \ldots{} \textit{-nai} `no one' & \textit{dare-mo-ga} `everyone' & \textit{dare-demo} `anyone' & `who?'\\
\textit{nani} `what' & \textit{nani-ka} `something' & \textit{nani-mo} \ldots{} \textit{-nai} `nothing' & (gap) & \textit{nan-demo} `anything' & `what?'\\
\textit{doko} `where' & \textit{doko-ka} `somewhere' & \textit{doko-mo} \ldots{} \textit{-nai} `nowhere' & \textit{doko-mo} `everywhere' & \textit{doko-demo} `anywhere' & `where?'\\
\textit{itu} `when' & \textit{itu-ka} `some time' & \textit{itu-mo} \ldots{} \textit{-nai} `never' & \textit{itu-mo} `always' (\S\ref{sec:unselective}) & \textit{itu-demo} `any time' & `when?'\\
\textit{dono} N `which N' & (rare) & \textit{dono} N-\textit{mo} \ldots{} \textit{-nai} `no N' & \textit{dono} N-\textit{mo} `every N' & \textit{dono} N-\textit{demo} `any N' & `which N?'\\
\bottomrule
\end{tabularx}
\end{center}

The glosses matter less than the shape of the paradigm: its rows are restrictions and its columns are forces. Korean has a parallel system, with the indeterminates \textit{nwukwu} `who', \textit{mwues} `what', \textit{eti} `where' and \textit{encey} `when', the particles \textit{-nka}, \textit{-na} and \textit{-to}, and a dedicated indeterminate \textit{amwu} that occurs only with the last two:

\begin{exe}
\ex \label{ex:korean-indet}
\begin{xlist}
\ex \label{ex:nwukwunka}
\gll Nwukwu-nka-ka o-ass-ta.\\
     who-\gl{nka}-\gl{nom} come-\gl{pst}-\gl{decl}\\
\glt `Someone came.'
\ex \label{ex:amwuna}
\gll Amwu-na o-l swu iss-ta.\\
     any-\gl{na} come-\gl{adn} way exist-\gl{decl}\\
\glt `Anyone can come.'
\ex \label{ex:amwuto-a}
\gll Amwu-to o-ci anh-ass-ta.\\
     any-\gl{to} come-\gl{ci} \gl{neg}-\gl{pst}-\gl{decl}\\
\glt `No one came.'
\end{xlist}
\end{exe}

One root, several forces, and the force sits on the particle (\citealt{GillHarlowTsoulas2004}; for the cross-linguistic picture, \citealt{Haspelmath1997}). Korean is the comparison language throughout the seminar, and \textit{amwu} is its most interesting item: an indeterminate that cannot stand alone, and that takes only the free-choice and the polarity particles. Laid out by type, the Korean inventory looks as follows:

\begin{center}\small
\begin{tabularx}{\linewidth}{@{}P{3.6cm}LL@{}}
\toprule
Item & Type & Where force sits\\
\midrule
\textit{motun} N, \textit{motwu} & a nominal universal (`every', `all') & on the nominal, as English \textit{every}\\
\textit{nwukwu}/\textit{mwues} + \textit{-nka}, \textit{-na}, \textit{-to} & indeterminate + force particle & on the particle\\
\textit{amwu}-(N)-\textit{to} & polarity item, clause-bound & on the particle, licensed by Neg\\
\textit{-ssik} & overt distance distributivity & on the share; operator on the spine (\S\ref{sec:zutu})\\
\textit{-man} `only' & exhaustive focus particle & on the spine (Foc); an intervener (\S\ref{sec:path})\\
\bottomrule
\end{tabularx}
\end{center}

Korean, in other words, marks quantificational force overtly almost everywhere, and that is exactly why the steps of the computation are visible there: the particle pronounces retrieval. \textit{Motun} is the control. It is a nominal universal standing next to the particle series, and it is what makes the hierarchy of \S\ref{sec:hierarchy} testable inside Korean, as \textit{subete-no} N makes it testable inside Japanese.

It is worth noticing which particles these are:

\begin{center}\small
\begin{tabularx}{\linewidth}{@{}P{2.4cm}LLP{1.2cm}@{}}
\toprule
Particle & Also \ldots & \ldots{} so the quantifier is & \\
\midrule
Japanese \textit{ka} & disjunction `or'; the question particle & existential: a generalised disjunction & $\exists = \bigvee$\\
Japanese \textit{mo} & additive `also'; repeated, `both \ldots{} and' & universal: a generalised conjunction & $\forall = \bigwedge$\\
Korean \textit{-na} & disjunction & existential, free choice & \\
Korean \textit{-to} & additive & universal, polarity-sensitive & \\
\bottomrule
\end{tabularx}
\end{center}

Japanese \textit{ka} is also the disjunction `or' and the question particle; \textit{mo} is also additive `also' and, repeated, conjunctive `both \ldots{} and'; Korean \textit{-na} is a disjunction and \textit{-to} an additive. The existential is built from `or' and the universal from `and', which is just what a logician would expect if $\exists$ is a generalised disjunction and $\forall$ a generalised conjunction (\citealt{GillHarlowTsoulas2004}; for the typology, \citealt{Gil2013} and \citealt{Szabolcsi2015}). The particle, in other words, does not denote a quantifier. It denotes a connective, and quantification arises when the connective applies to the set of alternatives that the indeterminate supplies. This is the idea on which everything else in this section rests, and it is older than any of the analyses that implement it: Kuroda saw it, and \citet{Szabolcsi2015} gives it its cleanest algebraic statement, as join and meet. The caveat of \S\ref{sec:nandomo}, \textit{nan-do-mo}, is the same idea pushed one step further.

\pIIneedspace{6cm}
\begin{excursus}[Particles as join and meet]
\begin{center}
\begin{tikzpicture}[every node/.style={font=\small}, e/.style={gtaccent!60, line width=0.5pt}]
  \node (top) at (0,2.2) {$p \vee q \vee r$};
  \node[text=gtghost, font=\scriptsize, anchor=west] at (1.0,2.2) {\textit{ka}: join};
  \node (p) at (-1.6,1.1) {$p$};
  \node (q) at (0,1.1) {$q$};
  \node (r) at (1.6,1.1) {$r$};
  \node (bot) at (0,0) {$p \wedge q \wedge r$};
  \node[text=gtghost, font=\scriptsize, anchor=west] at (1.0,0) {\textit{mo}: meet};
  \draw[e] (top)--(p); \draw[e] (top)--(q); \draw[e] (top)--(r);
  \draw[e] (bot)--(p); \draw[e] (bot)--(q); \draw[e] (bot)--(r);
  \node[text=gtghost, font=\scriptsize] at (0,-0.6) {the alternatives $\{p,q,r\}$ supplied by the indeterminate};
\end{tikzpicture}
\end{center}
\citet{Szabolcsi2015} proposes that the quantifier particles of Japanese and Hungarian (\textit{ka} and \textit{mo}; \textit{vagy} and \textit{is}) are the join and meet operators of a lattice, applied to whatever set of alternatives they are given, whether propositions, individuals or questions. `Or' and `some', `and' and `every', are then the same operation applied to different sets, and that is why one and the same morpheme can be a connective, a quantifier particle and a question marker. David Gil's \citeyearpar{Gil2013} typology points the same way: across languages, the universal particle is overwhelmingly built from `and' or `also', and the existential from `or'.

The lattice is the standard picture. What matters for us is that force is not in the particle itself but in the operation it names, applied at a position. This is the semantic face of the idea that force is assigned at retrieval, which runs through the rest of the seminar.
\end{excursus}

\pIIneedspace{5cm}
\begin{excursus*}[The typology of indefinites, in one map]
\begin{center}
\begin{tikzpicture}[x=0.98cm, every node/.style={font=\scriptsize, draw=gtaccent!50, rounded corners=2pt, inner sep=3pt, text width=1.45cm, align=center}, e/.style={gtaccent!50, line width=0.5pt}]
  \node (spec) at (0,1) {specific known};
  \node (spu) at (2.6,1) {specific unknown};
  \node (irr) at (5.3,1) {irrealis non-specific};
  \node (q) at (8.0,2.1) {question};
  \node (cond) at (8.0,-0.1) {conditional};
  \node (ind) at (10.6,1) {indirect negation};
  \node (comp) at (10.6,-1.0) {comparative};
  \node (dn) at (13.2,1) {direct negation};
  \node (fc) at (13.2,-1.0) {free choice};
  \draw[e] (spec)--(spu)--(irr); \draw[e] (irr)--(q); \draw[e] (irr)--(cond); \draw[e] (q)--(ind); \draw[e] (cond)--(ind); \draw[e] (cond)--(comp); \draw[e] (ind)--(dn); \draw[e] (comp)--(fc); \draw[e] (ind)--(fc);
\end{tikzpicture}
\end{center}
\posscitet{Haspelmath1997} semantic map of indefinite functions, redrawn schematically. The functions an indefinite series can cover are contiguous regions of this graph. Japanese \textit{-ka} covers the left of it; \textit{-mo}, with negation, the direct-negation region; \textit{-demo} the free-choice and irrealis region; Korean \textit{amwu-to} covers direct negation, and \textit{amwu-na} free choice. For our purposes the map is a \emph{description} of where the combinations of particle and indeterminate land. The theory of this lecture is what is supposed to make the landing places follow from the particle's operation and its position; it is meant to explain the map, not to replace it.
\end{excursus*}

That last idea, the connective applied to a set of alternatives, is the core of the two standard analyses. \citet{Nishigauchi1990} treated indeterminates as variables, unselectively bound by the particle in the manner of \citet{Lewis1975} and \citet{Heim1982}, and explained the island-insensitivity we will see in \S\ref{sec:distance} by pied-piping the whole island at LF: the island moves, rather than being moved out of. \citet{Shimoyama2001,Shimoyama2006} and \citet{KratzerShimoyama2002} dispensed with movement and with variables altogether, in a semantics of alternatives after \citet{Hamblin1973}. Both are answers to the restrictor problem (\ref{ex:restrictor-problem}) and to the binding problem (\ref{ex:binding-problem}), and the difference between them is exactly the difference on which this seminar turns: a syntactic relation (binding, movement) against a semantic one (alternatives). The question of Lecture~3 is which locality each predicts. Both deserve credit. Nishigauchi's LF pied-piping was the first serious syntactic treatment, and it is still the one that predicts island-sensitivity in the right places; the Hamblin treatment is the standard now. The box will be a third option: storage, but in the syntax.

The Hamblin semantics runs as follows:

\begin{exe}
\ex \label{ex:hamblin} A Hamblin semantics for indeterminates (after \citealt{KratzerShimoyama2002}):
\begin{xlist}
\ex $\sem{\textit{dare}} = \{\, x : x \text{ is a person} \,\}$
\ex \label{ex:hamblin-pfa} $\sem{\alpha\ \beta} = \{\, a(b) : a \in \sem{\alpha},\ b \in \sem{\beta} \,\}$ \src{(pointwise functional application)}
\ex \label{ex:hamblin-ka} $\sem{[X]\ \textit{ka}} = \{\, \lambda w.\ \exists p \in \sem{X}\,[\,p(w) = 1\,] \,\}$
\ex \label{ex:hamblin-mo} $\sem{[X]\ \textit{mo}} = \{\, \lambda w.\ \forall p \in \sem{X}\,[\,p(w) = 1\,] \,\}$
\ex \label{ex:hamblin-q} $\sem{[X]\ \textit{ka}_{\text{Q}}} = \sem{X}$ \src{(a question denotes its set of answers)}
\end{xlist}
\end{exe}

Every expression denotes a set: ordinary words denote singletons, and indeterminates denote sets with more than one member. The indeterminate introduces a set of alternatives, which expand up the tree by pointwise composition, (\ref{ex:hamblin-pfa}), until a particle catches them: \textit{ka} closes them existentially, \textit{mo} universally, and the question operator returns them as the question. Above the particle the alternatives are gone, and what remains is a singleton.\footnote{(\ref{ex:hamblin-ka})--(\ref{ex:hamblin-q}) are the propositional operators. Kratzer and Shimoyama also define closure at the level of DP; I keep to the propositional version here, since nothing below depends on the difference.} Selectivity follows from closeness: an indeterminate is caught by the first particle above it, and an indeterminate that should be caught by a higher particle cannot be, if a lower one intervenes. That last point is what will become `path locality' in \S\ref{sec:path}. Figure~\ref{fig:hamblin-deriv} works through the simplest case, \textit{Dono gakusei-mo kita} `Every student came' (\Hx{1}{ex:dono-gakusei}).

\pIIneedspace{9cm}
\begin{keiofig*}{fig:hamblin-deriv}{Alternatives expanding up the tree in \textit{Dono gakusei-mo kita}. The indeterminate denotes the set of students; the verb denotes a singleton, and (\ref{ex:hamblin-pfa}) applies it pointwise, giving one proposition per student; \textit{mo} closes the set by (\ref{ex:hamblin-mo}), and the result is a singleton. The particle is pronounced on the nominal, but the operator it spells out applies to the propositional alternatives at the clause.}
\SOtree{[{\scriptsize $\{\lambda w.\,\forall p\in\sem{\text{TP}}\,[p(w)=1]\}$\\ \scriptsize = one proposition}, SOtree, xslim, align=center
  [{TP\\ \scriptsize $\{\lambda w.\,\mathit{came}(a)(w),$\\ \scriptsize $\lambda w.\,\mathit{came}(b)(w),\ \ldots\}$}, draw=gtaccent, rounded corners=2pt
    [{DP\\ \scriptsize $\{a, b, c, \ldots\}$}, draw=gtaccent, rounded corners=2pt
      [{\LI{dono gakusei}\\ \scriptsize $\{x : \mathit{student}(x)\}$}, SOlexical]]
    [{VP\\ \scriptsize $\{\lambda x\lambda w.\,\mathit{came}(x)(w)\}$} [{\LI{kita}}, SOlexical]]]
  [{\LI{mo} \quad \scriptsize $\forall$: (\ref{ex:hamblin-mo})}, SOhead]]}
\end{keiofig*}

The two boxed nodes are the point of the picture. The set grows at DP and remains a set at TP until the operator applies. Nothing has moved, and nothing has been bound; the only locality the system knows is `the first particle up'. Keep in mind, too, the gap between where \textit{mo} is pronounced and where its operator applies, on the nominal and at the clause respectively. In Lecture~3 it becomes the claim that Japanese pronounces the retrieval step.

Two observations before going on. First, on this view nothing moves. The restriction stays where it is pronounced, and the alternatives it introduces are carried up the tree by the semantics until an operator is reached. That is a storage theory, done in the semantics: the set of alternatives is the store, and the particle is the retrieval operator; nothing in it is a quantifier in the sense of \citet{BarwiseCooper1981} until the particle has applied. Lecture~3 makes the comparison with Cooper storage explicit (\Hs{3}{sec:box-cooper}), and asks whether the Japanese facts favour a semantic store over a syntactic one, starting from the observation that \citet{KratzerShimoyama2002} is a storage theory done semantically. Second, what the particle contributes is not a quantifier in Barwise and Cooper's sense but an operation on a set, which is why, as \S\ref{sec:nandomo} shows, the force it yields depends on what kind of set it is given.

% =====================================================================
\section{Force at a distance}\label{sec:distance}

This section has two aims. The first is to show that the particle can fix force at a distance from the restriction, across a boundary that movement could not cross. The second is to separate two relations that are easily run together: the association between a particle and the indeterminates it binds, and the scope of the quantifier that results. The second aim is the one that pays off in Lecture~3, where association turns out to have the profile of a semantic store, and scope and licensing the profile of the box. Here I only draw the distinction.

\subsection{An indeterminate inside an island}\label{sec:island}

\begin{exe}
\ex \label{ex:shimoyama}
\gll [[Dono gakusei-ga syootaisita] sensei]-mo odotta.\\
     which student-\gl{nom} invite.\gl{pst} teacher-\gl{mo} dance.\gl{pst}\\
\glt `For every student $x$, the teacher(s) that $x$ invited danced.' \src{(after \citealt{Shimoyama2006})}
\end{exe}

The indeterminate \textit{dono gakusei} sits inside a relative clause, a complex NP and an island for movement; \textit{mo} attaches to the whole nominal, outside it. Yet \textit{mo} quantifies over students: (\ref{ex:shimoyama}) says that for every student, the teacher that student invited danced. The association between the particle and the indeterminate is not sensitive to the island. \citet{Nishigauchi1990} captured this by pied-piping the whole complex NP at LF, so that the island moves rather than being moved out of. On Shimoyama's analysis nothing moves at all: the alternatives introduced by \textit{dono gakusei} (one for each student) are composed through the relative clause like anything else, yielding one teacher-alternative per student, and \textit{mo} at the top closes them universally. Figure~\ref{fig:island} sets the structure beside the alternatives.

\pIIneedspace{8cm}
\begin{keiofig*}{fig:island}{The island in (\ref{ex:shimoyama}). Left: the structure; the dashed box is the relative clause, an island for movement, and \textit{mo} attaches outside it. Right: the alternatives, step by step. $R_x$ is `$x$ invited \_\_', $\mathit{teach}_x$ the teacher $x$ invited, and $p_x$ the proposition that $\mathit{teach}_x$ danced; \textit{mo}, by (\ref{ex:hamblin-mo}), then yields $\{\lambda w.\,\forall p\in\{p_a,p_b,\ldots\}\,[p(w)=1]\}$. For readability the relative clause is treated as a predicate abstract and the head as definite.}
\begin{twocol}{0.40}
\centering
\SOtree{[TP, SOtree, tiny, flat
  [{DP-\textit{mo}}, draw=gtaccent, rounded corners=2pt
    [DP
      [{CP \quad \ghost{\scriptsize island}}, draw=gtghost, dashed, rounded corners=2pt [{\LI{dono gakusei-ga syootaisita}}, SOlexical]]
      [NP [{\LI{sensei}}, SOlexical]]]
    [{\LI{-mo}}, SOhead]]
  [VP [{\LI{odotta}}, SOlexical]]]}
\midcol
\centering
\SOtree{[{TP \quad \scriptsize $\{p_a, p_b, \ldots\}$}, SOtree, tiny, flat, align=center
  [{DP \quad \scriptsize $\{\mathit{teach}_a, \mathit{teach}_b, \ldots\}$}, draw=gtaccent, rounded corners=2pt
    [{CP \quad \scriptsize $\{R_a, R_b, \ldots\}$ \quad \ghost{island}}, draw=gtghost, dashed, rounded corners=2pt
      [{\LI{dono gakusei-ga} \quad \scriptsize $\{a,b,\ldots\}$}, SOlexical]
      [{\LI{syootaisita}}, SOlexical]]
    [{\LI{sensei}}, SOlexical]]
  [{\LI{odotta}}, SOlexical]]}
\end{twocol}
\end{keiofig*}

Read the right-hand tree from the bottom. The indeterminate denotes the students, $\{a,b,\ldots\}$. Pointwise composition carries one alternative per student through the relative clause, one `invited by $x$' predicate per student: the island is not crossed by anything, and the set simply grows inside it. The head noun composes pointwise too, giving one teacher per student; at the clause there is one proposition per student, and \textit{mo} closes them universally. At no step does anything leave the island. The alternatives are inherited upward by (\ref{ex:hamblin-pfa}) like any other semantic value, which is why the semantic store is island-insensitive by design, and why \S\ref{sec:localities} has to ask what it is sensitive to instead.

Notice, now, what \textit{mo} actually quantifies over in (\ref{ex:shimoyama}):

\begin{center}\small
\begin{tabularx}{\linewidth}{@{}P{4.4cm}LP{1.6cm}@{}}
\toprule
Node & Alternatives & Indexed by\\
\midrule
\textit{dono gakusei} & $\{a, b, c\}$ & students\\
{}[\textit{dono gakusei-ga syootaisita}] \textit{sensei} & $\{$the teacher $a$ invited, the teacher $b$ invited, \ldots$\}$ & students\\
{}[\ldots]-\textit{mo odotta} & $\{\lambda w.\,\forall p \in \{\ldots\}\,[p(w)=1]\}$: one proposition & --\\
\bottomrule
\end{tabularx}
\end{center}

Not teachers, and not students directly, but a set of alternatives indexed by students. This, more than anything else, is what tells against treating \textit{dono} N-\textit{mo} as a lexical universal `every N'. A lexical `every' would have to quantify over the restriction it is attached to, here the teachers, and \textit{mo} does not. It is, I think, the strongest single argument for the configurational view of quantifiers in \Hx{1}{ex:Qconfig}.

\pIIneedspace{6cm}
\begin{excursus}[Nishigauchi's alternative: pied-piping the island at LF]
\begin{center}
\SOtree{[CP, SOtree, tiny
  [{DP-\textit{mo}}, SOmoved [{\LI{[[dono gakusei-ga \ldots] sensei]-mo}}, SOlexical]]
  [TP [{DP-\textit{mo}}, SOtrace [{$\langle$\ldots$\rangle$}]] [VP [{\LI{odotta}}, SOlexical]]]]}
\end{center}
On \posscitet{Nishigauchi1990} analysis the whole complex NP raises at LF, and the indeterminate inside it is bound by the operator from within the moved phrase. The island is respected, since nothing is extracted from it: the island itself moves, and the indeterminate rides along. The virtue of the analysis is that it predicts island-sensitivity exactly where \textit{wh}-in-situ shows it, and that it is a \emph{syntactic} theory of association, which the alternative semantics is not. Its cost, in the architecture of Lecture~1, is threefold: a level of LF, a movement that is never pronounced, and a binding relation into the pied-piped phrase that still needs an index. Each of these is something that the architecture gave up (\Hs{1}{sec:architecture}). That, and the fact that the path effects of \S\ref{sec:path} are hard to state on a movement theory, is my answer to the syntactician who asks why pied-piping is not enough.
\end{excursus}

\subsection{The particle pins scope}\label{sec:pins}

If the association is so free, one might expect the scope of the resulting quantifier to be free too. It is not. \citet{Mineshima2005} argued that a \textit{mo}-phrase cannot take scope outside a complex NP that contains it, whereas \textit{subete-no} N can:
\begin{exe}
\ex \label{ex:mineshima}
\begin{xlist}
\ex \label{ex:mineshima-a}
\gll Hanako-ga [[dono gakusei-mo kaita] ronbun]-o yonda.\\
     Hanako-\gl{nom} which student-\gl{mo} write.\gl{pst} paper-\gl{acc} read.\gl{pst}\\
\glt `Hanako read the paper that every student wrote.' \src{(the paper does not vary with the students)}
\ex \label{ex:mineshima-b}
\gll Hanako-ga [[subete-no gakusei-ga kaita] ronbun]-o yonda.\\
     Hanako-\gl{nom} all-\gl{gen} student-\gl{nom} write.\gl{pst} paper-\gl{acc} read.\gl{pst}\\
\glt `Hanako read the paper(s) that all the students wrote.' \src{(also: for each student, the paper that student wrote)}
\end{xlist}
\end{exe}
The generalisation is that the scope of a \textit{mo}-phrase is pinned to the position of \textit{mo}, and does not leave the complex NP, whereas a nominal universal can be understood outside it. The source is a talk rather than a published paper, and the wide reading of (\ref{ex:mineshima-b}) has to be told apart from a cumulative one, in which the students wrote the papers between them, which \textit{subete-no} N, like English \textit{all}, also allows; the judgement sheet asks both questions (\exref{sheet}, items 3, 4 and 11). And \Hx{1}{ex:dareka-dono} showed the same rigidity in a simple clause: \textit{dono hon-mo} in object position cannot scope over the subject unless it is scrambled above it. So we must distinguish two relations:

\begin{keybox}
\begin{exe}
\ex \label{ex:assoc-scope}
\begin{xlist}
\ex \label{ex:association} \textit{Association}: the relation between a particle and the indeterminates it binds. It is insensitive to islands but sensitive to closeness, in that the first particle up catches the alternatives.
\ex \label{ex:scope} \textit{Scope}: the relation between the quantifier formed by a particle and the operators of the clause. It is fixed at the particle's attachment site.
\end{xlist}
\end{exe}
\end{keybox}

The particle pins scope. That is the first sense in which force in Japanese is fixed at a position other than the restriction's: force is introduced where the particle is, and scope is computed from there. Two relations, two localities, and in the end two theories: association behaves like alternatives, that is, like a semantic store, while scope behaves like a position, which is to say syntactically. That is the whole of the argument of Lecture~3 in miniature, and I will not try to resolve it here.

\subsection{Conjunctive quantifiers and a covert D}\label{sec:conjunctive}

The particles also carry polarity. Quantifiers built from an indeterminate and a conjunction, \textit{-mo} in Japanese, \textit{-to} in Korean, \textit{-um} in Malayalam \citep{Jayaseelan2001}, are universal and polarity-sensitive: without a further element, they need a clausemate negation (or, in some languages, a particular mood):

\begin{exe}
\ex \label{ex:daremo}
\begin{xlist}
\ex \label{ex:daremo-neg}
\gll Dare-mo ko-nakat-ta.\\
     who-\gl{mo} come-\gl{neg}-\gl{pst}\\
\glt `No one came.'
\ex \label{ex:daremo-aff}
\gll \llap{*}Dare-mo kita.\\
     who-\gl{mo} come.\gl{pst}\\
\glt (on the intended reading)
\ex \label{ex:daremoga}
\gll Dare-mo-ga kita.\\
     who-\gl{mo}-\gl{nom} come.\gl{pst}\\
\glt `Everyone came.'
\end{xlist}
\ex \label{ex:amwuto}
\begin{xlist}
\ex \label{ex:amwuto-neg}
\gll Amwu-to o-ci anh-ass-ta.\\
     any-\gl{to} come-\gl{ci} \gl{neg}-\gl{pst}-\gl{decl}\\
\glt `No one came.' \src{(= (\ref{ex:amwuto-a}))}
\ex \label{ex:amwuto-aff}
\gll \llap{*}Amwu-to o-ass-ta.\\
     any-\gl{to} come-\gl{pst}-\gl{decl}\\
\glt (ungrammatical without negation)
\end{xlist}
\end{exe}

The judgement in (\ref{ex:daremo-aff}) is a judgement on the polarity reading, and for \textit{dare-mo} the accent matters: the universal and the polarity uses are distinguished prosodically as well as by case. (\ref{ex:daremoga}), with case, is a plain universal; case, as we shall see in a moment, bleeds the sensitivity.

In \citet{GilTsoulas2013} we argued that the polarity sensitivity is not a property of the indeterminate, nor of the conjunction, nor of their combination. It lives in a covert D above them:

\begin{exe}
\ex \label{ex:covertD} \lbr{DP} D[uNeg] / D[uMood] \lbr{QP} Indeterminate \ conj ]]\\
{\small D covert, licensed by clause-bound Agree with Neg (or Mood)}
\end{exe}

\begin{keiofig}{fig:covertD}{The covert D of (\ref{ex:covertD}) in \textit{Dare-mo ko-nakat-ta} `No one came', checked under Agree with a clausemate Neg. The tree is schematic: the verb and negation are drawn as separate heads, for the sake of the arrow.}
\SOtree{[NegP, SOtree, xslim, tikz+={\draw[->, gtaccent, line width=0.6pt] (dd.west) to[out=160, in=200] node[midway, left, text=gtghost] {\scriptsize Agree} (neg.west);}
  [{Neg \quad \LI{-nai}}, SOhead, name=neg]
  [\textit{v}P
    [DP, draw=gtaccent, rounded corners=2pt, name=dd
      [{D$_{[u\mathrm{Neg}]}$ \quad \ghost{\scriptsize covert}}, SOhead]
      [QP [{\LI{dare}}, SOlexical] [{\LI{-mo}}, SOhead]]]
    [VP [{\LI{ko-}}, SOlexical]]]]}
\end{keiofig}

The D bears an uninterpretable negative (or modal) feature, checked under Agree with a clausemate negation, which is why the dependency is clause-bound. We also proposed a principle, paraphrased in (\ref{ex:principle30}), to say when such a feature can be uninterpretable at all:

\begin{exe}
\ex \label{ex:principle30} A head may bear an uninterpretable feature [uF] only if [F] is expressible configurationally and is not lexicalised (\citealt{GilTsoulas2013}: (30), paraphrased; cf.\ \citealt{Ladusaw1996}).
\end{exe}

When D is filled by case, the sensitivity is bled, which is why (\ref{ex:daremoga}) is a plain universal. Two things in this analysis matter for Lecture~3. The force-provider (the particle) is kept apart from the element that needs licensing (the D), and the licensing is a syntactic relation with an operator on the spine, bounded by the clause. That is exactly the licensing layer the box theory will need for \textit{amwu-to} and \textit{dare-mo}, and it was in place a decade before the box. The split between store and retrieval is at least that old; what the box adds is the rationale that (\ref{ex:principle30}) could only stipulate, namely that force and licensing sit high because that is where retrieval happens. Lecture~3 inherits the principle in that form (\Hs{3}{sec:precedents}). The correspondence is close enough to tabulate:

\begin{center}\small
\begin{tabularx}{\linewidth}{@{}LL@{}}
\toprule
\citet{GilTsoulas2013} & The box theory (Lecture~3)\\
\midrule
indeterminate = restricted variable & indeterminate = stored restriction (low)\\
conj = QP head (force) & particle = retrieval operator (force)\\
covert D[uNeg/uM], needs licensing & high force position, consulted at the phase\\
Agree with clausemate Neg/Mood & consultation and discharge (clause-bound)\\
force configurational, not lexical (\ref{ex:principle30}) & force at retrieval, not a lexical quantifier\\
\bottomrule
\end{tabularx}
\end{center}

For now the table is a promissory note.

\subsection{A caveat: \textit{nan-do-mo}}\label{sec:nandomo}

Let me raise the obvious objection myself, before the reader does. If force sits on the particle, \textit{mo} should always yield a universal. It does not:

\begin{exe}
\ex \label{ex:nandomo}
\gll Taroo-wa sono eiga-o nan-do-mo mita.\\
     Taro-\gl{top} that film-\gl{acc} what-time-\gl{mo} see.\gl{pst}\\
\glt `Taro saw that film many times.'
\ex \label{ex:itidomo}
\gll Taroo-wa sono eiga-o iti-do-mo mi-nakat-ta.\\
     Taro-\gl{top} that film-\gl{acc} one-time-\gl{mo} see-\gl{neg}-\gl{pst}\\
\glt `Taro did not see that film even once.'
\end{exe}

With a numeral-classifier indeterminate, \textit{mo} gives `many' in the affirmative, not `every' (there is no sense in which Taro saw the film every number of times), and with the numeral `one' it gives `not even once'. The natural conclusion is that \textit{mo} is not a universal quantifier at all, but an additive and scalar operator, `also' and `even', and that the universal reading of \textit{dono} N-\textit{mo} is what conjunction over a set of individual alternatives amounts to. \citet{Nakanishi2006} and \citet{KobuchiPhilip2009} give the Japanese-internal arguments for this, and \citet{Szabolcsi2015} the general one. So the claim that force sits on the particle needs a qualification:

\begin{keybox}
\noindent\textit{Before.} Force sits on the particle.\par\smallskip
\noindent\textit{After.} The particle contributes an \emph{operation}, and force is the result of applying it to whatever has been stored; the same particle yields different forces for different stores.
\end{keybox}

This does not weaken the case for separating force from restriction. If anything it strengthens it: force is not a lexical property of any item, but something computed at the point of retrieval. The point will come back twice, in the fourth question for the discussion (\S\ref{sec:readings}), and in the question with which \S\ref{sec:hierarchy} ends: at which step, storage or retrieval, is force assigned, and is that step pronounced?

% =====================================================================
\section{Unselective binding}\label{sec:unselective}

\subsection{Lewis's problem}\label{sec:lewis}

The modern study of quantification on the clause begins with \citet{Lewis1975}:

\begin{exe}
\ex \label{ex:lewis}
\begin{xlist}
\ex \label{ex:quadratic} A quadratic equation usually has two solutions.
\ex Riders on the Thirteenth Avenue line seldom find seats.
\ex \label{ex:donkey} If a farmer owns a donkey, he always beats it.
\end{xlist}
\end{exe}

In (\ref{ex:quadratic}) \textit{usually} does not quantify over times (a given equation does not have two solutions on most occasions and fewer on others), and the indefinite does not have existential force (the sentence is not made true by one equation with two solutions). The natural paraphrase is \textit{most quadratic equations have two solutions}. The adverb has, so to speak, taken over the quantificational force that the generalised quantifier tradition locates in the determiner, and the indefinite has been demoted to a restriction. Lewis's proposal was that adverbs of quantification are unselective binders: they bind all the free variables in their scope at once, quantifying over cases, that is, over assignments to those variables that satisfy the restriction:

\begin{exe}
\ex \label{ex:always} $\text{\textsc{always}}_{x,y}\;[\mathrm{farmer}(x) \wedge \mathrm{donkey}(y) \wedge \mathrm{own}(x, y)]\;[\mathrm{beat}(x, y)]$
\end{exe}

\pIIneedspace{4cm}
\begin{keiofig*}{fig:tripartite}{Tripartite structure: the operator relates two open formulas, the restrictor and the nuclear scope, and binds every free index in both. This is the structure that Heim, Kamp, Diesing and Kratzer all inherit.}
\begin{tikzpicture}[every node/.style={font=\small, draw=gtaccent, rounded corners=2pt, inner sep=5pt, align=center}, >=Stealth]
  \node (op) at (0,0) {\textbf{Operator}\\ \textsc{always}$_{x,y}$};
  \node (res) at (4.9,0) {\textbf{Restrictor}\\ farmer$(x)$, donkey$(y)$, own$(x,y)$};
  \node (nuc) at (10.0,0) {\textbf{Nuclear scope}\\ beat$(x,y)$};
  \draw[->, gtaccent, line width=0.6pt] (op) -- (res);
  \draw[->, gtaccent, line width=0.6pt] (res) -- (nuc);
\end{tikzpicture}
\end{keiofig*}

The word to underline is `unselective': the adverb does not choose which variables to bind; it binds whatever is free. Three ingredients of this analysis became foundational, and each is still argued about. The first is tripartite structure: an operator, a restrictor and a nuclear scope, in place of the bipartite quantifier-plus-argument of generalised quantifier theory. The second is the treatment of indefinites as restricted variables, developed into a full theory by \citet{Kamp1981} and \citet{Heim1982}, with existential closure supplying force when nothing else does; this is the indefinite as the limiting case of the Q-configuration, as in \Hs{1}{sec:indefinites}. The third is that the restrictor is assembled non-locally: the \textit{if}-clause in (\ref{ex:donkey}) supplies the restriction although it is not the adverb's sister. That is the restrictor problem (\ref{ex:restrictor-problem}), answered, and it is the ingredient that Japanese \textit{-te-mo} makes visible (\S\ref{sec:tripartite}).

Unselectivity, taken literally, overgenerates. Suppose there are ten farmers, nine of whom own one donkey each and beat it, while the tenth owns a hundred donkeys and beats none:

\begin{exe}
\ex \label{ex:proportion} Most farmers who own a donkey beat it.
\end{exe}

(\ref{ex:proportion}) is true in this situation, but quantification over farmer--donkey pairs would make it false: there are nine beating pairs and a hundred non-beating ones, so that most of the 109 pairs are non-beating \citep{Kadmon1987}. Whatever binds in (\ref{ex:proportion}) is selective, and the binding problem (\ref{ex:binding-problem}) is not solved by binding everything. The responses define the space in which the syntax will eventually have to choose: situation-based selectivity, on which adverbs quantify over minimal situations and reach individuals only through them \citep{Berman1987,Heim1990,vonFintel1994,Elbourne2005}, and dynamic binding, on which indefinites keep their binding potential beyond their syntactic scope \citep{GroenendijkStokhof1991,Chierchia1992,Chierchia1995}. Quantificational variability, where an adverb quantifies over the parts of a plural or the answers to a question, belongs here too; in \citet{GillHarlowTsoulas2004} we argued that it arises with indeterminates and connectives as well.

The three excursions that follow are for readers from the semantic traditions concerned. Each of them, as will become clear, \emph{moves} the variable rather than removing it.

\pIIneedspace{5cm}
\begin{excursus}[For readers from DRT: the donkey sentence as a box]
\begin{center}
\begin{tikzpicture}[every node/.style={font=\small}]
  \node[draw=gtaccent, rounded corners=1pt, inner sep=5pt, align=left] (K1) at (0,0) {$x\ \ y$\\[2pt] farmer$(x)$\\ donkey$(y)$\\ own$(x,y)$};
  \node at (1.55,0) {$\Rightarrow$};
  \node[draw=gtaccent, rounded corners=1pt, inner sep=5pt, align=left] (K2) at (3.1,0) {\phantom{$x$}\\[2pt] beat$(x,y)$};
  \node[draw=gtaccent!50, rounded corners=1pt, inner sep=8pt, fit=(K1)(K2), label={[font=\scriptsize, text=gtghost]above:the main DRS, with no discourse referents of its own}] {};
\end{tikzpicture}
\end{center}
A discourse representation structure introduces \emph{discourse referents} ($x$, $y$) and \emph{conditions} on them \citep{Kamp1981,KampReyle1993}. An indefinite introduces a referent, and has no force of its own. The conditional $K_1 \Rightarrow K_2$ is verified if and only if every embedding of $K_1$ extends to an embedding of $K_2$, and that is where universal force comes from. The referents of $K_1$ are accessible in $K_2$, so that \textit{he} and \textit{it} pick up $x$ and $y$. This is unselective binding, drawn: $\Rightarrow$ quantifies over every referent of $K_1$ at once, and the proportion problem is the same problem here, with the same answers (duplex conditions). For readers from LFG and HPSG, DRT is the semantics that those frameworks have often been paired with. But the discourse referent is a variable by another name, and the question of Lecture~4 applies to it unchanged.
\end{excursus}

\begin{excursus*}[For readers from dynamic semantics: binding beyond scope]
In Dynamic Predicate Logic \citep{GroenendijkStokhof1991} a formula does not denote a truth value but a \emph{relation between assignments}: an input $g$ and the outputs $g'$ it can leave behind. Most formulas are \emph{tests}: they let $g$ through unchanged, or not at all. An existential is not a test. It \emph{resets} its variable and passes the new assignment on:
\begin{center}
\begin{tikzpicture}[>=Stealth, font=\footnotesize,
    st/.style={text=gttext},
    fm/.style={draw=gtaccent, rounded corners=2pt, inner sep=4pt, align=center},
    capt/.style={font=\scriptsize, text=gtghost, align=center},
    fl/.style={->, gtaccent, line width=0.7pt}]
  \node[st] (g) at (0.3,0) {input $g$};
  \node[fm] (ex) at (2.6,0) {$\exists x\,\mathit{farmer}(x)$};
  \node[st] (mid) at (5.8,0) {$g[x\mapsto a],\ g[x\mapsto b],\ \ldots$};
  \node[fm] (test) at (9.0,0) {$\mathit{whistles}(x)$};
  \node[st] (out) at (11.6,0) {$g[x\mapsto a]$, if $a$ whistles};
  \draw[fl] (g) -- (ex); \draw[fl] (ex) -- (mid); \draw[fl] (mid) -- (test); \draw[fl] (test) -- (out);
  \node[capt, below=2pt of ex] {resets $x$:\\ one output per farmer};
  \node[capt, below=2pt of mid] {the assignments flow on};
  \node[capt, below=2pt of test] {a test on $x$:\\ keeps the outputs that pass};
\end{tikzpicture}
\end{center}
The assignment flows left to right, and the second conjunct sees the $x$ that the first conjunct set, although it is outside the syntactic scope of $\exists x$. That is the whole system, summed up in one equivalence:
\[\exists x\,\varphi \wedge \psi \ \equiv\ \exists x\,(\varphi \wedge \psi)\]
`A farmer walks in. He whistles.' The second sentence's \textit{he} is bound by the first sentence's \textit{a farmer}, and that is what the equivalence says: conjoining after the existential is the same as conjoining inside it. The reach of a binder is measured by how far the change it makes to the assignment is carried, not by where its brackets close.

The donkey sentence (\ref{ex:donkey}) comes out as $\exists x\exists y\,[\mathit{farmer}(x)\wedge\mathit{donkey}(y)\wedge\mathit{own}(x,y)] \to \mathit{beat}(x,y)$. The antecedent resets $x$ and $y$, one output per farmer--donkey pair, and the dynamic conditional is defined so that \emph{every} output of the antecedent must satisfy the consequent. That is where the universal force comes from, with no universal quantifier anywhere in the sentence: the conditional does the quantifying, as in DRT, and the pronouns are tests on the assignments passed to them. If every reset variable is quantified over, the proportion problem (\ref{ex:proportion}) follows, and \citeauthor{Chierchia1995}'s dynamic binding \citep{Chierchia1992,Chierchia1995} is the answer: a determiner or an adverb binds \emph{selectively}, so that \textit{most farmers} quantifies over farmers only, while the donkey-value is carried dynamically into the nuclear scope and picked up there. The indefinite keeps its binding potential beyond its syntactic scope: it binds \textit{it} without being quantified over.

What it means for us is this. Dynamic binding is the purest form of the idea that a variable persists beyond the point at which it was introduced: a store of values indexed by variable \emph{names}, carried from one conjunct to the next by the semantics. The store is what this seminar wants; the names are what Inclusiveness forbids, and the store inherits in full the question of where the names come from. If a dynamic semanticist is reading, the question I would put is whether their assignments are objects that the syntax has to supply. Lecture~4 will say that the syntax supplies relations, not names.
\end{excursus*}

\begin{excursus}[For readers from situation semantics: minimal situations]
On the situation-based view \citep{Berman1987,Heim1990,Elbourne2005}, \textit{always} quantifies over \emph{minimal situations} satisfying the restrictor: for (\ref{ex:donkey}), the smallest situations containing a farmer, a donkey he owns, and nothing else, each of which extends to one in which he beats it. Individuals are reached only through the situations they inhabit. What this buys is selectivity: (\ref{ex:proportion}) quantifies over farmers, and the donkey is picked up by a situation pronoun in the nuclear scope (Elbourne's \textit{it} is `the donkey in $s$'; Heim's E-type pronouns are its precursor).

What it owes, from our point of view, is that a situation variable is still a variable. \textit{It} is a definite description with a situation argument, bound by the quantifier over situations, and that is an index again, now on situations. Each of the three traditions moves the variable rather than removing it: to a discourse referent, to a dynamic assignment, to a situation argument. The claim of this seminar is not that any of them is wrong, but that none of them answers the question of where the index comes from; that is the point of \S\ref{sec:bites}.
\end{excursus}

\subsection{Existential closure and the restrictor: Diesing and Rooth}\label{sec:diesing}

\citet{Diesing1992keio} gave the tripartite structure a syntax. On her Mapping Hypothesis, material inside VP is mapped onto the nuclear scope, where existential closure applies, and material outside VP onto the restrictor. German shows it on the surface:

\begin{exe}
\ex \label{ex:german}
\begin{xlist}
\ex \label{ex:german-a}
\gll \ldots{} weil ja doch Linguisten Kammermusik spielen.\\
     \ldots{} since \gl{prt} \gl{prt} linguists chamber.music play\\
\glt `\ldots{} since there are linguists playing chamber music.' \src{(existential)}
\ex \label{ex:german-b}
\gll \ldots{} weil Linguisten ja doch Kammermusik spielen.\\
     \ldots{} since linguists \gl{prt} \gl{prt} chamber.music play\\
\glt `\ldots{} since linguists (in general) play chamber music.' \src{(generic)}
\end{xlist}
\end{exe}

The particles \textit{ja doch} mark the edge of VP. An indefinite inside it is existentially closed; an indefinite outside it restricts a generic operator. German shows on the surface what English shows only through intonation and interpretation: scrambling the indefinite across the particles changes its force. Figure~\ref{fig:diesing} draws the two structures.

\pIIneedspace{8cm}
\begin{keiofig*}{fig:diesing}{Diesing's mapping as a tree. Left, (\ref{ex:german-a}): the indefinite inside the edge marked by \textit{ja doch}, existentially closed. Right, (\ref{ex:german-b}): the indefinite scrambled outside the edge, in the restrictor of the generic operator. The boxed node moves from \textit{v}P to TP as the indefinite scrambles.}
\begin{twocol}{0.47}
\centering
\SOtree{[CP, SOtree, tiny, flat
  [{C \quad \LI{weil}}, SOhead]
  [TP
    [{\LI{ja doch}\\ \ghost{\scriptsize edge}}, SOlexical]
    [{\textit{v}P\\ \ghost{\scriptsize nuclear scope:}\\ \ghost{\scriptsize $\exists$-closure}}, draw=gtaccent, rounded corners=2pt
      [{\LI{Linguisten}}, SOlexical] [{\LI{Kammermusik}\\ \LI{spielen}}, SOlexical]]]]}
\midcol
\centering
\SOtree{[CP, SOtree, tiny, flat
  [{C \quad \LI{weil}}, SOhead]
  [{TP\\ \ghost{\scriptsize restrictor: \textsc{gen}}}, draw=gtaccent, rounded corners=2pt
    [{\LI{Linguisten}}, SOmoved]
    [TP
      [{\LI{ja doch}}, SOlexical]
      [{\textit{v}P\\ \ghost{\scriptsize nuclear scope}} [{$\langle$\LI{Linguisten}$\rangle$}, SOtrace] [{\LI{Kammermusik}\\ \LI{spielen}}, SOlexical]]]]]}
\end{twocol}
\end{keiofig*}

Existential closure at the edge of \textit{v}P is, I think, the first retrieval-site proposal in the literature: a designated position on the clausal spine at which force-less material receives force. It is the first place, that is, where a position rather than an item supplies force. It is also, avant la lettre, a phase-theoretic idea, since the edge of \textit{v}P is where Transfer applies; the box's discharge at the phase is its descendant. Lecture~3 argues the point, but the parallel can be set out now:

\begin{center}\small
\begin{tabularx}{\linewidth}{@{}P{2.8cm}LL@{}}
\toprule
 & \citet{Diesing1992keio} & The box (Lecture~3)\\
\midrule
What is force-less & an indefinite inside VP & a boxed quantifier, pronounced at its token\\
Where force arrives & the edge of VP (\textit{v}P): existential closure & the phase edge: discharge of the verb's $\theta$-feature\\
What marks the site & \textit{ja doch} in German; nothing in English & a particle in Japanese and Korean; nothing in English\\
Locality & the mapping is clause-internal & $\theta$-discharge is clause-bound\\
\bottomrule
\end{tabularx}
\end{center}

Diesing's closure is a position on the spine that supplies force to force-less material. The box generalises the idea: every phase edge is such a position, and what is retrieved there is whatever was stored below. The parallel is the reason why the phase heads of (\ref{ex:spine}), \textit{v} and C, will turn out to be the box's retrieval sites, and the other heads the positions it consults.

\citet{Rooth1985,Rooth1995} showed that focus can do the same job:

\begin{exe}
\ex \label{ex:rooth}
\begin{xlist}
\ex \label{ex:rooth-a} Mary always takes \textsc{John} to the movies.\\
`Whenever Mary takes someone to the movies, it is John.'
\ex Mary always takes John to the \textsc{movies}.\\
`Whenever Mary takes John somewhere, it is to the movies.'
\end{xlist}
\end{exe}

The focused constituent falls in the nuclear scope, and the background supplies the restrictor: the same words, the same structure, and the partition is read off prosody. Diesing and Rooth are usually presented as rivals, the one reading the partition off position, the other off prosody. I would rather see them as two externalisations of one partition. Every language splits the clause into a restrictor and a nuclear scope; languages differ in whether the split is read off position (German scrambling around \textit{ja doch}, Japanese \textit{wa} and \textit{ga}, as in \citealt{Kuroda1972} and \Hx{1}{ex:wa-ga}) or off prosody (English focus). This is where the externalisation hypothesis of \Hx{1}{ex:variation} does its first piece of work in this lecture. It is also the same kind of variation as the one the rigidity generalisation of Lecture~3 turns on, whether a given step of the computation is pronounced; Lecture~3 states that generalisation in one direction only.

\begin{excursus}[Focus semantics in one definition]
On \posscitet{Rooth1992} theory, every expression has an ordinary semantic value $\sem{\alpha}^{o}$ and a \emph{focus value} $\sem{\alpha}^{f}$: the set of alternatives obtained by replacing the focused part by something of the same type. For the predicate of (\ref{ex:rooth-a}):
\begin{align*}
&\sem{\textit{takes }[\textsc{John}]_F\textit{ to the movies}}^{f}\\
&\qquad = \{\,\lambda x.\ x \text{ takes } y \text{ to the movies} : y \in D_e\,\}
\end{align*}
The squiggle operator $\sim$ introduces a variable $C$ whose value must be a subset of the focus value, $C \subseteq \sem{\alpha}^{f}$, and a focus-sensitive operator such as \textit{always} or \textit{only} takes $C$ as its restrictor. In (\ref{ex:rooth-a}), $C$ is the set of propositions of the form `Mary takes $y$ to the movies', and \textit{always} quantifies over the cases in $C$: whenever one of them holds, it is the John one.

Two things are worth noticing. Focus values are Hamblin alternatives, the very objects of (\ref{ex:hamblin}), so that \S\ref{sec:indeterminates} and this section run on the same machinery. And $\sim$ introduces a variable, $C$, anchored in context: an index once more. Beck's intervention effects (\S\ref{sec:path}) are what happens when $\sim$ evaluates alternatives that were meant for a higher operator.
\end{excursus}

\subsection{Japanese pronounces the tripartite structure}\label{sec:tripartite}

\begin{wide}
\begin{exe}
\ex \label{ex:tripartite}
\begin{xlist}
\ex \label{ex:itumo}
\gll Taroo-wa itu-mo okureru.\\
     Taro-\gl{top} when-\gl{mo} be.late.\gl{npst}\\
\glt `Taro is always late.'
\ex \label{ex:itukitemo}
\gll [Itu kite]-mo, Taroo-wa inai.\\
     when come-\gl{te}-\gl{mo} Taro-\gl{top} be.\gl{neg}.\gl{npst}\\
\glt `Whenever one comes, Taro isn't there.'
\ex \label{ex:dare-nani-temo}
\gll [Dare-ga nani-o itte]-mo, Hanako-wa kika-nai.\\
     who-\gl{nom} what-\gl{acc} say-\gl{te}-\gl{mo} Hanako-\gl{top} listen-\gl{neg}.\gl{npst}\\
\glt `No matter who says what, Hanako doesn't listen.'
\end{xlist}
\end{exe}
\end{wide}

In (\ref{ex:itumo}) the adverb `always' is itself assembled from an indeterminate over times, \textit{itu} `when', and the particle: the composition of `always' made visible. (\ref{ex:itukitemo}) and (\ref{ex:dare-nani-temo}) are unconditionals, the Japanese counterparts of Lewis's \textit{if}-clauses, with the operator pronounced at their edge. The \textit{-te} clause supplies the restrictor; \textit{mo}, attached at its edge, is the operator; the main clause is the nuclear scope. Compositionally, the \textit{-te} clause denotes a set of alternatives, one for each value of its indeterminates (one for each pair of a person and a thing said, in (\ref{ex:dare-nani-temo})); \textit{mo} at its edge closes the set universally, so that the main clause is asserted for every alternative, which is the unconditional reading; and when the alternatives are ordered by likelihood rather than enumerated, as in (\ref{ex:concessive}) below, the same closure is heard as `even if'. Figure~\ref{fig:unconditional} draws (\ref{ex:dare-nani-temo}); (\ref{ex:itukitemo}) is the same with one indeterminate.

\begin{keiofig}{fig:unconditional}{The unconditional (\ref{ex:dare-nani-temo}): the operator at the restrictor's edge. Compare the three boxes of Figure~\ref{fig:tripartite}: operator, restrictor and nuclear scope, here in that order, each with its own morphology.}
\SOtree{[CP, SOtree, tiny
  [{CP-\textit{mo} \quad \ghost{\scriptsize restrictor + operator}}, draw=gtaccent, rounded corners=2pt
    [{TP \quad \ghost{\scriptsize restrictor}} [{\LI{dare-ga}}, SOlexical] [{\LI{nani-o}}, SOlexical] [{\LI{itte}}, SOlexical]]
    [{\LI{-mo} \quad \ghost{\scriptsize operator}}, SOhead]]
  [{TP \quad \ghost{\scriptsize nuclear scope}}, draw=gtaccent, rounded corners=2pt [{\LI{Hanako-wa kika-nai}}, SOlexical]]]}
\end{keiofig}

In (\ref{ex:dare-nani-temo}) a single \textit{mo} binds both indeterminates: the alternatives are indexed by pairs $\langle x, y\rangle$ such that $x$ says $y$, and \textit{mo} closes over all of them, so that whatever the pair, Hanako does not listen. One particle, several variables, at the restrictor's edge: this is Lewis's unselective binder, pronounced (cf.\ \citealt{Nishigauchi1990}). Notice, finally, that without an indeterminate inside it the same \textit{-te-mo} clause is simply concessive:

\begin{exe}
\ex \label{ex:concessive}
\gll Ame-ga hutte-mo iku.\\
     rain-\gl{nom} fall-\gl{te}-\gl{mo} go.\gl{npst}\\
\glt `Even if it rains, I will go.'
\end{exe}

Here \textit{mo} is `even', exactly as in (\ref{ex:itidomo}), \textit{iti-do-mo}. It is the same particle in every case; what differs is what it finds in the store: a set of alternatives indexed by times, a set indexed by pairs of individuals, or a singleton with a scalar presupposition. That is the thread of this lecture. Force is computed at the point where the operator meets what has been stored, and three times now (\textit{nan-do-mo}, \textit{-te-mo}, \textit{dono} N-\textit{mo}) the same \textit{mo} has yielded a different force for a different store.

\subsection{Inclusiveness again: what does \textit{always} bind?}\label{sec:bites}

Unselective binding is index binding. In Heim's system the adverb binds every free index in its restrictor, and the indefinites \emph{are} those indices, restricted by their descriptive content. In a grammar without indices, what does \textit{always} bind? Three answers are available:

\begin{center}\small
\begin{tabularx}{\linewidth}{@{}P{3.0cm}LL@{}}
\toprule
Answer & How it binds & The cost\\
\midrule
A situation variable & reaches individuals through the situations they inhabit & a situation variable is still a variable\\
A set of alternatives, as in (\ref{ex:hamblin}) & no variables at all & the set is computed by the semantics: the store is semantic\\
A syntactic relation & between the operator and the positions it binds, built by the operations the syntax already has & to be developed in Lecture~4\\
\bottomrule
\end{tabularx}
\end{center}

This is the binding problem (\ref{ex:binding-problem}) restated for a grammar that cannot write \emph{x}, and the three rows are the three answers of \Hs{1}{sec:three-answers} in their clausal form. I leave the question open until Lecture~4, where the third option is developed and where it turns out that the box already depends on such a relation.

% =====================================================================
\section{Force on the clause}\label{sec:clause}

\citet{Partee1995} distinguished \emph{D-quantification}, where force is lexicalised in the nominal extended projection (determiners: \textit{every}, \textit{most}, \textit{no}), from \emph{A-quantification}, where it is expressed by adverbs, affixes, auxiliaries and argument-structure adjusters on the verbal side. \S\ref{sec:unselective} dealt with A-quantification in its most familiar forms, the adverbs of English, and existential closure and the generic operator in German. This section turns to languages in which the verbal side carries more of the load, and ends by stating the spine on which force can be fixed. The point of the section is that the A side is not a marginal option, but a set of positions that every language has.

\subsection{Salish and Mohawk}\label{sec:salish}

\citet{Jelinek1995} argued that Straits Salish has no D-quantifiers at all. Its arguments are pronominal, and quantificational force is expressed by elements that behave as predicates or adverbs, so that all quantification is A-quantification. \citet{Baker1995} argued that Mohawk lacks quantified NPs of the \textit{everybody} and \textit{nobody} type. On his analysis the overt NPs of a polysynthetic language are adjuncts linked to pronominal arguments, and an adjunct cannot be a quantifier that binds an argument position, so universal force has to be expressed in other ways. \citet{Bittner1995} made related arguments for Kalaallisut (West Greenlandic).

\citet{Matthewson2001} disputed the Salish half of this. In St'\'at'imcets, a Salish language of the Interior branch, there are quantifiers in the nominal after all, but they combine with full DPs, [Q [D NP]], as in \Hx{1}{ex:statimcets}, and she argued that this is the structure of quantified nominals in every language. I want to take her result constructively. If she is right, then what distinguishes Salish from English is not whether force can be nominal, but how much of the nominal's structure is pronounced, which is precisely the kind of variation that the externalisation hypothesis (\Hx{1}{ex:variation}) leads us to expect.

\subsection{Kratzer (2005): indefinites and the operators they depend on}\label{sec:kratzer}

\citet{Kratzer2005} is the paper that connects this section to \S\ref{sec:indeterminates}, and it deserves to be stated in full rather than cited. Its thesis is that an indefinite is a \emph{restriction} with no force of its own. Its force comes from an operator elsewhere in the clause: a particle, an adverb of quantification, or existential closure at a designated position. Languages differ in which of these operators they have and where they sit; they do not differ in the indefinite.

\begin{center}\small
\begin{tabularx}{\linewidth}{@{}P{2.0cm}LLL@{}}
\toprule
& The restriction & The operator that supplies force & Where it sits\\
\midrule
Japanese & the indeterminate, \textit{dare}, \textit{dono} N & the particle, \textit{ka} or \textit{mo} & at the edge of the nominal or of the clause (\S\S\ref{sec:indeterminates}--\ref{sec:distance})\\
German & the indefinite, \textit{Linguisten} & existential closure, or \gl{gen} & the edge of VP (\S\ref{sec:diesing})\\
St'\'at'imcets & a DP with the determiners \textit{i \ldots{} a}, which mark a domain & existential closure & outside the nominal, at the clause\\
\bottomrule
\end{tabularx}
\end{center}

The mechanism is the one of \S\ref{sec:indeterminates}: the indefinite introduces alternatives, and the operator closes them where it stands. One analysis covers a particle language and a language with determiners, and the dependence of \textit{dare} on \textit{ka} and of a Salish indefinite on its closure operator comes out as one dependence.

Three things make the paper matter for this seminar. First, it is the semantic precedent for the configuration this section ends on, (\ref{ex:config}): force on the clause, restriction in the nominal, and an association between them, stated as a thesis about indefinites in general and not as a fact about Japanese. Second, it ties \S\ref{sec:indeterminates} to the languages of this section. Without it the indeterminate system looks like an exotic particle morphology; with it, Japanese is the \emph{pronounced} case of something every language does, and the particle is where the dependence of restriction on operator can be seen. Third, it makes the variation a variation in operators and positions. If the indefinite is the same everywhere, what differs between languages is which operators sit on the spine and whether they are pronounced; that is the premise of the spine (\ref{ex:spine}) and of the hierarchy of \S\ref{sec:hierarchy}, where it becomes P5.

What the paper leaves open is locality. Kratzer's operators close alternatives wherever the alternatives have expanded to, and nothing in the account says how far an alternative may travel before it is closed, or what may stand in its way. That is the question of \S\ref{sec:localities}, which shows that it has two answers, not one. An early version of the paper was given, as `DPs without quantificational force', at the Strategies of Quantification conference in York in 2004, which is where the line of work I pursued with Gil and Harlow first met Kratzer's.

\subsection{Mandarin \textit{dou}}\label{sec:dou}

\begin{exe}
\ex \label{ex:dou}
\begin{xlist}
\ex \label{ex:tamen-dou}
\gll Tamen dou lai le.\\
     they \gl{dou} come \gl{prt}\\
\glt `They all came.'
\ex \label{ex:mei-dou}
\gll Mei-ge xuesheng *(dou) lai le.\\
     every-\gl{cl} student \gl{dou} come \gl{prt}\\
\glt `Every student came.'
\ex \label{ex:shei-dou}
\gll Shei dou lai le.\\
     who \gl{dou} come \gl{prt}\\
\glt `Everyone came.'
\end{xlist}
\end{exe}

On \posscitet{Cheng1995} analysis \textit{dou} is a universal, distributive operator on the clausal spine, which associates with a plural or \textit{wh} element to its left. With a \textit{wh}-phrase it yields a universal, as in (\ref{ex:shei-dou}), and it is required with a preverbal \textit{mei}-phrase, as in (\ref{ex:mei-dou}). The obligatory \textit{dou} of (\ref{ex:mei-dou}) is the crucial datum: the `universal' \textit{mei} cannot do without the operator on the spine. \citet{Lin1998} treated \textit{dou} as a distributivity operator over covers; \citet{Cheng2009} argued that \textit{mei} contributes a maximal plurality and that the universal force of \textit{mei} \ldots{} \textit{dou} is supplied by \textit{dou}, which is to say that the force is \textit{dou}'s and not \textit{mei}'s (see also \citealt{GiannakidouCheng2006}; \citealt{Xiang2008}). The convergent picture is that the Mandarin nominal supplies the restriction and that force, distributive universal quantification, is contributed by a head on the spine, plausibly the overt counterpart of Beghelli and Stowell's Dist (\Hx{1}{ex:scopepositions}).

\begin{keiofig}{fig:dou}{\textit{Dou} as the overt counterpart of Dist, in the Beghelli--Stowell spine RefP $>$ CP $>$ AgrSP $>$ DistP $>$ ShareP $>$ NegP $>$ AgrOP $>$ VP. Whether the restriction moves to Spec-DistP or is merged there is not at issue; what matters is that the surface order shows the configuration.}
\SOtree{[DistP, SOtree, xslim
  [{DP \quad \ghost{\scriptsize restriction}}, SOmoved [{\LI{mei-ge xuesheng}}, SOlexical]]
  [Dist$'$
    [{Dist \quad \hl{\LI{dou}}}, SOhead]
    [\textit{v}P [DP, SOtrace [{$\langle$\LI{mei-ge xuesheng}$\rangle$}]] [VP [{\LI{lai le}}, SOlexical]]]]]}
\end{keiofig}

The requirement that the restriction be to the left of \textit{dou} then reduces to the requirement that it reach the head's specifier, overtly, where English \textit{every} would get there covertly. In a sentence: Mandarin pronounces one row of the Beghelli--Stowell spine that was posited for English on covert evidence.

\subsection{\textit{-zutu} and \textit{-ssik}}\label{sec:zutu}

\begin{exe}
\ex \label{ex:zutu}
\begin{xlist}
\ex \label{ex:zutu-a}
\gll Kodomo-tati-ga ringo-o mittu-zutu tabeta.\\
     child-\gl{pl}-\gl{nom} apple-\gl{acc} three.\gl{cl}-\gl{zutu} eat.\gl{pst}\\
\glt `The children ate three apples each.' \src{(= \Hx{1}{ex:spot-zutu})}
\ex \label{ex:zutu-b}
\gll Gakusei-ga hitori-zutu heya-ni haitta.\\
     student-\gl{nom} one.person-\gl{zutu} room-\gl{dat} enter.\gl{pst}\\
\glt `The students entered the room one at a time.'
\end{xlist}
\ex \label{ex:ssik}
\gll Haksayng-tul-i chayk-ul twu-kwen-ssik sa-ss-ta.\\
     student-\gl{pl}-\gl{nom} book-\gl{acc} two-\gl{cl}-\gl{ssik} buy-\gl{pst}-\gl{decl}\\
\glt `The students bought two books each.'
\end{exe}

These are distance-distributive markers (the anti-quantifiers of \citealt{Choe1987}; \citealt{Gil1995}; \citealt{Zimmermann2002}). The marker sits on the share, the thing distributed, while the sorting key, the thing distributed over, is elsewhere: a plural argument in (\ref{ex:zutu-a}) and (\ref{ex:ssik}), and a series of events or occasions in (\ref{ex:zutu-b}). Force is pronounced, then, but neither on the restriction it ranges over nor on a head of the spine; Choe's analysis posits a covert distributive operator relating key and share, which I place on the spine, with the marker as its dependent exponent; \citet{Zimmermann2002} puts the operator inside the share DP, and nothing here turns on the difference. Figure~\ref{fig:keyshare} shows the three-way split: the key in one place, the share marked in another, the operator on the spine.

\begin{keiofig}{fig:keyshare}{Key, share and a covert operator on the spine, for (\ref{ex:zutu-a}). The tree is schematic: on \posscitet{Zimmermann2002} analysis the operator is inside the share DP.}
\SOtree{[DistP, SOtree, xslim, flat
  [{DP \quad \ghost{\scriptsize sorting key}}, draw=gtaccent, rounded corners=2pt [{\LI{kodomo-tati-ga}}, SOlexical]]
  [Dist$'$
    [{Dist \quad \ghost{\scriptsize covert}}, SOhead]
    [\textit{v}P
      [{DP \quad \ghost{\scriptsize share, marked}}, draw=gtaccent, rounded corners=2pt [{\LI{ringo-o}}, SOlexical] [{\LI{mittu-\hl{zutu}}}, SOlexical]]
      [V, SOhead [{\LI{tabeta}}, SOlexical]]]]]}
\end{keiofig}

(\ref{ex:zutu-b}) is worth pausing on, because its key is not an argument at all: it is the clause's own event, which is to say that distribution, here, is quantification over occasions. That is what connects \textit{-zutu} to the adverbs of \S\ref{sec:unselective}.

The same configuration turns up in an English contact variety. \citet{YeoTsoulas2013} analysed Singapore English \textit{all} as a distributive element next to the spine, associating leftward with a plural restriction, that is, with the syntax of \textit{dou} and English lexical material. It suggests that the \textit{dou} configuration is a stable point in the space of grammars, reachable under contact.

\subsection{The spine}\label{sec:spine}

Putting these together gives a spine of positions at which force can be fixed:

\begin{exe}
\ex \label{ex:spine} C $>$ Foc $>$ Neg $>$ Dist $>$ Asp $>$ \textit{v}
\end{exe}

\begin{center}\small
\begin{tabularx}{\linewidth}{@{}P{1.0cm}LL@{}}
\toprule
Head & Exponents & What is fixed there\\
\midrule
C & \textit{ka}, \textit{no} (Japanese questions) & the force of a question; its scope site\\
Foc & \textit{dake} `only', \textit{sae} `even', \textit{mo} `also'; \textit{-man} `only' (Korean) & the association of a focus particle with its alternatives\\
Neg & \textit{-nai} (Japanese); \textit{an}, \textit{-ci anh-} (Korean) & the licensing of a polarity item\\
Dist & \textit{dou} (Mandarin), \textit{all} (Singapore English); the share marked by \textit{-zutu}, \textit{-ssik} & distributive force\\
Asp & \gl{gen} (habitual and generic) & generic force\\
\textit{v} & existential closure \citep{Diesing1992keio} & existential force\\
\bottomrule
\end{tabularx}
\end{center}

The claim is that each head in (\ref{ex:spine}) but the lowest is pronounced in some language; whether existential closure at \textit{v} ever is, is the first question of \S\ref{sec:readings}. The spine, then, is a series of places where force can be fixed. Lecture~3 relates it to the box in two ways. The phase heads on it, \textit{v} and C, are the sites at which a boxed quantifier is discharged, and so the sites at which its scope is fixed; the other heads relate to the box, where they do, by consultation rather than discharge, as Neg does when it licenses a polarity item, which is why that relation is bounded by the clause (\Hs{3}{sec:wedge}). It is not a spine of scope positions in the sense of \citet{BeghelliStowell1997}, with a dedicated landing site for each type of quantifier. What it keeps from their proposal is the claim that the positions are there in every language, pronounced or not (\S\ref{sec:hierarchy}); what it drops is the movement into them.

\subsection{Interim synthesis}\label{sec:synthesis}

Across \S\S\ref{sec:indeterminates}--\ref{sec:clause} one configuration recurs:

\begin{keybox}
\begin{exe}
\ex \label{ex:config} [ \ldots{} Op\textsubscript{force} \ldots{} [ \ldots{} restriction \ldots{} ] \ldots{} ]: force on the clausal spine, restriction in the nominal, and an association between them, subject to locality.
\end{exe}
\end{keybox}

The languages differ only in what they pronounce, and where:

\begin{center}\footnotesize
\begin{tabularx}{\linewidth}{@{}P{2.6cm}P{2.4cm}P{2.8cm}LL@{}}
\toprule
Type & Restriction & Force & Association & What is pronounced\\
\midrule
English, D-quantification & NP & D (\textit{every}) & sisterhood & force, but not the scope site (QR)\\
English, A-quantification & indefinite, bare plural & adverb, \gl{gen} & mapping and focus & the adverb, when there is one; \gl{gen} never\\
Straits Salish (Jelinek) & pronominal argument & adverbial or predicative Q & predication & force, on the spine\\
St'\'at'imcets (Matthewson) & DP & Q over DP & sisterhood to DP & both; D non-quantificational\\
Japanese, Korean & indeterminate & particle (\textit{mo}, \textit{ka}; \textit{-to}, \textit{-na}, \textit{-nka}) & alternatives, or Agree & both\\
Mandarin & \textit{mei}-N, \textit{wh}, plural & \textit{dou} (Dist) & leftward association & both; the operator on the spine\\
\bottomrule
\end{tabularx}
\end{center}

The generalised quantifier picture of Lecture~1 survives as one row of this table, not as the theory of quantification. A general theory has instead to answer, for every language, the three questions of (\ref{ex:problems}): how the restriction is assembled, how binding works, and where force is assigned. The answers, I suggest, pattern parametrically, and \S\ref{sec:hierarchy} tries to say how. Before that, the two localities.

% =====================================================================
\section{Two localities, previewed}\label{sec:localities}

This section states the divide that the argument of Lecture~3 opens between a syntactic and a semantic store. The data of this lecture show two kinds of locality, and they measure different things. Domain locality (\S\ref{sec:domain}) counts edges: does a clause or phase boundary intervene? Path locality (\S\ref{sec:path}) counts interveners: does an operator of a certain kind stand on the path? Let me give the conclusion first, so that the data can be read as evidence for or against it: which of the two governs a dependency tells a syntactic store from a semantic one, and \S\ref{sec:wedge} says what follows if the two come apart. The two-by-two at the end of the section is the design of the judgement sheet.

\subsection{Domain locality}\label{sec:domain}

\begin{wide}
\begin{exe}
\ex \label{ex:npi-jp}
\begin{xlist}
\ex \label{ex:npi-jp-a}
\gll Taroo-wa [Hanako-ga nani-mo tabe-nakat-ta to] itta.\\
     Taro-\gl{top} Hanako-\gl{nom} what-\gl{mo} eat-\gl{neg}-\gl{pst} \gl{comp} say.\gl{pst}\\
\glt `Taro said that Hanako ate nothing.'
\ex \label{ex:npi-jp-b}
\gll \llap{*}Taroo-wa [Hanako-ga nani-mo tabeta to] iwa-nakat-ta.\\
     Taro-\gl{top} Hanako-\gl{nom} what-\gl{mo} eat.\gl{pst} \gl{comp} say-\gl{neg}-\gl{pst}\\
\glt (`Taro didn't say that Hanako ate anything.') \src{(\citealt{Muraki1978}; \citealt{Kato1985})}
\end{xlist}
\ex \label{ex:korean-npi}
\gll \llap{*}Chelswu-nun [Yenghi-ka amwu-to manna-ss-ta-ko] sayngkakha-ci anh-ass-ta.\\
     Chelswu-\gl{top} Yenghi-\gl{nom} any-\gl{to} meet-\gl{pst}-\gl{decl}-\gl{comp} think-\gl{ci} \gl{neg}-\gl{pst}-\gl{decl}\\
\glt (`Chelswu didn't think that Yenghi met anyone.') \src{(after \citealt{GilTsoulas2013})}
\end{exe}
\end{wide}

NPI licensing by negation is clause-bound: the polarity item and its negation must be clausemates (\citealt{Muraki1978}; \citealt{Kato1985}; for Korean, \citealt{Sohn1995} and \citealt{ChungPark1998}). Nothing stands between the two in (\ref{ex:npi-jp-b}) or (\ref{ex:korean-npi}) except a clause boundary, and a clause boundary is enough. The matrix verbs are chosen for that reason: a bridge verb, `say' or `think', puts no focus operator on the path, so that the only thing between the polarity item and its negation is the boundary itself. That is what makes these the clean cell of the two-by-two in \S\ref{sec:wedge}.

Call this \emph{domain locality}: what matters is whether a boundary intervenes, not what else does (Figure~\ref{fig:domain}). On the analysis of \S\ref{sec:conjunctive} it is the locality of Agree between D[uNeg] and Neg; on the box analysis of Lecture~3 it will be the locality of $\theta$-discharge, since a boxed element must be discharged at the phase of its own clause, or the derivation crashes. \citet{Shimoyama2011} gives a semantic statement of the same clause-boundedness for the Japanese indeterminate polarity items, in terms of the scope they must take with respect to their clausemate negation. Her statement is accurate, and it does not close the divide: a semantic store can state the clausemate condition, but it has to state it, as a condition added to the store, whereas a syntactic store inherits it from the phase, and that difference between stating and inheriting is exactly what (\Hx{3}{ex:locality-claim}) is about. Two syntactic theories, then, and one prediction: the dependency fails at the edge, regardless of what else is in the way. That is the signature of a store whose discipline comes from phases.

\begin{keiofig}{fig:domain}{Domain locality: licensing fails at a clause boundary, whatever else is, or is not, in the way.}
\begin{tikzpicture}[every node/.style={font=\small}, >=Stealth]
  \node[draw=gtaccent, rounded corners=2pt, inner sep=4pt] (neg) at (0,0) {Neg};
  \node[draw=gtghost, dashed, rounded corners=2pt, inner sep=6pt, minimum width=4.4cm, minimum height=1.1cm] (cp) at (4.4,0) {};
  \node[text=gtghost, font=\scriptsize] at (4.4,0.75) {clause boundary};
  \node[draw=gtaccent, rounded corners=2pt, inner sep=4pt] (npi) at (4.4,0) {\textit{nani-mo} / \textit{amwu-to}};
  \draw[->, gtalert, line width=0.8pt] (neg) -- node[midway, above, font=\scriptsize, text=gtalert] {* licensing} (cp.west);
\end{tikzpicture}
\end{keiofig}

\subsection{Path locality}\label{sec:path}

\begin{wide}
\begin{exe}
\ex \label{ex:path-jp}
\begin{xlist}
\ex[?*]{\label{ex:sika-a}
\gll Hanako-sika nani-o yoma-nakat-ta no?\\
     Hanako-only what-\gl{acc} read-\gl{neg}-\gl{pst} \gl{q}\\
\glt (`What did only Hanako read?')}
\ex[]{\label{ex:sika-b}
\gll Nani-o Hanako-sika yoma-nakat-ta no?\\
     what-\gl{acc} Hanako-only read-\gl{neg}-\gl{pst} \gl{q}\\
\glt `What did only Hanako read?' \src{(\citealt{Hoji1985}; \citealt{Takahashi1990}; \citealt{Tomioka2007})}}
\end{xlist}
\ex \label{ex:path-disj}
\begin{xlist}
\ex[?*]{\label{ex:disj-a}
\gll [John ka Bill]-ga nani-o katta no?\\
     John or Bill-\gl{nom} what-\gl{acc} buy.\gl{pst} \gl{q}\\
\glt (`What did John or Bill buy?')}
\ex[]{\label{ex:disj-b}
\gll Nani-o [John ka Bill]-ga katta no?\\
     what-\gl{acc} John or Bill-\gl{nom} buy.\gl{pst} \gl{q}\\
\glt `What did John or Bill buy?' \src{(\citealt[264]{Hoji1985}; \citealt{Hagstrom1998})}}
\end{xlist}
\ex \label{ex:path-kor}
\begin{xlist}
\ex[*]{\label{ex:man-a}
\gll Minswu-man nwukwu-lul po-ass-ni?\\
     Minswu-only who-\gl{acc} see-\gl{pst}-\gl{q}\\
\glt (`Who did only Minswu see?')}
\ex[]{\label{ex:man-b}
\gll Nwukwu-lul Minswu-man po-ass-ni?\\
     who-\gl{acc} Minswu-only see-\gl{pst}-\gl{q}\\
\glt `Who did only Minswu see?' \src{(\citealt{BeckKim1997}; \citealt{Beck2006})}}
\end{xlist}
\end{exe}
\end{wide}

An in-situ \textit{wh}-phrase cannot be c-commanded by a focus-sensitive element such as \textit{-sika} or \textit{-man} `only', or by a disjunction: the question is degraded, and scrambling the \textit{wh}-phrase above the intervener rescues it. Here no clause boundary is involved at all. The disjunction datum in (\ref{ex:path-disj}) goes back to \citet[264]{Hoji1985}; \citet[ch.~2]{Hagstrom1998} gives it an analysis in terms of the movement of the Q-particle \textit{ka}, which on his account moves to the periphery of the clause (not necessarily to C), and is blocked by an intervening \textit{ka} `or'.

Call this \emph{path locality}: what matters is what stands on the path between the \textit{wh}-phrase and its question particle, not how many boundaries the path crosses (Figure~\ref{fig:path}). On \posscitet{Beck2006} analysis, interveners are focus operators that evaluate the alternatives passing through them, so that the alternatives introduced by the \textit{wh}-phrase never reach Q. Intervention, on this view, is what happens when a semantic store meets another operator that reads alternatives: the alternatives are consumed by the intervener (the squiggle of the excursus on focus in \S\ref{sec:diesing}), and Q finds a singleton. Scrambling the \textit{wh}-phrase above the intervener simply takes it off the path. \citet{Tomioka2007} offers an information-structural account of the same facts. Either way, (\ref{ex:path-jp})--(\ref{ex:path-kor}) are what a semantic store predicts.

\pIIneedspace{3.5cm}
\begin{keiofig*}{fig:path}{Path locality: the alternatives of the \textit{wh}-phrase are evaluated by an intervening focus operator, and never reach Q.}
\begin{tikzpicture}[every node/.style={font=\small}, >=Stealth]
  \node[draw=gtaccent, rounded corners=2pt, inner sep=4pt] (q) at (0,0) {Q: \textit{no} / \textit{ni}};
  \node[draw=gtalert, rounded corners=2pt, inner sep=4pt, fill=gtalert!10] (int) at (6.2,0) {\textit{-sika}, \textit{-man}, \textit{ka} `or'};
  \node[draw=gtaccent, rounded corners=2pt, inner sep=4pt] (wh) at (11.2,0) {\textit{nani}, \textit{nwukwu}};
  \draw[->, gtaccent!60, line width=0.8pt] (wh) -- node[midway, above, font=\scriptsize, text=gtghost] {alternatives} (int);
  \draw[->, gtalert, line width=0.8pt, dashed] (int) -- node[midway, above, align=center, font=\scriptsize, text=gtalert] {* evaluated here;\\never reach Q} (q);
\end{tikzpicture}
\end{keiofig*}

\subsection{The divide}\label{sec:wedge}

\begin{keybox}
\begin{exe}
\ex \label{ex:wedge}
\begin{xlist}
\ex \label{ex:domain-loc} \textit{Domain locality} counts edges: a dependency fails if a clause or phase boundary intervenes, whatever else does.
\ex \label{ex:path-loc} \textit{Path locality} counts interveners: a dependency fails if an operator of a certain kind intervenes, however many boundaries there are.
\end{xlist}
\end{exe}
\end{keybox}

A syntactic store whose discipline is inherited from phases predicts (\ref{ex:domain-loc}); a semantic store of alternatives predicts (\ref{ex:path-loc}). The two measures cut across each other, which gives a two-by-two: same clause or across a clause boundary, with or without a focus intervener on the path.

\begin{center}\small
\renewcommand{\arraystretch}{1.3}
\begin{tabularx}{\linewidth}{@{}P{2.6cm}>{\centering\arraybackslash}X>{\centering\arraybackslash}X@{}}
\toprule
 & \textbf{No focus intervener} & \textbf{Focus intervener on the path}\\
\midrule
\textbf{Same clause} & syntactic \jok\quad semantic \jok\newline{\footnotesize\ghost{both predict good}} & syntactic \jok\quad semantic *\newline{\footnotesize\ghost{they diverge: the open test}}\\
\textbf{Across a clause boundary} & syntactic *\quad semantic \jok\newline{\footnotesize\ghost{they diverge: observed * for NPIs, (\ref{ex:npi-jp-b}), (\ref{ex:korean-npi})}} & syntactic *\quad semantic *\newline{\footnotesize\ghost{both predict bad}}\\
\bottomrule
\end{tabularx}
\end{center}

Two cells decide. The first is the cell across a clause boundary with no intervener: an embedded indeterminate with its conjunction, a matrix negation, a bridge verb and nothing else. For NPI licensing it is observed to fail, as in (\ref{ex:npi-jp-b}) and (\ref{ex:korean-npi}), and that matches the syntactic store, which measures domains, rather than the semantic one. The second is the cell within a single clause, with a focus intervener on the path. That is the decisive open cell: if the dependency survives the intervener, the store is syntactic; if it degrades, in the manner of the intervention effects that \citet{Linebarger1987} found for English polarity items, it is semantic.

If the two kinds of locality come apart in a single language, with a single family of particles, then there is a syntactic reflex that a semantic store cannot capture, and the box survives as an explanation rather than a notation. If they coincide, then, for this domain at least, a store in the syntax explains nothing that a store in the semantics does not. One caution: clause-crossing and the presence of an intervener have to be varied independently, or a domain effect can masquerade as a path effect. The judgement sheet in the appendix has items in each of the two cells that decide, for the association of \textit{ka} and \textit{mo} with an indeterminate and for the licensing of \textit{nani-mo} separately; Lecture~3 (\Hs{3}{sec:keio2x2}) takes those two cells, crosses them with the two relations, association and licensing, and enters the judgements of the literature.

Let me float the conjecture I will defend in Lecture~3, so that readers can hold it against their own judgements as they go:

\begin{keybox}
\noindent\textit{One particle, two localities.} Japanese seems to show both kinds of locality, and \textit{mo} shows both on its own: its association with an indeterminate crosses an island, as in (\ref{ex:shimoyama}), just as the association of \textit{ka} shows path effects, as in (\ref{ex:path-jp})--(\ref{ex:path-disj}), while the licensing of \textit{nani-mo} by negation shows domain effects, as in (\ref{ex:npi-jp}).
\end{keybox}

If that is right, it points not to the box everywhere but to a division of labour between a syntactic and a semantic store, of the kind the Greek emphatic and non-emphatic n-words suggested in a single paradigm \citep{Giannakidou1998,Giannakidou2000}. It is the first of the three conjectures of Lecture~3 (\Hx{3}{ex:conj1}), and I hedge it accordingly: Japanese \emph{seems} to show both. The judgement sheet (\exref{sheet}) has items in each cell, and Lecture~3 returns to them.

% =====================================================================
\section{The hierarchy}\label{sec:hierarchy}

The variation of this lecture and of \Hp{1} can now be gathered into a single structure, and its format is \posscitet{Roberts2019}: a parameter hierarchy. Parameters, on this view, are not independent switches. They are organised into hierarchies, in which each question arises only given a particular answer to the one above it, and in which the options further down are the more marked ones, adopted by a learner only on more specific evidence (Figure~\ref{fig:hierarchy-notion}). Three things follow from the format, and they matter for what comes next. A lower question presupposes an answer to the one above it: there is no asking `is Q on a particle?' of a language that has no Q in the nominal. The ordering is itself a hypothesis about acquisition: the learner meets the higher question first, and takes a lower branch only on more specific evidence. And the leaves need not be languages. A language can sit on two leaves at once, if two of its items answer the questions differently; that is not a defect of the format but a prediction about the items, and it is the reason Japanese and Korean will appear twice below.

\begin{keiofig}{fig:hierarchy-notion}{The format of a parameter hierarchy, after \citet{Roberts2019}: each question arises only given an answer to the one above it, and markedness increases downwards.}
\begin{tikzpicture}[
  font=\footnotesize,
  q/.style={draw=gttext!55, rounded corners=2pt, inner sep=3pt},
  leaf/.style={draw=gtaccent!60, rounded corners=2pt, inner sep=3pt, fill=gtexbg},
  arr/.style={-{Stealth[length=1.6mm]}, semithick, gttext!55},
  lab/.style={font=\scriptsize\itshape, text=gttext!65, inner sep=1pt}
]
  \node[q] (q1) at (0,0) {Question 1?};
  \node[leaf] (a1) at (-1.9,-1.0) {unmarked answer};
  \node[q] (q2) at (0.7,-1.0) {Question 2?};
  \node[leaf] (a2) at (-1.2,-2.0) {less marked};
  \node[q] (q3) at (1.4,-2.0) {Question 3?};
  \node[text=gtghost] at (1.4,-2.7) {$\vdots$};
  \draw[arr] (q1) -- node[lab, above left] {no} (a1);
  \draw[arr] (q1) -- node[lab, right] {yes} (q2);
  \draw[arr] (q2) -- node[lab, above left] {no} (a2);
  \draw[arr] (q2) -- node[lab, right] {yes} (q3);
  \draw[-{Stealth[length=1.6mm]}, gtghost] (3.4,0.1) -- node[lab, right, align=left, text=gtghost] {more\\ marked} (3.4,-2.6);
\end{tikzpicture}
\end{keiofig}

The hierarchy takes the form of five questions, four about the nominal and one about the clause:\toexercise{hierarchy}

\begin{exe}
\ex \label{ex:parameters}
\begin{tabular}[t]{@{}lp{0.80\linewidth}@{}}
P1 & Is [Q] lexicalised inside the nominal extended projection?\\
P2 & If so, is it on a particle, or is it a nominal quantifier?\\
P3 & For a nominal quantifier, is there a general classifier layer between Q and N?\\
P4 & Do D and Q co-occur?\\
P5 & (clausal) Is force fixed on the clausal spine, and is that step pronounced?
\end{tabular}
\end{exe}

P1--P4 are the nominal half, built from \Hs{1}{sec:nominal}; P5 is the clausal half, built from \S\S\ref{sec:unselective}--\ref{sec:clause} above. Two questions test any arrangement of the answers. Do \textit{subete-no} N and \textit{dono} N-\textit{mo} land on different branches, as Korean \textit{motun} and \textit{nwukwu} do? And where do Mandarin \textit{mei} \ldots{} \textit{dou} and Straits Salish go? P3 is asked because a pronounced classifier layer is where division is realised (\Hs{1}{sec:Qlayer}), and a language that pronounces division there has a reason to leave force to the spine, as Mandarin does with \textit{dou} (\citealt{Cheng2009}: \textit{mei}-N supplies a maximal plurality, \textit{dou} the force).

\pIIneedspace{17cm}
\subsection*{Building the nominal half}

Figure~\ref{fig:hierarchy} builds the nominal half one question at a time, and then asks the clausal question of every leaf. At P1, Straits Salish and Mohawk, if Jelinek and Baker are right, answer no: the nominal has no Q, and force is a matter for the clause alone. Every other language answers yes, and meets P2. At P2, given a Q in the nominal, the question is whether it is a particle or a nominal quantifier: Japanese \textit{dono} N-\textit{mo} and \textit{dare-ka}, Korean \textit{nwukwu-to}, \textit{-na} and \textit{-nka} take the particle branch, and that is the leaf this seminar is about. At P3, for a nominal quantifier, the question is whether there is a classifier layer between Q and N: Mandarin \textit{mei-ge} N answers yes, and its universal force comes from \textit{dou} on the spine, which P5 records. At P4, Greek \textit{o kathe} N and English \textit{all the} N answer yes: D and Q together, Q above D in English \textit{all the} N and D above Q in Greek \textit{o kathe} N (\Hx{1}{ex:statimcets} and the figure beside it). The last leaf is Q alone, syncretic with D: English \textit{every} N, Greek \textit{kathe} N, Japanese \textit{subete-no} N, Korean \textit{motun} N. The generalised quantifier of \Hp{1} is one leaf among five, and the most marked.

\begin{keiofig*}{fig:hierarchy}{One possible answer: four nominal parameters and one clausal one, after \citet{Roberts2019}. The leaves list items, not languages; the highlighted leaf is the one this seminar is about. P5 is asked of every leaf of the nominal tree.}
\begin{tikzpicture}[
  font=\small,
  node distance=6mm and 13mm,
  q/.style={draw=gttext!55, rounded corners=3pt, text width=5.0cm, inner sep=4pt, align=left},
  leaf/.style={draw=gttext!55, rounded corners=3pt, text width=5.0cm, inner sep=4pt, align=left},
  hlf/.style={leaf, draw=gtaccent, fill=gtaccent!6, thick},
  arr/.style={-{Stealth[length=2mm]}, semithick, gttext!65},
  lab/.style={font=\footnotesize\itshape, text=gttext!65, inner sep=2pt}
]
\node[q] (p1) {\textbf{P1}\\ Is [Q] lexicalised inside the nominal?};
\node[q, below=of p1] (p2) {\textbf{P2}\\ On a particle, or a nominal quantifier?};
\node[q, below=of p2] (p3) {\textbf{P3}\\ A general classifier layer between Q and N?};
\node[q, below=of p3] (p4) {\textbf{P4}\\ Do D and Q co-occur?};
\node[q, below=of p4] (l5) {\textbf{Q alone, syncretic with D}\\ English \textit{every} N; Greek \textit{kathe} N\\ Japanese \textit{subete-no} N; Korean \textit{motun} N};
\node[leaf, left=of p1] (l1) {\textbf{Force only on the clause}\\ Straits Salish \citep{Jelinek1995}\\ Mohawk \citep{Baker1995}};
\node[hlf, left=of p2] (l2) {\textbf{Force on a particle}\\ Japanese \textit{dono} N-\textit{mo}, \textit{dare-ka}\\ Korean \textit{nwukwu-to}, \textit{-na}, \textit{-nka}};
\node[leaf, left=of p3] (l3) {\textbf{Q with a classifier}\\ Mandarin \textit{mei-ge} N (force from \textit{dou})};
\node[leaf, left=of p4] (l4) {\textbf{D and Q together}\\ Greek \textit{o kathe} N; English \textit{all the} N};
\draw[arr] (p1) -- node[lab, above] {no} (l1);
\draw[arr] (p2) -- node[lab, above] {particle} (l2);
\draw[arr] (p3) -- node[lab, above] {yes} (l3);
\draw[arr] (p4) -- node[lab, above] {yes} (l4);
\draw[arr] (p1) -- node[lab, right] {yes} (p2);
\draw[arr] (p2) -- node[lab, right] {nominal quantifier} (p3);
\draw[arr] (p3) -- node[lab, right] {no} (p4);
\draw[arr] (p4) -- node[lab, right] {no} (l5);
\coordinate (mid) at ($(l1.west)!0.5!(p1.east)$);
\node[draw=gttext!55, fill=gtexbg, rounded corners=3pt, text width=11.0cm, inner sep=5pt, align=left, anchor=north]
  (p5) at ([yshift=-9mm] mid |- l5.south)
  {\textbf{P5 (clausal): is force fixed on the spine, and is that step pronounced?}\\
   Yes: Mandarin \textit{dou}; Japanese \textit{-zutu}, \textit{-te-mo}; Korean \textit{-ssik}; Straits Salish; Mohawk\\
   No, the spine is used but left silent: English, Greek};
\draw[arr] (l5.south) -- node[lab, right] {then, for every leaf} (l5.south |- p5.north);
\end{tikzpicture}
\end{keiofig*}

P5 is asked for every leaf of the nominal tree: a language's items are placed on the nominal branches, and then the language is asked whether the clausal step is pronounced. The items listed under `yes' in the figure are those whose exponent sits on the spine itself; the particles \textit{mo}, \textit{ka}, \textit{-to} and \textit{-nka} count as well, since they spell out an operator whose force is fixed on the clause (\S\ref{sec:distance}), although they are pronounced on the nominal, and that is why Japanese and Korean answer `yes' in the table that closes this lecture. Mandarin therefore sits on two levels at once, and Figure~\ref{fig:hierarchy-languages} draws the consequence for all seven languages: the tree laid flat, with each language linked to every leaf one of its items occupies. The lines cross, and they are meant to. A language is a set of answers, not one answer, and P5 is a property of the language rather than of a leaf, which is why it is shown there by shading rather than by a branch.

\begin{keiofig*}{fig:hierarchy-languages}{The whole tree, with the languages placed. Each language is linked to the leaves its items occupy; Japanese and Korean sit on two leaves, and so do Greek and English. Shaded languages answer `yes' to P5: force is fixed on the spine and the step is pronounced (\textit{dou}; \textit{-zutu}, \textit{-te-mo}; \textit{-ssik}; the Salish and Mohawk predicates). Greek and English use the spine but leave it silent. The dashed line is Straits Salish at P4, if Matthewson's result for St'\'at'imcets carries over to it.}
\begin{tikzpicture}[
  font=\footnotesize,
  q/.style={draw=gttext!55, rounded corners=2pt, text width=2.6cm, inner sep=3pt, align=left},
  leaf/.style={draw=gtaccent!60, rounded corners=2pt, text width=2.6cm, inner sep=3pt, align=left, fill=gtexbg},
  hlf/.style={leaf, draw=gtaccent, thick},
  lang/.style={draw=gttext!55, rounded corners=2pt, inner sep=3.5pt, minimum height=5.5mm},
  yes/.style={fill=gtaccent!20},
  arr/.style={-{Stealth[length=1.6mm]}, semithick, gttext!55},
  lk/.style={gtaccent!75, line width=0.6pt, looseness=0.45},
  lab/.style={font=\scriptsize\itshape, text=gttext!65, inner sep=1pt, fill=gtbg}
]
  \node[q] (p1) at (0,0)    {\textbf{P1}\ \ Q lexicalised inside the nominal?};
  \node[q] (p2) at (3.2,0)  {\textbf{P2}\ \ On a particle, or a nominal quantifier?};
  \node[q] (p3) at (6.4,0)  {\textbf{P3}\ \ A classifier layer between Q and N?};
  \node[q] (p4) at (9.6,0)  {\textbf{P4}\ \ Do D and Q co-occur?};
  \node[leaf] (l1) at (0,-2.1)    {\textbf{Force only on the clause}};
  \node[hlf]  (l2) at (3.2,-2.1)  {\textbf{Force on a particle}\\ \textit{dono} N-\textit{mo}, \textit{dare-ka}; \textit{nwukwu-to}, \textit{-na}, \textit{-nka}};
  \node[leaf] (l3) at (6.4,-2.1)  {\textbf{Q with a classifier}\\ \textit{mei-ge} N (force from \textit{dou})};
  \node[leaf] (l4) at (9.6,-2.1)  {\textbf{D and Q together}\\ \textit{o kathe} N; \textit{all the} N};
  \node[leaf] (l5) at (12.8,-2.1) {\textbf{Q alone, syncretic with D}\\ \textit{every} N; \textit{kathe} N; \textit{subete-no} N; \textit{motun} N};
  \draw[arr] (p1) -- node[lab, above=19pt] {yes} (p2);
  \draw[arr] (p2) -- node[lab, above=19pt] {nominal quantifier} (p3);
  \draw[arr] (p3) -- node[lab, above=19pt] {no} (p4);
  \draw[arr] (p1) -- node[lab, left] {no} (l1);
  \draw[arr] (p2) -- node[lab, left] {particle} (l2);
  \draw[arr] (p3) -- node[lab, left] {yes} (l3);
  \draw[arr] (p4) -- node[lab, left] {yes} (l4);
  \draw[arr] (p4.east) -- node[lab, above, sloped] {no} (l5.north);
  \node[lang, yes] (sal) at (-0.2,-4.3) {Straits Salish};
  \node[lang, yes] (moh) at (2.0,-4.3)  {Mohawk};
  \node[lang, yes] (ja)  at (4.1,-4.3)  {Japanese};
  \node[lang, yes] (ko)  at (5.9,-4.3)  {Korean};
  \node[lang, yes] (zh)  at (7.9,-4.3)  {Mandarin};
  \node[lang] (gr) at (10.3,-4.3) {Greek};
  \node[lang] (en) at (12.4,-4.3) {English};
  \draw[lk] (sal.north) to[out=90, in=-90] (l1.south);
  \draw[lk, dashed, gtghost] (sal.north) to[out=90, in=-90] (l4.south);
  \draw[lk] (moh.north) to[out=90, in=-90] (l1.south);
  \draw[lk] (ja.north) to[out=90, in=-90] (l2.south);
  \draw[lk] (ja.north) to[out=90, in=-90] (l5.south);
  \draw[lk] (ko.north) to[out=90, in=-90] (l2.south);
  \draw[lk] (ko.north) to[out=90, in=-90] (l5.south);
  \draw[lk] (zh.north) to[out=90, in=-90] (l3.south);
  \draw[lk] (gr.north) to[out=90, in=-90] (l4.south);
  \draw[lk] (gr.north) to[out=90, in=-90] (l5.south);
  \draw[lk] (en.north) to[out=90, in=-90] (l4.south);
  \draw[lk] (en.north) to[out=90, in=-90] (l5.south);
\end{tikzpicture}
\end{keiofig*}

Three things about this tree are worth saying, and they are where a tree built from \exref{hierarchy} and mine are most likely to differ; the differences deserve to be taken seriously.

First, the leaves list items, not languages. Japanese and Korean each sit on two branches: \textit{dono} N-\textit{mo} and \textit{nwukwu-to} on the particle branch, \textit{subete-no} N and \textit{motun} N with the nominal quantifiers. That makes the hierarchy testable within a single language rather than only across languages. If it is right, the items on the particle branch should have their scope pinned to the particle, and should show path effects in their association with the indeterminate, while the items on the other branch should behave like English \textit{every}, with whatever flexibility the language allows. Mineshima's contrast (\S\ref{sec:pins}; \Hs{1}{sec:jp-universals}) is the first test, and the judgement sheet has two more.

Second, Mandarin sits on two levels at once. \textit{Mei-ge} N is a nominal quantifier with a classifier (P3), but its universal force comes from \textit{dou} on the spine (P5). That is the configuration of (\ref{ex:config}) in its purest form: the nominal supplies a maximal restriction, the spine supplies the force. Straits Salish, if Jelinek is right, answers `no' to P1 and is decided entirely at P5; if Matthewson's result for St'\'at'imcets carries over to it, it belongs instead with Greek and English \textit{all the} at P4. Both placements are defensible, and the fact that the hierarchy can host the disagreement is, I think, a point in its favour. The Salish placement is a genuinely open question, and one on which others know the Salish literature better than I do.

Third, P5 is where English is most unlike the rest. English answers `no', but not because it lacks the positions of (\ref{ex:spine}); Beghelli and Stowell's whole argument was that it has them. It answers `no' because it does not pronounce them. That answers the inversion of \S\ref{sec:aims}: English uses the clausal spine as much as Japanese or Mandarin does, and simply leaves it silent. Lecture~3 argues that this is why Japanese scope is rigid; it does not follow, and Lecture~3 does not claim, that silence is what makes English scope flexible.

Let me end with the question that, in Lecture~4, I will suggest organises the whole hierarchy:

\begin{keybox}
\noindent At which step of the derivation, storage or retrieval, is force assigned, and is that step pronounced?\par\smallskip
{\small
\begin{tabularx}{\linewidth}{@{}P{3.0cm}LP{2.4cm}@{}}
\toprule
 & Force assigned at & Pronounced?\\
\midrule
Japanese, Korean & retrieval, in the particle & yes\\
Mandarin & retrieval, on the spine (\textit{dou}) & yes\\
English & on the stored item itself (\textit{every}) & neither step\\
\bottomrule
\end{tabularx}}
\end{keybox}

Japanese and Korean pronounce retrieval, in the particle; Mandarin pronounces it on the spine; English pronounces neither, and puts force on the stored item itself. If that is the right question, the hierarchy is not a list of facts about determiners but a map of which steps of one derivation the languages of the world choose to pronounce (\Hs{4}{sec:revisited}).

% =====================================================================
\section{Questions and readings}\label{sec:readings}

Of the questions below, the sixth is the one on which Lecture~3 depends, and the second is the one that speakers of Japanese are best placed to answer.

\begin{questions}
\begin{enumerate}
\item Is existential closure at \textit{v} ever pronounced? If not, why is it the one silent head of (\ref{ex:spine})?
\item What does case do to force? Compare \textit{dare-mo} with \textit{dare-mo-ga}, and \textit{amwu-to} with its case-marked counterparts.
\item Is the \textit{-te-mo} unconditional a case of unselective binding or of Hamblin alternatives, and is there anything that could tell the two apart?
\item If \textit{mo} is `even', what fixes the universal reading of \textit{dono} N-\textit{mo}? Does \textit{nan-do-mo} show that force is computed at the point of retrieval?
\item Where do Mandarin \textit{mei} \ldots{} \textit{dou} and Straits Salish go in the hierarchy of \S\ref{sec:hierarchy}, and where would a language you know best go?
\item Is the clausemate condition on NPIs a syntactic locality (Agree, phases) or a semantic one (the strength of the licensor)?
\end{enumerate}
\end{questions}

\paragraph{Key readings.} \citet{KratzerShimoyama2002}, the Hamblin treatment; \citet{Shimoyama2006}, its application to islands; \citet{Diesing1992keio}, the Mapping Hypothesis.

\paragraph{Further readings.} \citet{Kuroda1965}; \citet{Lewis1975}; \citet{Heim1982}; \citet{Nishigauchi1990}; \citet{Haspelmath1997}; \citet{Jelinek1995}; \citet{Baker1995}; \citet{Matthewson2001}; \citet{Kratzer2005}; \citet{Cheng1995,Cheng2009}; \citet{GillHarlowTsoulas2004}; \citet{GilTsoulas2013}; \citet{Beck2006}; \citet{Szabolcsi2015}; \citet{Roberts2019}.


% ======================================================================
%  PART 3
% ======================================================================
\keiolecture{3}{Scope in the box theory}{Scope in the box theory}%
  {Day 2, Lecture 3 \textperiodcentered{} Sunday 8 November 2026, 13:30--15:30}
\keioepigraph{Same mechanism, different explanation.}{}
\lecturemap

% =====================================================================
\section{Aims and plan}\label{sec:aims}

This lecture has three aims. The first is to give the box theory of QR in full: what the box is, how a quantifier gets into it and out of it, and how a grammar with a box gets covert scope without covert movement and a variable without an index. The second is to make the comparison with Cooper storage explicit, so that those who use storage every day can tell me whether I have simply reinvented it. I should say at once that the first, honest answer to that question is yes, in mechanism; the interest lies in what the mechanism is made to explain. The third is to put the result to the Japanese data of Lectures 1 and 2, and to derive scope rigidity, and the scope of scrambled phrases, from the order and the locality of discharge.

The question of the lecture is whether scope can come from storage, and whether the store is syntactic. Let me turn it around first. The usual question is whether a grammar needs covert movement in order to have inverse scope. The inverse question is, I think, more revealing: why should a grammar that has every ingredient of inverse scope ever lack it? Japanese has quantifiers, a clausal spine, phases and a store (if I am right), and yet \Hx{1}{ex:dareka-dono}, \textit{Dare-ka-ga dono hon-mo yonda}, cannot mean what \Hx{1}{ex:someone}, \textit{Someone read every book}, can. Whatever the box is, it has to explain that, and \S\ref{sec:rigidity} is where the explanation is cashed: a grammar lacks inverse scope when the order in which its quantifiers are discharged is not chosen but read off an order the derivation already has.

\subsection*{What this lecture takes from Lectures 1 and 2}

Three results from the first two lectures carry the argument of this one.

\begin{exe}
\ex \label{ex:recap}
\begin{xlist}
\ex \textit{The desideratum} (\Hx{1}{ex:desideratum}). Quantificational force must be separable from the restriction; the grammar must be able to fix force and scope away from the restriction, and to relate the two without adding anything to the workspace.
\ex \textit{The spine} (\Hx{2}{ex:spine}). The clausal spine offers a series of positions at which force can be fixed, C $>$ Foc $>$ Neg $>$ Dist $>$ Asp $>$ \textit{v}, all but the lowest of them pronounced in some language.
\ex \textit{Two localities} (\Hx{2}{ex:wedge}). Domain locality counts edges and is what a syntactic store predicts; path locality counts interveners and is what a semantic store predicts.
\end{xlist}
\end{exe}

The three results are the premises of the argument of this lecture, and they divide the work between them. The box is the theory that is meant to satisfy the first, to find on the second the heads that discharge it and the heads that consult it, and to be tested by the third. Nothing else from the first two lectures is presupposed, apart from the vocabulary of \Hs{1}{sec:architecture}: MERGE in a workspace, Minimal Yield, Determinacy, Form Copy and phases.

The plan is as follows. \S\ref{sec:wh-to-QR} goes from \textit{wh} to QR, by way of our earlier work on scope assignment. \S\ref{sec:cooper} sets out Cooper storage properly. \S\ref{sec:box} gives Chomsky's box at the level of ideas, and \S\ref{sec:boxQR} the box theory of QR, with two derivations in full and four objections. \S\ref{sec:worksheet} puts the proposal to work on a harder case, and lists the problems the derivation raises. \S\ref{sec:korean} turns to Korean and Japanese, which pronounce the step of retrieval. \S\ref{sec:box-cooper} asks whether the box is Cooper storage, and \S\ref{sec:wedge} tests the question on the 2\,$\times$\,2 from Lecture~2. \S\ref{sec:rigidity} derives rigidity and the scope of scrambling, and \S\ref{sec:conjectures} states three conjectures and the problems that remain. Along the way, excursions take up points that matter to readers from particular traditions (Cooper's own formulation, Keller's nesting, the CCG saxophonists, Pesetsky's \textit{ittai}, Bo\v{s}kovi\'c and Takahashi's lowering, and sketches of antecedent-contained deletion and inverse linking); the argument does not depend on them.

% =====================================================================
\section{From \textit{wh} to QR}\label{sec:wh-to-QR}

\subsection{LF movement and the copy theory do not mix}\label{sec:LF-copy}

This section has two aims: to recall why, in a grammar without a level of LF, scope cannot be assigned by covert movement, and to show how far one can get instead with overt movement and a choice of which copy to pronounce. It is the route Norman Yeo and I took in \citet{TsoulasYeo2017}, and with one more step it leads to the box. The paper's last section already tried to extend the idea to QR, and ran into precisely the problem of the lower copy that the box is meant to solve; that is why it is the point of departure.

The argument is short, and it is the single-output model of \Hs{1}{sec:cyclicity} applied. If there is no level of LF, there is no stage of the derivation after Spell-Out at which covert movement could apply; and if movement leaves copies, then what used to be called covert movement is simply overt movement with the lower copy pronounced. So we pursued the idea that all movement is overt, and that the difference between \textit{wh}-movement and \textit{wh}-in-situ lies in which copy is spelled out:

\begin{exe}
\ex \label{ex:tripartite-Q} \textit{The tripartite structure of questions} \citep{TsoulasYeo2017}: a clause-typing interrogative C, a Q particle, and a \textit{wh}-phrase. Q and the \textit{wh}-phrase start out together, as [Q \textit{wh}]; Q fixes scope; and the copies created by Internal Merge are spelled out selectively, so that Q and the \textit{wh}-phrase may be pronounced in different positions.
\end{exe}

Japanese supplies the clearest case. On \posscitet{Hagstrom1998} analysis the question particle starts out next to the \textit{wh}-phrase and moves to the clause periphery, where it is pronounced as \textit{ka} or \textit{no}:

\begin{exe}
\ex \label{ex:hagstrom}
\gll Taroo-ga nani-o katta no?\\
     Taro-\gl{nom} what-\gl{acc} buy.\gl{pst} \gl{q}\\
\glt `What did Taro buy?'\\[2pt]
\lbr{CP} \lbr{TP} Taroo-ga [ $\langle$Q$\rangle$ nani-o ] katta ] Q ] \src{Q pronounced at C, the \textit{wh}-phrase in situ}
\end{exe}

Hagstrom's Q-movement is the empirical anchor here: the particle is pronounced at the top, the \textit{wh}-phrase at the bottom, and the two are related by movement of Q alone. On this view a Japanese question is derived by exactly the same movement as an English one, and the two differ only in what is pronounced where. English pronounces the \textit{wh}-phrase at the top and Q nowhere; Japanese pronounces Q at the top and the \textit{wh}-phrase at the bottom. That is the externalisation hypothesis of \Hx{1}{ex:variation} at work, and Figure~\ref{fig:two-externalisations} draws it: one derivation, two externalisations, with the copies that are not pronounced drawn in grey.

\begin{keiofig*}{fig:two-externalisations}{One derivation, two externalisations. Left: Japanese pronounces Q at the top (\textit{no}) and the \textit{wh}-phrase at the bottom. Right: English pronounces the \textit{wh}-phrase at the top and Q nowhere. Grey, dashed nodes are copies that are not pronounced.}
\begin{minipage}[t]{0.47\linewidth}\centering
{\footnotesize Japanese: Q pronounced at the top, \textit{wh} at the bottom}\par\medskip
\begin{forest}
[CP, SOtree, xslim
  [TP [{\LI{Taroo-ga}}, SOlexical] [VP [DP [{$\langle$Q$\rangle$}, SOtrace] [{\LI{nani-o}}, SOlexical]] [{\LI{katta}}, SOlexical]]]
  [{C\quad\hl{\LI{no}}}, SOhead, SOmoved]]
\end{forest}
\end{minipage}\hfill
\begin{minipage}[t]{0.49\linewidth}\centering
{\footnotesize English: \textit{wh} pronounced at the top, Q nowhere}\par\medskip
\begin{forest}
[CP, SOtree, xslim
  [{\hl{\LI{what}}\quad\ghost{$\langle$Q$\rangle$}}, SOmoved]
  [C$'$ [C, SOhead [{\LI{did}}, SOlexical]] [TP [{\LI{Taro}}, SOlexical] [VP [{\LI{buy}}, SOlexical] [{$\langle$Q what$\rangle$}, SOtrace]]]]]
\end{forest}
\end{minipage}
\end{keiofig*}

\citet{Pesetsky1987} had already shown that the in-situ strategy is not uniform: \textit{wh}-phrases modified by \textit{ittai} `the hell' are island-sensitive in a way that bare in-situ \textit{wh}-phrases are not, which is what one expects if the bare ones are associated with Q without movement and the \textit{ittai} ones must move.

\begin{excursus*}[Pesetsky's \textit{ittai}: in situ is not uniform]
The contrast Pesetsky drew is between two ways of relating an in-situ \textit{wh}-phrase to the question particle within one language. A bare \textit{wh}-phrase inside a complex NP island can associate with a matrix \textit{no}; the same phrase modified by \textit{ittai} `the hell' resists it. The pair he discusses, with Nishigauchi's Japanese, is of this kind:
\begin{exe}
\ex \label{ex:ittai}
\begin{xlist}
\ex[]{\label{ex:ittai-a}
\gll Mary-wa [John-ni nani-o ageta hito]-ni atta no?\\
     Mary-\gl{top} John-\gl{dat} what-\gl{acc} give.\gl{pst} person-\gl{dat} meet.\gl{pst} \gl{q}\\
\glt `What did Mary meet the person who gave to John?'}
\ex[?*]{\label{ex:ittai-b}
\gll Mary-wa [John-ni ittai nani-o ageta hito]-ni atta no?\\
     Mary-\gl{top} John-\gl{dat} the.hell what-\gl{acc} give.\gl{pst} person-\gl{dat} meet.\gl{pst} \gl{q}\\
\glt (`What the hell did Mary meet the person who gave to John?') \src{(after \citealt{Pesetsky1987}; \citealt{Nishigauchi1990})}}
\end{xlist}
\end{exe}
The point does not depend on the details of the pair. Japanese shows two association mechanisms side by side: association by consultation, which involves no movement and is island-insensitive, and association by movement, which is island-sensitive. The box's consultation (\S\ref{sec:box}) is the first of these, made syntactic.
\end{excursus*}

In that last section we extended the idea, more speculatively, to QR: a quantifier would carry a silent Q, and its scope would be the position at which that Q is spelled out, while the quantifier itself is pronounced low. That was as far as we got in 2017, and it leaves one problem untouched: the lower copy. Selective spell-out decides which copy is pronounced; it does not tell the semantics how to interpret the copy that is not interpreted as the operator. For a \textit{wh}-phrase the problem can perhaps be finessed, since its lower copy is an indefinite, a set of alternatives. For a quantifier it cannot, and we are back with Trace Conversion and \Hs{1}{sec:copy}.

It is worth noting that scope without QR has a history at Keio. \citet{Kitahara1992,Kitahara1996} argued that the scope of quantifiers can be derived from movement that the grammar needs independently, for feature checking, rather than from a dedicated rule of Quantifier Raising (compare \citealt{Hornstein1995}); in the 1996 paper, scope is read off the chains that feature checking creates. The difference from what follows is exactly the lower copy. On that account there is one, and it is interpreted; on mine there is none: the box goes one step further, and there is no lower copy at all.

\subsection{The step to the box}\label{sec:step}

What we needed in 2017, I now think, was a way of relating a quantifier to two positions, one where it is pronounced and one where it takes scope, without either position being occupied by a copy that the semantics then has to convert. That is precisely what a store does: the element is set aside, a placeholder stands in its position, and the element is applied later.

Stated that way, the step from selective spell-out to a store is small, and it is worth seeing exactly what it changes. Selective spell-out keeps two copies and chooses which to pronounce; the problem was the copy it does not pronounce. A store keeps one occurrence of the quantifier, outside the structure, and puts in its place not a copy but a placeholder, something that was never a quantifier and so never needs converting. The three questions a syntactic store must then answer are the three that organise the rest of the lecture: what the placeholder is, if it is not an indexed variable (\S\ref{sec:boxQR}: the verb's own $\theta$-feature); what sets the quantifier aside and brings it back, if the only operation is MERGE (\S\ref{sec:boxQR}: Self-Merge, and External Merge at a phase edge); and where the locality of the whole arrangement comes from, if it is not stipulated (\S\S\ref{sec:box-cooper}--\ref{sec:wedge}: from phases). Before turning to a syntactic store, let us look carefully at the semantic one, because the measure of the syntactic store is what it does that the semantic one cannot.

% =====================================================================
\section{Cooper storage, properly}\label{sec:cooper}

The intuition behind Cooper storage \citep{Cooper1975,Cooper1983} is simple. Semantic interpretation proceeds bottom-up. When it meets a quantifier it does not try to interpret it in place; it sets it aside in a store and leaves a placeholder, a variable, in its position. When enough structure has been built for the quantifier to be interpreted, that is, at a node of type $t$, the quantifier can be retrieved and applied to the meaning built so far, abstracting over the placeholder. Cooper regarded scope as a semantic matter: there is one syntactic structure, and the ambiguity lives in the order of retrieval. Figure~\ref{fig:cooper-drawn} is the whole mechanism: set aside, leave a variable, apply later.

\begin{keiofig*}{fig:cooper-drawn}{Cooper storage, drawn. The quantifier goes into the store under an index; a placeholder bearing the same index stands in its position; at a node of type $t$ the quantifier is retrieved and abstracts over the placeholder.}
\begin{tikzpicture}[every node/.style={font=\footnotesize\rmfamily}, >=Stealth]
  \node[draw=gtaccent, rounded corners=2pt, inner sep=3pt] (qp) at (0,0) {\textit{every book}};
  \node[draw=gtaccent, rounded corners=2pt, inner sep=3pt, fill=gtexbg] (store) at (3.9,1.0) {store: $\{\langle\textit{every book}, 1\rangle\}$};
  \node[draw=gtghost, rounded corners=2pt, inner sep=3pt] (ph) at (3.9,0) {placeholder $x_1$ in object position};
  \node[draw=gtaccent, rounded corners=2pt, inner sep=3pt, text width=3.6cm, align=center] (t) at (8.9,0) {node of type $t$: retrieve, abstract over $x_1$};
  \draw[->, gtaccent, line width=0.6pt] (qp) -- node[above, sloped, font=\scriptsize, text=gtghost] {storage} (store.west);
  \draw[->, gtghost, line width=0.6pt] (qp) -- (ph);
  \draw[->, gtaccent, line width=0.6pt] (store.east) to[bend left=15] node[above right, font=\scriptsize, text=gtghost] {retrieval} (t.north);
\end{tikzpicture}
\end{keiofig*}

Each node is assigned a pair, a meaning and a store:

\begin{exe}
\ex \label{ex:storage} \textit{Storage:}\quad $\langle Q,\ S\rangle \ \Rightarrow\ \langle \lambda P.\,P(x_i),\ S \cup \{\langle Q, i\rangle\}\rangle$ \quad ($i$ fresh)
\ex \label{ex:retrieval} \textit{Retrieval} ($\varphi$ of type $t$):\quad $\langle \varphi,\ S \cup \{\langle Q, i\rangle\}\rangle \ \Rightarrow\ \langle Q(\lambda x_i.\,\varphi),\ S\rangle$
\end{exe}

By (\ref{ex:storage}), the quantifier $Q$ goes into the store under a fresh index $i$, and in its place stands the lifted placeholder $\lambda P.\,P(x_i)$ (or simply $x_i$, of type $e$). By (\ref{ex:retrieval}), at a node of type $t$, $Q$ is taken out of the store, the meaning built so far is abstracted over $x_i$, and $Q$ is applied to the result. Note that the index is written twice, once when the quantifier is stored and once when it is retrieved; that will be the first of the three features to which I return below. (This is a simplified modern statement; \citet{Cooper1983} works with assignment sequences and NP-intensions. The excursus that follows gives his own formulation, for those who want it.)

\begin{excursus}[Cooper's own formulation]
In \citet{Cooper1983} the store is defined over assignment sequences, and it is in the sequences that the index lives. The storage rule maps the standard value of a noun phrase to a pair:

\begin{exe}
\ex \label{ex:cooper-own-storage} \textit{NP-storage:}\quad $\langle\mspace{-3mu}\langle \text{NP}\rangle\mspace{-3mu}\rangle \ \to\ \langle\, \sigma \mapsto \{X \subseteq E : \sigma_i \in X\},\ \text{NP}\dbr{i} \,\rangle$, for any $i$
\end{exe}

The first member is the content that enters composition in the NP's place: relative to an assignment $\sigma$, the set of sets that contain the $i$-th member of $\sigma$, that is, an individual lifted to a generalised quantifier. The second member is the NP's own meaning, put on hold in the store under the index $i$. Retrieval selects the $k$-th stored quantifier and applies it to the meaning of the sentence, leaving the rest in the store to be retrieved higher:

\begin{exe}
\ex \label{ex:cooper-own-retrieval} \textit{NP-retrieval:}\quad $\langle\mspace{-3mu}\langle \text{S}\rangle\mspace{-3mu}\rangle \ \to\ \langle\, \text{S}'_k(\text{S}^c_0),\ \text{S}^c_1,\ldots,\text{S}^c_{k-1},\text{S}^c_{k+1},\ldots,\text{S}^c_n \,\rangle$, for any $k$, $0 < k \le n$
\end{exe}

The truth definition then says that the sentence is true at a world $w$ under $\sigma$ if and only if the set of individuals $a$ that make $\varphi$ true when assigned to the index $i$ belongs to the NP-intension at $w$. (\ref{ex:storage})--(\ref{ex:retrieval}) are the same thing with the sequences suppressed. The formulation is worth seeing once, if only because it shows where the index is in Cooper's own system: not in the syntax, which has none, but in the assignment sequences over which the semantic values are defined.
\end{excursus}

Here is how the system derives the two readings of \Hx{1}{ex:someone}, \textit{Someone read every book}:

\begin{exe}
\ex \label{ex:cooper-deriv}
\begin{minipage}[t]{\linewidth}\small%
\begin{tabularx}{\linewidth}[t]{@{}P{2.3cm}L@{}}
\textit{every book} & $\langle \lambda P.\,P(x_1),\ \{\langle \textit{every book}, 1\rangle\}\rangle$\\
\textit{read every book} & $\langle \lambda y.\,\mathrm{read}(y, x_1),\ \{\langle \textit{every book}, 1\rangle\}\rangle$\\
\textit{someone read every book} & $\langle \mathrm{read}(x_2, x_1),$\newline $\{\langle \textit{every book}, 1\rangle, \langle \textit{someone}, 2\rangle\}\rangle$
\end{tabularx}
\end{minipage}
\end{exe}

\begin{center}\small
\begin{tabularx}{\linewidth}{@{}P{3.6cm}Lc@{}}
\toprule
Order of retrieval & Result & Reading\\
\midrule
\textit{someone} first, then \textit{every book} & $\textit{every book}(\lambda x_1.\ \textit{someone}(\lambda x_2.\ \mathrm{read}(x_2, x_1)))$ & $\forall > \exists$\\
\textit{every book} first, then \textit{someone} & $\textit{someone}(\lambda x_2.\ \textit{every book}(\lambda x_1.\ \mathrm{read}(x_2, x_1)))$ & $\exists > \forall$\\
\bottomrule
\end{tabularx}
\end{center}

The quantifier retrieved last takes widest scope. Storage solves both of QR's problems without moving anything. The type mismatch never arises, because the object position is filled by a placeholder of type $e$ (or its lift), and inverse scope is simply the order of retrieval. The two results in the table are exactly the two that the box will deliver in \S\ref{sec:deriv1}, from the order of discharge rather than the order of retrieval. Three features of the system will matter when we compare it with the box.

First, it needs an index. The placeholder is $x_i$, and retrieval abstracts over exactly that index. Storage relocates the variable of \Hs{1}{sec:problem} into the semantics; it does not dispense with it.

Second, its locality is not built in. In the basic system retrieval may apply at any node of type $t$, so nothing prevents a quantifier stored in an embedded clause from being retrieved in the matrix clause, and the clause-boundedness of QR has to be imposed by hand.

Third, it needs a constraint on nesting. Consider \textit{a representative of every city}. If \textit{every city} is stored inside the object and the containing nominal is then stored as well, nothing in (\ref{ex:storage})--(\ref{ex:retrieval}) prevents the container from being retrieved in a position where the inner quantifier can no longer bind the placeholder it left, so that a variable ends up free. \citet{Keller1988} proposed nested stores: a stored nominal carries its own store with it, and its internal quantifiers are retrieved in a way that keeps them bound. It works, but it is a further stipulation on the store.

These three features will reappear as rows of the comparison in \S\ref{sec:box-cooper}. The index row hides the deepest difference between the store and the box; the rows for locality and nesting are where the box claims to explain what the store stipulates.

\begin{excursus*}[The nesting problem, drawn]
\begin{minipage}[t]{0.50\linewidth}\centering
{\small\textit{a representative of every city}}\par\medskip
\begin{forest}
[{DP\quad\scriptsize $\langle \lambda P.P(x_2),\ \{\langle\text{DP},2\rangle\}\rangle$}, SOtree, xslim, draw=gtaccent, rounded corners=2pt
  [{\LI{a representative of}}, SOlexical]
  [{QP\quad\scriptsize $\langle \lambda P.P(x_1),\ \{\langle\textit{every city},1\rangle\}\rangle$}, draw=gtaccent, rounded corners=2pt [{\LI{every city}}, SOlexical]]]
\end{forest}
\end{minipage}\hfill
\begin{minipage}[t]{0.46\linewidth}
\setlength\parskip{0.45\baselineskip}
\textit{Every city} is stored first, under the index 1, leaving $x_1$ inside the nominal headed by \textit{representative}. Then the whole nominal, which contains $x_1$, is stored too, under the index 2.

If the outer nominal is retrieved before \textit{every city} has been, its meaning, with $x_1$ free inside it, is placed under a $\lambda x_2$ that does not bind $x_1$. The later retrieval of \textit{every city} then abstracts over an $x_1$ that is no longer in its scope, and the result has a free variable.
\end{minipage}

\medskip
Keller's repair is to let a stored nominal carry its own store with it, so that its internal quantifiers are retrieved when it is, and in an order that keeps them bound. Nested stores are, in effect, a stack rather than a set. The nesting problem is also the problem of inverse linking, which returns as item~\ref{open:inverse} of the open list in \S\ref{sec:open}; the HPSG stores are Keller's, which is one reason to take it seriously here. The box's version of the question is whether phase structure gives nesting for free (\S\S\ref{sec:box-cooper}, \ref{sec:open}).
\end{excursus*}

Cooper storage, in one form or another, has been the usual way of scoping a quantifier without moving it. HPSG kept quantifiers in a store feature of the sign and retrieved them at designated points \citep{PollardSag1994}, and LFG's glue semantics reaches the same results through resource-sensitive deduction, with no store at all \citep{Dalrymple2001}; both were set out in the excursions of \Hp{1}. CCG goes another way. \citet{Steedman2012} gets the scope of universals from the derivation itself, since an object can combine with a derived \textit{S/NP} built by type-raising and composition, and treats indefinites as generalised Skolem terms. The derivation then predicts which readings exist, rather than generating all of them and filtering, as in the classic case of \textit{Every boy admires, and every girl detests, some saxophonist}, which has two readings, not four. That is locality and nesting derived rather than stipulated, and it is the standard the box has to meet. The methodological point is simple: a syntactic store is worth having only if its locality comes from the syntax.

\begin{excursus}[Why CCG gets two readings, not four]
Take \textit{Every boy admires, and every girl detests, some saxophonist}. On the first reading there is one saxophonist, admired by every boy and detested by every girl. On the second the saxophonists vary, with the boys and with the girls. What is not available is a mixed reading: one saxophonist for the boys but varying ones for the girls, or the reverse.

In CCG the coordinated constituent \textit{every boy admires, and every girl detests} is a single \textit{S/NP}, built by type-raising and composition, and \textit{some saxophonist} combines with it once, as one argument. The indefinite is a Skolem term, specified once for the whole coordination: either as a constant, which gives the wide reading, or as a function of the variables bound by both universals, which gives the narrow one. There is no derivation on which it is specified differently for the two conjuncts. A store with free retrieval, by contrast, generates four readings and must filter two out, and so does QR with across-the-board movement. The combinatorics derive the restriction, and that is the sense in which CCG explains what storage stipulates. It is also the sense in which the box must do the same from phases. For categorial grammarians this is the best card, and I would rather play it for them than have it played against me; the point of convergence is taken up in \S\ref{sec:box-cooper}.
\end{excursus}

So here is the measure for the rest of the lecture. Cooper storage does what QR did without movement; it does it in the semantics; it needs an index; and its locality is not its own. Of the box I will ask three questions: does it need an index, where does its locality come from, and is it syntactic? \S\ref{sec:box-cooper} answers them, row by row.

% =====================================================================
\section{The box}\label{sec:box}

This section sets out Chomsky's proposal at the level of ideas, as I understand it, without claiming to reproduce its details. Keio is where the proposal was first aired, and those who heard the Keio--EMU lectures in 2023 will know better than I do where my paraphrase slips.

\subsection{Chomsky's proposal}\label{sec:box-proposal}

\citet{Chomsky2024}, first aired in the 2023 Keio--EMU lectures given with Seely and Kitahara, proposes that an element internally merged to the edge of a phase is placed in a box, written \BX{X}. Four ideas, informally:

\begin{exe}
\ex \label{ex:box-ideas}
\begin{xlist}
\ex \textit{Storage.} The element is set aside.
\ex \textit{Segregation.} It is no longer part of the ongoing derivation. It is not at a phase edge and not in a scope position; it has no position at all.
\ex \textit{Accessibility.} It remains available: phase heads at later cycles consult it for instructions, without merging it again. There is no movement.
\ex \textit{Retrieval.} It is discharged later, and that is where its scope is fixed and its externalisation determined.
\end{xlist}
\end{exe}

The motivating case is A-bar movement. A \textit{wh}-phrase is merged in its $\theta$-position, internally merged to the edge of the lowest phase, and boxed there. From then on there is no successive-cyclic movement through intermediate edges, since every later phase head, and C[Q] in particular, can consult the box. The intermediate positions disappear, and with them the copies they would leave and the labelling problems they would raise; what reconstruction and successive-cyclic morphology show becomes a matter of consultation rather than of movement. Figure~\ref{fig:motivating} draws the picture the published applications assume: boxing at the first edge, consultation thereafter, retrieval at C[Q]. The dashed arrows are the point. Nothing is merged at the intermediate phases.

\begin{keiofig}{fig:motivating}{A-bar movement without successive cyclicity: boxed at the first edge, consulted (dashed), discharged at C[Q].}
\begin{tikzpicture}[every node/.style={font=\footnotesize\rmfamily}, >=Stealth]
  \node[draw=gtaccent, rounded corners=2pt, inner sep=3pt] (vp) at (0,0) {\vs P phase: \textit{wh} at the edge};
  \node[draw=gtghost, dashed, rounded corners=2pt, inner sep=3pt] (cp1) at (0,1.2) {embedded CP phase};
  \node[draw=gtghost, dashed, rounded corners=2pt, inner sep=3pt] (vp2) at (0,2.4) {matrix \vs P phase};
  \node[draw=gtaccent, rounded corners=2pt, inner sep=3pt] (cq) at (0,3.6) {matrix C[Q]: consults, retrieves};
  \node[draw=gtaccent, thick, rounded corners=2pt, inner sep=4pt, fill=gtexbg] (box) at (3.6,1.8) {\BX{wh}};
  \draw[->, gtaccent, line width=0.7pt] (vp.east) to[bend right=15] node[right, font=\scriptsize, text=gtghost] {boxed here} (box.south);
  \draw[->, gtghost, line width=0.5pt, dashed] (cp1.east) -- node[pos=0.45, sloped, above, font=\scriptsize, text=gtghost] {consult} (box.west);
  \draw[->, gtghost, line width=0.5pt, dashed] (vp2.east) -- node[pos=0.45, sloped, above, font=\scriptsize, text=gtghost] {consult} (box.west);
  \draw[->, gtaccent, line width=0.7pt] (box.north) to[bend right=15] node[right, font=\scriptsize, text=gtghost] {discharged here} (cq.east);
\end{tikzpicture}
\end{keiofig}

The published applications treat \textit{wh}-movement, across-the-board extraction \citep{Lee2024} and parasitic gaps; the literature also applies the box to \textit{that}-trace effects, remnant movement and parsing. As far as I have been able to find, nobody has yet treated QR, quantifier scope or Cooper storage in these terms, though the search could not be made exhaustive and the claim stays hedged. The nearest thing is on the formal side, and it is older than the box: \citet{Kobele2006,Kobele2018} interprets the movement of minimalist grammars by Cooper storage. That is storage as the semantics of movement, where what follows is storage in the place of movement; \Hs{4}{sec:nb-delay} takes up the comparison. There are also computational implementations: Ginsburg's has a Box operation \citep{Ginsburg2025}, and \posscitet{FongOishi2024} parser for the Strong Minimalist Thesis uses box notation.

\subsection{What needs to be said}\label{sec:needs}

For a box theory of QR, two questions need answers that the published proposal does not, as far as I can tell, supply:

\begin{exe}
\ex \label{ex:box-questions}
\begin{xlist}
\ex \label{ex:box-q-a} What does it mean to be put in a box?
\ex \label{ex:box-q-b} What operation puts it there?
\end{xlist}
\end{exe}

My answer to (\ref{ex:box-q-a}) is that a boxed item's interpretation (its structural position for the purposes of scope) and its locus of externalisation are both determined at a later stage of the derivation. The box is, in that sense, a store. My answer to (\ref{ex:box-q-b}) is more contentious. If MERGE is the only operation there is, boxing must be an instance of MERGE. The suggestion is that boxing is one of the outcomes of Self-Merge, MERGE($\alpha$, $\alpha$), an operation proposed in various places and for various purposes \citep{Guimaraes2000,Kayne2005,Adger2013,Tsoulas2017}. Most of those who discuss it take it to create the first projection line, so to speak. It is reasonable to suggest that boxing is one of its outcomes too, and \S\ref{sec:objections} says exactly what Self-Merge yields and why it does not violate Minimal Yield (\Hx{1}{ex:RR}). (\ref{ex:box-q-b}) is where readers will press, and rightly.

Note the difference from Chomsky's version, and let me name it before anyone else does. In his, the element is first merged in its $\theta$-position and boxed at the phase edge, and there is no independent operation that puts an element in a box \citep[204]{Kitahara2024}. In mine, for QR, the quantifier is boxed as soon as it is built, by an operation that does exactly that, and it is never merged in its $\theta$-position at all; the $\theta$-position is filled by something else, namely the verb's own $\theta$-feature. That is the step that yields a variable without an index, and it is the step readers should press hardest.

\begin{center}\small
\begin{tabularx}{\linewidth}{@{}P{2.6cm}LL@{}}
\toprule
 & Chomsky's box & The box of this lecture, for QR\\
\midrule
What is boxed & an element internally merged to a phase edge & a quantified phrase bearing [uQ], as soon as it is built\\
By what & no independent operation & Self-Merge, MERGE($\alpha$, $\alpha$) (\S\ref{sec:objections})\\
The $\theta$-position & filled by the element, by External Merge & filled by the verb's own $\theta$-feature, \uth\ (\S\ref{sec:proposal})\\
Applications & \textit{wh}, ATB, parasitic gaps & QR, and through it rigidity and scrambling\\
\bottomrule
\end{tabularx}
\end{center}

% =====================================================================
\section{The box theory of QR}\label{sec:boxQR}

This section states the proposal in five clauses, derives one English and one Japanese sentence with it, and then goes back over the clauses most likely to be challenged.

\subsection{The five clauses}\label{sec:proposal}

\begin{keybox}
\begin{exe}
\ex \label{ex:proposal}
\begin{xlist}
\ex \label{ex:box} \textit{Box.} A quantified phrase in an internal-argument relation, bearing an unvalued quantificational feature, [uQ], is boxed by Self-Merge. MERGE(QP, QP, WS) yields the singleton \{QP\}, written \bx{QP}, which remains a member of the workspace, outside the main syntactic object.
\ex \label{ex:fission} \textit{Fission.} The verb is split by two applications of MERGE. Self-Merge, MERGE(V, V, WS), yields the singleton \{V\}; Internal Merge of V with it yields \{V, \{V\}\}, and Form Copy construes the two occurrences of V as copies. The singleton \{V\}, in the internal argument position, is a token of the verb in which only the internal $\theta$-feature is visible. Call it \uth. It is of type $e$. It is the variable, and nothing has been added to the workspace to make it.
\ex \label{ex:discharge} \textit{Discharge.} At a phase, if the phase head's Minimal Search finds an undischarged \uth\ in its complement, the boxed QP is externally merged at the edge and discharges it; if none has done so by the time the last phase of the clause is complete, the derivation crashes. Discharge fixes scope: the QP takes scope over the complement of the phase head.
\ex \label{ex:case} \textit{Case.} A phase head with an unvalued Case feature consults the box and values Case on the boxed QP. This relation is anchored to the \uth\ in the phase head's domain, and it fixes where the QP will be pronounced.
\ex \label{ex:externalisation} \textit{Externalisation.} A boxed element is pronounced at the highest position at which it enters a relation visible to the sensorimotor interface (Case, or an edge feature). Discharge is not such a relation.
\end{xlist}
\end{exe}
\end{keybox}
The three clauses (\ref{ex:box})--(\ref{ex:discharge}) are stated here in their final form; the reasons for each, and for fission the formulation it replaces, are given in \S\ref{sec:objections}.

The five clauses divide into two pairs and a hinge. (\ref{ex:box})--(\ref{ex:fission}) build the two things that QR needed, an operator set aside and a variable in its place; (\ref{ex:discharge}) relates them, and is at once the theory of scope and the theory of clause-boundedness; (\ref{ex:case})--(\ref{ex:externalisation}) decide where the quantifier is heard. The whole account of covertness is in the contrast between two relations, one visible to the sensorimotor interface and one not:

\begin{center}\small
\begin{tabular}{@{}lp{4.9cm}c@{}}
\toprule
 & Fixed by & Visible to SM?\\
\midrule
Scope & discharge at a phase edge (\ref{ex:discharge}) & no\\
Pronunciation & Case (\ref{ex:case}), or an edge feature (\ref{ex:externalisation}) & yes\\
\bottomrule
\end{tabular}
\end{center}

(\ref{ex:externalisation}) generalises the formulation I used at the International Circle of Korean Linguistics meeting in Sheffield in July (ICKL), on which a chain containing a $\theta$-position and a Case position is externalised at the Case position. The generalisation is needed for long-distance scrambling (\S\ref{sec:scrambling}), and I flag it as new.

Let me take the clauses in turn.

\paragraph{On (\ref{ex:box}): boxing is not optional.} Boxing is forced by [uQ], so nothing is chosen for scope, which answers the worry about Determinacy raised in \Hx{1}{ex:determinacy}. A force-less indefinite lacks [uQ] and is never boxed, which is why indefinites take their exceptional scope by other means (\Hs{1}{sec:indefinites}). In Japanese and Korean, as \S\ref{sec:korean} shows, the feature sits on the particle. Figure~\ref{fig:boxing} shows the step in the notation of the workspace: once the box is filled, the workspace carries it as a subscript.

\begin{keiofig*}{fig:boxing}{Boxing is forced by [uQ]. The quantified phrase is built by External Merge; Self-Merge then yields the singleton \{QP\}, which stays in the workspace, outside the main object.}
\setcounter{minsyn@step}{0}%
\begin{minipage}[t]{0.40\linewidth}\vspace{0pt}
\begin{workspace}[EM]
  \footnotesize\WSeqTree{\{\LI{every},\LI{book}\}_{[u\mathrm{Q}]},\allowbreak\ \ldots}
    {[{QP$_{[u\mathrm{Q}]}$}, SOtree, tiny [Q, SOhead [{\LI{every}}, SOlexical]] [NP [{\LI{book}}, SOlexical]]]}
\end{workspace}
\end{minipage}\hfill
\begin{minipage}[t]{0.13\linewidth}
\vspace{0pt}\vspace*{9ex}\WSright[\SelfMergeIII{QP}]
\end{minipage}\hfill
\begin{minipage}[t]{0.40\linewidth}\vspace{0pt}
\begin{workspacebox}[\BoxOp{QP}]{QP}
  \footnotesize\WSeqBox{QP}{\BX{QP},\allowbreak\ \ldots}
  \par\smallskip\footnotesize MERGE(QP, QP) $= \{\mathrm{QP}\}$: a singleton, in the workspace, outside the main object.
\end{workspacebox}
\end{minipage}
\end{keiofig*}

For the present I restrict boxing to internal arguments. An external argument is externally merged in Spec-\vs P, as \textit{a child} is at W7 of Derivation~1 below, and takes its scope from the positions it occupies; the duality of semantics is then respected in its most literal form for subjects. The two arguments of a transitive verb are thus treated asymmetrically: the subject is merged in its $\theta$-position, the object never is. There is a symmetric alternative, on which \vs\ fissions too, so that the external $\theta$-role also has its token, both arguments may be boxed, and boxed quantifiers are paired with tokens by Case. I list it as an open question (\S\ref{sec:open}, item~\ref{open:symmetric}), since it bears directly on the problems the worksheet raises (\S\ref{sec:five-problems}).

How the grammar knows, at W3 of Derivation~1, that the QP it is boxing will be an internal argument is a fair question, and I would rather state my assumption than hide it. A quantified subject carries [uQ] too. What satisfies it is External Merge in Spec-\vs P: the subject has a $\theta$-position to be merged in, and it takes its scope from that position and from Spec-TP. The object has no such position, because the verb's own token occupies it (\ref{ex:fission}), so for the object boxing is the only way [uQ] can be satisfied. Boxing, on this reading, is a last resort rather than a lexical property, and the restriction is read off what the derivation can do with a phrase, not off the phrase itself. \textit{Dono gakusei-mo sono hon-o yonda} `every student read that book' is then derived like (\ref{ex:kodomo}) below, with the quantified subject in the place of \textit{kodomo} and nothing in the box.

\paragraph{On (\ref{ex:fission}): fission, and the variable without an index.}
The verb carries its internal $\theta$-role as an unvalued feature. Fission is two applications of MERGE, each of which adds exactly one object to the workspace (Figure~\ref{fig:fission}). Self-Merge, MERGE(V, V), yields the singleton \{V\}, V's own projection; Internal Merge of V with that singleton yields \{V, \{V\}\}. Form Copy construes the two occurrences of V as copies of one object (they are identical, one c-commands the other, and they are in the same phase), so nothing new has entered the workspace. What distinguishes the two occurrences is which features each makes visible. The singleton in complement position exposes only \uth: it is the verb's unsaturated argument slot, a token of type $e$ with a lexical source.
Semantically, the token is the open argument of the verb. One way to carry it to TP is function composition, Jacobson's Geach rule \citep{Jacobson1999}, on which the slot stays open through \vs\ and the subject and what reaches TP is the predicate $\lambda x$.\ \textit{a child read} $x$, of type $\langle e,t\rangle$; categorial grammarians will recognise it. That is not the route I take. Nothing is passed up the tree: the phase head finds the token where it is (\S\ref{sec:objections}), and the $\lambda$ of (\ref{ex:CI-B}) below is what discharge contributes. What kind of relation that is, is the question of \Hp{4}.

\begin{keiofig*}{fig:fission}{Fission as two applications of MERGE. Self-Merge on V yields the singleton \{V\} (left); Internal Merge of V with it yields \{V, \{V\}\} (right). Form Copy identifies the two occurrences of V; the singleton in complement position is the token \uth, of type $e$.}
\setcounter{minsyn@step}{0}%
\begin{minipage}[t]{0.36\linewidth}\vspace{0pt}
\begin{workspacebox}[\SelfMergeI\ on V]{QP}
  \footnotesize\WSeqBoxTree{QP}{\BX{QP},\allowbreak\ \Atom{\LI{read}},\allowbreak\ \ldots}
    {[V, SOtree, tiny, SMroot [V, SMcopy, SM1edge [{\LI{read}}, SOlexical]]]}
\end{workspacebox}
\end{minipage}\hfill
\begin{minipage}[t]{0.18\linewidth}
\vspace{0pt}\vspace*{11ex}\WSright[\IntMerge{\text{V}}{\Atom{\text{V}}}]
\end{minipage}\hfill
\begin{minipage}[t]{0.40\linewidth}\vspace{0pt}
\begin{workspacebox}[two occurrences; Form Copy]{QP}
  \footnotesize\WSeqBoxTree{QP}{\BX{QP},\allowbreak\ \{\LI{read},\{\LI{read}\}\},\allowbreak\ \ldots}
    {[VP, SOtree, tiny [V, SOhead [{\LI{read}}, SOlexical]] [{\{V\} $= $ V$_{\uthr}$\quad\ghost{type $e$}}, SMcopy, SOtrace]]}
\end{workspacebox}
\end{minipage}
\end{keiofig*}

\paragraph{On (\ref{ex:discharge}): two phases, two scope sites, and clause-boundedness.} There are two phases in a simple clause, so there are two scope sites, \vs P and CP, which correspond to the two landing sites of QR, \textit{v}P and TP (\Hx{1}{ex:QR-LFs}). Clause-boundedness follows. If \uth\ has not been discharged when the clause's last phase is complete, Transfer sends an unvalued feature to the interfaces and the derivation crashes. Two sites and one condition, then, as in Figure~\ref{fig:two-sites}. \textit{Wh}-phrases are different, because their scope is fixed by consultation from C[Q], and C[Q] can be in a higher clause. What is clause-bound is the discharge of \uth, not the box. A long-distance \textit{wh}-phrase must still have the $\theta$-feature it is related to discharged in its own clause, and its consultation by the matrix C is a separate relation. Two assumptions are packed into that sentence. The first is that a \textit{wh}-phrase bears [uQ] too, so that it boxes like any quantifier and has a \uth\ to discharge; what sets it apart is that its force is fixed by C[Q], which may be in a higher clause, and that in English its highest visible relation is the edge of that C, where it is pronounced, by (\ref{ex:externalisation}). The second is that discharge does not close the box. A boxed element remains consultable after its \uth\ has been discharged, and the box closes at the Transfer of the clause in which the element's highest visible relation is fixed; in Derivations 1 and 2 below the two coincide, which is why their last lines can say that the box is empty.

\begin{keiofig*}{fig:two-sites}{Two phases, two scope sites, one condition. Discharge at the \vs\ phase gives the surface reading of \textit{A child read every book}; discharge at the C phase gives the inverse one; if \uth\ is still undischarged when the clause is transferred, the derivation crashes.}
\begin{tikzpicture}[every node/.style={font=\footnotesize\rmfamily}, >=Stealth]
  \node[draw=gtaccent, rounded corners=2pt, inner sep=4pt] (v) at (0,0) {\vs\ phase: discharge $= \exists > \forall$};
  \node[draw=gtaccent, rounded corners=2pt, inner sep=4pt] (c) at (5.0,0) {C phase: discharge $= \forall > \exists$};
  \node[draw=gtalert, rounded corners=2pt, inner sep=4pt, fill=gtalert!10] (x) at (10.2,0) {Transfer, \uth\ undischarged: crash};
  \draw[->, gtghost] (v) -- (c); \draw[->, gtghost] (c) -- (x);
\end{tikzpicture}
\end{keiofig*}

\paragraph{On (\ref{ex:case}) and (\ref{ex:externalisation}): why QR is `covert'.} \vs\ consults the box exactly as C[Q] does in Chomsky's version, and values Case on the boxed quantifier. That relation is with the position of \uth, and since Case is visible to the sensorimotor interface and discharge is not, the quantifier is pronounced where its variable is, although it was never merged there. Compare A-movement: there the $\theta$-position is low and the Case position high, and the element is pronounced high. In QR the scope position is high and the Case position low, and it is pronounced low.

\begin{center}\small
\begin{tabularx}{\linewidth}{@{}P{2.4cm}LLP{2cm}@{}}
\toprule
 & $\theta$ or scope relation & Case or edge relation & Pronounced\\
\midrule
A-movement & low ($\theta$-position) & high (Spec-TP) & high\\
QR & high (discharge at the edge) & low (Case at \uth) & low\\
\bottomrule
\end{tabularx}
\end{center}

Covertness, then, is not a property of an operation. Nothing moves covertly: A-movement and QR are the same two relations with their heights reversed. The table is the crucial point of the proposal, generalised by (\ref{ex:externalisation}); \S\ref{sec:scrambling} extends it to four cases, and it is also the answer to the fourth of the objections in \S\ref{sec:objections}.

\subsection{Derivation 1: \textit{A child read every book}}\label{sec:deriv1}

\begin{exe}
\ex \label{ex:deriv-child} Workspace states for \textit{A child read every book}:
\end{exe}

\begin{center}\small
\renewcommand{\arraystretch}{1.12}
\begin{tabularx}{\linewidth}{@{}lLP{5.6cm}@{}}
\toprule
 & Workspace & Operation\\
\midrule
W1 & [\,every, book, read, \vs, a, child, T, C\,] & \\
W2 & [\,\{every, book\}, read, \vs, a, child, T, C\,] & MERGE(every, book) = QP\\
W3 & [\,\bx{QP}, read, \vs, a, child, T, C\,] & \textsc{box}: MERGE(QP, QP), forced by [uQ]\\
W4 & [\,\bx{QP}, \{read, \{read\}\uthr\}, \vs, a, child, T, C\,] & \textsc{fission}: MERGE(read, read) $=$ \{read\}; then MERGE(read, \{read\})\\
W5 & [\,\bx{QP}, \{\vs, VP\uthr\}, a, child, T, C\,] & MERGE(\vs, VP)\\
W6 & [\,\bx{QP}, \{\vs, VP\uthr\}, \{a, child\}, T, C\,] & MERGE(a, child) = DP\\
W7 & [\,\bx{QP}, \{DP, \{\vs, VP\uthr\}\}, T, C\,] & MERGE(DP, \vs P): external $\theta$ by EM\\
\multicolumn{3}{@{}p{\linewidth}@{}}{\itshape \vs\ phase: \vs\ consults \bx{QP} and values its Case, fixing its pronunciation at \uthr. Option A: discharge here ($\exists > \forall$). Option B: leave \uth\ undischarged.}\\
W8 & [\,\bx{QP}, \{T, \vs P\uthr\}, C\,] & MERGE(T, \vs P)\\
W9 & [\,\bx{QP}, \{DP, \{T, \vs P\uthr\}\}, C\,] & IM of the subject to Spec-TP\\
W10 & [\,\bx{QP}, \{C, TP\uthr\}\,] & MERGE(C, TP)\\
\multicolumn{3}{@{}p{\linewidth}@{}}{\itshape C phase: C's Minimal Search finds an undischarged \uth\ in TP, so the boxed QP is externally merged at the edge of CP.}\\
W11 & [\,\{QP, \{C, TP\}\}\,] & discharge; the box is empty; Transfer\\
\bottomrule
\end{tabularx}
\end{center}

Two conventions in the table need a word. The subscript on VP\uthr, \vs P\uthr\ and TP\uthr\ marks that the object contains an undischarged \uth; it is a reading convention, not a feature of VP, and nothing percolates (\S\ref{sec:objections}). And from W3 on, the workspace contains a boxed member, which in the figures that follow is written as a subscript on the workspace itself, $W^n_{[\mathrm{QP}]}$, the notation of the figures for a workspace with a box.
W4 compresses the two applications of MERGE that fission consists of (\ref{ex:fission}): Self-Merge yields \{read\}, and Internal Merge of \textit{read} with it yields \{read, \{read\}\}.

Figures~\ref{fig:deriv1-a}--\ref{fig:deriv1-c} give the same derivation in the workspace notation, step by step.

\begin{keiofig*}{fig:deriv1-a}{Derivation 1, W1--W3: the lexical items are selected; \textit{every} and \textit{book} are externally merged; the QP, bearing [uQ], is boxed.}
\setcounter{minsyn@step}{0}%
\begin{minipage}[t]{0.29\linewidth}\vspace{0pt}
\begin{workspace}[Select]
  \footnotesize\WSeq{\LI{every},\allowbreak\LI{book},\allowbreak\LI{read},\allowbreak\vs,\allowbreak\LI{a},\allowbreak\LI{child},\allowbreak\text{T},\allowbreak\text{C}}
\end{workspace}
\end{minipage}\hfill
\begin{minipage}[t]{0.34\linewidth}\vspace{0pt}
\begin{workspace}[\ExtMerge{\LI{every}}{\LI{book}}]
  \footnotesize\WSeqTree{\{\LI{every},\LI{book}\},\allowbreak\LI{read},\allowbreak\vs,\allowbreak\LI{a},\allowbreak\LI{child},\allowbreak\text{T},\allowbreak\text{C}}
    {[{QP$_{[u\mathrm{Q}]}$}, SOtree, tiny [Q, SOhead [{\LI{every}}, SOlexical]] [NP [{\LI{book}}, SOlexical]]]}
\end{workspace}
\end{minipage}\hfill
\begin{minipage}[t]{0.33\linewidth}\vspace{0pt}
\begin{workspacebox}[\BoxOp{QP}]{QP}
  \footnotesize\WSeqBox{QP}{\BX{QP},\allowbreak\LI{read},\allowbreak\vs,\allowbreak\LI{a},\allowbreak\LI{child},\allowbreak\text{T},\allowbreak\text{C}}
  \par\smallskip\raggedright\footnotesize \textsc{box}: MERGE(QP, QP), forced by [uQ]. The singleton stays in the workspace, outside the main object.
\end{workspacebox}
\end{minipage}
\end{keiofig*}

\begin{keiofig*}{fig:deriv1-b}{Derivation 1, W4--W7. W4: fission, by two applications of MERGE; the singleton in complement position exposes only \uth. W5: \vs\ is merged, and \uth\ is in the domain of the phase head. W6--W7: the subject is built and externally merged in its $\theta$-position. The quantifier, meanwhile, is still in the box: it has never been merged anywhere.}
\setcounter{minsyn@step}{3}%
\begin{minipage}[t]{0.48\linewidth}\vspace{0pt}
\begin{workspacebox}[\SelfMergeI; \IntMerge{\text{V}}{\Atom{\text{V}}}]{QP}
  \footnotesize\WSeqBoxTree{QP}{\BX{QP},\allowbreak\{\LI{read},\{\LI{read}\}_{\uthr}\},\allowbreak\vs,\allowbreak\LI{a},\allowbreak\LI{child},\allowbreak\text{T},\allowbreak\text{C}}
    {[VP, SOtree, tiny [V, SOhead [{\LI{read}}, SOlexical]] [{V$_{\uthr}$}, SMcopy, SOtrace]]}
\end{workspacebox}
\end{minipage}\hfill
\begin{minipage}[t]{0.48\linewidth}\vspace{0pt}
\begin{workspacebox}[\ExtMerge{\vs}{\text{VP}}]{QP}
  \footnotesize\WSeqBoxTree{QP}{\BX{QP},\allowbreak\{\vs,\text{VP}_{\uthr}\},\allowbreak\LI{a},\allowbreak\LI{child},\allowbreak\text{T},\allowbreak\text{C}}
    {[\vs P, SOtree, tiny [\vs, SOhead] [VP [V, SOhead [{\LI{read}}, SOlexical]] [{V$_{\uthr}$}, SOtrace]]]}
\end{workspacebox}
\end{minipage}

\medskip
\begin{minipage}[t]{0.40\linewidth}\vspace{0pt}
\begin{workspacebox}[\ExtMerge{\LI{a}}{\LI{child}}]{QP}
  \footnotesize\WSeqBox{QP}{\BX{QP},\allowbreak\{\vs,\text{VP}_{\uthr}\},\allowbreak\{\LI{a},\LI{child}\},\allowbreak\text{T},\allowbreak\text{C}}
\end{workspacebox}
\end{minipage}\hfill
\begin{minipage}[t]{0.56\linewidth}\vspace{0pt}
\begin{workspacebox}[\ExtMerge{\text{DP}}{\vs\text{P}}]{QP}
  \footnotesize\WSeqBoxTree{QP}{\BX{QP},\allowbreak\{\text{DP},\{\vs,\text{VP}_{\uthr}\}\},\allowbreak\text{T},\allowbreak\text{C}}
    {[\vs P, SOtree, tiny [DP [{\LI{a child}}, SOlexical]] [\vs$'$ [\vs, SOhead] [VP [V, SOhead [{\LI{read}}, SOlexical]] [{V$_{\uthr}$}, SOtrace]]]]}
\end{workspacebox}
\end{minipage}
\end{keiofig*}

At W7 the external $\theta$-role is assigned by External Merge, as the duality of semantics (\Hx{1}{ex:duality}) requires. Contrast the two arguments: the subject is externally merged in its $\theta$-position; the object never is. That asymmetry is the restriction of boxing to internal arguments (\S\ref{sec:proposal}), and it is the first place the symmetric alternative would differ.

Then comes the \vs\ phase, and with it a choice. \vs\ consults the box and values Case on \bx{QP}, anchoring the relation to \uth\ in its domain: the QP will be pronounced at \uth, whatever happens next. Now there are two ways to go. On option A, \vs's Minimal Search finds the undischarged \uth, and the boxed QP is externally merged at the edge of \vs P and discharges it, taking scope over the complement of \vs; the subject, which moves on to Spec-TP, ends up above it, and the reading is $\exists > \forall$. On option B, \uth\ is left undischarged, the box is still full, and the derivation continues to the C phase. The two options are two determinate derivations, not a choice inside one, and Figure~\ref{fig:deriv1-optA} draws the first. The QP node at the edge is drawn in the style of a moved element because it is at an edge, but it got there by discharge, not by Internal Merge.

\begin{keiofig}{fig:deriv1-optA}{Derivation 1, option A: discharge at the \vs\ phase. The QP is interpreted at the edge and pronounced at \uth: scope high, Case low.}
\begin{forest}
[\vs P, SOtree, tiny
  [{QP}, SOmoved [{\LI{every book}}, SOlexical]]
  [\vs P [DP [{\LI{a child}}, SOlexical]] [\vs$'$ [\vs, SOhead] [VP [V, SOhead [{\LI{read}}, SOlexical]] [{V$_{\uthr}$\\\ghost{\scriptsize pronounced: \LI{every book}}}, SOtrace]]]]]
\end{forest}
\end{keiofig}

On option B, T is merged, the subject is internally merged to Spec-TP, and C is merged; TP still contains an undischarged \uth. C's Minimal Search finds it, and the boxed QP is externally merged at the edge of CP and discharges it. The box is empty, and the phase is transferred. The scope is over TP, and the reading $\forall > \exists$ (Figure~\ref{fig:deriv1-c}).

\begin{keiofig*}{fig:deriv1-c}{Derivation 1, option B. W8--W10 (compressed): T is merged, the subject is internally merged to Spec-TP, C is merged; TP still contains an undischarged \uth. W11: the boxed QP is externally merged at the edge of CP and discharges \uth; the box is empty; Transfer. Scope over TP: $\forall > \exists$.}
\setcounter{minsyn@step}{9}%
\begin{minipage}[t]{0.38\linewidth}\vspace{0pt}
\begin{workspacebox}[EM T; IM the subject; EM C]{QP}
  \footnotesize\WSeqBox{QP}{\BX{QP},\allowbreak\{\text{C},\text{TP}_{\uthr}\}}
  \par\smallskip\footnotesize W8--W10: T; the subject to Spec-TP; C. TP still contains an undischarged \uth.
\end{workspacebox}
\end{minipage}\hfill
\begin{minipage}[t]{0.58\linewidth}\vspace{0pt}
\begin{workspace}[discharge at C; the box is empty]
  \footnotesize\WSeqTree{\{\text{QP},\{\text{C},\text{TP}\}\}}
    {[CP, SOtree, tiny
      [QP, SOmoved [{\LI{every book}}, SOlexical]]
      [C$'$ [C, SOhead] [TP [DP, SOmoved [{\LI{a child}}, SOlexical]] [\vs P [DP, SOtrace [{$\langle$\LI{a child}$\rangle$}]] [VP [{\LI{read}}, SOlexical] [{V$_{\uthr}$}, SOtrace]]]]]]}
\end{workspace}
\end{minipage}
\end{keiofig*}

The outputs at the two interfaces are these:

\begin{exe}
\ex \label{ex:outputs}
\begin{xlist}
\ex \label{ex:SM} SM: \textit{a child read every book}. The quantifier is pronounced at its Case site, which is the position of \uth.
\ex \label{ex:CI-B} CI, option B ($\forall > \exists$):\\ $\sem{\textit{every book}}\big(\lambda x.\ \exists y\,[\,\mathrm{child}(y) \wedge \mathrm{read}(y, x)\,]\big)$
\ex \label{ex:CI-A} CI, option A ($\exists > \forall$):\\ $\exists y\,\big[\,\mathrm{child}(y) \wedge \sem{\textit{every book}}\big(\lambda x.\ \mathrm{read}(y, x)\big)\,\big]$
\end{xlist}
\end{exe}

In words, (\ref{ex:CI-B}) says that for every book, some child or other read it, and (\ref{ex:CI-A}) that there is one child who read every book: the two readings of \Hx{1}{ex:someone}, reached without moving anything. The derivation contains no covert movement and no lower copy of the quantifier. The object position is occupied by the verb's own \uth\ token, and the $\lambda$ in (\ref{ex:CI-B}) and (\ref{ex:CI-A}) is what discharge contributes. Compare Cooper's (\ref{ex:cooper-deriv}): the same two results, from the order of discharge rather than the order of retrieval. What discharge is, if it is not abstraction over an index, is the question of \Hp{4}. The $\lambda$ is, for now, an honest debt; \Hp{4} pays it with a relation between tokens.

The two readings correspond to two derivations, which differ in whether \vs\ or C discharges the box. Each derivation is determinate; the ambiguity lies in which derivation we have, not inside either of them. Of the three ways out of the Determinacy problem (ambiguity resolved at the interface; two determinate derivations; scope as the one corner of the grammar where determinacy fails), this is the second, and it is the one I will defend. On either derivation Determinacy is preserved: no operation ever faces two targets.

What, then, chooses between the two derivations in English? It is tempting to restate \posscitet{Fox2000} Scope Economy inside the derivation: discharge at the first phase at which it can apply, unless a later discharge yields a different interpretation. I now think it is better not to. A condition of that form looks ahead, to an interpretation not yet computed, which is exactly what a cyclic derivation cannot do (\S\ref{sec:open}, item~\ref{open:ambiguity}). Both derivations are generated, neither is filtered by the syntax, and the preference for the surface reading that English speakers show out of the blue is left to processing, which is where the hesitations speakers show over \Hx{1}{ex:someone} suggest it belongs. What this costs is the work Scope Economy did elsewhere, above all the parallelism of scope under ellipsis \citep{Fox2000}, which the box theory now owes. In Japanese, something else chooses, as \S\ref{sec:rigidity} argues, and it chooses in the grammar.

\subsection{Derivation 2: \textit{Kodomo-ga dono hon-mo yonda}}\label{sec:deriv2}

\begin{exe}
\ex \label{ex:kodomo}
\gll Kodomo-ga dono hon-mo yonda.\\
     child-\gl{nom} which book-\gl{mo} read.\gl{pst}\\
\glt `A child read every book.' / `The children read every book.' \src{(subject $>$ object only)}
\ex \label{ex:deriv-kodomo} Workspace states for (\ref{ex:kodomo}):
\end{exe}

\begin{center}\small
\renewcommand{\arraystretch}{1.12}
\begin{tabularx}{\linewidth}{@{}lLP{6.0cm}@{}}
\toprule
 & Workspace & Operation\\
\midrule
W1 & [\,\textit{dono}, \textit{hon}, \textit{mo}, \textit{yom-}, \vs, \textit{kodomo}, T, C\,] & \\
W2 & [\,\{\textit{dono}, \textit{hon}\}, \textit{mo}, \textit{yom-}, \vs, \textit{kodomo}, T, C\,] & MERGE(\textit{dono}, \textit{hon}): the restriction\\
W3 & [\,\{\{\textit{dono}, \textit{hon}\}, \textit{mo}\}, \textit{yom-}, \vs, \textit{kodomo}, T, C\,] & MERGE(R, \textit{mo}): \textit{mo} closes the alternatives; [uQ] on \textit{mo}\\
W4 & [\,\bx{QP}, \textit{yom-}, \vs, \textit{kodomo}, T, C\,] & \textsc{box}, forced by [uQ]\\
W5 & [\,\bx{QP}, \{\textit{yom-}, \{\textit{yom-}\}\uthr\}, \vs, \textit{kodomo}, T, C\,] & \textsc{fission}: Self-Merge, then Internal Merge\\
W6 & [\,\bx{QP}, \{\vs, VP\uthr\}, \textit{kodomo}, T, C\,] & MERGE(\vs, VP)\\
W7 & [\,\bx{QP}, \{\textit{kodomo}, \{\vs, VP\uthr\}\}, T, C\,] & MERGE(\textit{kodomo}, \vs P)\\
\multicolumn{3}{@{}p{\linewidth}@{}}{\itshape \vs\ phase: \vs\ consults \bx{QP} and values its Case; \textit{mo} is pronounced in the slot of \textup{-}\textit{o}. Discharge order is read off an existing order (\S\ref{sec:rigid}): the $\theta$-hierarchy and the order of externalisation both put \textit{kodomo} above \textit{dono hon-mo}, so the box is discharged here, below the subject's final position.}\\
\ldots & T, the subject to Spec-TP, C & nothing left in the box; Transfer\\
\bottomrule
\end{tabularx}
\end{center}

Derivation 2 has the shape of Derivation 1, again in seven steps, with the particle carrying [uQ]. The difference is at the \vs\ phase, drawn in Figure~\ref{fig:deriv2-v}, head-final. The box is discharged at \vs, not at C, and not by choice; why not by choice is the subject of \S\ref{sec:rigidity}. After that, T is merged, the subject moves to Spec-TP, and C is merged, with nothing left in the box; the phase is transferred, and there is one reading, subject $>$ object.

\begin{keiofig}{fig:deriv2-v}{Derivation 2, the \vs\ phase. \vs\ consults the box and values Case; the QP is discharged at the edge of \vs P, below the subject's final position, and pronounced at \uth, with \textit{mo} in the slot of the case particle.}
\begin{forest}
[\vs P, SOtree, tiny
  [{QP}, SOmoved [{\LI{dono hon-mo}}, SOlexical]]
  [\vs P [DP [{\LI{kodomo-ga}}, SOlexical]] [\vs$'$ [VP [{V$_{\uthr}$\\\ghost{\scriptsize pronounced: \LI{dono hon-\hl{mo}}}}, SOtrace] [V, SOhead [{\LI{yonda}}, SOlexical]]] [\vs, SOhead]]]]
\end{forest}
\end{keiofig}

What changes when force is a particle? Three things. First, the trigger for boxing, [uQ], sits on a pronounced element, \textit{mo}. Second, \textit{mo} is pronounced exactly in the slot where the structural case particle would be: \textit{dono hon-mo}, not *\textit{dono hon-o-mo}; it displaces \textit{ga} and \textit{o}, though not postpositional \textit{ni}, as in \textit{dono hito-ni-mo} `to everyone'. This is the small fact flagged in \Hs{1}{sec:vote}, that after a full nominal the force particle and the case particle compete for one slot; the lexicalised \textit{dare-mo-ga} of \Hx{2}{ex:daremoga}, where \textit{ga} follows \textit{mo}, is the exception that Lecture~2 treated as case filling D. On (\ref{ex:case}) it is what one would expect: the boxed quantifier is externalised at its Case site, and in Japanese the particle that carries its force is what fills the slot. Third, discharge is not free. (\ref{ex:kodomo}) has one reading, and the reason, I will argue in \S\ref{sec:rigidity}, is that in Japanese the order of discharge is not chosen but read off an order the derivation already has.

\subsection{Four objections, and lines of answer}\label{sec:objections}

\textit{Is Self-Merge an instance of MERGE as MSMT defines it, and what does it add to the workspace?} Self-Merge has three uses in this system: it can build a first projection line, it can supply the token of a fissioned verb, the first of fission's two steps, and it can box. For boxing the answer is simple. Since sets have no repeated members, MERGE($\alpha$, $\alpha$) yields $\{\alpha, \alpha\} = \{\alpha\}$, a singleton. It adds exactly one new accessible object, the singleton itself; $\alpha$ was accessible before and remains so. Minimal Yield is respected, and Inclusiveness too, since the singleton contains nothing but $\alpha$. What I have to stipulate is that a singleton root is inert, not integrated into the main object until it is discharged. I do not think that follows from anything yet; it is the price of (\ref{ex:box}), and I would rather own it than hide it. (MSMT says that Merge targets two and only two elements, and \citet[18]{Chomsky2024} leaves singleton phrases as an open question; on the other hand, the UCLA lectures merge an element with itself for the successor function, \citealt[25]{Chomsky2019}. Self-Merge is, at least, not foreign to the system.)

Fission is the harder case. On my first formulation, one application of Self-Merge yielded two tokens of the verb, identical except in which of their features were visible. That was inconsistent with the singleton argument just given: MERGE(V, V) yields \{V\}, one object, not two tokens, and if it did yield two it is not obvious what the second would add to the workspace, or how Minimal Yield could be met. The answer, I now think, is that fission is not one application of MERGE but two (Figure~\ref{fig:fission}). The first, Self-Merge, MERGE(V, V) $= \{\mathrm{V}\}$, yields a singleton, V's own projection: one new accessible object. The second, Internal Merge, MERGE(V, $\{\mathrm{V}\}$) $= \{\mathrm{V},\{\mathrm{V}\}\}$, yields one more. Form Copy construes the two occurrences of V as copies. The singleton \{V\} in complement position is the token \uth: of type $e$, with a lexical source, and with nothing added. Each application of MERGE adds exactly one object, so Minimal Yield is met at each step. What distinguishes fission from boxing is that the singleton is not left as an inert root: the very next application of MERGE integrates it, as the complement of V. Note, finally, that what is internally merged at the second step is V, not the token. The token is the singleton formed by Self-Merge, a degenerate case of External Merge, so the $\theta$-position is still filled by External Merge, as MSMT requires, and there is no movement into a $\theta$-position.

\textit{Isn't `$\theta$ high, Case low' backwards?} Not quite. The $\theta$-role is assigned low, by the verb, to its own \uth\ token, in a configuration created by Self-Merge, which is a degenerate case of External Merge. Discharge is not $\theta$-assignment; it is the identification of the quantifier with the variable, which is to say binding, which is what scope is. In that sense the duality of semantics (\Hx{1}{ex:duality}) survives: the $\theta$-relation is still the verb's, made by (External) Merge, and the scope relation is made by the box.
But what operation is discharge? It is External Merge of the boxed singleton to the phase edge, triggered by the phase head's Minimal Search finding an undischarged token in its domain. It is not a third operation: the boxed singleton is a member of the workspace, and merging a member of the workspace with the phase is External Merge by definition. Two things follow. First, the duality of semantics holds in the form the box theory of \textit{wh} already gives it: $\theta$-positions are filled by External Merge, and edges yield scope, by Internal Merge or by consultation and discharge. Second, nothing percolates. The phase head finds the token by Minimal Search, so the revised Inclusiveness condition of \Hp{4} can list MERGE (with discharge as a case of it), Agree and Form Copy as the only operations that build relations.

\textit{Isn't \uth\ percolation exactly the percolation of indices that \citet[114]{Chomsky2000keio} warned against?} Yes, if \uth\ percolates. Chomsky allows that indices and the like could be innocuous, but questions arise if they function within the computation, and his example is the percolation of indices. On the version just given, \uth\ does not percolate: the phase head finds the undischarged \uth\ by Minimal Search, just as C[Q] consults the box. \Hp{4} takes up what remains of the objection: what kind of relation, if not an index, ties the boxed quantifier to the token it discharges.

\textit{With no level of LF, what makes QR covert?} Nothing moves covertly. The quantifier is pronounced at its Case site and interpreted where it discharges, by (\ref{ex:case}) and (\ref{ex:externalisation}); the single-output literature \citep{Bobaljik2002} is the ally here, not the opponent.

% =====================================================================
\section{A harder case: derive it, break it}\label{sec:worksheet}

Now let us put the proposal to work on a harder case: the first Japanese sentence of Lecture~1, with two quantifiers and one reading.\toexercise{worksheet}

\begin{exe}
\ex \label{ex:ws-sentence}
\gll Dare-ka-ga dono hon-mo yonda.\\
     who-\gl{ka}-\gl{nom} which book-\gl{mo} read.\gl{pst}\\
\glt `Someone read every book.' \src{(someone $>$ every book only; \Hx{1}{ex:dareka-dono})}
\end{exe}

The exercise (\exref{worksheet}) is to derive (\ref{ex:ws-sentence}) in the notation of (\ref{ex:deriv-child}) and (\ref{ex:deriv-kodomo}), marking every point at which the derivation has a choice (whether the subject boxes, and at which phase the boxed object is discharged), and to test each choice against Determinacy, Resource Restriction and Minimal Yield, as set out in \Hx{1}{ex:conditions}. (\ref{ex:ws-deriv}) gives the derivation with its choice points marked, and what follows is what it turns up.

\begin{exe}
\ex \label{ex:ws-deriv} Workspace states for (\ref{ex:ws-sentence}), with the choice points:
\end{exe}

\begin{center}\small
\renewcommand{\arraystretch}{1.12}
\begin{tabularx}{\linewidth}{@{}lLP{4.2cm}P{3.0cm}@{}}
\toprule
 & Workspace & Operation & A choice here?\\
\midrule
W1 & [\,\textit{dare}, \textit{ka}, \textit{dono}, \textit{hon}, \textit{mo}, \textit{yom-}, \vs, T, C\,] & & no\\
W2 & [\,\{\textit{dono}, \textit{hon}\}, \textit{dare}, \textit{ka}, \textit{mo}, \textit{yom-}, \vs, T, C\,] & MERGE(\textit{dono}, \textit{hon}): the restriction & no\\
W3 & [\,\{\{\textit{dono}, \textit{hon}\}, \textit{mo}\}, \textit{dare}, \textit{ka}, \textit{yom-}, \vs, T, C\,] & MERGE(R, \textit{mo}): \textit{mo} closes the alternatives; [uQ] on \textit{mo} & no\\
W4 & [\,\bx{QP}, \textit{dare}, \textit{ka}, \textit{yom-}, \vs, T, C\,] & \textsc{box}, forced by [uQ] & no: forced, so Determinacy holds\\
W5 & [\,\bx{QP}, \{\textit{yom-}, \{\textit{yom-}\}\uthr\}, \textit{dare}, \textit{ka}, \vs, T, C\,] & \textsc{fission}: Self-Merge, then Internal Merge & no: one object per step, so Minimal Yield holds\\
W6 & [\,\bx{QP}, \{\vs, VP\uthr\}, \textit{dare}, \textit{ka}, T, C\,] & MERGE(\vs, VP) & no\\
W7 & [\,\bx{QP}, \{\textit{dare}, \textit{ka}\}, \{\vs, VP\uthr\}, T, C\,] & MERGE(\textit{dare}, \textit{ka}): the subject & yes, unless boxing is restricted to internal arguments: does \textit{dare-ka} box?\\
W8 & [\,\bx{QP}, \{\{\textit{dare}, \textit{ka}\}, \{\vs, VP\uthr\}\}, T, C\,] & MERGE(DP, \vs P): the external argument, by External Merge & yes, unless the order is read off (\S\ref{sec:rigid}): discharge at \vs, or at C?\\
\bottomrule
\end{tabularx}
\end{center}

The two choice points are the derivation's whole interest. At W7, with boxing restricted to internal arguments, \textit{dare-ka} does not box and is externally merged in Spec-\vs P, so the derivation is (\ref{ex:deriv-kodomo}) again; on the symmetric alternative it would box, and the box would have two members. At W8, \vs\ consults the box, values Case on the QP and discharges it below the subject: the order of discharge is read off an order the derivation already has (\S\ref{sec:rigid}), so the discharge at C that would give inverse scope is not available, and there is one reading, subject $>$ object.

\subsection{Five problems the derivation raises}\label{sec:five-problems}

The derivation raises at least five problems, and a careful derivation finds most of them. I set them out here in the order in which they arise.

\textit{Does dare-ka box?} \textit{Ka} supplies existential force, so by (\ref{ex:box}) the phrase might be expected to carry [uQ], and the inventory of \S\ref{sec:inventory} treats it so. But if a \textit{ka}-phrase is closed at its own edge, as a choice function in \posscitet{Kratzer2005} sense, its force is fixed locally and no phase has to fix it: it lacks [uQ] in the relevant sense, does not box, and is merged in Spec-\vs P exactly like \textit{kodomo} in (\ref{ex:deriv-kodomo}); the derivation is then (\ref{ex:deriv-kodomo}) again, and the problem disappears. If it does box, the box has two members, and the next three problems arise.

\textit{The pairing problem.} With two boxed quantifiers and two $\theta$-slots (the subject's, on \vs, and the object's \uth, on V), what determines which quantifier discharges which slot? Determinacy demands a unique answer. Candidates: Case (the quantifier whose Case \vs\ values discharges the \uth\ in \vs's domain), or some other featural match between quantifier and slot.

\textit{The order problem.} In what order are the two discharged, and does the order give the single reading reported for (\ref{ex:ws-sentence})? If discharge simply follows Minimal Search from the phase head, the closest slot, the subject's, is found first, and the first quantifier discharged ends up lower, so that inverse scope would be the default: the wrong result for Japanese. Something has to fix the order, and that is what conjecture~3 (\S\ref{sec:rigid}) is for.

\textit{The identity problem.} After fission there are two occurrences of \textit{yom-}. Any later operation that targets the verb, such as V-to-\vs\ or V-to-T, faces two candidates, unless the visibility of \uth\ distinguishes them. Does it, for the purposes of Determinacy?

\textit{The box as a set.} When C consults a box with two members, is the search determinate? A box with two members is a set, and sets are unordered; Determinacy requires either a feature that distinguishes the two (a Case value, a type of [uQ]) or an order, a stack, which is what \posscitet{Keller1988} nesting amounts to.

Minimal Yield, by contrast, causes no trouble: each application of Self-Merge adds one object, and discharge, as External Merge of a workspace member, adds one more. The real problems are pairing and the box as a set; order is what the rigidity conjecture addresses.

With boxing restricted to internal arguments (\S\ref{sec:proposal}), the first problem has an answer that does not depend on the semantics of \textit{ka}: \textit{dare-ka-ga} is the external argument, and is externally merged in Spec-\vs P whatever its force, so that the derivation of (\ref{ex:ws-sentence}) is (\ref{ex:deriv-kodomo}) again. The box then never has more than one member in a simple transitive clause, and the pairing problem and the problem of the box as a set arise, in a simple transitive clause, only on the symmetric alternative, in which \vs\ fissions too; with two quantified internal arguments they return in any case (\S\ref{sec:open}, item~\ref{open:ditransitives}). The order problem survives in a reduced form: the question is no longer in which order two boxed quantifiers are discharged, but at which of the two phases the one boxed object is, below the subject or above it. That is still the question conjecture~3 answers. The identity problem is unaffected.

\S\ref{sec:open} maps these problems onto a list of open questions: the trigger, ambiguity and Determinacy, $\theta$ high and Case low, the symmetric alternative, ditransitives, phrase or feature, weak crossover, antecedent-contained deletion, and inverse linking.

% =====================================================================
\section{Korean and Japanese pronounce the retrieval step}\label{sec:korean}

This section has two aims: to show that Korean and Japanese pronounce a step of storage and retrieval that English keeps silent, and to use them to answer a question the box theory leaves open, namely whether what goes into the box is a phrase or a feature. It is the argument I made for Korean at ICKL, compressed and, in one respect, narrowed.

\subsection{The claim}\label{sec:claim}

\begin{exe}
\ex \label{ex:externalised-box} Indeterminate quantification in Korean and Japanese pronounces the retrieval step. The indeterminate is the stored restriction, in situ; the particle is the retrieval operator, which fixes force and scope.
\end{exe}
At ICKL I said that Korean externalises the box. That was too strong, and (\ref{ex:externalised-box}) says less. What Korean and Japanese pronounce is the step of retrieval, and that step would be pronounced whether the store were syntactic or semantic. Which kind of store is at work depends, I will argue, on the relation: association between a particle and an indeterminate has the profile of a semantic store, while licensing and scope have the profile of the box (\S\ref{sec:wedge}, and conjecture~1 in \S\ref{sec:conjectures}).

In \textit{nwukwu-nka} `someone', the restriction stays where the indeterminate root is, and only the particle carries force (Figure~\ref{fig:nwukwunka}). That is `store the restriction, retrieve through the operator', done overtly. Across languages, scope is mostly silent; a language that wears its scope-taking on its morphology is a window onto an architecture that English keeps in the dark.

\begin{keiofig}{fig:nwukwunka}{\textit{Nwukwu-nka} `someone': the restriction stays in situ; the particle retrieves it, fixing force and scope.}
\begin{tikzpicture}[every node/.style={font=\footnotesize\rmfamily}, >=Stealth]
  \node[draw=gtaccent, rounded corners=2pt, inner sep=3pt, text width=2.3cm, align=center] (root) at (0,0) {\textit{nwukwu}\\\ghost{\scriptsize restriction, stored, in situ}};
  \node[draw=gtaccent, rounded corners=2pt, inner sep=3pt, fill=gtexbg, text width=2.6cm, align=center] (prt) at (4.4,0) {\textit{-nka}\\\ghost{\scriptsize retrieval operator: force $\exists$, scope}};
  \draw[->, gtaccent, line width=0.7pt] (prt) -- node[above, font=\scriptsize, text=gtghost] {retrieves} (root);
\end{tikzpicture}
\end{keiofig}

\subsection{The inventory, with Japanese counterparts}\label{sec:inventory}

\begin{center}\small
\begin{tabularx}{\linewidth}{@{}P{2.9cm}P{3.1cm}P{3.0cm}L@{}}
\toprule
Korean & Japanese & Type & Relation to the box\\
\midrule
\textit{motun} N & \textit{subete-no} N & universal in the nominal & a clean case of a boxed phrase\\
\textit{nwukwu-nka}, \textit{-na}, \textit{-to} & \textit{dare-ka}, \textit{dare-demo}, \textit{dare-mo} & indeterminate + particle & the retrieval step pronounced: restriction stored, particle retrieves\\
\textit{amwu-to} \ldots{} \gl{neg} & \textit{dare-mo}, \textit{nani-mo} \ldots{} \textit{-nai} & conjunctive quantifier, NPI & clause-bound, like $\theta$-discharge\\
\textit{-ssik} & \textit{-zutu} & distance distributive & scope made visible on the share\\
\textit{-man} & \textit{-dake} & `only' & retrieval with exhaustification\\
\bottomrule
\end{tabularx}
\end{center}

The first row is the control: a nominal universal that boxes as a phrase, like English \textit{every}. The generalisation is that Korean and Japanese mark quantificational force overtly almost everywhere, which is exactly why the step of retrieval is visible in them.

The rows below the first are where the two languages earn their place in the argument. The second row is the heart of it: one root, several particles, several forces, with the restriction left in place and the force carried by the particle, which is the configuration that (\ref{ex:externalised-box}) calls the retrieval step pronounced. The third row is the same configuration under negation, and it supplies the relation on which the 2\,$\times$\,2 of \S\ref{sec:wedge} turns, since licensing by negation is bounded by the clause in a way that association is not (\Hx{2}{ex:korean-npi}; \Hx{2}{ex:npi-jp-b}). The fourth row is the one case in which the morphology sits on neither the restriction nor the operator but on the share: \textit{-zutu} and \textit{-ssik} attach to the distributed quantity, while the distributive force is fixed on the spine (\Hs{2}{sec:zutu}), so that both ends of a scope relation are visible at once. The fifth row, `only', is retrieval with something added: the particle fixes force and scope, and also exhaustifies over the alternatives to its host, which is why \textit{-man} and \textit{-dake} reappear in \S\ref{sec:wedge} as interveners, the operators that consume alternatives on their way to a question particle (\Hs{2}{sec:localities}).

\subsection{Phrase or feature?}\label{sec:phrase-feature}

The box theory boxes maximal projections. In my version, Self-Merge applies to the whole QP, and operator and restriction are set aside as a unit. Korean and Japanese do not seem to do that. In \Hx{2}{ex:shimoyama} the restriction, \textit{dono gakusei}, is deep inside a relative clause and the particle at the edge of the whole nominal; in a question the restriction is in situ and the particle at C (Figure~\ref{fig:phrase-feature}). Only the force is high. Korean and Japanese therefore look like boxing a feature (Move-F, in the sense of \citealt[ch.~4]{Chomsky1995keio}) rather than boxing a phrase. They are a natural experiment on a question the box theory leaves open, and they seem to answer it: feature.

\begin{keiofig}{fig:phrase-feature}{Phrase or feature? In a Japanese question the restriction stays low and only the force, the particle, is high.}
\begin{forest}
[CP, SOtree, xslim
  [TP [{\LI{Taroo-ga}}, SOlexical] [VP [{\LI{nani-o}\\\ghost{\scriptsize restriction, low}}, SOlexical] [{\LI{katta}}, SOlexical]]]
  [{C\quad\hl{\LI{no}}\\\ghost{\scriptsize force, high}}, SOhead]]
\end{forest}
\end{keiofig}

There is, I think, a way of having it both ways, which I offer as a suggestion: \emph{phrase at storage, feature at retrieval}. What goes into the box is a phrase; what a phase head consults at retrieval is a feature, [uQ]. English pronounces the stored phrase and nothing at retrieval. Korean and Japanese pronounce the operator that does the retrieving, and so what we see looks like Move-F.

\begin{center}\small
\begin{tabularx}{\linewidth}{@{}P{2.6cm}LL@{}}
\toprule
 & Pronounces at storage & Pronounces at retrieval\\
\midrule
English & the stored phrase (\textit{every book}) & nothing\\
Korean, Japanese & the restriction (\textit{nwukwu}, \textit{dono gakusei}) & the operator that retrieves (\textit{-nka}, \textit{mo}, \textit{no})\\
\bottomrule
\end{tabularx}
\end{center}

One caution: `high' here means `at the edge of the domain the particle closes', which is C for a question particle and the edge of the nominal for \textit{mo} and \textit{ka} as quantifiers. Whether those are the same kind of retrieval is the question of \S\ref{sec:wedge}, and the caution is the seed of conjecture~1: the question particle at C and \textit{mo} at the edge of a nominal may be different kinds of retrieval, semantic association in the one case and syntactic licensing in the other.

\subsection{No lexical triggers}\label{sec:triggers}

On a lexical-trigger view, there is a class of `boxers', the genuine generalised quantifiers (\textit{every}, \textit{no}, \textit{most}), while indefinites take their scope in situ; that is \posscitet{Reinhart2006} division of labour, restated as bearing a feature or not. It is the natural lexical-trigger theory, and Korean does not fit it. \textit{Nwukwu-nka} ($\exists$) and \textit{nwukwu-na} ($\forall$, free choice) share one root. Force is compositional, a matter of the particle and its height, not of the lexical class of the restriction, and there is no class of generalised quantifiers to carry the boxing feature. Korean relocates the feature from the quantified phrase to the particle, that is, to the retrieval step. That is what (\ref{ex:box}) amounts to when [uQ] sits on the particle, and it is the same argument from a one-root paradigm that supported the Q-configuration hypothesis in Lecture~1 (\Hx{1}{ex:Qconfig}).

\subsection{Two precedents}\label{sec:precedents}

The architecture is older than the box. In \citet{GilTsoulas2013} we analysed conjunctive quantifiers with a covert D bearing [uNeg] or [uMood] above the indeterminate and the conjunctive particle (\Hs{2}{sec:conjunctive}; \Hx{2}{ex:covertD}). Set side by side with the box theory, the two are close to notational variants. The table was shown in \Hp{2} as a promissory note; here it is the payment.

\begin{center}\small
\begin{tabularx}{\linewidth}{@{}LL@{}}
\toprule
\citet{GilTsoulas2013} & The box theory\\
\midrule
the indeterminate is a restricted variable & the indeterminate root is the stored restriction, low\\
the conjunctive particle heads QP and supplies force & the particle is the retrieval operator, and supplies force\\
a covert D[uNeg] or D[uMood] needs licensing & a high force position, consulted at the phase\\
Agree with a clausemate Neg or Mood & consultation and discharge, clause-bound\\
force is configurational, not lexical (Gil and Tsoulas's principle (30); \Hx{2}{ex:principle30}) & force is fixed at retrieval, not by a lexical quantifier\\
case in D bleeds polarity sensitivity & Case marks the low position of externalisation\\
\bottomrule
\end{tabularx}
\end{center}

The second precedent is \citet{Tsoulas2003}, on floating quantifiers as overt scope markers (\Hs{1}{sec:FQ}). There a DP carries a scope feature, valued by long-distance Agree with a scopal head, and the valuing site is spelled out as the marker, so that scope freezes there. It is the same instinct, but it is not the box:

\begin{center}\small
\begin{tabularx}{\linewidth}{@{}LcL@{}}
\toprule
\citet{Tsoulas2003} & & The box\\
\midrule
the scope feature & $\approx$ & the boxing feature, [uQ]\\
the marker, spelled out at the valuing site & $\approx$ & the externalised particle\\
freezing & $\approx$ & rigidity\\
scope set by Agree & & scope set by discharge\\
no storage, no segregation & $\neq$ & storage and segregation\\
a Beghelli--Stowell spine of scope positions & $\neq$ & phase edges, and no spine of scope positions\\
\bottomrule
\end{tabularx}
\end{center}

Scope is syntactic and set by Agree, but there is no storage and no segregation, and the spine of scope positions is stipulated where the box has phase edges. Of the two precedents it is the one I would most like readers to see, because its two weaknesses, no store and a stipulated spine, are exactly what the box repairs.

The point, then, is that the split between storage and retrieval is at least a decade old in this work. Gil and Tsoulas already separated the force-provider from the element that needed licensing. What the box adds is the rationale that Gil and Tsoulas's (30) (\Hx{2}{ex:principle30}) only stipulated: force and licensing sit high because that is where retrieval happens, and Agree is consultation. And D[uNeg] is exactly the licensing layer the box needs for \textit{amwu-to} and \textit{dare-mo}. There is a third precedent, for discharge itself: \citet{Tsoulas1999}, on long-distance scrambling in Korean. It belongs with scrambling, and \S\ref{sec:scrambling} takes it up.

% =====================================================================
\section{The box and Cooper's store}\label{sec:box-cooper}

To a first approximation, the boxed, set-aside element is Cooper's store. Here is the comparison, row by row, and with it the answers to the three questions of \S\ref{sec:cooper}: does the box need an index (no), where does its locality come from (from phases), and is it syntactic (yes). The row for nesting is honestly open.

\begin{center}\small
\begin{tabularx}{\linewidth}{@{}P{1.9cm}LLP{2.6cm}@{}}
\toprule
 & Cooper storage & The box & What it does\\
\midrule
Storage & the NP-storage rule (\ref{ex:storage}) & \textsc{box}, by Self-Merge (\ref{ex:box}) & delays interpretation\\
Placeholder & an indexed variable, $x_i$ & the fissioned \uth, of type $e$, with no index & the bound position\\
Retrieval & at any node of type $t$ (\ref{ex:retrieval}) & discharge at a phase (\ref{ex:discharge}) & fixes scope, and force\\
Ambiguity & free order of retrieval & which phase discharges: two derivations & yields the readings\\
Locality & stipulated, or absent & inherited from phases: discharge is clause-bound & bounds scope\\
Nesting & Keller's nested stores, stipulated & open: does phase structure give it for free? & keeps variables bound\\
Where & in the semantics & in the syntax & \\
\bottomrule
\end{tabularx}
\end{center}

So, is the box simply Cooper's store? In mechanism, yes. The box is a store: there is one syntactic object, there is no covert stratum, and scope is read off the order of discharge, just as in Cooper's system it is read off the order of retrieval. If that is all the box is, the single-cycle grammar has bought a store, not a second level, and the minimalist who rejects LF should have no quarrel on this point with anyone who has used a store.

But the rows for locality and nesting are not the same. The discipline of the box is inherited from phase syntax: clause-boundedness follows from the fact that \uth\ must be discharged before its clause is transferred, and, if the question in the nesting row can be answered positively, the constraint Keller had to add would follow as well. A store in the semantics has to assume both. That is the claim I want to defend about the two stores:

\begin{keybox}
\begin{exe}
\ex \label{ex:locality-claim} \textit{Same mechanism, different explanation.} The box is Cooper's store in mechanism. It differs in where its locality comes from: a store in the syntax inherits its locality from phases, and a store in the semantics has to assume it.
\end{exe}
\end{keybox}

Two remarks. First, (\ref{ex:locality-claim}) is an empirical claim, not a matter of taste. If the locality of scope-taking turned out to be the locality of paths (interveners) rather than of domains (edges), the box's explanation would be wrong and a semantic store right. That is what \S\ref{sec:wedge} tests. Second, the placeholder row hides the deepest difference. Cooper's placeholder is an indexed variable; the box's is the verb's own $\theta$-feature, and the relation between it and the boxed quantifier carries no index. That does not mean the box has dispensed with what an index encodes. It encodes it as a relation between two tokens, and \Hp{4} is about what kind of relation that is.

The claim is about these two stores, and not about frameworks. Cooper's is not the only way of scoping a quantifier without moving it, and some of the others inherit their locality as the box does. For categorial grammarians, then, a point of convergence rather than of contrast. Steedman's CCG derives locality and nesting from the combinatorics of the derivation; the box derives them from phases. Both explain what a store stipulates. One might say that phases play for the box the role that combinators play in CCG, and I would welcome a discussion of whether the two make the same predictions. The saxophonist sentence of \S\ref{sec:cooper} is the obvious case on which to compare them. \Hs{4}{sec:neighbours} sets the box beside the rest of its neighbours.

% =====================================================================
\section{The divide, and the \texorpdfstring{2\,$\times$\,2}{2x2} from Lecture~2}\label{sec:wedge}

\subsection{Hamblin alternatives are a storage theory}\label{sec:KS-storage}

A problem for (\ref{ex:locality-claim}) comes from Japanese. \citet{KratzerShimoyama2002} is a storage theory, done in the semantics: the alternatives introduced by an indeterminate expand up the tree until a particle catches them (\Hx{2}{ex:hamblin}). That is Cooper, not the box. If Japanese particle binding is nothing but the propagation of alternatives, then a store in the semantics is all that Japanese needs, and the box collapses into its semantic shadow. Shimoyama's and Kratzer's work is the best version of the view the box has to beat.

It is also, like the box and unlike Cooper's store, free of indices. An alternative set is not a variable: nothing in the propagation is abstracted over, and the particle that catches the alternatives does not need to know which indeterminate introduced them. On the index row of the comparison in \S\ref{sec:box-cooper}, then, the Hamblin store and the box come out even, and the contest between them is fought entirely on the rows for locality and nesting. That is why the test that follows is a test of locality, and of nothing else.

\subsection{The test is locality}\label{sec:test-locality}

The two theories make different predictions about locality (\Hs{2}{sec:localities}). Intervention effects \citep{BeckKim1997,Beck2006} mark where the propagation of alternatives is blocked, and are argued to be semantic: they count interveners on a path. The locality of $\theta$-discharge marks where boxing is bounded, and is syntactic: it counts edges. If the two come apart, there is a syntactic reflex that the semantic store cannot capture, and the box survives. If they coincide, then in this domain a store in the syntax explains nothing that a store in the semantics does not. This is the empirical crux.

At ICKL I put the question as a two-by-two for Korean:

\begin{center}\small
\renewcommand{\arraystretch}{1.3}
\begin{tabularx}{\linewidth}{@{}P{3.2cm}LL@{}}
\toprule
 & \textbf{No focus intervener} & \textbf{A focus intervener on the path}\\
\midrule
\textbf{Same clause} & box \jok, store \jok\newline{\footnotesize\ghost{both predict good}} & box \jok, store \jbad\newline{\footnotesize\ghost{they diverge: the open test}}\\
\textbf{Across a clause boundary} & box \jbad, store \jok\newline{\footnotesize\ghost{they diverge; observed \jbad}} & box \jbad, store \jbad\newline{\footnotesize\ghost{both predict bad}}\\
\bottomrule
\end{tabularx}
\end{center}

The cross-clause cell without an intervener (an embedded indeterminate with \textit{-to}, a matrix negation, the bridge verb \textit{sayngkakha-} `think', and no focus operator anywhere) is observed to fail (\citealt{GilTsoulas2013}: (5); \Hx{2}{ex:korean-npi}). That matches the box, which counts domains, and not the semantic store, which should let the alternatives through. The decisive cell is the open one: same clause, with a focus intervener such as \textit{-man} `only' or another quantifier on the path. If licensing survives, the effect is syntactic; if it degrades, in the way \citet{Linebarger1987} described for English NPIs, the effect is semantic. One caution applies throughout: clause-crossing and the presence of an intervener must be varied independently, or a domain effect can masquerade as a path effect.

\subsection{The \texorpdfstring{2\,$\times$\,2}{2x2} from Lecture~2}\label{sec:keio2x2}

Japanese lets us do something the Korean version could not: hold the particle family constant and vary the relation. \textit{Mo} enters into two relations, association with an indeterminate and licensing by a clausemate negation, and \textit{ka} into association with a \textit{wh}-phrase.

Here are the cells, with the judgements the literature reports entered in them and, in each cell, the items of the judgement sheet in the appendix (\exref{sheet}) that test it:

\begin{center}\footnotesize
\renewcommand{\arraystretch}{1.3}
\begingroup
\begin{tabularx}{\linewidth}{@{}P{2.3cm}LL@{}}
\toprule
 & \textbf{Across a boundary (clause or island), no intervener} & \textbf{Same clause, a focus intervener on the path}\\
\midrule
\textit{Association} (particle and indeterminate) & \jok\ good: \textit{mo} reaches into a relative clause (\Hx{2}{ex:shimoyama}; item 10); a \textit{wh}-phrase in an embedded clause associates with a matrix question particle (item 7) & \jqquest\ degraded: ?*[\textit{John ka Bill}]-\textit{ga nani-o katta no?} (\Hx{2}{ex:disj-a}; item 5), rescued by scrambling (\Hx{2}{ex:disj-b}; item 6)\\
\textit{Licensing} (NPI and negation) & \jbad\ bad: the clausemate condition, \Hx{2}{ex:npi-jp-b} (items 8 and 9, the clausemate control and the cross-clause case) and, for Korean, \Hx{2}{ex:korean-npi} & \hl{the open cell}: the design in (\ref{ex:design}) (item 12)\\
\bottomrule
\end{tabularx}\endgroup
\end{center}


Read the table by rows. Association has no domain effects and does have path effects: that is the profile of a semantic store. Licensing has domain effects: that is the profile of the box. The open cell decides whether licensing also has path effects. If it does not, the two localities come apart within a single particle family, and conjecture~1 (\S\ref{sec:conjectures}) holds: not the box everywhere, but a division of labour between a semantic store for association and a syntactic one for licensing and scope, as the Greek emphatic and non-emphatic n-words suggested within a single paradigm (\Hs{2}{sec:wedge}; \citealt{Giannakidou1998,Giannakidou2000}). If it does, licensing has both profiles, and the case for a syntactic store has to be made on other grounds. Both outcomes are stated before any results are in; that is the point of stating them.

\subsection{The open cell}\label{sec:open-cell}

The design, which I would like to build with colleagues at Keio and run properly, is schematically:

\begin{exe}
\ex \label{ex:design}
\begin{xlist}
\ex \mbox{[ \ldots{} NPI \ldots{} [ \ldots{} X-\textit{dake} \ldots{} ] \ldots{} Neg ]}\\ \hspace*{\fill}{\small(the focused phrase is embedded and does not c-command the NPI: no intervener on the path)}
\ex \mbox{[ \ldots{} X-\textit{dake} \ldots{} [ \ldots{} NPI \ldots{} ] \ldots{} Neg ]}\\ \hspace*{\fill}{\small(the focused phrase c-commands the NPI below negation)}
\end{xlist}
\end{exe}

with \textit{-sae}, \textit{-mo} and a quantifier in place of \textit{-dake}, both orders, matched fillers, acceptability on a seven-point scale and a yes-or-no for each reading. The judgement sheet in the appendix is a pilot at best, and I would much rather design the real study with people who know how to do it, in Language Faculty Science and on the experimental side, than report a pilot as a result.

% =====================================================================
\section{Rigidity as a consequence, and scrambling}\label{sec:rigidity}

\subsection{Rigidity}\label{sec:rigid}

Recall the sentences with which Lecture~1 opened: \textit{Dare-ka-ga dono hon-mo yonda} has one reading, \textit{Dono hon-mo dare-ka-ga yonda} two (\Hx{1}{ex:dareka-dono} and \Hx{1}{ex:dono-dareka}); one reading for SOV, two for OSV. The same pattern shows up with negation:

\begin{exe}
\ex \label{ex:zenin}
\begin{xlist}
\ex \label{ex:zenin-a}
\gll Zen'in-ga sono tesuto-o uke-nakat-ta.\\
     all-\gl{nom} that test-\gl{acc} take-\gl{neg}-\gl{pst}\\
\glt `All did not take that test.' \src{(all $>$ not; *not $>$ all)}
\ex \label{ex:zenin-b}
\gll Sono tesuto-o zen'in-ga uke-nakat-ta.\\
     that test-\gl{acc} all-\gl{nom} take-\gl{neg}-\gl{pst}\\
\glt `That test, all did not take.' \src{(all $>$ not; not $>$ all; \citealt{Miyagawa2001})}
\end{xlist}
\ex \label{ex:motun-neg}
\gll Toli-ka motun khwukhi-lul an mek-ess-ta.\\
     Toli-\gl{nom} every cookie-\gl{acc} \gl{neg} eat-\gl{pst}-\gl{decl}\\
\glt `Toli didn't eat every cookie.' \src{($\forall > \neg$; $\neg > \forall$ for one group of speakers only)}
\end{exe}

On \posscitet{Miyagawa2001} analysis, the subject of (\ref{ex:zenin-a}) raises to Spec-TP, above negation, while in (\ref{ex:zenin-b}) the scrambled object satisfies the requirement on Spec-TP and the subject can stay below negation. \citet{Goro2007} brings acquisition evidence to the Japanese pattern. Note that \textit{zen'in} is a nominal universal, not an item marked by a particle; that will matter for conjecture~3. In Korean the picture is more mixed: \citet{HanLidzMusolino2007} found that speakers split into two groups on sentences like (\ref{ex:motun-neg}), an object universal under short-form negation, and argued that the input does not settle whether the verb raises, so that each speaker settles it for themselves (see also \citealt{HanMusolinoLidz2016}). \citet{BobaljikWurmbrand2012} put the cross-linguistic pattern in terms of a soft preference for matching orders at the two interfaces (Scope Transparency) and its characteristic signature: of the four combinations of two orders and two scopes, three occur, since the base order has only surface scope and the derived order has both.

\begin{center}\small
\renewcommand{\arraystretch}{1.25}
\begin{tabular}{@{}lcc@{}}
\toprule
 & S $>$ O & O $>$ S\\
\midrule
base order (SOV) & \jok & \jbad\\
derived order (OSV) & \jok & \jok\\
\bottomrule
\end{tabular}
\end{center}

The three-out-of-four signature is what conjecture~3 is built to derive. Bobaljik and Wurmbrand state it as a soft preference at the interface; the box states it as a property of the order of discharge. If scope is the height of discharge, then what fixes scope is whatever fixes the order of discharge, and in a rigid language that order is not free. More precisely:

\begin{keybox}
\begin{exe}
\ex \label{ex:conj-rigidity} \textit{Conjecture 3: rigidity reads an existing order.} In a rigid language the order of discharge is not chosen; it copies an order that the derivation already has, either the $\theta$-hierarchy or the order of externalisation. Where the two agree, as in SOV, there is one reading; where they differ, as in OSV, there are two. In a flexible language the order of discharge floats free, and the preference for surface scope is left to processing (\S\ref{sec:deriv1}).
\end{exe}
\end{keybox}

\begin{center}\small
\begin{tabularx}{\linewidth}{@{}lLLL@{}}
\toprule
 & $\theta$-hierarchy & Order of externalisation & Orders of discharge available\\
\midrule
SOV & S $>$ O & S $>$ O & one: S $>$ O\\
OSV & S $>$ O & O $>$ S & two: S $>$ O, O $>$ S\\
\bottomrule
\end{tabularx}
\end{center}

(\ref{ex:conj-rigidity}) derives the three-out-of-four signature directly: in the base order the $\theta$-hierarchy and the order of externalisation coincide, and only one order of discharge is available; in the derived order they come apart, and both are. With boxing restricted to internal arguments, `order of discharge' comes down to the height of the one boxed object: in SOV it is discharged below the subject, at \vs, as in Derivation~2; in OSV it may be discharged at \vs\ or at C.

A word on when the order is read. The $\theta$-hierarchy is available at the \vs\ phase; the order of externalisation of a scrambled object is fixed only when the edge feature that moves it is satisfied, at C, after the phase at which discharge may first apply. The conjecture is therefore a condition on the completed clause: both derivations are generated, and at Transfer the one whose order of discharge copies one of the two orders is the one that survives. That is not look-ahead of the kind that made me reject Scope Economy (\S\ref{sec:deriv1}), since what is read is a fact about the syntactic object, not an interpretation.

What (\ref{ex:conj-rigidity}) does not do is say why Japanese is rigid and English is not, and I no longer think the answer I gave at ICKL is right. There I said that pronounced retrieval explains rigidity: spell the operator out, as Korean does, and the language is rigid; leave it covert, as English does, and it is flexible. The second half fails. \textit{Zen'in} in (\ref{ex:zenin}) is rigid although its force is not on a particle, and German is rigid without pronouncing anything at retrieval at all: \textit{Mindestens ein Student hat jedes Buch gelesen} `at least one student read every book' has only the surface reading, and the fronted \textit{Jedes Buch hat mindestens ein Student gelesen} has both, the three-out-of-four signature again (see \citealt{BobaljikWurmbrand2012} for the German pattern and its literature). (\ref{ex:conj-rigidity}) is a property of the language, not of particle-marked items. What survives is a one-way conjecture, in the terms of \Hs{2}{sec:hierarchy}: a language that fixes force at retrieval and pronounces that step is rigid, at least where it pronounces it, because the learner then has evidence about where retrieval happens and can tie it to an order already present; but rigidity does not require it. The claim is about pronounced retrieval, not about pronounced force: English pronounces force, on \textit{every}, but fixes it at storage, and the conjecture says nothing about it. The rigidity parameter is still relocated, from the availability of movement towards the externalisation of the retrieval step, but only in one direction.

Two honest caveats. First, Korean is not uniformly rigid: scrambling and floating numerals allow flexibility exactly where scope is positional rather than marked by a particle, which supports the generalisation (flexibility survives where the morphology goes quiet) but also shows that it is not exceptionless. Second, the claim that Korean is rigid is itself contested, since the split on (\ref{ex:motun-neg}) looks like competition between grammars rather than a single setting; the residual variability lives in the covert, positional corner, not in the morphologically marked one. One last test is worth running: if \textit{subete-no} N can escape a complex NP \citep{Mineshima2005} but cannot take scope over a subject in SOV, then the two diagnostics measure different things, nesting and the order of discharge (sheet items 4 and 11).

\subsection{Scrambling}\label{sec:scrambling}

Short scrambling creates new scope, as \Hx{1}{ex:dono-dareka} showed. Long-distance scrambling does not:

\begin{exe}
\ex \label{ex:ld-scrambling}
\gll Dono hon-mo dare-ka-ga [Hanako-ga yonda to] omotte iru.\\
     which book-\gl{mo} who-\gl{ka}-\gl{nom} Hanako-\gl{nom} read.\gl{pst} \gl{comp} think-\gl{te} be.\gl{npst}\\
\glt `Someone thinks that Hanako read every book.' \src{(*every book $>$ someone; \citealt{Saito1992keio}; \citealt{Tada1993})}
\end{exe}

\citet{Saito1989} argued that long-distance scrambling is semantically vacuous, undone at LF; \citet{Saito1992keio} and \citet{Tada1993} distinguished clause-internal scrambling, which can create new scope and binding, from long-distance scrambling, which always reconstructs. \citet{BoskovicTakahashi1998} turned the picture upside down: a scrambled phrase is merged in its surface position and lowers at LF to its $\theta$-position, in order to check a $\theta$-feature. They are the closest precedent for `pronounced high, related to the argument position by $\theta$'. The insight I want to keep is that a $\theta$-relation is the glue between a phrase pronounced high and its argument position. The lowering I want to replace with discharge. The long-distance scrambled phrase is pronounced high, in the matrix clause, because that is where it enters a relation visible to the sensorimotor interface, an edge feature, by (\ref{ex:externalisation}); but the \uth\ it discharges belongs to the embedded verb, and discharge is clause-bound, so its scope is low.

\begin{excursus*}[Bo\v{s}kovi\'c and Takahashi's lowering, and the box's discharge]
\begin{minipage}[t]{0.48\linewidth}\centering
{\footnotesize Bo\v{s}kovi\'c and Takahashi: merge high, lower at LF to check $\theta$}\par\medskip
\begin{forest}
[CP, SOtree, tiny
  [{\LI{dono hon-mo}}, SOmoved]
  [TP [{\LI{dare-ka-ga}}, SOlexical] [VP [CP [{\LI{Hanako-ga}}, SOlexical] [VP [{$\theta$}, SOtrace] [{\LI{yonda to}}, SOlexical]]] [{\LI{omotte iru}}, SOlexical]]]]
\end{forest}
\par\smallskip{\scriptsize\ghost{LF lowering lands at the $\theta$-position}}
\end{minipage}\hfill
\begin{minipage}[t]{0.48\linewidth}\centering
{\footnotesize The box: pronounced at an edge, discharged in its own clause}\par\medskip
\begin{forest}
[CP, SOtree, tiny
  [{\LI{dono hon-mo}}, SOmoved]
  [TP [{\LI{dare-ka-ga}}, SOlexical] [VP [CP [{\LI{Hanako-ga}}, SOlexical] [VP [{V$_{\uthr}$}, SOtrace] [{\LI{yonda to}}, SOlexical]]] [{\LI{omotte iru}}, SOlexical]]]]
\end{forest}
\par\smallskip{\scriptsize\ghost{pronounced at the edge; \uth\ discharged at the embedded phase}}
\end{minipage}

\medskip
Two trees, one surface, two theories. The glue is the same in both, but on the second there is no lowering: the $\theta$-relation is the discharge of the embedded verb's \uth, which is clause-bound, so scope is low; pronunciation is high because the edge relation is visible to the sensorimotor interface, by (\ref{ex:externalisation}). The box's version needs no LF and no downward movement.

Discharge has a history, and I should give it. In \citet{Tsoulas1999} I argued, against lowering, that a long-distance scrambled phrase is merged where it is pronounced and never occupies its $\theta$-position, which is filled in the overt derivation by an empty pronoun that the scrambled phrase binds. The evidence was a set of binding facts in Korean. When \textit{caki} `self' is scrambled out of an embedded clause, the subject of that clause can no longer bind it, and only the matrix subject can; and when \textit{caki} and the embedded subject are both scrambled, the reading on which the one binds the other is lost. A phrase that was simply put back in its $\theta$-position at LF should show no such effect. Two empty pronouns in one clause, coindexed, violate Condition~B. The scrambled phrase depends on the pronoun for its $\theta$-role and for its scope, so that, as I put it then, `scope is determined thematically', and radical reconstruction is that dependence and not an operation. For clause-internal scrambling, where a pronoun would rule every sentence out, I adopted \posscitet{NeelemanReinhart1998} proposal that a $\theta$-role is a feature of the assigning head which, if it cannot be discharged at first Merge, passes up the verbal projection until it can. That is where the discharge of a $\theta$-feature comes from. The box theory keeps three things from that paper: nothing lowers; the scope of a long-distance scrambled phrase is the scope of its $\theta$-position; and a $\theta$-feature may be discharged at a distance from the verb. It changes two: the placeholder is the verb's own token and not a pronoun, and nothing percolates (\Hs{4}{sec:percolation}). The price of the first change should be stated. A token of the verb is not a pronoun, Condition~B has nothing to apply to, and the Korean binding facts, which the pronoun derived, are not yet derived here. \citet{Tsoulas2004ypl} had already found the pronoun insufficient, since the local anaphor \textit{caki-casin} is bound by the nearest subject wherever it is scrambled to, and had begun to ask what an index is in a grammar that respects Inclusiveness, which is the question of \Hp{4}.
\end{excursus*}

\begin{keybox}
\begin{exe}
\ex \label{ex:conj-mirror} \textit{Conjecture 2: QR and long-distance scrambling are mirror images.} QR is externalised low (at its Case site) and scoped high (where it discharges). Long-distance scrambling is externalised high (at an edge) and scoped low, because discharge stays within the clause. Radical reconstruction then needs no lowering.
\end{exe}
\end{keybox}

This is why the earlier principle had to be generalised into (\ref{ex:externalisation}). With Case alone, the scrambled object, which gets its Case in the embedded clause, would be pronounced low. (\ref{ex:externalisation}) says instead that a boxed element is pronounced at its highest relation visible to the sensorimotor interface: for QR that is Case, which is low; for A-movement, Case or the EPP, which is high; for scrambling, an edge feature, which is high. Clause-internal scrambling falls out too: the scrambled object is pronounced at the edge of its own clause and may discharge either at \vs\ or at C of that clause, which is why OSV is ambiguous, in line with (\ref{ex:conj-rigidity}). The table generalises the comparison of A-movement and QR in \S\ref{sec:proposal} to four cases:

\begin{center}\small
\begin{tabularx}{\linewidth}{@{}P{3.2cm}P{2.3cm}P{2.6cm}L@{}}
\toprule
 & Pronounced & Scope & The visible relation (\ref{ex:externalisation})\\
\midrule
QR & low & high & Case, at \uth\\
A-movement & high & high & Case or the EPP, at Spec-TP\\
long-distance scrambling & high & low & an edge feature, in the matrix clause\\
clause-internal scrambling & high (own edge) & \vs\ or C of its own clause & an edge feature\\
\bottomrule
\end{tabularx}
\end{center}

One more step is worth taking, though I take it tentatively. If the edge feature that drives scrambling belongs to a phase head that consults the box, then long-distance scrambling is consultation without discharge, exactly like long-distance \textit{wh}: the matrix phase head consults the boxed phrase and fixes its pronunciation, while the discharge of \uth\ stays in the embedded clause. This requires what the remark on \textit{wh}-phrases in \S\ref{sec:proposal} already assumed: discharge fixes scope but does not close the box, which stays open to consultation until the clause in which the element's highest visible relation is fixed is transferred.

% =====================================================================
\section{Three conjectures, and what remains open}\label{sec:conjectures}

\subsection{The conjectures}\label{sec:conj-list}

Three conjectures came out of preparing this lecture. All three are new, as far as I know, and all three can be tested with speakers of Japanese: the first with the open cell of the 2\,$\times$\,2 from Lecture~2, the second with judgements on long-distance scrambling, and the third with the SOV and OSV patterns of nominal universals.

\begin{keybox}
\begin{exe}
\ex \label{ex:conj1} \textit{Conjecture 1: one particle, two localities.} The association of a particle with an indeterminate (\textit{mo} with \textit{dono} N, \textit{ka} with a \textit{wh}-phrase) shows path effects and no domain effects; the licensing of an NPI by negation shows domain effects. That is the decorrelation the caution of \Hs{2}{sec:wedge} asked for, and it points not to the box everywhere but to a division of labour: a semantic store for association, a syntactic store for licensing and scope.
\ex \label{ex:conj2} \textit{Conjecture 2: mirror images} (= (\ref{ex:conj-mirror})). QR is externalised low and scoped high; long-distance scrambling is externalised high and scoped low, because discharge is clause-bound. Radical reconstruction needs no lowering.
\ex \label{ex:conj3} \textit{Conjecture 3: rigidity reads an existing order} (= (\ref{ex:conj-rigidity})). In a rigid language the order of discharge copies the $\theta$-hierarchy or the order of externalisation; hence one reading for SOV and two for OSV.
\end{exe}
\end{keybox}

\subsection{What remains open}\label{sec:open}

Here is the list onto which I map the problems of \S\ref{sec:worksheet}, each with the line of answer I can currently offer. The first five are the ones the derivation of \S\ref{sec:worksheet} turns up, and the worksheet (\exref{worksheet}) with it, item~\ref{open:ditransitives} once a second object is added; item~\ref{open:phrase} comes from \S\ref{sec:phrase-feature}; the last three are the semanticists' questions.

\begin{enumerate}
\item \label{open:trigger} \textit{The trigger.} Boxing is forced by [uQ] (\ref{ex:box}). Open: is [uQ] the same feature that, in \Hx{1}{ex:Qconfig}, makes a nominal a QP?
\item \label{open:ambiguity} \textit{Ambiguity and Determinacy.} Two readings are two determinate derivations (\S\ref{sec:deriv1}). Open: what selects the higher discharge in English? Restated inside the derivation, Scope Economy would have to look ahead; left to processing, as I now propose, it owes an account of scope parallelism under ellipsis.
\item \label{open:theta} \textit{$\theta$ high, Case low.} $\theta$ is assigned low, to \uth; discharge is binding (\S\ref{sec:objections}). Open: is discharge a third relation beside External and Internal Merge, and if so, is that compatible with the Strong Minimalist Thesis? On the version adopted in \S\ref{sec:objections}, discharge is External Merge of the boxed singleton, and the third relation is gone; the question then becomes whether the inertness of a singleton root, stipulated in (\ref{ex:box}), can be derived.
\item \label{open:symmetric} \textit{Internal arguments only, or both?} Boxing is restricted here to internal arguments (\S\ref{sec:proposal}). The symmetric alternative lets \vs\ fission too, so that the external $\theta$-role also has a token and a quantified subject may be boxed, with boxed quantifiers paired with tokens by Case. It would treat the two arguments alike, but it brings back the pairing problem and the box as a set (\S\ref{sec:five-problems}).
\item \label{open:ditransitives} \textit{Ditransitives.} The restriction to internal arguments keeps the box to one member only in a simple transitive clause. With two quantified internal arguments both carry [uQ] and both box, and the pairing problem, the box as a set and the order problem of \S\ref{sec:five-problems} return. Pairing has an answer here that it lacks on the symmetric alternative: the two members differ in Case, overtly so in Japanese, where \textit{-ni} survives under \textit{mo} while \textit{-o} is displaced (\textit{dono gakusei-ni-mo} beside \textit{dono hon-mo}), so that a quantifier can be paired with the token in the domain of the head that values its Case. Order is where the data are known. Scope between the two objects is frozen in the double object construction and free in the prepositional dative \citep{Bruening2001}, and in Japanese the base order is rigid while scrambling the direct object over the indirect one creates ambiguity \citep{Hoji1985}. Conjecture~3, applied to two boxed objects, predicts exactly that: in the double object construction and in the Japanese base order the $\theta$-hierarchy and the order of externalisation agree, so one order of discharge is available; in the dative and under scrambling they diverge, so two are. If that holds, frozen scope is a dividend, and English flexibility is a matter of the box against an unboxed subject rather than of two members against each other. Open: the second token. Fission as stated puts a token only in a complement position, and on every current analysis one of the two objects is a specifier, of Appl or of the lower shell; either the second $\theta$-assigning head fissions too, with its token in a specifier, or the specifier object is externally merged in place like a subject and does not box, in which case the frozen reading needs the lower head to be a discharge site.
\item \label{open:phrase} \textit{Phrase or feature.} Phrase at storage, feature at retrieval (\S\ref{sec:phrase-feature}). Open: whether \textit{mo} as a quantifier and \textit{ka} in a question retrieve in the same sense.
\item \label{open:wco} \textit{Weak crossover.} With no movement, weak crossover cannot be a constraint on movement. QR never cured it anyway; \Hp{4} takes it up.
\item \label{open:acd} \textit{Antecedent-contained deletion.} In \textit{John read every book that Mary did}, the box never puts the quantifier inside the verb phrase that serves as antecedent, which contains only the verb's \uth\ token; if that is right, antecedent containment never arises, and ACD may turn out to be a dividend rather than a debt, with the relative clause merged inside the box, in the spirit of late merger \citep{FoxNissenbaum1999,Fox2002}.
\item \label{open:inverse} \textit{Inverse linking.} In \textit{Someone from every city despises it}, a quantifier inside a quantified nominal takes scope over it \citep{May1985,Sauerland2005}. That needs nesting, which \citet{Keller1988} had to stipulate. Does phase structure give it for free? The Japanese probe is Mineshima's \textit{dono sensei-ga susumeta dono hon-mo}, two indeterminates and one particle: is retrieval free or nested?
\end{enumerate}

Of the last three, ACD is the one with most to gain and inverse linking the one with most to lose, since it is where nesting is tested. The excursus that follows sketches both; neither is more than a sketch.

\begin{excursus*}[Two sketches: antecedent-contained deletion and inverse linking]
\begin{minipage}[t]{0.48\linewidth}
\setlength\parskip{0.45\baselineskip}
\textbf{ACD:} \textit{John read every book that Mary did}. The antecedent VP is [\textit{read} V$_{\uthr}$]: it contains the verb's token, not the quantifier, so the elided VP, also [\textit{read} V$_{\uthr}$], is not contained in its own antecedent, and no regress arises. The relative clause \textit{that Mary did} is merged inside the box, with the QP, before discharge: late merger \citep{FoxNissenbaum1999}, in the box.

\begin{center}
{\footnotesize\BOX{\text{every book [that Mary did [read V$_{\uthr}$]]}}}\\[2pt] \ldots\ [\textit{read} V$_{\uthr}$]
\end{center}
\end{minipage}\hfill
\begin{minipage}[t]{0.48\linewidth}
\setlength\parskip{0.45\baselineskip}
\textbf{Inverse linking:} the classic example is a subject, \textit{Someone from every city despises it}; since boxing is restricted to internal arguments and a force-less indefinite never boxes, take an object with a quantified container, \textit{Mary visited most museums in every city}. \textit{Every city} is boxed inside \textit{most museums in every city}, which is itself boxed: a box within a box, on the assumption that P fissions as V does, so that \textit{every city} has a token to discharge. For \textit{every city} to take scope over the containing nominal and bind \textit{it}, the inner box must be discharged after the outer one is introduced, at a phase that c-commands both. That is Keller's nesting, stated as a constraint on the order of discharge. Whether phases give it for free is open.

\begin{center}
\BOX{\text{most museums in \BOX{\text{every city}}}}
\end{center}
\end{minipage}

\medskip
If a semanticist asks about ACD, box-internal late merger is the line of answer; if about inverse linking, nested boxes are the obvious representation, and the question is whether phase structure orders their discharge. Both are to be worked through.
\end{excursus*}

% =====================================================================
\section{Questions, readings}\label{sec:readings}\label{sec:expect}

Some questions are predictable, and it is better to have the answers on the page.

\begin{questions}[Likely questions, with short answers]
\begin{enumerate}
\item \textit{How does your box differ from Chomsky's, and from \posscitet{KitaharaSeely2025}?} In three ways: boxing by Self-Merge, an operation that puts an element in a box where Chomsky's version has none; the fission of \uth; and the application to QR, where the quantifier never occupies its $\theta$-position.
\item \textit{Isn't this Cooper storage with extra steps?} Same mechanism, different explanation: the box inherits its locality from phases, the store stipulates it, and the 2\,$\times$\,2 from Lecture~2 is the test (\S\S\ref{sec:box-cooper}--\ref{sec:wedge}).
\item \textit{Shimoyama's mo reaches into islands; doesn't Japanese favour the semantic store?} For association, yes; licensing is clause-bound. One particle, two localities.
\item \textit{Are the judgements robust?} The judgement sheet in the appendix is a pilot at best, and the proper study is an invitation.
\item \textit{Indefinites scope out of islands; where do they fit?} Force-less indefinites never box; their exceptional scope is choice-functional (\Hs{1}{sec:indefinites}).
\item \textit{Why restrict boxing to objects?} Because the subject is externally merged in its $\theta$-position anyway, and the restriction removes the pairing problem; but the symmetric alternative is open (\S\ref{sec:open}, item~\ref{open:symmetric}).
\end{enumerate}
\end{questions}

\subsection*{Key readings}

\citet{Cooper1983}; \citet{Chomsky2024}; \citet{TsoulasYeo2017}; \citet{BeckKim1997}; \citet{Miyagawa2001}; \citet{BobaljikWurmbrand2012}. Between them they cover the store, the box, our 2017 paper, intervention, the Japanese data on negation, and Scope Transparency.

\subsection*{Further readings}

\citet{Keller1988}; \citet{Kobele2018}; \citet{Kitahara1996}; \citet{KitaharaSeely2025}; \citet{Steedman2012}; \citet{Hoji1985}; \citet{Hagstrom1998}; \citet{Saito1992keio}; \citet{BoskovicTakahashi1998}; \citet{Tomioka2007}; \citet{Goro2007}; \citet{HanLidzMusolino2007}.


% ======================================================================
%  PART 4
% ======================================================================
\keiolecture{4}{Indices after Inclusiveness}{Indices after Inclusiveness}%
  {Day 2, Lecture 4 \textperiodcentered{} Sunday 8 November 2026, 15:45--17:45}
\keioepigraph{The ban on indices was right about values, and too strong about relations.}{}
\lecturemap

% =====================================================================
\section{Aims and plan}\label{sec:aims}

We started, in Lecture~1, with a question: how does a grammar that cannot write \emph{x} quantify? Lecture~3 gave a partial answer. A quantifier can be stored, by Self-Merge, and discharged at a phase, binding the verb's own $\theta$-feature, so that nothing moves covertly and no copy needs converting (\Hx{3}{ex:proposal}). But the answer left a debt. Discharge relates a boxed quantifier to a \uth\ token, and the semantics reads that relation as $\lambda$-abstraction (\Hx{3}{ex:CI-B}). What is that relation, in a grammar without indices? Figure~\ref{fig:debt} redraws the last stage of the derivation of Lecture~3 (\Hx{3}{ex:deriv-child}) so as to show only what matters for this question: the boxed quantifier at the edge, the verb's own token in the complement, and nothing in between but structure.

\begin{keiofig}{fig:debt}{The debt of Lecture~3. The dashed token is the one the variable relates to the quantifier; this lecture is about the relation between the boxed QP and $u$.}
\begin{forest}
[CP, SOtree, xslim
  [QP, SObox]
  [C$'$ [C, SOhead]
    [TP [\ldots]
      [VP [V, SOhead [{\LI{read}}, SOlexical]] [{V$_{\uthr}$ $= u$}, SOtrace]]]]]
\end{forest}
\end{keiofig}

This lecture answers the opening question directly: the variable is a relation over tokens. It is not an item in the workspace, and it carries no value. The box already relies on one such relation, and the first task is to say what it is. The account is then tested on Japanese bound-variable anaphora, with \textit{kare}, \textit{soitu}, \textit{zibun} and \textit{pro}, after which I return to the architecture of Lecture~1 and put the argument of the four lectures together. There is also a reason to care that goes beyond the box. Every working syntactician and semanticist uses indices; if a grammar without them is to be taken seriously, it owes an account of what it puts in their place.

This lecture has three aims. First, to say exactly what is wrong with an index, and what is not. Second, to state a revised Inclusiveness condition under which the relations the grammar needs are admissible and the values it does not need are excluded. Third, to show that the \uth\ of the box is an index of the admissible kind, a relation between tokens that carries no value. Diagnosis, condition, application: the third is the one that pays the debt of Lecture~3.

The usual question is how a grammar without indices can express binding at all. Let me turn it around. Indices are extremely useful, and nobody does without them in practice. Why should the grammar be denied them? What exactly is objectionable about an index? I will suggest that the objection is not to relations, which the grammar computes all the time (sisterhood, c-command, Agree, the copy relation), but to values, which it has no means of producing. This is, I think, a genuine shift of perspective rather than a rhetorical move: the ban on indices was a ban on a notation, and what it should have been is a ban on values. Once that is clear, much of the apparent conflict between Inclusiveness and the semantics of binding dissolves. Not all of it, though: percolation (\S\ref{sec:percolation}) and weak crossover (\S\ref{sec:wco}) remain.

The plan is as follows. \S\ref{sec:features-indices} distinguishes features from indices and states a trilemma. \S\ref{sec:revised} revises Inclusiveness. \S\ref{sec:formcopy} reads Form Copy as a theory of indices. \S\ref{sec:testbed} is the Japanese test bed: \textit{kare}, \textit{soitu}, \textit{zibun} and \textit{pro}. \S\ref{sec:box-index} shows that the box has an index too, and faces the percolation problem. \S\ref{sec:wco} is an open problem, weak crossover. \S\ref{sec:revisited} revisits the architecture of Lecture~1 and states three open questions. \S\ref{sec:neighbours} sets the proposal beside its neighbours in other frameworks, and \S\ref{sec:conclusion} concludes. The excursions set in grey boxes (the logics of trees, sharing in unification grammars, Kayne's derivation, \textsc{slash}, Dependent Type Semantics, attention as a store) are for readers from particular frameworks, and the argument does not depend on them.

% =====================================================================
\section{Features type, indices identify}\label{sec:features-indices}

In an unpublished paper from 2004, \textit{Indices and the theory of grammar}, Murat Kural and I drew a distinction that I think is still the right place to start. Features \emph{type} tokens: they sort lexical items and phrases into classes (third person, plural, nominal, interrogative), and many tokens share each value; the inventory of values is finite and fixed by the lexicon. Indices \emph{identify} tokens: they say which tokens count as the same (coreference, binding, the links of a chain), and their only content is identity and distinctness; their inventory, if there is one, is not finite. In a line, a feature answers the question \emph{what kind of thing is this?}, and an index the question \emph{is this the same thing as that?} Everything that follows is about what happens when a grammar has only the first kind of object and needs the second kind of information.

\subsection{What an index does in a textbook}\label{sec:textbook-index}

It helps to be explicit about what the index does in the analyses we all teach. One notation does at least three jobs:

\begin{exe}
\ex \label{ex:index-jobs}
\begin{xlist}
\ex \label{ex:job-coref} \textit{Coreference:}\quad John$_1$ thinks that he$_1$ is clever.
\ex \label{ex:job-binding} \textit{Binding:}\quad Every boy$_1$ loves his$_1$ mother.
\ex \label{ex:job-chain} \textit{Chains:}\quad Who$_1$ did you see $t_1$?
\end{xlist}
\end{exe}

And it does a fourth, which is the one the semantics cannot do without: the index that Predicate Abstraction reads (\Hx{1}{ex:PA}), the numeral that becomes the $\lambda$-bound variable. Figure~\ref{fig:textbook-index} gives the textbook tree for (\ref{ex:job-binding}), with Heim and Kratzer's binder index.

\begin{keiofig}{fig:textbook-index}{The textbook tree for (\ref{ex:job-binding}). The numeral 1 occurs three times and does three different things.}
\begin{forest}
[S, SOtree, xslim
  [{DP$_1$} [{\LI{every boy}}, SOlexical]]
  [{$\lambda$1}
    [S [{$t_1$}, SOtrace]
      [VP [{\LI{loves}}, SOlexical] [{\LI{his$_1$ mother}}, SOlexical]]]]]
\end{forest}\\[4pt]
{\footnotesize\color{gtghost} Predicate Abstraction at $\lambda$1:\ \ $\lambda x.\,\sem{\mathrm{S}}^{g[1\mapsto x]}$}
\end{keiofig}

The numeral 1 appears three times in the tree, and it does three different things: it marks the moved phrase, it is the variable that $t_1$ and \textit{his$_1$} stand for, and it is the argument of the abstraction. The question for a grammar without indices is which of these it can do without, and which it must reconstruct. The rest of this lecture sorts the three into relations, which are admissible, and values, which are not.

\subsection{Three properties, and what is objectionable}\label{sec:three-properties}

In the 2004 paper we noted three properties that set indices apart from features \citep{KuralTsoulas2004}:

\begin{exe}
\ex \label{ex:index-props}
\begin{xlist}
\ex \label{ex:prop-meaning} Indices have no intrinsic meaning. The index 7 means nothing; all that matters is whether two indices are the same or different.
\ex \label{ex:prop-unbounded} Indices are unbounded. A sentence may need as many distinct indices as it has independent referents and binders, so the inventory cannot be finite, unlike any inventory of feature values.
\ex \label{ex:prop-distance} Indices act at a distance. Coindexation relates tokens arbitrarily far apart, whereas the relations that features enter (selection, Agree) are local.
\end{xlist}
\end{exe}

Each of (\ref{ex:index-props}a--c) is a reason why Inclusiveness should exclude indices, and it is worth keeping the three apart: (\ref{ex:prop-meaning}) is about \emph{content}, (\ref{ex:prop-unbounded}) about the \emph{inventory}, and (\ref{ex:prop-distance}) about \emph{locality}. The repair in \S\ref{sec:revised} answers each of them separately.

\citet[114]{Chomsky2000keio} drew the relevant line. Indices and the like `could be innocuous'; the questions arise when they function within the computation, and his example is the `percolation of indices'. What is objectionable, in other words, is not the notation but indices that do work in the computation. Yet the semantics needs exactly the work they do. Predicate Abstraction (\Hx{1}{ex:PA}) abstracts over an index, and binding is a relation between an index and a binder. So the conflict is real, and it is not a conflict about notation. A grammar that must express identity between tokens, but has only features, faces a trilemma:

\begin{keybox}
\begin{exe}
\ex \label{ex:trilemma} \textit{The trilemma.}
\begin{xlist}
\ex \label{ex:bound-domain} \textit{Bound the domain.} Allow only finitely many index values. The cost is a fixed ceiling on the number of simultaneous dependencies; whether natural language ever exceeds a small ceiling is, interestingly, an empirical question.
\ex \label{ex:unbound-values} \textit{Unbound the values.} Allow features with unboundedly many values: indices proper. That abandons the finiteness of the feature inventory, and Inclusiveness with it.
\ex \label{ex:share} \textit{Share instances.} Let Agree make two locations share a single feature instance, rather than copying a value from one to the other \citep{PesetskyTorrego2007}. Instances are distinct objects, and sharing one instance is what it is to be coindexed.
\end{xlist}
\end{exe}
\end{keybox}

Each horn is in the literature, and each costs something. Figure~\ref{fig:trilemma} sets them out; the heavier arrow is where we are going.

\begin{keiofig*}{fig:trilemma}{The trilemma (\ref{ex:trilemma}). The third horn is the one I keep, and \S\ref{sec:revised} generalises it.}
\begin{tikzpicture}[>=Stealth, every node/.style={font=\footnotesize\rmfamily, align=center}]
  \node[draw=gtaccent, rounded corners=2pt, inner sep=5pt, fill=gtexbg] (q) at (0,0) {A grammar with only features\\must express identity between tokens};
  \node[draw=gtaccent, rounded corners=2pt, inner sep=4pt, text width=4.4cm] (a) at (-5.3,-2.2) {(\ref{ex:bound-domain}) \textbf{Bound the domain}\\finitely many index values\\{\scriptsize\color{gtghost}cost: a fixed ceiling on dependencies}};
  \node[draw=gtaccent, rounded corners=2pt, inner sep=4pt, text width=4.4cm] (b) at (0,-2.2) {(\ref{ex:unbound-values}) \textbf{Unbound the values}\\features with unboundedly many values\\{\scriptsize\color{gtghost}cost: Inclusiveness, and finiteness}};
  \node[draw=gtaccent, rounded corners=2pt, inner sep=4pt, fill=gtexbg, text width=4.4cm] (c) at (5.3,-2.2) {(\ref{ex:share}) \textbf{Share instances}\\two locations, one feature instance\\{\scriptsize\color{gtghost}cost: only where Agree can reach}};
  \draw[->, gtaccent] (q) -- (a);
  \draw[->, gtaccent] (q) -- (b);
  \draw[->, gtaccent, line width=1.2pt] (q) -- (c);
\end{tikzpicture}
\end{keiofig*}

The first horn is less absurd than it sounds, and it has a logic behind it worth a moment even for readers who do not work formally. With a bound on the number of indices, structures of the Government and Binding kind stay within the regular tree languages; without one, they do not \citep{Rogers1998}. \citet{GrafAbner2012} asked the natural follow-up, whether \emph{binding} with a bounded index set is regular. Boundedness is not, then, an arbitrary restriction but a hypothesis with a measurable consequence: how many distinct bound variables does any sentence actually need at once? As far as I know, nobody has shown that a sentence needs more than three or four simultaneous dependencies, though I would not want to rest anything on that.

\begin{excursus}[A page of logic: what can be said about a tree]
The logics in question differ in what they quantify over, and so in what they can define on trees \citep{ThatcherWright1968,Doner1970,Rogers1998}:

\begin{center}\footnotesize
\begin{tabular}{@{}P{2.0cm}P{2.6cm}P{4.6cm}@{}}
\toprule
Logic & What it quantifies over & What it can define on trees\\
\midrule
Propositional & nothing: fixed atomic facts & local label checks\\
First-order (FO) & \emph{nodes}: $\exists x\,\forall y\,(\ldots)$ & dominance, c-command, `some N dominates some V'; \textbf{not} an unbounded chain of dependencies\\
Monadic second-order (MSO) & nodes \emph{and sets of nodes}: $\exists X\,(\ldots)$ & every \textbf{regular} tree language, and only those (Thatcher \& Wright; Doner): paths, closure, `all the nodes between $x$ and $y$'\\
Full second-order & sets, relations, functions & far too much: arbitrary relations between arbitrary nodes\\
\bottomrule
\end{tabular}
\end{center}

Read: a first-order formula says `there is a node $x$ such that for every node $y$ \ldots'; a monadic second-order formula may also say `there is a \emph{set} of nodes $X$ such that \ldots', and it is the sets that let it talk about paths and closure, which is what the regular tree languages need. Full second-order logic may quantify over arbitrary relations, and can then say anything about any two nodes.

Rogers's result is that GB theory, minus free indexation, is MSO-definable, and that free indexation, with unboundedly many distinct indices each relating arbitrary nodes, is what pushes it out. So (\ref{ex:prop-unbounded}) and (\ref{ex:prop-distance}) together are what make indices expensive in the formal sense: unbounded \emph{and} at a distance. Either alone would be fine. The cost of an index, in other words, is formal as well as conceptual.
\end{excursus}

The second horn is, I think, the honest name for what most of us do. The third is \posscitet{PesetskyTorrego2007} feature sharing. After Agree there is one instance of the feature in two places, and that is already a coindexation, since what the two locations share is not a value but an identity. Figure~\ref{fig:sharing} contrasts it with valuation as copying.

\begin{keiofig*}{fig:sharing}{Valuation as copying (left) and as sharing (right), after \citet{PesetskyTorrego2007}. The circled 1 is not a value but an identity.}
\begin{tikzpicture}[>=Stealth, every node/.style={font=\scriptsize\rmfamily}]
  \node[draw=gtghost, rounded corners=2pt, inner sep=3pt] (p1) at (0,0) {T [$\varphi$: \underline{\ \ \ }]};
  \node[draw=gtghost, rounded corners=2pt, inner sep=3pt] (g1) at (4.6,0) {DP [$\varphi$: 3sg]};
  \draw[->, gtghost] (g1) -- node[above]{copy the value} (p1);
  \node[text=gtghost] at (2.3,-0.6) {valuation as copying};
  \node[draw=gtaccent, rounded corners=2pt, inner sep=3pt] (p2) at (8.2,0) {T [$\varphi$: \tikz[baseline]\node[draw=gtaccent, circle, inner sep=1pt]{1};]};
  \node[draw=gtaccent, rounded corners=2pt, inner sep=3pt] (g2) at (13.6,0) {DP [$\varphi$: \tikz[baseline]\node[draw=gtaccent, circle, inner sep=1pt]{1}; 3sg]};
  \draw[<->, gtaccent, line width=0.8pt] (g2) -- node[above, align=center]{one instance,\\two locations} (p2);
  \node[text=gtaccent] at (10.9,-0.6) {sharing};
\end{tikzpicture}
\end{keiofig*}

The last option is, I think, an index honestly described: a relation between occurrences that carries no value of its own. Set against the three properties of (\ref{ex:index-props}), it fares as follows:

\begin{center}\small
\begin{tabularx}{\linewidth}{@{}P{3.0cm}LP{3.4cm}@{}}
\toprule
Property of indices & What a shared instance does about it & Verdict\\
\midrule
(\ref{ex:prop-meaning}) no intrinsic meaning & an instance has no meaning beyond being itself & respected\\
(\ref{ex:prop-unbounded}) unbounded & instances come \emph{with tokens}, not from a list of values; no infinite inventory & respected, without the cost\\
(\ref{ex:prop-distance}) at a distance & sharing is built by Agree, so only where Agree can reach & met only \textbf{locally}\\
\bottomrule
\end{tabularx}
\end{center}

Note the third row. Locality is not respected here but imposed, and that is a prediction rather than a cost; \S\ref{sec:partition} cashes it as the claim that there is no unbounded coindexation in syntax. The rest of this lecture argues that (\ref{ex:share}), generalised, is what the grammar has always been doing, and that Inclusiveness, properly stated, already allows it.

\begin{excursus}[Sharing instances is what unification grammars have always done]
Readers who work in HPSG or LFG will have been thinking it already \citep{PollardSag1994,Bresnan2001}. In a feature structure two paths can be \emph{reentrant}: they lead to one and the same object, marked by a tag such as \fbox{1}. Agreement, control and binding are stated as reentrancy, not as copying:

\begin{center}\footnotesize\sffamily
\begin{tabular}{@{}l@{}}
[ \textsc{subj} \ \ [ \textsc{index} \ \fbox{2} ] \\
\ \ \textsc{head} \ \ [ \textsc{agr} \ \ \fbox{2} ] \\
\ \ \textsc{comps} \ $\langle$ [ \textsc{anaphor} \ \fbox{2} ] $\rangle$ ]
\end{tabular}
\end{center}

The tag \fbox{2} is not a value. It says that three paths end at one object, which is (\ref{ex:share}) exactly, and it is older than Pesetsky and Torrego. The difference from the minimalist case lies in locality. Unification is unbounded within the feature structure (any two paths may be identified), so (\ref{ex:prop-distance}) is not imposed: the index of HPSG \emph{can} act at a distance, which is why HPSG needs its Binding Theory stated separately. The minimalist bet, which \S\ref{sec:partition} makes explicit, is that the locality of Agree and Form Copy \emph{is} the Binding Theory, or a good part of it.
\end{excursus}

% =====================================================================
\section{Inclusiveness revised}\label{sec:revised}

\subsection{Tokens}\label{sec:tokens}

\citet{CollinsStabler2016}, in their formalisation of minimalist syntax (their Definition~5), distinguish lexical items from their tokens. A lexical item token is a pair $\langle \mathrm{LI}, k\rangle$, where $k$ distinguishes separate selections of the same item from the lexicon, such as the two tokens of \textit{the} in \textit{the cat saw the dog} (Figure~\ref{fig:tokens-ws}). Is $k$ an index in the forbidden sense? I do not think it is. It is not added in the course of the computation; it is fixed when the item enters the workspace, and nothing the computation does ever changes it. It is a bookkeeping device of the formalisation, not something Merge consults: the 1 and the 2 do no work in the derivation, and are there only so that we can tell the two selections apart. If pressed on whether $k$ is a value in disguise, the honest answer is that it is a value that never does any work, and what the revised condition below excludes is values that do.

\begin{keiofig}{fig:tokens-ws}{The tokens of \textit{the cat saw the dog} at Select. The two tokens of \textit{the} are distinct objects that look alike.}
\begin{derivation}
\WSsetstep{0}
\begin{workspace}[Select]
  \footnotesize\WSeq{\langle\LI{the},1\rangle,\allowbreak\ \langle\LI{cat},1\rangle,\allowbreak\ \langle\LI{saw},1\rangle,\allowbreak\ \langle\LI{the},2\rangle,\allowbreak\ \langle\LI{dog},1\rangle}
\end{workspace}
\end{derivation}
\end{keiofig}

External Merge of two tokens yields two objects; Internal Merge of one token yields two occurrences of one object:

\begin{exe}
\ex \label{ex:rep-copy}
\begin{xlist}
\ex \label{ex:rep} \textit{Repetitions}: two tokens of one lexical item, $\langle \textit{the}, 1\rangle$ and $\langle \textit{the}, 2\rangle$. They are distinct objects that happen to look alike.
\ex \label{ex:copy} \textit{Copies}: two occurrences of one token, created by Internal Merge. They are one object.
\end{xlist}
\end{exe}

Figure~\ref{fig:rep-copy} shows the two cases side by side. On the left, two selections, two objects: repetitions. The subscripts there are the $k$ of the token pair, shown only so that the two DPs can be told apart. On the right, one selection, moved, two occurrences: copies, with one subscript because there is one token.

\begin{keiofig*}{fig:rep-copy}{Repetitions and copies (\ref{ex:rep-copy}). Left: two External Merges, two tokens, two objects. Right: Internal Merge of one DP, one object with two occurrences.}
\WSsetstep{0}
\begin{derivation}
\begin{twocol}{0.47}
\begin{workspace}[two External Merges]
  \scriptsize\WSeqTree{\{\LI{the}_1,\LI{cat}\},\allowbreak\ \LI{saw},\allowbreak\ \{\LI{the}_2,\LI{dog}\}}
    {[{}, SOtree, tiny, phantom, s sep=18pt
      [DP [{\LI{the}$_1$}, SOlexical] [{\LI{cat}}, SOlexical]]
      [DP [{\LI{the}$_2$}, SOlexical] [{\LI{dog}}, SOlexical]]]}
\end{workspace}
\midcol
\begin{workspace}[Internal Merge of DP with TP]
  \scriptsize\WSeqTree{\{\text{DP},\{\text{T},\{\ldots\text{DP}\ldots\}\}\}}
    {[TP, SOtree, tiny, s sep=14pt, SMroot
      [{DP \LI{the}$_1$ \LI{cat}}, SMcopy, SOmoved]
      [T$'$ [T, SOhead] [\vs P [{$\langle$DP \LI{the}$_1$ \LI{cat}$\rangle$}, SMcopy, SOtrace] [\ldots]]]]}
\end{workspace}
\end{twocol}
\end{derivation}
\end{keiofig*}

\citet{Chomsky2021} builds the same distinction into the grammar, with Form Copy deciding which identical inscriptions are copies and which are repetitions (\S\ref{sec:formcopy}). The set-theoretic foundations of the distinction have been questioned: \citet{Gartner2022} asks whether copies can be had from standard set theory at all, as they are claimed to be. They have also been given an algebraic formulation \citep{MarcolliChomskyBerwick2025}. And the problem itself has a history at Keio: Kitahara's work on Inclusiveness, trace features and multiple occurrences of a single item \citep{Kitahara1997,Kitahara1998,Kitahara2000} asked many of these questions first, before Form Copy existed.

Why does this matter for us? If the grammar can tell copies from repetitions \emph{without} indices, by the history of the tokens and by the relations built among them, then it already has the resources for an index of the admissible kind. The revised condition says what those resources are.

\subsection{The revised condition}\label{sec:revised-condition}

\begin{keybox}
\begin{exe}
\ex \label{ex:incl-revised} \textit{Inclusiveness (revised).}
\begin{xlist}
\ex \label{ex:incl-tokens} \textit{Tokens.} The computation operates on tokens of lexical items, $\langle \mathrm{LI}, k\rangle$. Token identity is fixed by the input to the computation, not created by it.
\ex \label{ex:incl-relations} \textit{Relations.} The computation may establish relations among tokens, but only those built by MERGE (including discharge), Agree and Form Copy.
\ex \label{ex:incl-values} \textit{No new values.} The computation introduces no feature values that are not present in the lexical items that enter it.
\end{xlist}
\end{exe}
\end{keybox}

There are three clauses, one for each of the three properties in (\ref{ex:index-props}). Clause (\ref{ex:incl-tokens}) answers (\ref{ex:prop-meaning}) and (\ref{ex:prop-unbounded}) by making identity part of the input: two inscriptions are the same or different because of where they came from, not because the computation has labelled them. Clause (\ref{ex:incl-relations}) answers (\ref{ex:prop-distance}) by making every relation inherit the locality of the operation that builds it. It admits relations: sisterhood and c-command, which MERGE builds, discharge included; Agree (including the sharing of a single instance, (\ref{ex:share})); and the copy relation that Form Copy assigns. None of these is an object added to the workspace; they are facts about the structures the computation builds. Clause (\ref{ex:incl-values}) is where the condition keeps its teeth. It excludes exactly what makes an index objectionable: a value, such as 7, that has no lexical source and does work in the computation.

Clause (\ref{ex:incl-relations}) names discharge explicitly, and the reason is worth stating. Lecture~3 took discharge to be External Merge of the boxed singleton to the phase edge, triggered by the phase head's Minimal Search for an undischarged token (\Hs{3}{sec:objections}). If that is right, discharge is not a fourth relation-builder beside MERGE, Agree and Form Copy but a case of the first, and the reconstruction of coindexation in (\ref{ex:coindexation}) below names nothing that (\ref{ex:incl-relations}) does not already admit. The duality of semantics then keeps its original form, and the three-clause structure of (\ref{ex:incl-revised}) is unchanged. Should discharge turn out to need a relation of its own, it would be a fourth builder, and reducing it to MERGE would be the open line.

The original condition (\Hx{1}{ex:inclusiveness}) bans indices outright, together with bar levels and traces. Under (\ref{ex:incl-revised}) traces and bar levels are still excluded, since each would be a new object with no lexical source. But an index is admissible if, and only if, it is a relation among tokens built by MERGE, Agree or Form Copy, carrying no value of its own. Nothing that the original condition banned for good reason comes back; only the relations do:

\begin{center}\small
\begin{tabularx}{\linewidth}{@{}LP{3.2cm}P{4.6cm}@{}}
\toprule
 & Original condition (\Hx{1}{ex:inclusiveness}) & Revised condition (\ref{ex:incl-revised})\\
\midrule
traces & banned & banned (a new object)\\
bar levels & banned & banned (a new object)\\
an index as a \textbf{value} & banned & banned by (\ref{ex:incl-values})\\
an index as a \textbf{relation} & banned & admitted by (\ref{ex:incl-relations})\\
\bottomrule
\end{tabularx}
\end{center}

In a nutshell: the ban on indices was right about values, and too strong about relations. This is idea (b) of \Hs{1}{sec:architecture}, revised, and nothing else in the architecture of Lecture~1 changes because of it. Two consequences are worth stating now, both of them predictions. First, since every admissible index is built by MERGE, Agree or Form Copy, every admissible index inherits the locality of the operation that built it; there is no unbounded coindexation in syntax (\S\ref{sec:partition}), and \S\ref{sec:testbed} tests this on long-distance \textit{zibun}. Second, the \uth\ relation of discharge is already one of the three: discharge is External Merge (\Hs{3}{sec:objections}), and (\ref{ex:incl-relations}) lists it under MERGE. \S\ref{sec:box-index} says what that buys, and it is where the debt of Lecture~3 is paid.

\subsection{Assignments from tokens: a remark on the semantics}\label{sec:assignments}

(\ref{ex:incl-revised}) has a direct semantic counterpart. In \citet{HeimKratzer1998} an assignment is a partial function from the natural numbers to individuals, $g : \mathbb{N} \to D$, and a pronoun or trace with index $n$ denotes $g(n)$; Predicate Abstraction modifies $g$ at $n$, giving $g[n \mapsto x]$, the assignment that is like $g$ except that it sends $n$ to $x$. But an assignment function need not map numerals to individuals; it can map tokens to individuals, $g : \mathrm{Tok} \to D$. If what the assignment maps are tokens, with the verb's own token $u$ in the place the textbook gives to the variable's index (\S\ref{sec:box-index}), then the domain of the assignment is a set of objects that the syntax already has, and nothing has been added to interpret them (Figure~\ref{fig:assignments}). Nothing in the Heim and Kratzer machinery depends on the variables being numerals. What it depends on is that they are distinct, and that abstraction can target one of them; tokens give both. I return to this in (\ref{ex:discharge-sem}).

\begin{keiofig}{fig:assignments}{Assignments from numerals, which are values with no lexical source, and from tokens, which the syntax already has.}
\begin{tikzpicture}[>=Stealth, every node/.style={font=\footnotesize\rmfamily}]
  \node (n) at (0,1.2) {$g$: \ 1 $\mapsto$ Taroo, \ 2 $\mapsto$ Hanako, \ 7 $\mapsto$ \ldots};
  \node[text=gtghost, font=\scriptsize\rmfamily] at (0,0.75) {numerals: values with no lexical source};
  \node (t) at (0,-0.3) {$g$: \ $u_{\mathrm{read}} \mapsto$ Taroo, \ $u_{\mathrm{see}} \mapsto$ Hanako};
  \node[text=gtaccent, font=\scriptsize\rmfamily] at (0,-0.75) {tokens: objects the syntax already has};
  \draw[->, gtaccent, line width=0.8pt] (0,0.48) -- (0,0.02);
\end{tikzpicture}
\end{keiofig}

% =====================================================================
\section{Form Copy as a theory of indices}\label{sec:formcopy}

\subsection{What Form Copy computes}\label{sec:fc-computes}

\begin{keybox}
\begin{exe}
\ex \label{ex:FC} \textit{Form Copy} (\citealt{Chomsky2021}; as summarised in \citealt{ChomskyEtAl2023}). At the phase level, two inscriptions are assigned the relation \emph{copy} if
\begin{xlist}
\ex \label{ex:FC-a} they are identical,
\ex \label{ex:FC-b} one c-commands the other, with no identical inscription between them (minimality), and
\ex \label{ex:FC-c} they are in the same phase.
\end{xlist}
\end{exe}
\end{keybox}

Identical inscriptions that fail (\ref{ex:FC}) are repetitions, distinct objects. A moved phrase and its source are copies; two separately selected tokens of \textit{John} are repetitions. Form Copy is an operation of the phase level: it applies to the structure that MERGE has built, and assigns a relation to it. It adds nothing to the workspace, which is why it passes (\ref{ex:incl-relations}). One caveat should be stated at once. If co-phasehood turns out not to be a condition after all, the locality of Form Copy is minimality alone, and the consequence drawn in \S\ref{sec:partition} has to be restated accordingly.

Figure~\ref{fig:fc-action} shows Form Copy at work on the two cases. On the left is the ordinary case, copies created by Internal Merge: the two inscriptions of \textit{which book} are identical, the higher c-commands the lower with nothing identical between them, and they are in the same phase. On the right are two tokens of \textit{John} which meet the first two conditions but fail the third. Grammar does not say whether they denote the same person: there may be two people called John, or one, but if one, not by grammar.

\begin{keiofig*}{fig:fc-action}{Form Copy in action. Left: Internal Merge yields copies. Right: two selections of \textit{John} fail (\ref{ex:FC-c}) and remain repetitions.}
\begin{twocol}{0.48}
\centering\footnotesize Internal Merge: copies\par\smallskip
\begin{forest}
[CP, SOtree, tiny
  [{\LI{which book}}, SOmoved]
  [C$'$ [C, SOhead]
    [TP [{\LI{Taroo}}, SOlexical]
      [\vs P [{\LI{bought}}, SOlexical] [{$\langle$\LI{which book}$\rangle$}, SOtrace]]]]]
\end{forest}\par\smallskip
\scriptsize (\ref{ex:FC-a}) identical \jok\quad (\ref{ex:FC-b}) c-command, nothing identical between \jok\quad (\ref{ex:FC-c}) same phase \jok\\ $\Rightarrow$ \textbf{copies}
\midcol
\centering\footnotesize Two selections: repetitions\par\smallskip
\begin{forest}
[TP, SOtree, tiny
  [{\LI{John}$_1$}, SOlexical]
  [\vs P [{\LI{thinks}}, SOlexical]
    [CP [{\LI{that}}, SOlexical]
      [TP [{\LI{John}$_2$}, SOlexical] [{\LI{is clever}}, SOlexical]]]]]
\end{forest}\par\smallskip
\scriptsize (\ref{ex:FC-a}) identical \jok\quad (\ref{ex:FC-b}) c-command \jok\quad (\ref{ex:FC-c}) same phase \jbad\\ $\Rightarrow$ \textbf{repetitions}
\end{twocol}
\end{keiofig*}

\citet{Chomsky2021} also proposed that obligatory control is Form Copy applied to two \emph{externally} merged inscriptions of the controller, as in \textit{John tried [John to leave]}. That brings Form Copy very close to binding, and it is what makes Form Copy a candidate theory of binding and not only of chains.

\subsection{Form Copy records a partition}\label{sec:partition}

Applied across a derivation, what Form Copy delivers is an equivalence relation over inscriptions, reflexive, symmetric and transitive: which inscriptions are occurrences of the same object. An equivalence relation is the same thing as a partition. The inscriptions fall into classes, and two are related if, and only if, they are in the same class (Figure~\ref{fig:partition}). That is exactly what a system of indices encodes, a partition of the nominal expressions into classes, coindexation being membership of the same class.

\begin{keiofig}{fig:partition}{The partition Form Copy computes over the inscriptions of Figure~\ref{fig:fc-action}: three classes. `Coindexed' means `in the same class'.}
\begin{tikzpicture}[every node/.style={font=\scriptsize\rmfamily}]
  \draw[gtaccent, rounded corners=6pt, fill=gtexbg] (-0.2,-0.2) rectangle (2.5,1.0);
  \node at (0.5,0.62) {\textit{which book}};
  \node at (1.6,0.18) {$\langle$\textit{which book}$\rangle$};
  \draw[gtaccent, rounded corners=6pt, fill=gtexbg] (2.8,-0.2) rectangle (4.1,1.0);
  \node at (3.45,0.4) {\textit{John}$_1$};
  \draw[gtaccent, rounded corners=6pt, fill=gtexbg] (4.4,-0.2) rectangle (5.7,1.0);
  \node at (5.05,0.4) {\textit{John}$_2$};
\end{tikzpicture}
\end{keiofig}

Form Copy, then, is a theory of indices, and one that respects (\ref{ex:incl-revised}): it computes the partition from structure (identity, c-command, the phase) rather than writing values onto tokens. Indices are not banned. They are computed, and only under the conditions of (\ref{ex:incl-relations}). Stated as a reconstruction of the textbook notion:

\begin{keybox}
\begin{exe}
\ex \label{ex:coindexation} \textit{Coindexation, reconstructed.} Two expressions are coindexed if, and only if, their tokens have been related by Form Copy, by Agree, or by discharge.
\end{exe}
\end{keybox}

This has a strong consequence. Any dependency that a grammar expresses by coindexation must be built by one of the relations in (\ref{ex:incl-relations}), and it will inherit that relation's locality:

\begin{center}\small
\begin{tabularx}{\linewidth}{@{}P{2.6cm}LP{4.2cm}@{}}
\toprule
Relation-builder & Its locality & The coindexations it builds\\
\midrule
Form Copy & minimality; the phase (\ref{ex:FC-b}--\ref{ex:FC-c}) & chains, control\\
Agree & the locality of Agree: a probe's domain & $\varphi$-sharing, minimal pronouns\\
MERGE, as discharge & the clause (\Hx{3}{ex:discharge}); the other cases of MERGE build sisterhood and c-command, which bind nothing by themselves & quantifier binding of \uth\\
\bottomrule
\end{tabularx}
\end{center}

There is, then, no unbounded coindexation in syntax. Long-distance dependencies that look like coindexation (long-distance \textit{zibun}, discourse anaphora) must be something else: a dependency of the interpretive kind (\S\ref{sec:testbed}), not a syntactic one. The consequence is a prediction, and \S\ref{sec:testbed} tests it: anything bound across a clause boundary, with no Agree or Form Copy path, cannot be syntactic binding.

\subsection{Pronouns}\label{sec:fc-pronouns}

Pronouns pose an obvious problem, since Form Copy needs identity and a pronoun is not identical to its antecedent. There are two routes, both in the literature (Figure~\ref{fig:pronouns}). On \posscitet{Kayne2002} analysis, the antecedent starts out inside the pronoun's phrase and moves out of it, so that antecedent and pronoun are related by Internal Merge after all. On \posscitet{Kratzer2009} analysis, a bound pronoun starts out as a minimal pronoun, without $\varphi$-features of its own, and acquires them from its binder through Agree, so that \textit{his} is simply the spell-out of the features the relation transmitted. Either way, the relation that binds is one of those in (\ref{ex:incl-relations}), and it inherits that relation's locality; neither route needs a value. I do not need to choose between them here.

\begin{keiofig*}{fig:pronouns}{Two routes from pronouns to (\ref{ex:incl-relations}). Left: \citet{Kayne2002}, binding as Internal Merge out of a doubling constituent. Right: \citet{Kratzer2009}, a minimal pronoun valued by Agree.}
\begin{twocol}{0.48}
\centering
\begin{forest}
[TP, SOtree, tiny
  [{\LI{John}}, SOmoved]
  [T$'$ [{T \LI{thinks}}, SOhead]
    [CP [{\LI{that}}, SOlexical]
      [TP [DP [{$\langle$\LI{John}$\rangle$}, SOtrace] [{\LI{he}}, SOlexical]] [{\LI{is clever}}, SOlexical]]]]]
\end{forest}\par\smallskip
\scriptsize antecedent and pronoun related by \textbf{Internal Merge}
\midcol
\centering
\begin{forest}
[TP, SOtree, tiny
  [{\LI{every boy}}, SOlexical]
  [T$'$ [{T$_{[\varphi]}$}, SOhead]
    [VP [{\LI{loves}}, SOlexical]
      [DP [{\textit{pro}$_{[\,\varphi:\ \_\ ]}$ $\to$ \LI{his}}, SOlexical] [{\LI{mother}}, SOlexical]]]]]
\end{forest}\par\smallskip
\scriptsize antecedent and pronoun related by \textbf{Agree}; \textit{his} spells out the transmitted $\varphi$-features
\end{twocol}
\end{keiofig*}

\begin{excursus*}[Kayne's derivation, in steps]
In \posscitet{Kayne2002} derivation the antecedent is first merged \emph{with} the pronoun, as a doubling constituent, and then moves out of it to the matrix subject position (the embedded CP layer is omitted below). A bound pronoun is thus the remnant of a doubling constituent, and Form Copy relates the two occurrences of the antecedent, just as (\ref{ex:FC}) says.

\medskip
\WSsetstep{0}
\begin{derivation}
\begin{twocol}{0.40}
\begin{workspace}[External Merge of \textit{John} and \textit{he}]
  \scriptsize\WSeqTree{\{\LI{John},\LI{he}\},\allowbreak\ \ldots}
    {[DP, SOtree, tiny [{\LI{John}}, SOlexical] [{\LI{he}}, SOlexical]]}
\end{workspace}
\midcol
\begin{workspace}[Internal Merge of \textit{John} with TP]
  \scriptsize\WSeqTree{\{\LI{John},\{\text{T},\ldots\}\}}
    {[TP, SOtree, tiny, flat, SMroot [{\LI{John}}, SMcopy, SOmoved] [T$'$ [{T \LI{thinks that}}, SOhead] [TP [DP [{$\langle$\LI{John}$\rangle$}, SMcopy, SOtrace] [{\LI{he}}, SOlexical]] [{\LI{is clever}}, SOlexical]]]]}
\end{workspace}
\end{twocol}
\end{derivation}
\medskip

The cost should be said openly: the derivation involves an A-movement out of a finite clause, which is why many prefer Kratzer's route. Either is fine for present purposes.
\end{excursus*}


% =====================================================================
\section{The Japanese test bed}\label{sec:testbed}

Japanese is the obvious place to test (\ref{ex:coindexation}), for three reasons. First, the distinction between a formal dependency, which needs c-command, and a looser, interpretive one, which does not, is a Japanese research topic: \citet{Ueyama1998} built a theory of bound-variable anaphora in Japanese on exactly that split. Second, Japanese has four different kinds of element that might be bound, \textit{kare}, \textit{soitu}, \textit{zibun} and \textit{pro}, and they do not behave alike. Third, the prediction of (\ref{ex:coindexation}), that syntactic coindexation is local and everything else interpretive, can be read directly off which readings each element allows, and across which boundaries. The data first: which bound readings does each allow?

\begin{wide}
\begin{exe}
\ex \label{ex:bound}
\begin{xlist}
\ex[*]{\label{ex:kare}
\gll Dono gakusei-mo kare-no sensei-o sonkeisite iru.\\
     which student-\gl{mo} he-\gl{gen} teacher-\gl{acc} respect-\gl{te} be.\gl{npst}\\
\glt (`Every student respects his teacher', on the bound reading)}
\ex[]{\label{ex:soitu}
\gll Dono gakusei-mo soitu-no sensei-o sonkeisite iru.\\
     which student-\gl{mo} that.guy-\gl{gen} teacher-\gl{acc} respect-\gl{te} be.\gl{npst}\\
\glt `Every student respects his teacher.' \src{(bound reading available)}}
\ex[]{\label{ex:zibun}
\gll Dono gakusei-mo zibun-no sensei-o sonkeisite iru.\\
     which student-\gl{mo} self-\gl{gen} teacher-\gl{acc} respect-\gl{te} be.\gl{npst}\\
\glt `Every student respects his own teacher.' \src{(bound reading available)}}
\ex[]{\label{ex:pro}
\gll Dono gakusei-mo [pro sensei-ni homerareta to] itta.\\
     which student-\gl{mo} pro teacher-\gl{dat} praise-\gl{pass}-\gl{pst} \gl{comp} say.\gl{pst}\\
\glt `Every student said that he had been praised by the teacher.' \src{(bound reading available)}}
\end{xlist}
\ex \label{ex:zibun-ld}
\gll Taroo-wa [Hanako-ga zibun-o semeta to] omotte iru.\\
     Taro-\gl{top} Hanako-\gl{nom} self-\gl{acc} blame.\gl{pst} \gl{comp} think-\gl{te} be.\gl{npst}\\
\glt `Taro thinks that Hanako blamed him / herself.' \src{(\textit{zibun} = Taro or Hanako)}
\end{exe}
\end{wide}

The classic observation is that \textit{kare}, although it is the pronoun of the textbooks, resists bound-variable readings \citep{SaitoHoji1983,Hoji1991}, while \textit{soitu}, a demonstrative-based expression, literally `that guy', allows them (\citealt{Hoji1995}). The star on (\ref{ex:kare}) is for the bound reading only: with a name such as \textit{Taroo} as antecedent, \textit{kare} is fine, and so is \textit{soitu}. \textit{Zibun} is bindable, subject-oriented and, as (\ref{ex:zibun-ld}) shows, can be bound across a clause boundary; \textit{pro} is bindable. (\ref{ex:zibun-ld}) is the case that (\ref{ex:coindexation}) must handle: a clause boundary, no Agree path, no Form Copy, and yet a dependency.

A word for readers who work in Language Faculty Science. The judgements here are reported from the literature, and the proper test is the one the programme has developed, in which what matters is whether bound readings are systematically absent where the structural condition is not met, not whether they are always present where it is \citep{Hoji2015}. Put in a table, the informative cells are in the right-hand column:

\begin{center}\small
\begin{tabularx}{\linewidth}{@{}L>{\centering\arraybackslash}p{3.4cm}>{\centering\arraybackslash}p{3.4cm}@{}}
\toprule
 & c-command met & c-command not met\\
\midrule
bound reading reported & required as a control; confirms nothing & \textbf{the refuting cell}\\
bound reading absent & noise is possible & \textbf{the confirming cell}\\
\bottomrule
\end{tabularx}
\end{center}

The prediction of (\ref{ex:coindexation}) for each element is a prediction about that column. For \textit{soitu}, \textit{zibun} and \textit{pro}, the bound reading should be systematically absent without c-command by an admissible relation; for \textit{kare}, it should be absent everywhere. I would welcome correction on any of (\ref{ex:bound})--(\ref{ex:zibun-ld}) on those terms. Here are the four elements side by side:

\begin{center}\small
\begin{tabularx}{\linewidth}{@{}l>{\centering\arraybackslash}X>{\centering\arraybackslash}X>{\centering\arraybackslash}X>{\centering\arraybackslash}X@{}}
\toprule
 & bound by a quantifier & subject-oriented & across a clause boundary & coreferent with a name\\
\midrule
\textit{kare} `he' & \jbad\ (\ref{ex:kare}) & no & (n/a) & \jok\\
\textit{soitu} `that guy' & \jok\ (\ref{ex:soitu}) & no & open & \jok\\
\textit{zibun} `self' & \jok\ (\ref{ex:zibun}) & yes & \jok\ (\ref{ex:zibun-ld}) & \jok\\
\textit{pro} & \jok\ (\ref{ex:pro}) & no & \jok\ (\ref{ex:pro}) & \jok\\
\bottomrule
\end{tabularx}
\end{center}

Two cells carry the argument. One element out of four refuses binding altogether, and the question is whether that is a lexical accident or a consequence of (\ref{ex:incl-revised}). Two elements are bound across a clause boundary, \textit{zibun} in (\ref{ex:zibun-ld}) and \textit{pro} in (\ref{ex:pro}), and the question is whether that is syntactic binding, which (\ref{ex:coindexation}) forbids, or something else. Here is the proposal.

\begin{keybox}
\begin{exe}
\ex \label{ex:kare-anchored} \textit{Kare} is anchored: like the indexed definites of \citet{JenksKonate2022}, it carries an index that must be valued by the context. \textit{Soitu}, \textit{zibun} and \textit{pro} carry no such index, and are bindable.
\end{exe}
\end{keybox}

Indexed definites, in a line: some languages distinguish a definite that picks out a familiar referent by an index the discourse supplies (an \emph{anaphoric} definite) from one that picks out a unique referent by its description. The former carries an index as part of its lexical entry, and the index wants a value. \textit{Kare}, on this view, is that kind of thing in a pronoun's clothing: a pronoun with the semantics of an anaphoric definite, which refers to whoever the context has made salient and can therefore be coreferent.

Why should an anchored expression resist binding? Because, under (\ref{ex:incl-revised}), a binding relation is a relation between tokens, and such relations carry no value. They can identify one token with another; they cannot supply a value to an index that demands one (Figure~\ref{fig:anchored}). The objection that Agree does transmit values, the $\varphi$-values of a minimal pronoun for instance, misses what kind of value \textit{kare}'s index wants: an individual, not a feature value drawn from the lexicon. No token carries an individual as a feature, so no relation of (\ref{ex:incl-relations}) can transmit one; only an assignment can, and that is what `the context' means here. Nor does a lexically present index without a value offend (\ref{ex:incl-values}), which bans values the computation introduces, not slots the lexicon supplies. \textit{Kare}'s index can therefore be valued only by the context, which is to say that \textit{kare} can be coreferent but not bound. \textit{Soitu}, \textit{zibun} and \textit{pro}, which demand no value, can be related to a binder by one of the relations of (\ref{ex:incl-relations}): Agree, in the manner of \posscitet{Kratzer2009} minimal pronouns, or something like Form Copy. On this view the lexical split in (\ref{ex:bound}) is not an arbitrary property of four words, but a direct reflection of the difference between relations and values that (\ref{ex:incl-revised}) draws.

\begin{keiofig*}{fig:anchored}{Why an anchored expression resists binding. A relation of (\ref{ex:incl-relations}) can identify a token with its binder; it cannot supply the value that \textit{kare}'s index demands, which only the context can.}
\begin{tikzpicture}[>=Stealth, every node/.style={font=\scriptsize\rmfamily, align=center}]
  \node[draw=gtaccent, rounded corners=2pt, inner sep=4pt] (q) at (0,0) {\textit{dono gakusei-mo}\\binder};
  \node[draw=gtaccent, rounded corners=2pt, inner sep=4pt] (s) at (6.6,1.3) {\textit{soitu} / \textit{zibun} / \textit{pro}\\no index to value};
  \node[draw=gtalert, rounded corners=2pt, inner sep=4pt] (k) at (6.6,-1.3) {\textit{kare}\\index$_{\,i}$: \underline{needs a value}};
  \node[draw=gtghost, rounded corners=2pt, inner sep=4pt] (c) at (12.6,-1.3) {the context\\supplies a value};
  \draw[->, gtaccent, line width=0.8pt] (q) -- node[above=1pt, sloped, pos=0.45]{a relation of (\ref{ex:incl-relations}) identifies} (s);
  \draw[->, gtalert, line width=0.8pt, dashed] (q) -- node[below=1pt, sloped, pos=0.4]{but cannot supply\\a value} (k);
  \draw[->, gtghost, line width=0.8pt] (c) -- node[above]{coreference} (k);
\end{tikzpicture}
\end{keiofig*}

Two questions remain. First, long-distance \textit{zibun} in (\ref{ex:zibun-ld}), and \textit{pro} in (\ref{ex:pro}), cross a clause boundary, and so cannot be Form Copy, which is phase-bound, or Agree, which is local. They must be the interpretive dependency (Ueyama's other type), or logophoric, as Kuno and Sells argued for long-distance reflexives more generally \citep{Kuno1987,Sells1987}. That is not an embarrassment for (\ref{ex:coindexation}); it is what (\ref{ex:coindexation}) predicts. A dependency with no syntactic builder is not a syntactic dependency. The critical case is a quantified antecedent across a finite boundary, \textit{Dono gakusei-mo [Hanako-ga zibun-o semeta to] omotte iru} `every student thinks that Hanako blamed him', and here the account makes a prediction that separates it from \posscitet{Ueyama1998} formal dependency, which is c-command-based and not clause-bound: if the long-distance dependency runs through the logophoric centre, a quantifier should bind embedded \textit{zibun} and \textit{pro} in attitude complements, where there is such a centre, and not in relative or adjunct clauses, where there is none, whereas a c-command condition alone draws no such line. I do not know the facts well enough to say which way they fall, and it is the first thing I would want tested. Second, if (\ref{ex:kare-anchored}) is right, \textit{kare} should behave like an indexed definite in other respects too, for instance in the conditions under which it can be anaphoric to a quantified antecedent across sentences, where an anaphoric definite needs a discourse referent and a bound variable has none. I leave that open, though it is a judgement that the Language Faculty Science methodology is well placed to design properly.

A third question, which I will not pursue, is whether \textit{soitu}'s demonstrative base is what makes it bindable, by giving it a \emph{descriptive} rather than an indexed semantics. If so, Japanese would be showing the split that Jenks and Konate find among definites within its pronominal system rather than its determiner system.

% =====================================================================
\section{The box has an index too}\label{sec:box-index}

\subsection{\texorpdfstring{u$\theta$}{u-theta} as a thematic index}\label{sec:thematic-index}

In the 2004 paper Kural and I proposed what we called thematic indices: an argument is coindexed with a slot in its predicate's $\theta$-grid, in the spirit of \citet{Stowell1981} and of \posscitet{Higginbotham1985} $\theta$-discharge, and the index does nothing but record that relation. Some background, for readers who did not grow up with it. A verb comes with a $\theta$-grid, \textit{read} $\langle 1, 2\rangle$. Higginbotham's $\theta$-discharge is the saturation of a slot by an argument, marked on the grid, so that \textit{read} $\langle 1, 2^{*}\rangle$ once the object is in; Stowell had the argument carry the slot's index. Figure~\ref{fig:thematic-2004} gives the 2004 picture. The index $i$ relates the argument to slot~2 and does nothing else: it has no meaning (\ref{ex:prop-meaning}), it is not drawn from a list (\ref{ex:prop-unbounded}), and it is local (\ref{ex:prop-distance}). It was already an index of the admissible kind, before we had a name for the kind. What was missing was a grammar in which the slot is a token and the relation is built by an operation.

\begin{keiofig}{fig:thematic-2004}{Thematic indices \citep{KuralTsoulas2004}: the index $i$ relates the argument to slot~2 of the verb's grid, and does nothing else.}
\begin{forest}
[VP, SOtree, slim
  [V [{\LI{read} $\langle 1, 2_{i}\rangle$}, SOlexical]]
  [{DP$_{i}$} [{\LI{every book}}, SOlexical]]]
\end{forest}
\end{keiofig}

Now look again at the box of Lecture~3. The fissioned \uth\ is a slot of the verb's $\theta$-grid, made visible as a token by Self-Merge, and discharge relates the boxed quantifier to that token. That is the 2004 thematic index, rebuilt under (\ref{ex:incl-revised}), and each clause of (\ref{ex:incl-revised}) is met by one ingredient of the box (Figure~\ref{fig:box-thematic} and the table beside it). The token is the verb's own (\ref{ex:incl-tokens}): a singleton $\{\mathrm{V}\}$, with a lexical source, nothing added (strictly, a singleton that MERGE builds from the verb's token; I use `token' for it as Lecture~3 does). The relation is built at discharge (\ref{ex:incl-relations}), which on the line taken in Lecture~3 is an application of MERGE. And nothing carries a value (\ref{ex:incl-values}): the relation identifies, and $u$ has no content beyond being the verb's token. The variable of QR, in other words, is an index of the admissible kind, and the debt of Lecture~3 is paid.

\begin{keiofig*}{fig:box-thematic}{The box's \uth\ is the thematic index, rebuilt under (\ref{ex:incl-revised}): token, relation, no value.}
\begin{twocol}{0.40}
\centering
\begin{forest}
[CP, SOtree, tiny
  [QP, SObox]
  [C$'$ [C, SOhead]
    [TP [\ldots]
      [VP [V, SOhead [{\LI{read}}, SOlexical]] [{V$_{\uthr}$ $= u$}, SOtrace]]]]]
\end{forest}
\midcol
\footnotesize
\begin{tabular}{@{}P{2.2cm}P{5.6cm}@{}}
\toprule
Clause & How the box meets it\\
\midrule
(\ref{ex:incl-tokens}) tokens & the token is \textbf{the verb's own}: a singleton $\{\mathrm{V}\}$, lexical source, nothing added\\
(\ref{ex:incl-relations}) relations & the relation is \textbf{built at discharge}, External Merge of the boxed singleton at the phase edge\\
(\ref{ex:incl-values}) no new values & \textbf{nothing carries a value}: the relation identifies, and $u$ is just the verb's token\\
\bottomrule
\end{tabular}
\end{twocol}
\end{keiofig*}

Its interpretation can be stated directly:

\begin{exe}
\ex \label{ex:discharge-sem} \textit{Discharge, interpreted.} If a boxed QP is discharged at the edge of XP, and discharge relates it to the token $u$ of \uth\ in XP, then, where $g$ is an assignment from tokens (not numerals) to individuals:\\[4pt]
$\sem{[\,\mathrm{QP}\ \mathrm{XP}\,]}^{g} \;=\; \sem{\mathrm{QP}}^{g}\big(\lambda x.\ \sem{\mathrm{XP}}^{g[u \mapsto x]}\big)$
\end{exe}

In words, (\ref{ex:discharge-sem}) says that the quantifier applies to the property of being an $x$ such that XP holds when the token $u$ stands for $x$. This is the $\lambda$ of \Hx{3}{ex:CI-B}, with the index replaced by a token. The modified assignment $g[u \mapsto x]$ changes its value at a token the syntax already has, so nothing has been added to interpret it. Compare Predicate Abstraction (\Hx{1}{ex:PA}): it is the same rule, with the trace's index $n$ replaced by the token $u$, and the binder index itself replaced by the discharge relation. Figure~\ref{fig:discharge-tree} walks through the composition: TP denotes a proposition with $u$ free, discharge abstracts over $u$, and the quantifier applies. Nothing in it should look new to a semanticist except the domain of $g$, and that is the point.

\begin{keiofig}{fig:discharge-tree}{Discharge interpreted, (\ref{ex:discharge-sem}): TP with $u$ free; abstraction over $u$ at the edge; the quantifier applies.}
\begin{forest}
[{CP}, SOtree, xslim
  [QP, SObox]
  [{TP} [\ldots [{$u$}, SOtrace]]]]
\end{forest}\\[4pt]
{\footnotesize
\begin{tabular}{@{}r@{\ \ }l@{}}
CP: & $\sem{\mathrm{QP}}^{g}(\lambda x.\,\sem{\mathrm{TP}}^{g[u\mapsto x]})$\\
TP: & $\sem{\mathrm{TP}}^{g}$, with $u$ free\\
$u$: & $g(u)$
\end{tabular}}
\end{keiofig}

One could also do without assignments altogether, since the discharge relation tells the semantics which argument slot to abstract over, which is what the combinators of a variable-free semantics do \citep{Jacobson1999}. On that formulation the open slot is passed up by function composition (the Geach rule, $\mathbf{g}$) and closed by the quantifier when it applies, and $u$ is not a variable but a \emph{marked slot}. The verb with its \uth\ token exposed is $\lambda y.\lambda x.\ \textit{read}(x,y)$ with $y$ flagged as the slot that discharge will target, and the \vs P built above it is $\lambda y.\ \textit{a child read } y$, of type $\langle e,t\rangle$, which is exactly the observation of Lecture~3 that what reaches TP is a predicate (\Hs{3}{sec:proposal}). For present purposes the two formulations come to the same thing. The assignment version makes the comparison with Predicate Abstraction transparent; the combinator version makes it transparent that no variable name was ever needed. There is a point here for the categorial grammarians. What CCG and Jacobson do with composition, the box does with a token and a relation. The difference is not in the semantics but in where the locality of the dependency comes from: the type system there, the phase here (\Hs{3}{sec:box-cooper}).

\subsection{The catch: percolation}\label{sec:percolation}

In the version of the box theory I presented at ICKL, \uth\ was passed up the tree until a phase head discharged it, much like the \textsc{slash} feature of GPSG and HPSG (Figure~\ref{fig:percolation}). The idea came from \citet{NeelemanReinhart1998}, by way of my own use of it for scrambling (\citealt{Tsoulas1999}; \Hs{3}{sec:scrambling}). That is a problem. As \S\ref{sec:three-properties} recalled, \citet[114]{Chomsky2000keio} gives the percolation of indices as his example of indices that are not innocuous, and \uth\ percolation looks exactly like that. Percolation is the one thing Chomsky names, and the earlier version did it.

\begin{keiofig}{fig:percolation}{The earlier picture: \uth\ written on every node between the token and the phase head.}
\begin{forest}
[{CP}, SOtree, xslim
  [C, SOhead]
  [{TP$_{\uthr}$\ {\color{gtalert}$\uparrow$}} [{\LI{a child}}, SOlexical]
    [{\vs P$_{\uthr}$\ {\color{gtalert}$\uparrow$}} [\vs, SOhead]
      [{VP$_{\uthr}$\ {\color{gtalert}$\uparrow$}} [V, SOhead] [{V$_{\uthr}$}, SOtrace]]]]]
\end{forest}
\end{keiofig}

\begin{excursus}[\textsc{slash}, for the analogy]
In GPSG and HPSG an unbounded dependency is a feature, \textsc{slash}, whose value is the category of the gap \citep{GazdarEtAl1985,PollardSag1994}. It is introduced at the gap, passed up by a principle (the Foot Feature Principle, the Nonlocal Feature Principle), and bound off by the filler:

\begin{center}
\begin{forest}
[{S}, SOtree, tiny
  [{NP} [{\LI{which book}}, SOlexical]]
  [{S[\textsc{slash} NP]} [{\LI{did}}, SOlexical]
    [{S[\textsc{slash} NP]} [{\LI{Taroo}}, SOlexical]
      [{VP[\textsc{slash} NP]} [{\LI{buy}}, SOlexical] [{NP[\textsc{slash} NP]}, SOtrace]]]]]
\end{forest}
\end{center}

The dependency is local at every step and unbounded overall, and its value is a \emph{category}, not an identity: two gaps of the same category have the same \textsc{slash} value. That is why HPSG also needs indices. That earlier \uth\ was a \textsc{slash} feature with identity built in, the slot of one particular verb. That is what makes it look like percolation of an index, and it is what the third exit below removes.
\end{excursus}

There are three ways out:

\begin{exe}
\ex \label{ex:exits}
\begin{xlist}
\ex \label{ex:exit-sharing} \textit{Feature sharing.} The \uth\ is a single instance shared by the nodes it passes through \citep{PesetskyTorrego2007}. This is allowed by (\ref{ex:incl-relations}), since Agree builds sharing relations, but it needs a chain of Agree relations up the spine with no independent motivation.
\ex \label{ex:exit-upward} \textit{Upward Agree.} The quantifier, once at the edge, enters an Agree relation with the \uth\ below it under c-command, so that the two share one instance \citep{Zeijlstra2012}. Nothing percolates, but the account is committed to a particular direction of Agree, and it requires the quantifier to c-command the token at discharge.
\ex \label{ex:exit-search} \textit{No percolation.} The phase head finds the undischarged \uth\ in its complement by Minimal Search, just as C[Q] consults the box, and discharge relates the quantifier introduced at the edge to the token that the search returned.
\end{xlist}
\end{exe}

\begin{keiofig*}{fig:exits}{The three exits of (\ref{ex:exits}), as mini-trees.}
\begin{minipage}[t]{0.31\linewidth}\centering
\begin{forest}
[C, SOtree, tiny [{T\ \tikz[baseline]\node[draw=gtaccent,circle,inner sep=1pt]{1};}]
  [{\vs\ \tikz[baseline]\node[draw=gtaccent,circle,inner sep=1pt]{1};} [{V\ \tikz[baseline]\node[draw=gtaccent,circle,inner sep=1pt]{1};}]]]
\end{forest}\par\smallskip
\scriptsize (\ref{ex:exit-sharing}) one instance, many nodes
\end{minipage}\hfill
\begin{minipage}[t]{0.31\linewidth}\centering
\begin{forest}
[CP, SOtree, tiny [QP, SObox] [TP [\ldots [{$u$\ $\uparrow$}, SOtrace]]]]
\end{forest}\par\smallskip
\scriptsize (\ref{ex:exit-upward}) the quantifier values $u$ from above
\end{minipage}\hfill
\begin{minipage}[t]{0.31\linewidth}\centering
\begin{forest}
[CP, SOtree, tiny [QP, SObox] [{C\ {\color{gtaccent}search\,$\searrow$}}, SOhead] [TP [\ldots [{$u$}, SOtrace]]]]
\end{forest}\par\smallskip
\scriptsize (\ref{ex:exit-search}) C searches, finds $u$, discharges
\end{minipage}
\end{keiofig*}

Which of the three would one take, and why?\toexercise{exits} In brief, the first is cheap but unmotivated, and it predicts that every head on the spine between token and edge takes part in the dependency, since each must Agree; the second is a commitment about the direction of Agree; the third uses nothing the box does not already use, and it predicts Determinacy effects wherever two undischarged tokens fall in one search domain (Figure~\ref{fig:two-tokens} below).

My own preference is (\ref{ex:exit-search}), for two reasons and a consequence. It uses nothing that the box does not already use: phase heads that search and consult. And it gives the duality of semantics its cleanest form, as the next paragraph shows. The consequence is a reinterpretation of the notation: on (\ref{ex:exit-search}), the notation TP\uthr\ in \Hx{3}{ex:deriv-child} is not a label on TP but a record of what Minimal Search from C finds, and the token $u$ in (\ref{ex:discharge-sem}) is simply the token the search returned (Figure~\ref{fig:minsearch}): nothing is written on TP, \vs P or VP.

\begin{keiofig}{fig:minsearch}{The preferred exit, (\ref{ex:exit-search}): C finds $u$ by Minimal Search; nothing is written on the nodes in between.}
\begin{forest}
[CP, SOtree, tiny
  [QP, SObox]
  [C$'$ [{C\ {\color{gtaccent}Minimal Search}}, SOhead]
    [TP [{\LI{a child}}, SOlexical]
      [\vs P [\vs, SOhead]
        [VP [{V \LI{read}}, SOhead] [{$u$\ {\color{gtaccent}found}}, SOtrace]]]]]]
\end{forest}
\end{keiofig}

The duality then reads as follows. The $\theta$-relation is the verb's, made by Self-Merge, a case of External Merge (of the verb with itself). The scope relation is made by External Merge too, of the boxed singleton at the phase edge, triggered by the search that found $u$. The third exit and the External Merge reading of discharge adopted in Lecture~3 are the same picture seen from two sides: Minimal Search is how discharge finds its token, and External Merge is what it then does. The revised condition accordingly lists MERGE (including discharge), Agree and Form Copy, and nothing else.

The third exit leaves one thing to be settled, which the derivation of Lecture~3 already raised (\Hs{3}{sec:worksheet}). If both arguments could be boxed, which is the symmetric alternative Lecture~3 leaves open (\Hs{3}{sec:open}), there would be two undischarged tokens in the complement of C, and Minimal Search would find the subject's first (Figure~\ref{fig:two-tokens}). The order in which a language discharges its boxed quantifiers would then be exactly what conjecture~3 (\Hx{3}{ex:conj-rigidity}) has to fix: the pairing problem in a box with two members. On the restriction adopted in Lecture~3, however, boxing applies to internal arguments only. The subject is never boxed, there is one token, and the question changes shape: rigidity then reads the order of externalisation, not of search.

\begin{keiofig}{fig:two-tokens}{Two undischarged tokens in the complement of C, if both arguments are boxed: the search finds the subject's first.}
\begin{forest}
[CP, SOtree, xslim
  [{C}, SOhead]
  [TP
    [\vs P [{$u_{\mathrm{subj}}$}, SOtrace]
      [\vs$'$ [\vs, SOhead] [VP [V, SOhead] [{$u_{\mathrm{obj}}$}, SOtrace]]]]]]
\end{forest}
\end{keiofig}

% =====================================================================
\section{An open problem: weak crossover}\label{sec:wco}

\begin{wide}
\begin{exe}
\ex \label{ex:wco-en}
\begin{xlist}
\ex[]{\label{ex:wco-en-a} Every boy loves his mother. \src{(bound reading available)}}
\ex[??]{\label{ex:wco-en-b} His mother loves every boy. \src{(bound reading degraded)}}
\ex[??]{\label{ex:wco-en-c} Who does his mother love? \src{(bound reading degraded)}}
\end{xlist}
\ex \label{ex:wco-jp}
\begin{xlist}
\ex[?*]{\label{ex:wco-jp-a}
\gll [Soitu-no hahaoya]-ga dono gakusei-mo aisite iru.\\
     that.guy-\gl{gen} mother-\gl{nom} which student-\gl{mo} love-\gl{te} be.\gl{npst}\\
\glt (`His mother loves every student', on the bound reading)}
\ex[]{\label{ex:wco-jp-b}
\gll Dono gakusei-mo [soitu-no hahaoya]-ga aisite iru.\\
     which student-\gl{mo} that.guy-\gl{gen} mother-\gl{nom} love-\gl{te} be.\gl{npst}\\
\glt (bound reading available)}
\end{xlist}
\end{exe}
\end{wide}

The name comes from (\ref{ex:wco-en-c}): the \textit{wh}-phrase has crossed over the pronoun on its way to the front, and the pronoun cannot then be bound by it. The effect is `weak' because it is milder than the \emph{strong} crossover of \textit{*Who$_1$ does he$_1$ love $t_1$?}, where the pronoun c-commands the gap. (\ref{ex:wco-en-b}) is the same configuration without any visible movement. On the QR theory \textit{every boy} raises covertly over \textit{his}, and the bound reading is degraded in the same way. That is where the phenomenon touches us: a grammar in which nothing moves covertly has to say what goes wrong in (\ref{ex:wco-en-b}). In (\ref{ex:wco-jp}) the pronoun is \textit{soitu}, not \textit{kare}, since \textit{kare} cannot be bound anyway (\S\ref{sec:testbed}), and Japanese adds a twist: clause-internal scrambling of the object over the subject repairs the effect, as in (\ref{ex:wco-jp-b}) \citep{SaitoHoji1983,Hoji1985}. There the quantifier is pronounced above the pronoun, and the bound reading returns. Any account of weak crossover for Japanese has to make the pronounced position matter.

Weak crossover \citep{Postal1971,Wasow1972} has had many explanations, and it is instructive to set them out with what each is stated over:

\begin{center}\small
\begin{tabularx}{\linewidth}{@{}P{3.4cm}LP{3.2cm}@{}}
\toprule
Account & The condition, in a line & Stated over\\
\midrule
Leftness \citep{Chomsky1976} & a variable cannot be the antecedent of a pronoun to its left & the LF variable (trace)\\
Bijection Principle \citep{KoopmanSportiche1982} & one operator binds exactly one variable; the pronoun and the trace would make two & operator--variable chains\\
A-binding \citep{Reinhart1983} & a pronoun is bound only from an argument position that c-commands it & A- vs $\bar{\text{A}}$-positions\\
Crossover situations \citep{Buring2004} & a configuration of binder, pronoun and gap defined on the structure & movement and its launching site\\
\bottomrule
\end{tabularx}
\end{center}

All of them are stated over movement, or over traces and the positions they occupy. The point is not that they are wrong but that their vocabulary is not available to us: in a grammar with no covert movement and no traces, weak crossover cannot be a constraint on movement. The question is whether it can fall out of the geometry of consultation. Of the four, Reinhart's is the one that survives translation. There are also accounts of crossover that never used the vocabulary of movement, in variable-free semantics, in continuation semantics and in LFG; \S\ref{sec:nb-nesting} sets them beside the suggestion that follows.

Here is a suggestion, which I put forward as a question rather than an answer. A boxed quantifier has three relations: to its \uth\ token, where it is $\theta$-related; to the position where it is pronounced, its Case site or an edge; and to its discharge site. Suppose that a pronoun can be bound by a quantifier only through one of the quantifier's tokens that the syntax can see, that the relation requires c-command from that token, and that it can proceed only from the first two of the three relations, never from the discharge site, the pronounced position counting only where it is distinct from the discharge site (Figure~\ref{fig:wco-relations}).

\begin{keiofig}{fig:wco-relations}{The three relations of a boxed quantifier. On the suggestion, binding of a pronoun proceeds only from the first two.}
\begin{tikzpicture}[>=Stealth, every node/.style={font=\scriptsize\rmfamily, align=center}]
  \node[draw=minsyn@box, very thick, inner sep=3pt] (q) at (0,1.5) {QP};
  \node[draw=gtaccent, rounded corners=2pt, inner sep=3pt] (th) at (-2.6,-0.2) {$u$: the $\theta$-token\\{\color{gtaccent}can bind}};
  \node[draw=gtaccent, rounded corners=2pt, inner sep=3pt] (pr) at (0,-0.2) {pronounced position\\{\color{gtaccent}can bind}};
  \node[draw=gtghost, rounded corners=2pt, inner sep=3pt] (ds) at (2.6,-0.2) {discharge site\\{\color{gtghost}cannot bind}};
  \draw[->, gtaccent] (q) -- (th); \draw[->, gtaccent] (q) -- (pr); \draw[->, gtghost, dashed] (q) -- (ds);
\end{tikzpicture}
\end{keiofig}

Then (\ref{ex:wco-en-a}) is fine, since the subject, which is never boxed, is externally merged in Spec-\vs P and pronounced in Spec-TP, and both positions c-command \textit{his}; (\ref{ex:wco-en-b}) is out, since the object's \uth\ and its Case site are both low, and the discharge site, which is high and would c-command the pronoun, does not count; and (\ref{ex:wco-jp-b}) is fine, since the scrambled object is pronounced at the clause edge, above \textit{soitu} (Figure~\ref{fig:wco-applied}). (\ref{ex:wco-en-c}), the \textit{wh}-case, is (\ref{ex:wco-en-b}) with Q in the box, with one difference the clause just added marks: \textit{who} is pronounced at the site where C[Q] discharges it, so its pronounced position is its discharge site and does not count, whereas the scrambled object of (\ref{ex:wco-jp-b}) is pronounced at an edge distinct from the site of its discharge, and that position does. Three cases, one condition. This is \posscitet{Reinhart1983} insight, binding from an argument position, restated for a grammar in which the only argument position a boxed quantifier has is the verb's token (or, with Case, the position the token's Case is checked in).

\begin{keiofig*}{fig:wco-applied}{The suggestion applied to (\ref{ex:wco-en-a}), (\ref{ex:wco-en-b}) and (\ref{ex:wco-jp-b}): binding from the $\theta$-token or the pronounced position, never from the discharge site.}
\begin{minipage}[t]{0.31\linewidth}\centering\footnotesize
(\ref{ex:wco-en-a}) \textit{every boy loves his mother}\par\smallskip
\begin{forest}
[TP, SOtree, tiny [{\LI{every boy}}, SOlexical] [VP [{\LI{loves}}, SOlexical] [{\LI{his mother}}, SOlexical]]]
\end{forest}\par\smallskip
\scriptsize the unboxed subject, in Spec-\vs P and Spec-TP, c-commands \textit{his}: \jok
\end{minipage}\hfill
\begin{minipage}[t]{0.31\linewidth}\centering\footnotesize
(\ref{ex:wco-en-b}) \textit{his mother loves every boy}\par\smallskip
\begin{forest}
[CP, SOtree, tiny [QP, SObox] [TP [{\LI{his mother}}, SOlexical] [VP [{\LI{loves}}, SOlexical] [{$u_{\mathrm{obj}}$}, SOtrace]]]]
\end{forest}\par\smallskip
\scriptsize the object's \uth\ and Case site are both low; the discharge site is high but cannot bind: \jbad
\end{minipage}\hfill
\begin{minipage}[t]{0.31\linewidth}\centering\footnotesize
(\ref{ex:wco-jp-b}) \textit{dono gakusei-mo soitu-no hahaoya-ga aisite iru}\par\smallskip
\begin{forest}
[CP, SOtree, tiny [{\LI{dono gakusei-mo}}, SOmoved] [TP [{\LI{soitu-no hahaoya-ga}}, SOlexical] [VP [{$u_{\mathrm{obj}}$}, SOtrace] [{\LI{aisite iru}}, SOlexical]]]]
\end{forest}\par\smallskip
\scriptsize the scrambled object is pronounced at the clause edge, above \textit{soitu}: \jok
\end{minipage}
\end{keiofig*}

The suggestion has an obvious problem, and the clause just added makes it worse rather than better. Long-distance scrambling also pronounces the object high, at a matrix edge distinct from the embedded site of its discharge, so the suggestion as stated predicts that it binds; yet it is reported not to repair weak crossover \citep{Saito1992keio}, so the pronounced position of a long-distance scrambled phrase must not count. Perhaps binding from a pronounced position requires that it be in the clause in which the \uth\ is discharged, which would tie the repair to conjecture~2 (\Hx{3}{ex:conj-mirror}): the long-distance scrambled phrase scopes low, and binds low. I do not know yet. A test suggests itself, which speakers could run: if what counts is the pronounced position and not the edge as such, clause-internal scrambling of the object to a position between the subject and a higher adverb should repair the effect just as well as scrambling to the clause edge.

One consolation: QR never cured weak crossover either. In (\ref{ex:wco-en-b}) QR puts \textit{every boy} above \textit{his}, and some further condition was always needed to stop the bound reading, so the box is no worse off than the theory it replaces. It may even be better off, since it offers a principled reason why one of the quantifier's relations should not count for binding: the discharge site is a scope relation, not an argument position. On the External Merge reading of discharge the boxed singleton is merged at the edge, but as a scope-taker and not as an argument, which is exactly the distinction that Reinhart's condition draws.

% =====================================================================
\section{The architecture revisited, and three open questions}\label{sec:revisited}

\subsection{What the argument changed}\label{sec:changed}

Lecture~1 began with five ideas (\Hs{1}{sec:architecture}). Here is what has happened to each of them:

\begin{center}\small
\begin{tabularx}{\linewidth}{@{}P{2.6cm}P{4.2cm}L@{}}
\toprule
Idea & What Lecture 1 assumed & What the argument changed\\
\midrule
(a) MERGE and its conditions & one operation, under Minimal Yield, Determinacy, Minimal Search and Form Copy & Self-Merge boxes and fissions within Minimal Yield; ambiguity is two determinate derivations; pairing in a box with two members is still open\\
(b) Inclusiveness and No Tampering & no indices, no traces & \textbf{revised}: tokens, relations, no new values; an index is a relation, and \uth\ is one\\
(c) The duality of semantics & External Merge for $\theta$, Internal Merge for scope & $\theta$ by External Merge, of the verb with itself (Self-Merge), assigned low; scope by External Merge of the box at the phase edge (discharge), or by Internal Merge at edges (scrambling)\\
(d) Cyclicity without LF & Transfer by phase; `covert' as lower-copy pronunciation & nothing moves covertly; the locality of QR is the clause-boundedness of discharge; \textit{wh} is consulted across phases\\
(e) Spines and externalisation & force on different heads; variation in externalisation & pronounced retrieval entails rigidity, though not conversely (\Hs{3}{sec:rigid}); a boxed element is pronounced at its highest relation visible to the sensorimotor interface (\Hx{3}{ex:externalisation})\\
\bottomrule
\end{tabularx}
\end{center}

Row (c) deserves a word. Had discharge been a third relation of its own, beside External and Internal Merge, (c) would have had to read: $\theta$ is assigned low by the verb, through Self-Merge, and scope is fixed by discharge, a third relation. On the External Merge reading of discharge adopted in Lecture~3 the duality survives in its original two-way form, with both halves of the boxed derivation built by MERGE.

What did not change is just as worth stating, especially for the minimalists: the workspace, Minimal Yield, Determinacy, phases, Form Copy, and the claim that variation is in externalisation. The box added two uses of MERGE the system already had, Self-Merge and External Merge of a singleton at a phase edge, and no new relation; the revised Inclusiveness condition added nothing, and removed a prohibition that was doing no work.

One question, I suggested in Lecture~2, organises the whole parameter hierarchy (\Hs{2}{sec:hierarchy}): at which step of the derivation, storage or retrieval, is force assigned, and is that step pronounced?

\begin{center}\small
\begin{tabularx}{\linewidth}{@{}P{2.8cm}P{3.6cm}P{2.6cm}L@{}}
\toprule
 & Force assigned at & Pronounced? & Scope\\
\midrule
English & storage (\textit{every}) & n/a: retrieval is silent & flexible\\
Japanese, Korean & retrieval (the particle) & yes & rigid\\
Mandarin & retrieval, on the spine (\textit{dou}) & yes & rigid in simple clauses \citep{Huang1982,AounLi1993}\\
Straits Salish, Mohawk & on the clause only \citep{Jelinek1995,Baker1995} & yes & not reported, as far as I know\\
\bottomrule
\end{tabularx}
\end{center}

English puts force on the stored item (\textit{every}) and pronounces nothing at retrieval, and its scope is flexible. Japanese and Korean pronounce the retrieving operator, the particle, and their scope is rigid. Mandarin pronounces retrieval on the spine (\textit{dou}), and Mandarin, too, has long been described as scope-rigid in simple clauses. Straits Salish and Mohawk, if Jelinek and Baker are right, assign force only on the clause.

That gives the typologists a prediction worth testing, which I state in one direction only: languages that fix force at retrieval, and pronounce that step, should be scope-rigid. What counts as pronouncing the step is an exponent of the retrieving operator, \textit{dou} on the spine or \textit{mo} on the nominal, and not an exponent of force on the stored item, which is what English \textit{every} is; that is why English does not bear on the prediction, and why the Salish and Mohawk row cannot be filled until their scope has been studied. The converse does not follow, since rigidity has other sources (German, and Japanese \textit{zen'in}, \Hx{3}{ex:zenin}), so a rigid language need not pronounce the step. What would refute the conjecture is a flexible language that pronounces it. The stronger claim, that pronounced retrieval explains rigidity, is withdrawn (\Hs{3}{sec:rigid}).

\subsection{Three open questions}\label{sec:open-questions}

\textit{Phrase or feature?} Does the pronounced retrieval step of Korean and Japanese (\Hx{3}{ex:externalised-box}) show that what is boxed is a feature rather than a maximal projection? I suggested phrase at storage, feature at retrieval (\Hs{3}{sec:phrase-feature}), but the question is not settled, and Korean and Japanese are where it will be.

\textit{Syntactic or semantic store?} The answer from the 2\,$\times$\,2 (\Hs{3}{sec:keio2x2}), if conjecture~1 holds, is both, divided by relation: association of a particle with an indeterminate runs on a semantic store, and licensing and scope on a syntactic one. That is less elegant than the box everywhere, and, I think, more likely to be true.

\textit{Free or nested retrieval?} \citeauthor{Cooper1983}'s retrieval is free; \citeauthor{Keller1988}'s is nested \citep{Cooper1983,Keller1988}; the box, if phases do the work, should be nested for free. Mineshima's example is the Japanese probe:

\begin{wide}
\begin{exe}
\ex \label{ex:nesting-probe}
\gll [Dono sensei-ga susumeta] dono hon-mo \ldots{}\\
     which teacher-\gl{nom} recommend.\gl{pst} which book-\gl{mo} {}\\
\glt `every book that any teacher recommended \ldots{}', or `for every teacher $x$ and every book $y$ that $x$ recommended \ldots{}' \src{(after \citealt{Mineshima2005})}
\end{exe}
\end{wide}

Two indeterminates, one particle. Does \textit{mo} bind both, as it binds both indeterminates of \Hx{2}{ex:dare-nani-temo}, \textit{[dare-ga nani-o itte]-mo}? And can the inner one be caught by a different, lower particle while the outer one is caught by \textit{mo}? Figure~\ref{fig:nesting} draws the two parses. If the inner indeterminate can be caught only by the outermost particle, retrieval is nested; if either order is possible, it is free. The answer decides between free and nested retrieval, and so between a store whose nesting is stipulated and one whose nesting is inherited.

\begin{keiofig*}{fig:nesting}{The nesting probe (\ref{ex:nesting-probe}). Left: \textit{mo} binds both indeterminates. Right: is the inner indeterminate caught lower, by another particle?}
\begin{twocol}{0.48}
\centering
\begin{forest}
[{DP-\textit{mo}}, SOtree, tiny
  [CP [{\LI{dono sensei-ga}}, SOlexical, name=in1] [{\LI{susumeta}}, SOlexical]]
  [{NP} [{\LI{dono hon}}, SOlexical, name=out1]]
  [{\textit{mo}}, SOhead, name=mo1]]
\draw[->, gtaccent, dashed] (mo1.south) .. controls +(0,-0.9) and +(0.2,-0.7) .. (out1.south);
\draw[->, gtaccent, dashed] (mo1.south) .. controls +(0,-1.5) and +(0,-1.3) .. (in1.south);
\end{forest}\par\smallskip
\scriptsize\color{gtaccent} \textit{mo} binds both indeterminates
\midcol
\centering
\begin{forest}
[{DP-\textit{mo}}, SOtree, tiny
  [{CP-\textit{X}}, name=CPX [{\LI{dono sensei-ga}}, SOlexical, name=in2] [{\LI{susumeta}}, SOlexical]]
  [{NP} [{\LI{dono hon}}, SOlexical, name=out2]]
  [{\textit{mo}}, SOhead, name=mo2]]
\draw[->, gtaccent, dashed] (mo2.south) .. controls +(0,-0.9) and +(0.2,-0.7) .. (out2.south);
\draw[->, gtalert, dashed] (CPX.south) .. controls +(0,-1.3) and +(0,-1.0) .. (in2.south);
\end{forest}\par\smallskip
\scriptsize\color{gtalert} the inner indeterminate caught lower, by another particle \textit{X}?
\end{twocol}
\end{keiofig*}

% =====================================================================
\section{The box among its neighbours: convergences and differences}\label{sec:neighbours}

The lectures have compared the proposal with other ways of doing the same job as they went along: with Cooper's store in \Hp{3}, with HPSG, glue semantics, CCG and variable-free semantics in the excursions of \Hp{1}, and with the semantic frameworks in those of \Hp{2}. This section gathers the comparisons, and adds some neighbours that the lectures passed over. It is organised by question and not by framework, because the frameworks do not line up as wholes: on each question the box turns out to have different neighbours. The questions are where the variable comes from, what delays the quantifier, where locality comes from, and what handles nesting and crossover. Before them, a word about the question I am no longer asking.

\subsection{Two labels, and a better map}\label{sec:labels}

In the announcement of these lectures, and at the meeting of the International Circle of Korean Linguistics in July (ICKL), I put the whole comparison as one question: is a monocyclic grammar with a store simply a monostratal one under another name? The lectures as given carried the answer as their thesis, monostratal in mechanism and not in explanation. I now think that the question was badly put, and that the answer, though I stand by what it says about Cooper's store (\Hx{3}{ex:locality-claim}), is a remark about where the proposal stands and not the proposal. It is worth saying why.

The two labels first. What I called monocyclic is better called \emph{single-cycle}: a grammar in which the overt and the covert cycles of the Y-model have collapsed into one, so that there is one derivation, handed to the interfaces phase by phase, and no level of LF at which operations invisible to the sensorimotor interface could apply \citep{Chomsky2004,Chomsky2005}. That was idea (d) of \Hs{1}{sec:cyclicity}. A \emph{monostratal} grammar is one with a single level of syntactic representation and no derivation relating one level to another; the term set GPSG, and after it HPSG and categorial grammar, against the grammars with several strata \citep{GazdarEtAl1985,PollardSag1994}.

There are three reasons why the labels decide little. First, they do not divide the field where the question needs them to. CCG is monostratal and derivational at once \citep{Steedman2000}. LFG has no derivation but two parallel syntactic structures, so that it is monostratal only in the sense that neither structure is derived from the other \citep{Bresnan2001}. \posscitet{Brody1995} minimalism is representational, with a single level. And the minimalist grammars of \citet{Stabler1997} are standardly studied through their derivation trees, one structure from which form and meaning are both read \citep{Kobele2006}. Second, a store is no longer what the monostratal frameworks have in common. HPSG has for the most part moved from Cooper storage to underspecified representations, glue semantics never had a store, and CCG has neither store nor movement (\Hs{1}{sec:box-preview}, and the excursions that follow it). Third, the question invites an answer about frameworks where the lectures have evidence only about constructions. The result of \Hp{3} is not that a store in the syntax is right and a store in the semantics wrong. It is that Japanese uses both, divided by relation (\Hx{3}{ex:conj1}).

A better map is \posscitet{Jacobson2002}, which sorts theories of scope by the way syntax and semantics are organised and not by the number of strata \citep[see also][3--5]{BarkerJacobson2007}. On one side are the directly compositional theories, in which every expression the syntax builds has a meaning as it is built: Hendriks's type-shifting and Cooper's store in the strong form, Montague's Quantifying-In in a weaker one. Next to them Jacobson puts a kind of theory that she compares to Generative Semantics, in which rules build structures and interpret them in tandem, and further rules adjust the structure before it is pronounced. On the other side is what she calls the surface-to-LF view, the view of \citet{May1977} and of \citet{HeimKratzer1998}: the syntax first builds a structure with no interpretation, then adjusts it into a Logical Form, and only then does the semantics apply. QR is the paradigm of the last kind.

The box theory, it seems to me, is of the middle kind, and that is the first convergence. A boxed quantifier is interpreted where it is discharged, at a phase, in the one cycle there is. What is adjusted is not its interpretation but its pronunciation, which Case fixes lower down (\Hx{3}{ex:externalisation}). QR says `pronounced here, interpreted later and higher'; the box says `interpreted here, pronounced lower'. It is, if one wants a historical name for it, the Quantifier Lowering of Generative Semantics \citep{McCawley1970,Lakoff1971} with externalisation in the place of the lowering transformation; \Hs{3}{sec:scrambling} said the same of \posscitet{BoskovicTakahashi1998} lowering of scrambled phrases. The single-cycle architecture, in other words, has moved to the directly compositional side of Jacobson's map. That is what was right in the old question, and it can be said without the word monostratal.

\pIIneedspace{8.5cm}
\subsection{Where the variable comes from}\label{sec:nb-variable}

This was the question of the lectures, and the neighbours give it six answers besides mine.

\begin{center}\small
\begin{tabularx}{\linewidth}{@{}P{3.2cm}P{4.6cm}L@{}}
\toprule
The variable is & Where & What stands in the argument position\\
\midrule
an index, valued by an assignment & QR with Trace Conversion; Cooper's store; the store of HPSG, and the variables of MRS & an indexed placeholder, $x_i$, over which the quantifier abstracts\\
a discourse referent & DRT and dynamic semantics (\Hs{2}{sec:unselective}) & a referent, which can be bound from outside its syntactic scope\\
a set of alternatives & Hamblin semantics (\Hx{2}{ex:hamblin}) & nothing: the alternatives are closed by the first operator that meets them\\
nothing at all & CCG; variable-free semantics & nothing: the open argument is handed on by the combinators\\
a hypothesis & glue semantics; type-logical grammar; Convergent Grammar & an assumption, made and later discharged in a proof\\
a proof term & Dependent Type Semantics & a term built from the proof of the context\\
a relation between tokens & the box & the verb's own token, \uth, which discharge relates to the quantifier\\
\bottomrule
\end{tabularx}
\end{center}

Three things about the table. First, the box's nearest neighbour is not the store. Cooper's placeholder is an indexed variable, and the store of HPSG inherits it; that was the placeholder row of \Hs{3}{sec:box-cooper}. The nearest neighbour is the hypothesis of the logical grammars. In glue semantics a quantifier is scoped by assuming a meaning for its argument position, building the clause on that assumption, and then withdrawing it: the hypothesis is discharged, and on the meaning side the discharge is abstraction over the variable that the hypothesis introduced \citep{Asudeh2023}. Type-logical grammar scopes quantifiers by the same hypothetical reasoning \citep{Moortgat1997,KubotaLevine2020}, and in \posscitet{Pollard2011} Convergent Grammar variables correspond to hypotheses by design. I said in \Hp{1} that this was the nicest analogue in any framework of what this lecture would claim. It is more than an analogue. The token \uth\ is a hypothesis: something the derivation puts in the argument position, builds on, and discharges at a phase edge, where the semantics reads the discharge as abstraction (\ref{ex:discharge-sem}). That the same word turns up is half an accident, since mine comes from \posscitet{Higginbotham1985} $\theta$-discharge and theirs from natural deduction, but the two uses name one move. That the variable is not an item is, then, common ground between the box and the logical grammars, and I claim no novelty for it.

Second, what is new is where the hypothesis lives. In glue semantics it is a step of a proof conducted over meanings, and the syntax knows nothing of it. In the box it is a syntactic object with a lexical source, the singleton that Self-Merge builds from the verb, and because it is in the syntax, syntactic relations can see it. Case is valued at it, which is why the quantifier is pronounced there (\Hx{3}{ex:case}); and a phase head searches for it, which is why its discharge is bounded (\Hx{3}{ex:discharge}). The logical grammars get the variable right and have nothing to say about why a quantifier is pronounced where its variable is. That is the part of the old QR picture that a token inherits and a hypothesis does not. The cost is on my side: a hypothesis is free, and the token has to be paid for with fission and with the stipulations that \Hs{3}{sec:objections} owned.

Third, the frameworks that do without a placeholder, CCG and variable-free semantics, are making a different bet: that the dependency can be coded in the types of the expressions that lie on its path. The semantics of discharge can be stated their way, with the Geach rule handing the open slot upwards (\S\ref{sec:thematic-index}), so on the semantic side the two are intertranslatable. They part in the syntax. On the exit I prefer, (\ref{ex:exit-search}), nothing lies on the path.

\subsection{What delays the quantifier}\label{sec:nb-delay}

Every theory of scope needs a way of interpreting a quantifier later, and so higher, than the place where its argument position is. Outside the grammar of LF nobody does it with a movement made for the purpose, and that is the convergence the old question was after. The devices differ.

\begin{center}\small
\begin{tabularx}{\linewidth}{@{}P{4.0cm}L@{}}
\toprule
The device & Where\\
\midrule
a second, covert cycle & QR at LF \citep{May1977,May1985,HeimKratzer1998}\\
a store & \citet{Cooper1983}; HPSG \citep{PollardSag1994,PollardYoo1998}; Convergent Grammar \citep{Pollard2011}; the semantics of minimalist grammars \citep{Kobele2006,Kobele2018}; the box\\
an underspecified description & MRS \citep{CopestakeEtAl2005} and the other underspecified semantics of HPSG \citep{RichterSailer1999,KoenigRichter2021}; the semantics of tree-adjoining grammar \citep{KallmeyerRomero2008}\\
the order of steps in a proof & glue semantics \citep{Dalrymple2001,Asudeh2023}; type-logical grammar \citep{Moortgat1997,KubotaLevine2020}\\
the order of combination & CCG, with type-raising, composition and Skolem terms \citep{Steedman2012}\\
a shift in the type of the predicate, or of the quantifier's context & \citet{Hendriks1993}; continuations \citep{Barker2002,BarkerShan2014}\\
\bottomrule
\end{tabularx}
\end{center}

On the store I owe an acknowledgement that the lectures did not make. Kobele gave Stabler's minimalist grammars a compositional semantics computed on derivations, in which the scope interactions of quantifiers are tied to movement and what is in transit is kept in a store \citep{Kobele2006,Kobele2012}; and \citet{Kobele2018} shows that Cooper storage has the structure of what is called a graded applicative functor, and that where the store is finite, the case that `corresponds to ideas in current transformational approaches to syntax', interpretation is a homomorphism on a regular set of derivation trees. A syntactic Cooper store is, in that work, a result about minimalist derivations, and it was there before the box was. What separates it from the proposal of \Hp{3} is what is stored and what is left behind. In Kobele's semantics storage interprets movement, so that constraints on movement become constraints on scope, which is the point of doing it in the syntax. In the box theory the quantifier is never in its argument position and never moves: the store takes the place of the movement, and what is left behind is not the trace of the stored phrase but a token of the verb. \posscitet{Pollard2011} Convergent Grammar is the same idea from the side of logical grammar, with Cooper's store reconstructed proof-theoretically: each stored operator is kept with the variable it will bind, and nothing happens in the syntax when it is retrieved. Glue semantics, for its part, has been connected to a Minimalist syntax directly, as an alternative to a Logical Form derived by covert movement \citep{Gotham2018}; of all the neighbours it is the one that has come closest to the single-cycle grammar from the other side.

The difference that matters on this row is not between the box and the store. It is between resolving scope and leaving it unresolved. The box, the stores, CCG and glue semantics deliver the readings one at a time. MRS and its relatives deliver one description that all the readings satisfy, on the principle that \citet[1264]{BorsleyMueller2021} put as `everything that is not ruled out works'. For a flexible language the two come to the same thing. For a rigid one they do not: a grammar that underspecifies has to add a constraint that removes a reading, and a grammar that resolves has to say why a derivation is missing. Conjecture~3 is an attempt at the second (\Hx{3}{ex:conj-rigidity}).

On rigidity the neighbours converge with conjecture~3 more closely than I had expected. The conjecture says that in a rigid language the order of discharge copies an order that the derivation already has, the $\theta$-hierarchy or the order of externalisation. \citet{Frey1993} derived the scope of German quantifiers from their surface positions together with their base positions. \citet{Kiss2001}, in HPSG and without traces, derives it from two orders that his title calls relational and configurational, the ranking of the arguments and their surface arrangement. \citet{BobaljikWurmbrand2012} state a preference for matching orders at the two interfaces. And in LFG the same two dimensions, rank on the hierarchy of grammatical functions and linear order, are what constrain the binding of a pronoun by an operator, with languages differing in which of them they impose \citep{Bresnan1995,DalrympleKaplanKing2001}. From several directions, then, one arrives at two orders, an argument hierarchy and a surface order, as what scope and binding are read off where they are rigid; and \citet{Gotham2019} shows that rigidity can be built into glue semantics as well, for German and for the frozen scope of the double object construction, by keeping track in the logic of the order in which quantifiers apply and fixing it in the lexical entries. I take the convergence as evidence that the two orders are real. The question it leaves for all of us is why some languages read scope off them and others do not, which conjecture~3 does not answer either (\Hs{3}{sec:rigid}).

\subsection{Where locality comes from}\label{sec:nb-locality}

This is the row on which \Hp{3} made its claim, and the comparison obliges me to state the claim more narrowly than I did at ICKL.

\begin{center}\small
\begin{tabularx}{\linewidth}{@{}P{3.6cm}LP{3.3cm}@{}}
\toprule
 & What bounds the scope of a quantifier & Stated, or inherited?\\
\midrule
Cooper's store; the store of HPSG & nothing in the mechanism: conditions on retrieval, or on the sites of retrieval & stated\\
MRS and its relatives & constraints on how a description may be resolved & stated\\
type-shifting & which lexical items may shift & stated\\
Hamblin alternatives & the first operator that can close the set & inherited, as a path effect\\
glue semantics & nothing in the logic; a constraint on the path to the scope level, or modes of the logic that a premise cannot cross & stated, or inherited from the logic\\
continuations & delimiters of the context a quantifier can take; a refinement of the semantic types & stated in the types\\
CCG & the combinatory rules and the lexical categories & inherited from the derivation\\
tree-adjoining grammar & the elementary tree of the verb, and the order of composition & inherited in part\\
minimalist grammars with a store & the constraints on movement & inherited from the syntax\\
QR & Subjacency, the ECP, phases; Scope Economy & inherited from movement\\
the box & the phase: \uth\ must be discharged before its clause is transferred & inherited from the syntax\\
\bottomrule
\end{tabularx}
\end{center}

Against a store in the semantics the claim of \Hp{3} stands. Cooper's mechanism lets a stored quantifier be retrieved at any node of the right type, and island conditions have to be added to it, as Cooper himself began to do \citep{Cooper1983}; \citet{RuysWinter2011} make the point about the standard scope mechanisms together, storage and type-shifting among them, that every one of them has to be made subject to the island constraints. The same holds of the propagation of alternatives, which is bounded by the first operator that can close them and by nothing else. A store in the syntax inherits a boundary from the syntax. That is the difference \Hp{3} pressed, and it is the content of (\Hx{3}{ex:locality-claim}).

Against the field as a whole the claim does not stand in the form I gave it, because several neighbours inherit locality too. CCG takes the scope of a universal from the derivation, so that it goes where syntactic combination goes \citep{Steedman2012}; I said as much in \Hs{3}{sec:cooper}. In tree-adjoining grammar the quantifier is composed with the elementary tree of its verb, and that tree sets the upper limit of its scope at the tensed clause \citep{KallmeyerRomero2008}. In the semantics of minimalist grammars, tying scope to movement brings the constraints on movement `into the purview of the grammar' \citep[103]{Kobele2012}. And glue semantics, which is by its practitioners' own account `towards the unconstrained end of the spectrum' \citep[147]{Gotham2021}, can be given scope islands in two ways. One states them, as a constraint on the path along which a quantifier finds its scope. The other builds them into the logic, as modes that a premise cannot cross, so that they follow from the rules of inference \citep{Gotham2021}. \citet[547]{Gotham2018} had already defined a scope island as a constituent whose premises must all be used up inside it, and remarked that the definition is `conceptually similar to the notion of a ``phase''\,'. `The syntax explains what a store has to assume' is true of the box and Cooper's store. It is just as true, with `the derivation' or `the logic' in the place of `the syntax', of several of the neighbours.

What distinguishes the frameworks, then, is not whether they inherit locality but what they inherit it from, and that has empirical content. A boundary inherited from phases is a domain, and it is the same for every quantifier. A boundary set by the first closing operator is a path. A boundary coded in types or modes may differ from one scope-taker to the next. The first two are the two localities of \Hs{2}{sec:localities}, and the open cell of the 2\,$\times$\,2 is the test between them (\Hs{3}{sec:open-cell}). The third is tested by a different kind of fact, and here the evidence is against me. \citet{Barker2022} shows that tensed clauses are not scope islands across the board. A universal can take scope out of the tensed complement of \textit{ensure}, as in \textit{An accomplice ensured that every prisoner escaped}, where the accomplices may vary with the prisoners, but not out of the complement of \textit{think}; and \textit{no prisoner} is trapped where \textit{every prisoner} escapes \citep[the examples are those of][151]{Gotham2021}. Barker's generalisation, the Scope Island Subset Constraint, is that the scope-takers trapped by one island are always a subset of those trapped by another, and a logic with ordered modes delivers it directly \citep{Gotham2021}.

These facts are a problem for the box as it stands, and I would rather say so here than have it said for me. I allowed in \Hs{1}{sec:QR}, with \citet{FarkasGiannakidou1996}, that the rule that scope is bounded by the clause has exceptions. But (\Hx{3}{ex:discharge}) draws the boundary at the last phase of the clause, and draws it for every quantifier alike. The phase account has a resource for the first of Barker's facts and none yet for the second. The resource is the size of the complement. If the complement of \textit{ensure} lacks the phasal layer that the complement of \textit{think} has, then the first phase to be completed above the quantifier is in the matrix clause, and discharge there is what (\Hx{3}{ex:discharge}) allows. \citet{Cecchetto2004} made phases the source of scope islands, and Barker himself sketches a syntactic account in which the strength of an island depends on the size of the complement that the embedding predicate selects. For the second fact one would need different quantifiers to be discharged at different heads of the spine, which is Beghelli and Stowell's picture, and which \Hs{2}{sec:spine} set aside. Whether the gradience is in the grammar at all is a further question: \citet{Wurmbrand2018} takes it to be a cost of processing, which is where these lectures left the preference for surface scope (\Hs{3}{sec:deriv1}). The question belongs with the open ones of \Hs{3}{sec:open}, and of all of them it is the one on which the neighbours bear most directly.

\subsection{Nesting and crossover}\label{sec:nb-nesting}

\Hp{3} left the nesting row of its comparison open, and \S\ref{sec:wco} left weak crossover as a problem. On both, the neighbours have results that the box has to match.

There are two facts about nesting. The first is the one that \Hs{3}{sec:cooper} drew: the inner quantifier must not be left with nothing to bind, and \posscitet{Keller1988} nested stores secure that by making the inner store travel with its container. The second is \posscitet{Larson1985}. In \textit{Two politicians spy on someone from every city}, \textit{every city} may take scope over \textit{someone}, but no third quantifier may come between the two: the sentence has no reading on which the cities vary, the politicians vary with the cities, and the person spied on varies with the politicians (see \citealt{MayBale2006}, who weigh a counterexample of May's and conclude that the generalisation probably holds). Larson's own account was a store: he made Cooper's store a push-down one, from which quantifiers are retrieved in the reverse of the order in which they were stored. \citet{May1985} accounts for it by adjoining the inner quantifier to the containing nominal instead of extracting it. \citet{JoshiKallmeyerRomero2008} derive it in tree-adjoining grammar: the two quantifiers are composed with each other before either is composed with the verb, so that they form a set which no other quantifier can split. \citet{Kobele2010} derives the readings of inverse linking with no semantic operations beyond function application and composition. In glue semantics the inverse reading is a choice of scope level for the premise of the inner quantifier \citep[148--149]{Gotham2021}. For the box, Larson's generalisation would follow if the inner quantifier had to be discharged at the same phase edge as its container, and next to it, and I do not yet see what would force that. So the open question of \Hs{3}{sec:open} has a sharper form than it had there: not whether phases give nesting, but whether they give Larson's generalisation.

On crossover, the four accounts in the table of \S\ref{sec:wco} are all stated over movement or traces, and I said that their vocabulary was not available to us. Three neighbours never had that vocabulary and have accounts of crossover all the same. In variable-free semantics the asymmetry is built into the binding rule \textsf{z}, which lets a higher argument of a predicate bind into a lower one and not the reverse \citep{Jacobson1999}. \citet{ShanBarker2006} derive crossover and superiority from a single assumption, that expressions are evaluated from left to right, in a semantics with continuations. And LFG, which has no traces, states the condition as prominence of the operator over the pronoun on one or both of two dimensions, rank on the hierarchy of grammatical functions and linear order, with languages differing in which they require \citep{Bresnan1995,DalrympleKaplanKing2001}.

The suggestion of \S\ref{sec:wco} is closest to the last of these. Binding from the $\theta$-token is binding by rank. Binding from the pronounced position is binding by linear order, and the repair by scrambling in (\ref{ex:wco-jp-b}) shows that Japanese allows it. The exclusion of the discharge site says that scope is not a third dimension. LFG reached the two dimensions from data like the Japanese; the box reaches them from the relations that a boxed quantifier enters, and they are the two orders that conjecture~3 reads (\S\ref{sec:nb-delay}). If that convergence is not an accident, the account to test against it is Shan and Barker's, on which it is the order of evaluation that matters; the two come apart wherever rank and linear order do.

Resumption belongs here as well. \citet{Asudeh2012} analyses the resumptive pronouns of Irish, Hebrew, Swedish and Vata in glue semantics as ordinary pronouns: the pronoun contributes a resource like any other, the dependency has no use for it, and a further resource, lexically supplied, removes it. The dependency is not interpreted through the pronoun. That is the point \S\ref{sec:conclusion} makes about resumptives, in another idiom: a resumptive pronounces a token related to the operator, and not the variable.

\subsection{Further afield}\label{sec:bridges}

Three bridges lead further off, and I mark them without crossing them.

\textit{To philosophy.} \posscitet{Fine2007} coordination, \posscitet{Recanati2012} mental files and \posscitet{Heim1982} file cards are three ways of saying that sameness of reference is a relation, which is the thought of \S\ref{sec:revised}.

\textit{To type theory.} \posscitet{Bekki2014} Dependent Type Semantics handles anaphora with proof terms instead of indices, a third architecture worth setting beside storage and relations.

\begin{excursus}[Dependent Type Semantics, in a line]
In DTS a sentence denotes a \emph{type}, and a discourse is true if the type is inhabited by a proof \citep{Bekki2014}. An indefinite is a $\Sigma$-type, a pair of a witness and a proof that it satisfies the restriction; a pronoun is an underspecified term, which the hearer resolves by proof search for an object of the right type in the context:

\begin{center}\footnotesize
\textit{A man walks. He whistles.}\\[3pt]
$u : \Sigma x{:}\mathbf{e}.\ (\mathrm{man}(x) \times \mathrm{walk}(x))$\\[2pt]
$\mathrm{whistle}(@\,:\,\mathbf{e}) \ \leadsto\ \mathrm{whistle}(\pi_1 u)$
\end{center}

Read: $u$ is a pair, an entity and a proof that it is a man who walks; the pronoun is a placeholder $@$ of type $\mathbf{e}$, resolved to $\pi_1 u$, the entity half of that pair. There is no index here: \textit{he} is resolved to $\pi_1 u$, the first projection of the proof of the first sentence. Anaphora is a relation to a proof term, computed, not stipulated. That is a further answer to the question of the seminar: the variable comes from the \emph{proof}, as here it comes from the \emph{token}. DTS, too, is a relations-based theory, with proof terms where we have tokens. Where it and the box differ is locality: proof search ranges over the whole context, discharge over the phase. The Japanese data of \S\ref{sec:testbed} are, in effect, a test of whether anaphora resolution and quantifier binding have the same locality; on (\ref{ex:coindexation}), they do not.
\end{excursus}

\textit{To computation.} Attention in a transformer \citep{VaswaniEtAl2017} keeps every token retrievable at every layer, which is storage with no phases. What does it lack that the box has?

\begin{excursus}[Attention as a store with no phases]
The analogy is deliberately loose. Attention is storage: every token is kept retrievable at every layer, and a query retrieves by similarity rather than by structure. It is storage with no phases: nothing is ever transferred, so nothing ever becomes inaccessible, and locality, where the models show it, is learned, not built in.

\begin{center}
\begin{tikzpicture}[>=Stealth, every node/.style={font=\scriptsize\rmfamily}]
  \foreach \i/\w in {0/every, 1/book, 2/a, 3/child, 4/read} {
    \node[draw=gtghost, rounded corners=2pt, inner sep=2pt, minimum width=0.9cm] (t\i) at (\i*1.1,0) {\w};
  }
  \node[draw=gtaccent, rounded corners=2pt, inner sep=2pt, minimum width=0.9cm, fill=gtexbg] (q) at (4.4,1.5) {read, layer $n$};
  \foreach \i in {0,1,2,3,4} { \draw[->, gtghost] (t\i.north) -- (q.south); }
  \node[text=gtghost] at (2.2,-0.55) {every token, every layer, no boundary};
\end{tikzpicture}
\end{center}

What the box has that attention lacks is therefore exactly what the four lectures argued for: locality inherited from the architecture, domain effects that do not have to be learned from the data. Whether the models learn them anyway is an empirical question, and a good one to argue about.
\end{excursus}

\subsection{What the comparison shows}\label{sec:nb-summary}

Let me draw it together. There are three convergences. Nobody outside the grammar of LF moves a quantifier in order to scope it, and the single-cycle architecture, with the box, has joined the theories in which a quantifier is interpreted as the structure is built. The variable is not an item for the logical grammars any more than it is for me: their hypothesis and my token are one move. And where scope and binding are rigid, several frameworks read them off the same two orders.

There are three differences, and each is a question of fact. The hypothesis lives in the syntax, as a token that Case and Minimal Search can see, or it lives in a proof over meanings; the first predicts that a quantifier is pronounced where its variable is, and the second predicts nothing about pronunciation. Locality is inherited from phases, from the combinatorics, from the modes of a logic, or from the first operator on a path; these predict domain effects, effects of the derivation, effects that vary with the scope-taker, and path effects, and the open cell of the 2\,$\times$\,2 and Barker's islands are where they can be told apart. And scope is resolved by the grammar or left to be resolved; the difference shows in what each has to say when a reading is missing.

The question I asked at ICKL has, then, a short answer and a long one. The short answer was given in \Hp{3}: in mechanism the box is a store, and Cooper would recognise every step of it. The long answer is this section. The box is one of a family of devices that have taken the place of the covert cycle. What distinguishes it within the family is that its placeholder is a syntactic object and its boundary a phase. Whether those are virtues is decided by Japanese, and not by the architecture.

% =====================================================================
\section{Conclusion: the variable is a relation}\label{sec:conclusion}

Let me put the argument of the four lectures together, in five steps. The first three belong to Lectures 1 and 2, and nothing since has changed them; the last two belong to Lectures 3 and 4.

First, natural language quantifies by operator and variable, yet, as far as I can tell, the restriction of a quantifier is always pronounced, its force often, its scope site sometimes, and its variable never; and a grammar that respects Inclusiveness has no means of writing the variable either (\Hp{1}). Second, the variable does not come from movement: copies need Trace Conversion, which adds a determiner and a variable that were never in the workspace (\Hs{1}{sec:problem}; \Hs{3}{sec:wh-to-QR}). Third, force is separable from restriction, and the clausal spine offers a series of places where it can be fixed, which Japanese and Korean pronounce (\Hs{1}{sec:nominal}; \Hp{2}). Fourth, scope comes from storage. A quantifier is boxed by Self-Merge, pronounced where it is Case-licensed, and interpreted where it discharges the verb's own $\theta$-feature. The box is Cooper's store in mechanism, but its locality is inherited from phases, and it predicts domain effects that a semantic store has to stipulate; Japanese shows both kinds of locality, divided by relation, and its rigidity and the scope of its scrambled phrases follow from the order and the locality of discharge (\Hp{3}). Fifth, the variable is a relation over tokens. The \uth\ that the box discharges is an index honestly described, a relation that carries no value, admitted by a revised Inclusiveness condition that excludes values and admits relations; Form Copy is a theory of such indices; and \textit{kare}, the one Japanese pronoun that needs a value, is the one that cannot be bound (this lecture). The fifth step is the one that answers the question with which the seminar opened.

\begin{keybox}
\textit{The answer.} The variable is not an item but a relation between tokens. Scope is storage in the syntax, and its locality is the locality of phases.
\end{keybox}

The two sentences are the two halves of the argument: the second is the result of Lecture~3, and the first that of this lecture, with the debt of the variable paid. How the answer stands to the other ways of doing the same job was the subject of \S\ref{sec:neighbours}, and it is a remark about the proposal, not a part of it. What the Japanese data added is two localities, divided by relation (association on a semantic store, licensing and scope on a syntactic one), and a language that pronounces the retrieval step and is rigid because of the order it reads.

And we can now go back to the sentence with which the seminar began. If the restriction is always pronounced and the variable never, it is not because the variable is a silent item that every language happens to leave unpronounced. It is because the variable is not an item at all, but a relation, and relations are not the kind of thing that gets pronounced. What is pronounced, on this view, is always a token: the restriction, the force, sometimes the scope site. What is never pronounced is the relation between the quantifier and the verb's token, because a relation has no phonology to give. And that is why a grammar that cannot write $x$ can still quantify. It never needed to write $x$. It needed a token to relate the quantifier to, and the verb supplies one.

The claim needs to be stated with some care, since bound pronouns and resumptives are the obvious objection to it. `The variable is never pronounced' (\Hx{1}{ex:generalisation}) is a claim about the foot of the quantifier's own chain: the position the quantifier is $\theta$-related to, the verb's token. That is what the box says is a relation, and that is what is never pronounced. Three things are pronounced and are not counterexamples. Bound pronouns (\textit{every boy loves his mother}) pronounce a token related to the quantifier by Agree or Form Copy (\S\ref{sec:fc-pronouns}), not the foot of its chain. Minimal pronouns \citep{Kratzer2009} are the same case, with the $\varphi$-features supplied by the relation. And the \textit{wh}-resumptives of Vata, Swedish and Irish pronounce a token at the foot of a \textit{wh}-dependency, which is a Form Copy relation, not a discharge relation (see \citealt{Asudeh2012} for all three, and \S\ref{sec:nb-nesting}). None of these is to be dismissed, but none is a counterexample either: each pronounces a \emph{token related to} the operator, which is exactly what (\ref{ex:incl-relations}) admits, and none pronounces the relation itself.

% =====================================================================
\section{Questions and readings}\label{sec:readings}

\begin{questions}
\begin{enumerate}
\item Is the line between values and relations stable? Is the $k$ of a token $\langle \mathrm{LI}, k\rangle$ a value in disguise?
\item Can Form Copy be extended to bound pronouns, and if so by which route, \posscitet{Kayne2002} or \posscitet{Kratzer2009}?
\item Is \textit{kare} anchored, and does it behave like an indexed definite elsewhere?
\item Which of the three exits in (\ref{ex:exits}) is preferable, and why?
\item Can weak crossover fall out of the geometry of consultation?
\item Of the sources of locality in \S\ref{sec:nb-locality}, which make the same predictions as phases, and where do they come apart?
\end{enumerate}
\end{questions}

The first question is the hardest, and the one a minimalist is most likely to ask. My answer is the one given in \S\ref{sec:tokens}: $k$ is a value that never does any work, and the revised condition bans values that do.

\subsection*{Key readings}

\citet{CollinsStabler2016}; \citet{Chomsky2021}; \citet{JenksKonate2022}. For the comparison of \S\ref{sec:neighbours}: \citet{Jacobson2002}; \citet{Kobele2018}; \citet{Gotham2021}.

\subsection*{Further readings}

\citet{KuralTsoulas2004}; \citet{GrafAbner2012}; \citet{PesetskyTorrego2007}; \citet{Ueyama1998}; \citet{Hoji1995}; \citet{Kratzer2009}; \citet{MarcolliChomskyBerwick2025}; \citet{Zeijlstra2012}; \citet{Fine2007}. On the neighbours: \citet{RuysWinter2011}; \citet{BorsleyMueller2021}; \citet{Mueller2013}; \citet{Asudeh2023}; \citet{BarkerShan2014}; \citet{Barker2022}; \citet{JoshiKallmeyerRomero2008}.


% ======================================================================
%  APPENDIX: THE EXERCISES
% ======================================================================
\keioappendix

These are the exercises that go with the four lectures, in the order in which the text reaches them. They are here for readers who want to work through them, on their own or in a class. The sentences are repeated with the numbers they have in the text, and each number is a link to it; the foot of each exercise links back to the place in the text where it belongs, and where what it turns up is discussed.

\section{Lecture 1: A grammar that cannot write \texorpdfstring{\emph{x}}{x}}

\begin{exercise}{vote}{Three judgements before any theory}{1}{sec:vote}
Before anything else, three judgements. For each of the three sentences below the question is the same: can the books vary with the readers? That is, can the sentence be true in a situation in which every book was read by someone, but not by the same someone in each case?

\begin{exe}
\exagain{1}{ex:someone} Someone read every book.
\exagain{1}{ex:dareka-dono}
\gll Dare-ka-ga dono hon-mo yonda.\\
     who-\gl{ka}-\gl{nom} which book-\gl{mo} read.\gl{pst}\\
\glt `Someone read every book.'
\exagain{1}{ex:dono-dareka}
\gll Dono hon-mo dare-ka-ga yonda.\\
     which book-\gl{mo} who-\gl{ka}-\gl{nom} read.\gl{pst}\\
\glt `Every book, someone read.'
\end{exe}

\textit{Context for the varying reading.} The Mita reading circle drew up a list of twelve novels for the summer. By September every novel on the list had been read by at least one member (Hanako read \textit{Kokoro}, Taroo read \textit{Sansiroo}, Kenta read \textit{Yukiguni}), but nobody had read all twelve. The convenor is now reporting to the department. Is (\ref{H1:ex:someone}) a true report of what happened? Is (\ref{H1:ex:dareka-dono})? Is (\ref{H1:ex:dono-dareka})?

The vignette is there so that the varying reading has a concrete situation to be true in. Scope judgements given out of the blue are unreliable, and out of context the surface reading is usually preferred; so take a moment with the context before answering. Lecture~3 returns to the three sentences (\Hs{3}{sec:rigid}).
\end{exercise}

\begin{exercise}{spot}{Spot the variable}{1}{sec:spot}
Take each of the sentences below and mark four things: the \emph{restriction} (what the quantification is over), the \emph{force} (what kind of quantification it is), the \emph{scope site} (where its scope is fixed), and the \emph{variable} (the position it binds). For each, say whether it is pronounced, and if so, by which morpheme. Disagreement is expected, particularly over the force of (\ref{H1:ex:spot-sannin}), and the disagreement is part of the lesson; the text's own answers follow the sentences in \Hs{1}{sec:spot}.

\begin{exe}
\exagain{1}{ex:spot-every} Every student left.
\exagain{1}{ex:spot-sannin}
\gll Gakusei-ga san-nin kita.\\
     student-\gl{nom} three-\gl{cl} come.\gl{pst}\\
\glt `Three students came.'
\exagain{1}{ex:spot-dareka}
\gll Dare-ka-ga kita.\\
     who-\gl{ka}-\gl{nom} come.\gl{pst}\\
\glt `Someone came.'
\exagain{1}{ex:spot-dono}
\gll Dono gakusei-mo kita.\\
     which student-\gl{mo} come.\gl{pst}\\
\glt `Every student came.'
\exagain{1}{ex:spot-zutu}
\gll Kodomo-tati-ga ringo-o mittu-zutu tabeta.\\
     child-\gl{pl}-\gl{nom} apple-\gl{acc} three.\gl{cl}-\gl{zutu} eat.\gl{pst}\\
\glt `The children ate three apples each.'
\exagain{1}{ex:spot-who}
\gll Dare-ga kita no?\\
     who-\gl{nom} come.\gl{pst} \gl{q}\\
\glt `Who came?'
\end{exe}

\begin{center}
\renewcommand{\arraystretch}{1.6}
\begin{tabularx}{\linewidth}{|l|L|L|L|L|}
\hline
 & \textbf{Restriction} & \textbf{Force} & \textbf{Scope site} & \textbf{Variable}\\ \hline
(\ref{H1:ex:spot-every}) & & & & \\ \hline
(\ref{H1:ex:spot-sannin}) & & & & \\ \hline
(\ref{H1:ex:spot-dareka}) & & & & \\ \hline
(\ref{H1:ex:spot-dono}) & & & & \\ \hline
(\ref{H1:ex:spot-zutu}) & & & & \\ \hline
(\ref{H1:ex:spot-who}) & & & & \\ \hline
\end{tabularx}
\end{center}
\end{exercise}

\section{Lecture 2: Quantification on the clausal spine}

\begin{exercise}{sheet}{The judgement sheet}{2}{sec:aims}
Twelve Japanese sentences, each with a reading described beneath it. For each, say whether the sentence has that reading: yes, no, or not sure. Judge the reading described, not whether you yourself would say the sentence, and if you cannot get the reading at all, say so. The sheet is meant for native speakers of Japanese, though anyone may fill it in; if you send me your judgements, say what your first language is. The items cover the cells of a small design that \Hs{2}{sec:wedge} explains and \Hs{3}{sec:keio2x2} uses. Items 1--4 and 11 bear on the two Japanese universals of \Hs{1}{sec:jp-universals}; items 5--7 and 10 on the association of \textit{ka} and \textit{mo} with an indeterminate; items 8, 9 and 12 on the licensing of \textit{nani-mo} by negation, of which item 12 is the open cell. Most of the items are from the literature; items 3, 4, 7, 11 and 12 are constructed after it, and they are the ones on which I would most like to be corrected.

\begin{exe}
\exi{{1}}
\gll Dare-ka-ga dono hon-mo yonda.\\
     who-\gl{ka}-\gl{nom} which book-\gl{mo} read.\gl{pst}\\
\glt `Someone read every book.' \src{\Hx{1}{ex:dareka-dono}}
\keioreading{the books vary with the readers: every book was read by someone, but not by the same someone in each case.}
\exi{{2}}
\gll Dono hon-mo dare-ka-ga yonda.\\
     which book-\gl{mo} who-\gl{ka}-\gl{nom} read.\gl{pst}\\
\glt `Every book, someone read.' \src{\Hx{1}{ex:dono-dareka}}
\keioreading{as in item 1: the books vary with the readers.}
\exi{{3}}
\gll Hanako-ga [[dono gakusei-mo kaita] ronbun]-o yonda.\\
     Hanako-\gl{nom} which student-\gl{mo} write.\gl{pst} paper-\gl{acc} read.\gl{pst}\\
\glt `Hanako read the paper that every student wrote.' \src{(after \citealt{Mineshima2005})}
\keioreading{the papers vary with the students: for each student, Hanako read the paper that student wrote.}
\exi{{4}}
\gll Hanako-ga [[subete-no gakusei-ga kaita] ronbun]-o yonda.\\
     Hanako-\gl{nom} all-\gl{gen} student-\gl{nom} write.\gl{pst} paper-\gl{acc} read.\gl{pst}\\
\glt `Hanako read the paper(s) that all the students wrote.' \src{(after \citealt{Mineshima2005})}
\keioreading{as in item 3: for each student, Hanako read the paper that student wrote.}
\exi{{5}}
\gll [John ka Bill]-ga nani-o katta no?\\
     John or Bill-\gl{nom} what-\gl{acc} buy.\gl{pst} \gl{q}\\
\glt `What did John or Bill buy?' \src{\Hx{2}{ex:disj-a}}
\keioreading{a question about what was bought, to be answered by naming the thing: `a book'.}
\exi{{6}}
\gll Nani-o [John ka Bill]-ga katta no?\\
     what-\gl{acc} John or Bill-\gl{nom} buy.\gl{pst} \gl{q}\\
\glt `What did John or Bill buy?' \src{\Hx{2}{ex:disj-b}}
\keioreading{as in item 5.}
\exi{{7}}
\gll Hanako-wa [Taroo-ga nani-o katta to] omotte iru no?\\
     Hanako-\gl{top} Taro-\gl{nom} what-\gl{acc} buy.\gl{pst} \gl{comp} think-\gl{te} be.\gl{npst} \gl{q}\\
\glt `What does Hanako think that Taro bought?'
\keioreading{a question about what Hanako thinks Taro bought, the question particle being in the main clause and the \textit{wh}-word in the embedded one.}
\exi{{8}}
\gll Taroo-wa [Hanako-ga nani-mo tabe-nakat-ta to] itta.\\
     Taro-\gl{top} Hanako-\gl{nom} what-\gl{mo} eat-\gl{neg}-\gl{pst} \gl{comp} say.\gl{pst}\\
\glt `Taro said that Hanako ate nothing.' \src{\Hx{2}{ex:npi-jp-a}}
\keioreading{as glossed: what Taro said is that Hanako ate nothing.}
\exi{{9}}
\gll Taroo-wa [Hanako-ga nani-mo tabeta to] iwa-nakat-ta.\\
     Taro-\gl{top} Hanako-\gl{nom} what-\gl{mo} eat.\gl{pst} \gl{comp} say-\gl{neg}-\gl{pst}\\
\glt (`Taro didn't say that Hanako ate anything.') \src{\Hx{2}{ex:npi-jp-b}}
\keioreading{Taro did not say, of anything, that Hanako ate it: the negation is in the main clause and \textit{nani-mo} in the embedded one.}
\exi{{10}}
\gll [[Dono gakusei-ga syootaisita] sensei]-mo odotta.\\
     which student-\gl{nom} invite.\gl{pst} teacher-\gl{mo} dance.\gl{pst}\\
\glt `For every student $x$, the teacher(s) that $x$ invited danced.' \src{\Hx{2}{ex:shimoyama}}
\keioreading{for every student, the teacher that student invited danced; the teachers may differ from student to student.}
\exi{{11}}
\gll Dare-ka-ga subete-no hon-o yonda.\\
     who-\gl{ka}-\gl{nom} all-\gl{gen} book-\gl{acc} read.\gl{pst}\\
\glt `Someone read all the books.'
\keioreading{the books vary with the readers, as in item 1. If yes: does the reading survive when the subject is \textit{dare-ka hitori-ga} `some one person'? The second question separates genuine inverse scope from the cumulative reading that \textit{subete-no} N, like English \textit{all}, also allows.}
\exi{{12}}
\gll Taroo-dake-ga nani-mo tabe-nakat-ta.\\
     Taro-only-\gl{nom} what-\gl{mo} eat-\gl{neg}-\gl{pst}\\
\glt `Only Taro ate nothing.'
\keioreading{Taro, and nobody else, ate nothing: \textit{nani-mo} is licensed by the negation of its own clause, with the focused subject between the two.}
\end{exe}
\end{exercise}

\begin{exercise}{hierarchy}{Build the hierarchy}{2}{sec:hierarchy}
Take the blank tree below, the five questions of \Hx{2}{ex:parameters}, and seven languages (English, Greek, Japanese, Korean, Mandarin, Straits Salish, Mohawk), plus any language you know well. Place each language where its quantifiers belong. Some will want to go in two places at once. That is part of the point, and a language may be split between two leaves.

Two questions to guide the placing: do \textit{subete-no} N and \textit{dono} N-\textit{mo} land on different branches, as Korean \textit{motun} and \textit{nwukwu} do? And where do Mandarin \textit{mei} \ldots{} \textit{dou} and Straits Salish go? The blank tree:

\begin{center}
\begin{tikzpicture}[
  font=\footnotesize,
  node distance=3mm and 13mm,
  q/.style={draw=gttext!55, rounded corners=2pt, text width=3.9cm, inner sep=3pt, align=left},
  leaf/.style={draw=gtaccent!60, rounded corners=2pt, text width=2.9cm, minimum height=7mm, inner sep=3pt, align=left, fill=gtexbg},
  arr/.style={-{Stealth[length=1.6mm]}, semithick, gttext!55},
  lab/.style={font=\scriptsize\itshape, text=gttext!65, inner sep=1.5pt}
]
  \node[q] (p1) {\textbf{P1}\ \ Is [Q] lexicalised inside the nominal?};
  \node[q, below=of p1] (p2) {\textbf{P2}\ \ On a particle, or a nominal quantifier?};
  \node[q, below=of p2] (p3) {\textbf{P3}\ \ A general classifier layer between Q and N?};
  \node[q, below=of p3] (p4) {\textbf{P4}\ \ Do D and Q co-occur?};
  \node[leaf, below=of p4] (l5) {\ghost{Q alone, syncretic with D: \ldots}};
  \node[leaf, left=of p1] (l1) {\ghost{force only on the clause: \ldots}};
  \node[leaf, left=of p2] (l2) {\ghost{force on a particle: \ldots}};
  \node[leaf, left=of p3] (l3) {\ghost{Q with a classifier: \ldots}};
  \node[leaf, left=of p4] (l4) {\ghost{D and Q together: \ldots}};
  \draw[arr] (p1) -- node[lab, above] {no} (l1);
  \draw[arr] (p2) -- node[lab, above] {particle} (l2);
  \draw[arr] (p3) -- node[lab, above] {yes} (l3);
  \draw[arr] (p4) -- node[lab, above] {yes} (l4);
  \draw[arr] (p1) -- node[lab, right] {yes} (p2);
  \draw[arr] (p2) -- node[lab, right] {nominal quantifier} (p3);
  \draw[arr] (p3) -- node[lab, right] {no} (p4);
  \draw[arr] (p4) -- node[lab, right] {no} (l5);
  \node[q, anchor=north west, text width=7.9cm, draw=gtaccent!60, fill=gtexbg] (p5) at ([yshift=-5mm] l1.west |- l5.south) {\textbf{P5} (clausal)\ \ Is force fixed on the spine, and is that step pronounced?\\[2pt] \ghost{yes: \ldots}\\ \ghost{no: \ldots}};
\end{tikzpicture}
\end{center}
\end{exercise}

\section{Lecture 3: Scope in the box theory}

\begin{exercise}{worksheet}{Worksheet: derive it, break it}{3}{sec:worksheet}
Derive the sentence below, the first Japanese sentence of Lecture~1 (two quantifiers, one reading), in the notation of \Hx{3}{ex:deriv-child} and \Hx{3}{ex:deriv-kodomo}, filling in the workspace from W2 to W8. Mark every point at which the derivation has a choice (whether the subject boxes, and at which phase, \vs\ or C, the boxed object is discharged), and test each choice against Determinacy, Resource Restriction and Minimal Yield, as set out in \Hx{1}{ex:conditions}. At the foot, write in one sentence the choice this derivation had to make. Lecture~3 gives the filled table (\Hx{3}{ex:ws-deriv}); fill this one in before looking at it. The text's own list of the problems the derivation raises is in \Hs{3}{sec:five-problems}.

\begin{exe}
\exagain{3}{ex:ws-sentence}
\gll Dare-ka-ga dono hon-mo yonda.\\
     who-\gl{ka}-\gl{nom} which book-\gl{mo} read.\gl{pst}\\
\glt `Someone read every book.' \src{(someone $>$ every book only)}
\end{exe}

\begin{center}\footnotesize
\renewcommand{\arraystretch}{2.1}
\begingroup
\begin{tabularx}{\linewidth}{|l|L|P{1.9cm}|P{1.3cm}|P{1.5cm}|}
\hline
\textbf{Step} & \textbf{Workspace} & \textbf{Operation} & \textbf{A choice here?} & \textbf{Determinacy, RR, MY}\\ \hline
W1 & [\,\textit{dare}, \textit{ka}, \textit{dono}, \textit{hon}, \textit{mo}, \textit{yom-}, \vs, T, C\,] & items accessed from the lexicon & & \\ \hline
W2 & & & & \\ \hline
W3 & & & & \\ \hline
W4 & & & & \\ \hline
W5 & & & & \\ \hline
W6 & & & & \\ \hline
W7 & & & & \\ \hline
W8 & & & & \\ \hline
\end{tabularx}\endgroup
\end{center}

\textbf{The choice this derivation had to make:}\par\vspace{1.4em}
\noindent\rule{\linewidth}{0.4pt}\par\vspace{1.2em}
\noindent\rule{\linewidth}{0.4pt}
\end{exercise}

\section{Lecture 4: Indices after Inclusiveness}

\begin{exercise}{exits}{Which exit?}{4}{sec:percolation}
Before reading my own preference (\Hs{4}{sec:percolation}): which of the three exits from the percolation problem, \Hx{4}{ex:exit-sharing}, \Hx{4}{ex:exit-upward} and \Hx{4}{ex:exit-search}, would you take, and why?
\end{exercise}

% ======================================================================
%  Back matter
% ======================================================================
\backmatter
\keioreferences

% ======================================================================
%  Terms and abbreviations (back matter)
% ======================================================================
\addchap{Terms and abbreviations}

{\small\noindent The working vocabulary of the volume, with the place where each term is introduced or does its main work. Definitions are given in the sense in which the text uses them, which is not always the only sense in the literature.\par}

\begingroup
\small
\newcommand\glossterm[2]{\item[#1\enspace{\mdseries\sffamily\scriptsize\color{gtghost}#2}]}
\begin{description}[style=nextline, leftmargin=1.2em, font=\normalfont\bfseries, itemsep=3pt]
\glossterm{Anchored}{\Hs{4}{sec:testbed}}
  Said of an expression whose index must be valued by the context, as an indexed definite's is, rather than by a binder. Japanese \textit{kare} is anchored, which is why relations that carry no value cannot bind it.
\glossterm{Association}{\Hs{2}{sec:pins}}
  The relation between a particle (\textit{mo}, \textit{ka}) and the indeterminate it quantifies over. It shows path effects: an intervening operator can block it, a clause boundary need not.
\glossterm{Box, boxing}{\Hs{3}{sec:box}}
  Chomsky's (2024) proposal that an internally merged element is set aside, segregated from the ongoing derivation but consulted by later phase heads. In the version developed here a quantifier is boxed by Self-Merge, and the box is, to a first approximation, a syntactic Cooper store.
\glossterm{Conservativity}{\Hs{1}{sec:kuroda}}
  $Q(A)(B) \Leftrightarrow Q(A)(A \cap B)$: the restriction of a natural-language determiner is obligatory and privileged.
\glossterm{Cooper storage}{\Hs{3}{sec:cooper}}
  A semantic mechanism (Cooper 1983) that puts a quantifier's meaning on hold, leaves an indexed placeholder in its place, and retrieves the quantifier at a later node, where it takes scope.
\glossterm{Determinacy}{\Hs{1}{sec:merge}}
  An operation applies only to a uniquely determined target; if the search returns two candidates, it cannot apply.
\glossterm{Discharge}{\Hs{3}{sec:boxQR}}
  The step at which a boxed quantifier is introduced at a phase edge and identifies the verb's \uth\ token, fixing its scope. In the revised formulation it is External Merge of the boxed singleton, triggered by the phase head's Minimal Search, not a third operation.
\glossterm{Domain locality, path locality}{\Hs{2}{sec:localities}}
  Domain locality counts edges (a dependency fails across a clause or phase boundary); path locality counts interveners (it fails across an operator of a certain kind). Their separation, the divide, is the test between a syntactic and a semantic store.
\glossterm{Duality of semantics}{\Hs{1}{sec:duality}}
  External Merge yields argument structure; Internal Merge yields scope and discourse-related properties.
\glossterm{EM, IM}{\Hs{1}{sec:merge}}
  External Merge (both objects are members of the workspace) and Internal Merge (one is a term of the other): two cases of one operation, MERGE.
\glossterm{Existential closure}{\Hs{2}{sec:diesing}}
  The default existential force supplied by the clause to free variables in its nuclear scope (Heim, Diesing).
\glossterm{Externalisation}{\Hs{1}{sec:spines}}
  The mapping of the narrow syntax to the sensorimotor interface, and, on the working hypothesis of the seminar, the locus of cross-linguistic variation: languages differ in which steps of the computation they pronounce.
\glossterm{FC, Form Copy}{\Hs{4}{sec:formcopy}}
  The procedure that construes identical inscriptions as copies of one object, rather than repetitions of two, under identity, c-command and co-occurrence in a phase. Read in Lecture~4 as a theory of indices: what it computes is which tokens count as the same.
\glossterm{Fission}{\Hs{3}{sec:boxQR}}
  The two applications of MERGE, Self-Merge and then Internal Merge of the verb, that leave the singleton \{V\}, its \uth, in the internal argument position, as the variable; Form Copy identifies the two occurrences of V.
\glossterm{Indeterminate}{\Hs{2}{sec:indeterminates}}
  A pronoun (Japanese \textit{dare}, \textit{nani}; Korean \textit{nwukwu}, \textit{mwues}) that contributes a restriction and no force of its own; force is supplied by a particle (Kuroda 1965).
\glossterm{Inclusiveness}{\Hs{1}{sec:inclusiveness}}
  The objects the computation forms are made of the features of the lexical items that enter it; no indices, bar levels and the like are added. Lecture~4 revises it to a condition on tokens, relations and values.
\glossterm{Index}{\Hs{4}{sec:features-indices}}
  What identifies, as opposed to what types: a relation of sameness between tokens. An index is innocuous as notation and objectionable only when the computation consults it as a value.
\glossterm{LFS}{\Hs{4}{sec:testbed}}
  Language Faculty Science (Hoji 2015): a programme that separates what the grammar excludes from what is merely hard to get, by testing for predicted refutations.
\glossterm{Minimal Search}{\Hs{1}{sec:merge}}
  MERGE, Agree and labelling find their targets by the least costly search: the first, structurally closest candidate is found and the search stops.
\glossterm{Monostratal}{\Hs{4}{sec:labels}}
  Said of a grammar with a single level of syntactic representation and no derivation relating one level to another, such as HPSG or categorial grammar. The label fits LFG only in the sense that neither of its two syntactic structures is derived from the other.
\glossterm{MSMT}{\Hs{1}{sec:architecture}}
  \textit{Merge and the Strong Minimalist Thesis} (Chomsky et al. 2023).
\glossterm{NTC}{\Hs{1}{sec:inclusiveness}}
  The No Tampering Condition: MERGE(X, Y) leaves X and Y unchanged.
\glossterm{PA}{\Hs{1}{sec:inclusiveness}}
  Predicate Abstraction (Heim \& Kratzer 1998): the rule that interprets an index adjoined below a moved phrase as $\lambda$-abstraction over the corresponding trace.
\glossterm{PIC}{\Hs{1}{sec:cyclicity}}
  The Phase Impenetrability Condition: the domain of a phase head is inaccessible to operations outside the phase; only the head and its edge remain.
\glossterm{Q-configuration hypothesis}{\Hs{1}{sec:nominal}}
  A QP is any nominal extended projection in which a [Q] feature in the D-region composes with the predicative core; where [Q] is pronounced and where its force is fixed are independent.
\glossterm{QR}{\Hs{1}{sec:QR}}
  Quantifier Raising (May 1977, 1985): covert adjunction of a quantified phrase to a clausal projection, leaving a variable.
\glossterm{Resource Restriction, Minimal Yield}{\Hs{1}{sec:merge}}
  An operation adds as few new accessible objects to the workspace as possible; MERGE adds exactly one.
\glossterm{Retrieval}{\Hs{3}{sec:cooper}}
  The step at which a stored quantifier is taken out of the store and applied to the meaning computed so far. Korean and Japanese pronounce it, in their particles.
\glossterm{Rigidity}{\Hs{3}{sec:rigidity}}
  The absence of inverse scope in a language or construction. Conjecture~3 derives it from the order of discharge reading an order the grammar already has.
\glossterm{Scope Economy}{\Hs{1}{sec:QR}}
  QR applies only if it makes a semantic difference (Fox 2000).
\glossterm{Self-Merge}{\Hs{3}{sec:boxQR}}
  MERGE of an object with itself, MERGE($\alpha$, $\alpha$), yielding the singleton \{$\alpha$\}. It boxes a quantifier and fissions a verb.
\glossterm{Single-cycle}{\Hs{1}{sec:cyclicity}}
  Said of a grammar with one derivation, transferred to the interfaces phase by phase, and no separate covert cycle or level of LF (Chomsky 2004). The announcement of the lectures says \emph{monocyclic}.
\glossterm{Thetic, categorical}{\Hs{1}{sec:kuroda}}
  Kuroda's (1972) distinction, after Brentano and Marty, between a judgement that presents a state of affairs and one that predicates something of a subject; Japanese \textit{ga} and \textit{wa}.
\glossterm{Token, repetition, copy}{\Hs{4}{sec:tokens}}
  A token is a lexical item paired with a number $k$ that distinguishes separate selections of the item from the lexicon, $\langle \mathrm{LI}, k\rangle$ (\Hs{4}{sec:tokens}). Two tokens of one item are repetitions; two occurrences of one token, related by Internal Merge and construed by Form Copy, are copies.
\glossterm{Trace Conversion}{\Hs{1}{sec:copy}}
  Fox's (2002) rules that turn the lower copy of a moved quantifier into a definite description containing a variable, by Variable Insertion and Determiner Replacement.
\glossterm{\uth}{\Hs{3}{sec:boxQR}}
  The singleton \{V\} created by fission, an occurrence of the verb's own token, in which only the internal $\theta$-feature is visible. It is of type $e$, and it is the variable; Lecture~4 argues that it is the thematic index of earlier work, rebuilt under the revised Inclusiveness condition.
\glossterm{Unselective binding}{\Hs{2}{sec:unselective}}
  Quantification by an adverb or operator over all the free variables in its scope at once (Lewis 1975), as opposed to binding a single variable.
\glossterm{UPP}{\Hs{1}{sec:FQ}}
  Uniquely predicational possessive: a possessive with a quantified possessor that is possible only in predicative position (Tsoulas \& Woods 2019).
\glossterm{Weak crossover}{\Hs{4}{sec:wco}}
  The degradation of a bound reading when a quantifier or \textit{wh}-phrase crosses a pronoun that it would otherwise bind and that does not c-command the position crossed from (where the pronoun does c-command it, the crossover is strong); an open problem for the token view.
\end{description}
\endgroup

% Keio's announcement, as issued.
\IfFileExists{2026110708B.pdf}{\includepdf[pages=1]{2026110708B.pdf}}{}

\end{document}
